\chapter{Convolutional Models (CNNs)}\label{ch:cnn}

\glspl{CNN} represent a paradigmatic case of equivariant neural architectures, where the underlying mathematical structure dictates both their capabilities and their limitations. This chapter develops the mathematical foundations of convolutional equivariance: from its historical and empirical origins to its rigorous formalization, and the systematic extensions that arise when asking what happens when equivariance is required with respect to symmetry groups more general than translations.

We start from classical \glspl{CNN}, whose success in computer vision is mathematically explained by translation equivariance, but whose restriction to this single symmetry group also constitutes their fundamental limitation. Once these limitations are formalized and the connection with harmonic analysis is established, we systematically extend the convolution operation to increasingly general symmetry groups: finite groups, compact Lie groups such as $SO(2)$ and $SO(3)$, non-compact groups such as $SE(3)$, Riemannian manifolds without global group structure, and finally fiber bundles with local gauge equivariance. Each extension broadens the class of geometric domains that the architecture can efficiently exploit.

This progression from the initial restriction toward full generality is precisely the one predicted by the unifying framework of \autoref{ch:gdl}: \glspl{CNN} were identified there as instances of the \gls{MPNN} framework over regular grids (\autoref{prop:cnn_mpnn}), and this chapter develops in depth the mathematical foundations of that connection.

\section{Classical Convolutional Neural Networks}\label{sec:cnn_clasicas}

\glspl{CNN} arise in the context of automatic pattern recognition in images. Their origin dates back to the late 1980s and early 1990s, when Yann LeCun demonstrated their potential for handwritten digit recognition~\citep{lecun1989}.

The idea of convolutional networks is born from the techniques used to detect patterns in images. The technique was relatively simple: it consisted of a convolution with a kernel that detected a specific pattern in the image. For example, a kernel that detected horizontal, vertical, diagonal edges, etc. Mathematically, this can be translated into the following operation, where \(I\) is the image and \(K\) is the kernel:

\begin{equation}
    (I \star K)(x,y) = \sum_{u=-k}^{k}\sum_{v=-k}^{k} I(x+u, y+v)K(u,v)
    \label{eq:convolution2d}
\end{equation}

Strictly speaking, this operation is a \textit{cross-correlation}, not a convolution, since mathematical convolution uses $I(x-u, y-v)K(u,v)$ (see \autoref{def:discrete_convolution}). However, in the deep learning literature it is conventional to call it ``convolution'' because, when learning the kernel weights, the difference between convolution and cross-correlation is irrelevant: the learned kernel is simply flipped. We follow this standard convention hereafter.

A classic example illustrates the type of kernels that were designed by hand for image processing:

\begin{example}[Sobel Kernel]\label{ex:kernel_sobel}
The \textbf{Sobel kernels} detect horizontal and vertical edges through approximations of the partial derivatives:

\textbf{Horizontal Sobel} (detects vertical edges):
\[
K_x = \begin{pmatrix}
-1 & 0 & 1 \\
-2 & 0 & 2 \\
-1 & 0 & 1
\end{pmatrix}
\]

\textbf{Vertical Sobel} (detects horizontal edges):
\[
K_y = \begin{pmatrix}
-1 & -2 & -1 \\
0 & 0 & 0 \\
1 & 2 & 1
\end{pmatrix}
\]

These kernels approximate the partial derivatives $\frac{\partial I}{\partial x}$ and $\frac{\partial I}{\partial y}$ respectively, where $I(x,y)$ represents the image intensity. The gradient magnitude is computed as:
\[
|\nabla I| \approx \sqrt{(I \star K_x)^2 + (I \star K_y)^2}
\]
\end{example}

Kernels such as Sobel's exhibit revealing symmetry properties: edge detectors are antisymmetric with respect to the detection direction and sum to zero (eliminating direct-current components), while smoothing kernels are symmetric. These symmetries directly reflect the geometric transformations that each kernel detects.

The problem with using kernels was finding the specific one for each pattern to be detected; moreover, such patterns had to be simple. The solution proposed by LeCun was to use a neural network to learn the kernels. This idea can be taken a step further, not only learning kernels to detect simple patterns, but also to detect more complex patterns by applying kernels to the results of other kernels. That is, one can learn a series of kernels in the first layer of the neural network and continue adding layers of kernels on top of the previous ones to detect more complex patterns. The first layers learn simple strokes such as horizontal and vertical lines, but each subsequent layer learns more abstract concepts, such as cubes or rectangles, and so on until it is capable of recognizing complex objects such as a cat or a dog.

The explosion of convolutional networks and deep learning took place in 2012 when AlexNet won the ImageNet competition by a significant margin~\citep{alexnet}. The three necessary and relevant factors for the success of deep learning came together: a large-scale dataset, a neural network architecture that used deep learning techniques, and the processing power of \glspl{GPU}.

\begin{figure}[htbp]
    \centering
    \includegraphics[width=0.95\textwidth]{figures/alexnet_architecture}
    \caption[AlexNet architecture]{AlexNet architecture: a convolutional network with 8 layers (5 convolutional and 3 fully connected) that achieved a Top-5 error of 15.3\% on ImageNet ILSVRC-2012, a significant improvement over the 26.2\% of the second-best model.}
    \label{fig:alexnet}
\end{figure}

From a mathematical perspective, the central architectural principle of \glspl{CNN} is \textit{weight sharing}: the same kernels are applied at all spatial positions, which is precisely the realization of translation equivariance as an inductive bias. The local connectivity of each kernel, combined with the hierarchical composition of layers (from simple edges to abstract representations), allows the network to progressively build representations of increasing complexity. However, the restriction to the translation group also constitutes the fundamental limitation of these architectures, motivating the geometric extensions developed in the later sections of this chapter.

Convolutions can be applied not only to images, but can be extrapolated to any dimension. For example, in the case of time series, we can apply convolutions in one dimension, and in the case of 3D volumes, we can apply convolutions in three dimensions. The formulation is similar:

\begin{equation}
    (I * K)(x) = \sum_{u=-k}^{k} I(x+u)K(u)
    \label{eq:convolution1d}
\end{equation}

\begin{equation}
    (I * K)(x,y,z) = \sum_{u=-k}^{k}\sum_{v=-k}^{k}\sum_{w=-k}^{k} I(x+u, y+v, z+w)K(u,v,w)
    \label{eq:convolution3d}
\end{equation}

\subsection{Mathematical Foundations and Limitations of Classical CNNs}\label{subsec:fundamentos_limitaciones_cnn}

\subsubsection{Convolution as an Equivariant Operation}\label{subsubsec:convolucion_equivariante}

The fundamental property that makes \glspl{CNN} successful is their translation equivariance: when the input is translated, the output is translated in the same way. This property, which we will explore with a simple example, is the basis of their efficiency in image processing.

\begin{example}[Translation Equivariance in 1D]\label{ex:equivarianza_traslaciones_1d}
Consider a 1D signal $f: \mathbb{Z} \to \mathbb{R}$ and a kernel $K = [-1, 0, 1]$ (edge detector).
The translation action by $\tau \in \mathbb{Z}$ is given by $(T_\tau f)(n) = f(n-\tau)$.

Let $f = [1, 2, 1, 0, 0]$ and $\tau = 1$. Using definition~\eqref{eq:convolution1d} with zero extension outside the domain:
\begin{align}
(T_1 f) &= [0, 1, 2, 1, 0] \quad \text{(translated signal)} \\
(f * K) &= [2, 0, -2, -1, 0] \quad \text{(original convolution)} \\
((T_1 f) * K) &= [1, 2, 0, -2, -1] \quad \text{(convolution of translated signal)} \\
T_1((f * K)) &= [1, 2, 0, -2, -1] \quad \text{(translation of convolution)}
\end{align}

We observe that $((T_1 f) * K) = T_1((f * K))$, verifying the equivariance.
\end{example}

This equivariance property is fundamental to understanding both the strengths and the limitations of \glspl{CNN}. The formal concepts of equivariance and their generalization to other transformation groups are developed in detail in \autoref{ch:gdl}.

\begin{figure}[htbp]
	\centering
	\includegraphics[width=0.85\textwidth]{figures/convolution_operation}
	\caption[2D convolution operation]{2D convolution operation: a kernel $\mathbf{W}$ of size $3 \times 3$ slides over the input image $\mathbf{X}$, computing the local inner product at each position to produce the output activation map $(\mathbf{X} \star \mathbf{W})$.}
	\label{fig:convolution_operation}
\end{figure}

\begin{definition}[Discrete Convolution]\label{def:discrete_convolution}
For functions \(f, g: \mathbb{Z}^d \to \mathbb{R}\), the \textit{discrete convolution} is:
\[
(f * g)(x) = \sum_{a \in \mathbb{Z}^d} f(a)\,g(x - a)
\]
where the sum ranges over all points \(a\) of the domain \(\mathbb{Z}^d\). In practice, if \(f\) has finite support (as in a digital image), the sum reduces to the indices where \(f(a) \neq 0\).
\end{definition}

\begin{definition}[Continuous Convolution]\label{def:continuous_convolution}
For functions \(f, g: \mathbb{R}^d \to \mathbb{R}\), the \textit{continuous convolution} is:
\[
(f * g)(x) = \int_{\mathbb{R}^d} f(a)\,g(x - a) \, \mathrm{d}a
\]
where the integral is taken with respect to the Lebesgue measure over all of \(\mathbb{R}^d\), and exists when \(f, g \in L^1(\mathbb{R}^d)\) or under compact support conditions.
\end{definition}

The discrete convolution can be seen as a discretization of the continuous convolution, where the integral is approximated by a Riemann sum. In the context of digital image processing, we work with the discrete version, while in analog signals or theoretical analysis the continuous version is used.

\begin{proposition}[Algebraic Properties of Convolution]\label{prop:propiedades_convolucion}
The discrete convolution satisfies the following properties for functions
\(f, g, h: \mathbb{Z}^d \to \mathbb{R}\) with finite support:
\begin{enumerate}
    \item \textbf{Commutativity}: \(f * g = g * f\)
    \item \textbf{Associativity}: \((f * g) * h = f * (g * h)\)
    \item \textbf{Distributivity}: \(f * (g + h) = f * g + f * h\)
    \item \textbf{Identity element}: \(f * \delta = f\), where \(\delta\) is the Kronecker delta
\end{enumerate}
\end{proposition}

\begin{proof}
\textbf{(1) Commutativity.} Through the change of variable \(b = x - a\):
\begin{align}
(f * g)(x) &= \sum_{a \in \mathbb{Z}^d} f(a)\,g(x - a) = \sum_{b \in \mathbb{Z}^d} f(x - b)\,g(b) = (g * f)(x)
\end{align}

\textbf{(2) Associativity.} By direct application of the definition and reordering of summation (discrete Fubini):
\begin{align}
((f * g) * h)(x) &= \sum_{b \in \mathbb{Z}^d} (f * g)(b)\,h(x - b) = \sum_{b \in \mathbb{Z}^d} \left[\sum_{a \in \mathbb{Z}^d} f(a)\,g(b - a)\right] h(x - b) \nonumber \\
                  &= \sum_{a \in \mathbb{Z}^d} f(a) \sum_{b \in \mathbb{Z}^d} g(b - a)\,h(x - b)
\end{align}
With the change \(c = b - a\):
\begin{align}
&= \sum_{a \in \mathbb{Z}^d} f(a) \sum_{c \in \mathbb{Z}^d} g(c)\,h(x - a - c) = \sum_{a \in \mathbb{Z}^d} f(a)\,(g * h)(x - a) = (f * (g * h))(x)
\end{align}

\textbf{(3) Distributivity.} By linearity of the sum:
\begin{align}
(f * (g + h))(x) &= \sum_{a \in \mathbb{Z}^d} f(a)\,(g + h)(x - a) \nonumber \\
                  &= \sum_{a \in \mathbb{Z}^d} f(a)\,g(x - a) + \sum_{a \in \mathbb{Z}^d} f(a)\,h(x - a) = (f * g)(x) + (f * h)(x)
\end{align}

\textbf{(4) Identity element.} The Kronecker delta satisfies \(\delta(a) = 1\) if \(a = 0\) and \(\delta(a) = 0\) otherwise:
\begin{align}
(f * \delta)(x) &= \sum_{a \in \mathbb{Z}^d} f(a)\,\delta(x - a) = f(x)
\end{align}
since \(\delta(x - a) \neq 0\) only when \(a = x\).
\end{proof}

\begin{theorem}[Translation Equivariance of Convolution]\label{theorem:convolution_translation_equivariance}
The convolution operation is equivariant to the translation group \(T = \mathbb{Z}^d\):
\[
T_b(f * g) = (T_b f) * g = f * (T_b g)
\]
where \(T_b f(x) = f(x - b)\) is the translation by \(b\).
\end{theorem}

\begin{proof}
For the first equality, through the index change $a' = a - b$:
\begin{align}
[(T_b f) * g](x) &= \sum_a (T_b f)(a) g(x - a) = \sum_a f(a - b) g(x - a) \\
                  &= \sum_{a'} f(a') g(x - a' - b) = (f * g)(x - b) = [T_b(f * g)](x)
\end{align}
For the second equality:
\begin{align}
[T_b(f * g)](x) &= (f * g)(x - b) = \sum_a f(a)g((x-b) - a) \\
                 &= \sum_a f(a)g((x - a) - b) = \sum_a f(a)[T_b g](x - a) \\
                 &= [f * (T_b g)](x)
\end{align}
\end{proof}

Translation equivariance implies that the kernel weights must be shared spatially, which drastically reduces the number of parameters compared to a fully connected network. This fundamental property is generalized in \autoref{sec:group_convolutions}.

\subsubsection{Limitations of Classical CNNs}

\begin{example}[Lack of Rotational Equivariance]
Consider an image of a digit ``6'' and its $180^{\circ}$ rotation that looks like a ``9''. A standard \gls{CNN} must learn these transformations from the data through data augmentation, rather than incorporating them structurally. If $R_{90}$ denotes a $90^{\circ}$ rotation:
\[
f(R_{90}(\text{image})) \neq R_{90}(f(\text{image}))
\]
This illustrates how standard \glspl{CNN} are only equivariant to translations, not to other geometric transformations.
\end{example}

Although \glspl{CNN} have revolutionized computer vision, their design presents fundamental limitations that restrict their applicability. These limitations, which arise directly from their restriction to translation equivariance, can be classified into four main categories:

\begin{enumerate}
    \item \textbf{Symmetry limitations}: They are only equivariant to discrete translations $\mathbb{Z}^d$, requiring data augmentation for transformations such as rotations, reflections, or scalings. This is inefficient both computationally and statistically.

    \item \textbf{Topological limitations}: Designed for regular Euclidean grids, they are inadequate for data with non-trivial topologies such as graphs, point clouds, Riemannian manifolds, or simplicial complexes.

    \item \textbf{Locality limitations}: Their dependence on local structures implies that global relationships require multiple layers to be captured, limiting efficiency in domains where long-range dependencies are critical.

    \item \textbf{Receptive field limitations}: The linear growth of the receptive field with depth requires deep architectures to capture global patterns, introducing vanishing gradient problems and computational inefficiency.
\end{enumerate}

These systematic limitations have motivated the development of geometric extensions that constitute the core of this chapter. Each extension addresses specific sets of these limitations through the controlled generalization of the preserved symmetries:

\begin{itemize}
    \item \textbf{Group Convolutions} (\autoref{sec:group_convolutions}): Generalize equivariance from translations toward finite groups such as $D_4$ (rotations and reflections), addressing the symmetry limitations.
    \item \textbf{Spherical/Steerable \glspl{CNN}}: Extend equivariance to continuous groups such as $SO(2)$ and $SO(3)$, enabling the processing of spherical and three-dimensional data.
    \item \textbf{\glspl{CNN} on Manifolds} (\autoref{sec:manifold_cnns}): Overcome the topological limitations through generalization to curved surfaces and Riemannian manifolds.
    \item \textbf{Gauge Convolutions} (\autoref{sec:gauge_cnns}): Introduce equivariance to gauge transformations on fiber bundles, the most abstract generalization possible.
\end{itemize}

This systematic progression from restricted equivariance toward full generality constitutes the central narrative of the chapter and exemplifies the principles of \gls{GDL} in action.

\section{Harmonic Analysis and Spectral Convolutions}\label{sec:analisis_armonico}

While the geometric extensions of \glspl{CNN} address the limitations of symmetry, topology, and locality through the direct generalization of convolutional operations, harmonic analysis provides a dual and complementary perspective that reveals the deep mathematical structure underlying all these architectures. This section develops the theoretical framework that connects classical convolutions with their spectral generalizations, from the foundations of Fourier analysis to modern applications in irregular domains.

The spatial-spectral duality constitutes one of the most fundamental principles in applied mathematics: while convolution operations are naturally performed in the spatial domain, their algebraic structure is best understood in the frequency domain. The convolution theorem, $\mathcal{F}[f * g] = \mathcal{F}[f] \cdot \mathcal{F}[g]$, reveals that the complex convolution operation in space reduces to a simple pointwise multiplication in frequency. This perspective not only provides computational advantages through algorithms such as the \gls{FFT}, but also offers deep insights into what features neural networks learn and how more efficient architectures can be designed.

In the context of \gls{GDL}, harmonic analysis acquires fundamental importance by providing the mathematical tools needed to generalize convolutions to non-Euclidean domains. Group representation theory, developed within the framework of harmonic analysis, allows the design of convolutional layers that automatically preserve the specific symmetries of the data domain. Likewise, spectral analysis on graphs through the Laplacian provides the theoretical basis for spectral \glspl{GCN}, while wavelets offer natural hierarchical representations that align with modern deep architectures.

\subsection{Foundations of Fourier Analysis}\label{subsec:fundamentos_fourier}

Fourier analysis constitutes one of the most powerful mathematical tools for the study of periodic phenomena and frequency analysis. Its application in the context of convolutional neural networks allows a deep understanding of the spectral properties of filters and convolutional operations.

\begin{definition}[Continuous Fourier Transform]\label{def:transformada_fourier_continua}
Let $f: \mathbb{R}^d \to \mathbb{C}$ be an integrable function. The \textit{Fourier transform} of $f$ is defined as:
\begin{equation}
\mathcal{F}[f](\omega) = \hat{f}(\omega) = \int_{\mathbb{R}^d} f(x) e^{-2\pi i \omega \cdot x} dx
\end{equation}
where $\omega \in \mathbb{R}^d$ represents the frequency variable and $\omega \cdot x = \sum_{j=1}^d \omega_j x_j$ is the Euclidean dot product.
\end{definition}

\textbf{Existence Conditions.} The Fourier transform is well defined for functions in $L^1(\mathbb{R}^d)$ (absolutely integrable functions). For functions in $L^2(\mathbb{R}^d)$, the transform is defined as a mean-square limit.

\begin{definition}[Inverse Fourier Transform]\label{def:transformada_fourier_inversa}
If $\hat{f} \in L^1(\mathbb{R}^d)$, the \textit{inverse Fourier transform} is defined as:
\begin{equation}
\mathcal{F}^{-1}[\hat{f}](x) = f(x) = \int_{\mathbb{R}^d} \hat{f}(\omega) e^{2\pi i \omega \cdot x} d\omega
\end{equation}
\end{definition}

\begin{theorem}[Fourier Inversion Theorem]\label{thm:inversion_fourier}
If $f, \hat{f} \in L^1(\mathbb{R}^d)$, then:
\begin{equation}
f(x) = \mathcal{F}^{-1}[\mathcal{F}[f]](x) = \int_{\mathbb{R}^d} \hat{f}(\omega) e^{2\pi i \omega \cdot x} d\omega
\end{equation}
for almost every $x \in \mathbb{R}^d$. For the proof, see \citep[Theorem~8.26]{Folland1999}.
\end{theorem}

\begin{definition}[Convolution on $\mathbb{R}^d$]\label{def:convolucion_continua}
For functions $f, g: \mathbb{R}^d \to \mathbb{C}$, the \textit{convolution} is defined as:
\begin{equation}
(f * g)(x) = \int_{\mathbb{R}^d} f(y) g(x - y) dy = \int_{\mathbb{R}^d} f(x - y) g(y) dy
\end{equation}
when the integral converges.
\end{definition}

\subsubsection{Concrete Examples of Fourier Transforms}

The following examples illustrate the fundamental properties of the Fourier transform in contexts relevant to the analysis of CNNs.

\begin{example}[Transform of the Gaussian Function]\label{ex:fourier_gaussian}
For the Gaussian function $f(x) = e^{-\alpha x^2}$ with $\alpha > 0$:
\begin{equation}
\mathcal{F}[e^{-\alpha x^2}](\omega) = \sqrt{\frac{\pi}{\alpha}} e^{-\pi^2 \omega^2/\alpha}
\end{equation}

\textbf{Notable properties}:
\begin{itemize}
    \item The transform of a Gaussian is another Gaussian
    \item Dual localization: narrow function $\Leftrightarrow$ wide transform
    \item Minimizes the time-frequency uncertainty principle
    \item Fundamental in smoothing filters for CNNs
\end{itemize}
\end{example}

\begin{example}[Transform of the Step Function]\label{ex:fourier_step}
For the unit step function $\mathbf{1}_{[0,1]}(x)$:
\begin{equation}
\mathcal{F}[\mathbf{1}_{[0,1]}](\omega) = \frac{e^{-2\pi i \omega} - 1}{-2\pi i \omega} = \frac{\sin(\pi\omega)}{\pi\omega} e^{-\pi i \omega}
\end{equation}

\textbf{Interpretation for CNNs}:
\begin{itemize}
    \item Represents averaging filters (constant kernels)
    \item The sinc function in frequency indicates low-pass behavior
    \item Oscillations in frequency due to discontinuities in the spatial domain
\end{itemize}
\end{example}

\begin{example}[Transform of the Sinc Function]\label{ex:fourier_sinc}
For the sinc function $f(x) = \frac{\sin(\pi x)}{\pi x}$:
\begin{equation}
\mathcal{F}[\operatorname{sinc}(x)](\omega) = \mathbf{1}_{[-1/2,1/2]}(\omega)
\end{equation}

This step $\leftrightarrow$ sinc duality is fundamental:
\begin{itemize}
    \item Ideal low-pass filter in frequency $\rightarrow$ sinc response in space
    \item Impossible to implement exactly (infinite support)
    \item Inspires the design of approximate filters in spectral CNNs
\end{itemize}
\end{example}

\begin{example}[Transform of the Dirac Delta]\label{ex:fourier_delta}
For the delta distribution $\delta(x)$:
\begin{equation}
\mathcal{F}[\delta](\omega) = 1 \quad \text{and} \quad \mathcal{F}[1](\omega) = \delta(\omega)
\end{equation}

\textbf{Meaning for CNNs}:
\begin{itemize}
    \item Spatial delta = flat spectrum (all frequencies present)
    \item Spatial constant = frequency impulse (zero frequency only)
    \item Localization principle in the spectral analysis of filters
\end{itemize}
\end{example}

\textbf{Computational Implementation.} In practice, these transforms are computed through the discrete \gls{FFT}. The spectral analysis of common kernels reveals important patterns:

\textbf{Analysis process}:
\begin{enumerate}
    \item \textbf{Padding}: The kernel is extended with zeros to improve frequency resolution
    \item \textbf{2D \gls{FFT}}: The two-dimensional fast Fourier transform is computed
    \item \textbf{Spectrum centering}: fftshift is applied to center the low frequencies
    \item \textbf{Magnitude and phase analysis}: The components are separated for interpretation
\end{enumerate}

\textbf{Examples of typical kernels}:
\begin{itemize}
    \item \textbf{Sobel X}: Edge detector with high-pass response in the horizontal direction
    \item \textbf{Gaussian}: Low-pass filter with smooth and symmetric response
    \item \textbf{Laplacian}: Isotropic edge detector with high-pass response
\end{itemize}

\textbf{Connection with Wavelets.} The previous transforms motivate the use of wavelets in CNNs:
\begin{itemize}
    \item \textbf{Gaussians}: Basis for Gabor wavelets (time-frequency analysis)
    \item \textbf{Steps}: Related to Haar wavelets (edge features)
    \item \textbf{Sinc}: Foundation of Daubechies wavelets (perfect reconstruction)
\end{itemize}
The ability to analyze multiple scales simultaneously makes wavelets especially attractive for hierarchical CNN architectures.

\begin{figure}[htbp]
	\centering
	\includegraphics[width=0.85\textwidth]{figures/fourier_pairs}
	\caption[Fourier transform pairs]{Fourier transform pairs: correspondence between signals in the time domain ($t$) and their representations in the frequency domain ($\omega$). The localization-delocalization duality that underlies the uncertainty principle is observed.}
	\label{fig:fourier_pairs}
\end{figure}

\subsection{Fundamental Properties of Harmonic Analysis}\label{subsec:propiedades_fundamentales}

\begin{theorem}[Convolution Theorem in the Fourier Domain]\label{thm:convolucion_fourier}
For functions \(f, g: \mathbb{R}^d \to \mathbb{C}\):
\[
\mathcal{F}[f * g] = \mathcal{F}[f] \cdot \mathcal{F}[g]
\]
where \(\mathcal{F}\) denotes the Fourier transform.
\end{theorem}

\begin{proof}
Let $f, g \in L^1(\mathbb{R}^d)$. Applying Fubini's theorem:
\begin{align}
\mathcal{F}[f * g](\omega) &= \int_{\mathbb{R}^d} (f * g)(x) e^{-2\pi i\omega \cdot x} dx \\
&= \int_{\mathbb{R}^d} \left[\int_{\mathbb{R}^d} f(y)g(x-y) \, dy\right] e^{-2\pi i\omega \cdot x} dx \\
&= \int_{\mathbb{R}^d} f(y) \left[\int_{\mathbb{R}^d} g(x-y) e^{-2\pi i\omega \cdot x} dx\right] dy \\
&= \int_{\mathbb{R}^d} f(y) e^{-2\pi i\omega \cdot y} \mathcal{F}[g](\omega) \, dy \\
&= \mathcal{F}[f](\omega) \cdot \mathcal{F}[g](\omega)
\end{align}
where in the fourth step the change of variable $z = x - y$ and the definition of $\mathcal{F}[g]$ are used.
\end{proof}

\begin{corollary}[Computational Efficiency via FFT]
The convolution can be computed in \(O(n \log n)\) using:
\[
f * g = \mathcal{F}^{-1}[\mathcal{F}[f] \cdot \mathcal{F}[g]]
\]
\end{corollary}

\begin{proof}
By the previous convolution theorem, $\mathcal{F}[f * g] = \mathcal{F}[f] \cdot \mathcal{F}[g]$, hence $f * g = \mathcal{F}^{-1}[\mathcal{F}[f] \cdot \mathcal{F}[g]]$. For discrete signals of length $n$, the \gls{FFT} computes $\mathcal{F}[f]$ and $\mathcal{F}[g]$ each in $O(n \log n)$ operations, the pointwise multiplication $\mathcal{F}[f] \cdot \mathcal{F}[g]$ requires $O(n)$ operations, and the inverse transform $\mathcal{F}^{-1}$ has the same cost $O(n \log n)$. The total cost is $O(n \log n)$, compared to the $O(n^2)$ of the direct computation of the convolution.
\end{proof}

\begin{theorem}[Plancherel Identity]\label{thm:plancherel}
For functions $f, g \in L^2(\mathbb{R}^d)$, the Fourier transform preserves the inner product:
\begin{equation}
\langle f, g \rangle_{L^2} = \langle \hat{f}, \hat{g} \rangle_{L^2}
\end{equation}
In particular, $\|\hat{f}\|_{L^2} = \|f\|_{L^2}$ (isometry). See \citep[Theorem~8.29]{Folland1999}.
\end{theorem}

\textbf{Importance in Neural Networks.} The Plancherel identity guarantees that operations in the frequency domain preserve the energy of signals, which is fundamental for the numerical stability of spectral algorithms in CNNs.

\begin{definition}[Discrete Fourier Transform (DFT)]\label{def:dft}
For a finite sequence $\{x_n\}_{n=0}^{N-1}$, the \textbf{discrete Fourier transform} is defined as:
\begin{equation}
X_k = \sum_{n=0}^{N-1} x_n e^{-2\pi i kn/N}, \quad k = 0, 1, \ldots, N-1
\end{equation}
\end{definition}

\begin{definition}[Fast Fourier Transform (FFT)]\label{def:fft}
The \gls{FFT} algorithm computes the DFT in time $O(N \log N)$ instead of $O(N^2)$ through recursive factorization:
\begin{equation}
X_k = \sum_{n=0}^{N/2-1} x_{2n} e^{-2\pi i kn/(N/2)} + e^{-2\pi i k/N} \sum_{n=0}^{N/2-1} x_{2n+1} e^{-2\pi i kn/(N/2)}
\end{equation}
\end{definition}

\begin{theorem}[Nyquist-Shannon Sampling Theorem]\label{thm:nyquist_shannon}
Let $f: \mathbb{R} \to \mathbb{C}$ be a band-limited function with $\hat{f}(\omega) = 0$ for $|\omega| > W$. Then $f$ is completely determined by its samples at intervals of $1/(2W)$:
\begin{equation}
f(t) = \sum_{n=-\infty}^{\infty} f\left(\frac{n}{2W}\right) \operatorname{sinc}(2W(t - \frac{n}{2W}))
\end{equation}
where $\operatorname{sinc}(x) = \sin(\pi x)/(\pi x)$. This fundamental result was established by \citet{Shannon1948}.
\end{theorem}

\textbf{Nyquist Frequency and Aliasing.} The Nyquist frequency $f_N = W$ is the maximum frequency that can be represented without aliasing. In CNNs, downsampling (pooling) can introduce aliasing if appropriate anti-aliasing filters are not applied.

\begin{definition}[Fourier Series]\label{def:serie_fourier}
For a periodic function $f: \mathbb{R} \to \mathbb{C}$ with period $T$, the Fourier series is expressed as:
\begin{equation}
f(t) = \sum_{n=-\infty}^{\infty} c_n e^{2\pi i n t/T}
\end{equation}
where the coefficients are $c_n = \frac{1}{T} \int_0^T f(t) e^{-2\pi i n t/T} dt$.
\end{definition}

\begin{definition}[Short-Time Fourier Transform (STFT)]\label{def:stft}
To analyze non-stationary signals, it is defined as:
\begin{equation}
\text{STFT}_f(t, \omega) = \int_{-\infty}^{\infty} f(\tau) w(\tau - t) e^{-2\pi i \omega \tau} d\tau
\end{equation}
where $w$ is a window function (typically Gaussian or Hann).
\end{definition}

\textbf{Time-Frequency Uncertainty Principle.} The STFT is subject to the uncertainty principle: $\Delta t \cdot \Delta \omega \geq \frac{1}{4\pi}$, which implies a fundamental trade-off between temporal and frequency resolution.

\subsubsection{Numerical Examples of the Convolution Theorem}

The convolution theorem is fundamental to understanding the computational efficiency of spectral CNNs. The following examples illustrate its practical application.

\begin{example}[Fast Convolution via FFT: Complexity Analysis]\label{ex:convolution_fft_complexity}
Consider the convolution of an image $f \in \mathbb{R}^{N \times N}$ with a kernel $g \in \mathbb{R}^{K \times K}$.

\textbf{Direct method}:
\begin{equation}
(f * g)[i,j] = \sum_{u=0}^{K-1} \sum_{v=0}^{K-1} f[i-u, j-v] \cdot g[u,v]
\end{equation}
Complexity: $O(N^2 K^2)$

\textbf{Spectral method}:
\begin{align}
\hat{f} &= \text{FFT2D}(f) \quad \text{[Complexity: } O(N^2 \log N)\text{]} \\
\hat{g} &= \text{FFT2D}(\text{pad}(g, N)) \quad \text{[Complexity: } O(N^2 \log N)\text{]} \\
\hat{h} &= \hat{f} \odot \hat{g} \quad \text{[Complexity: } O(N^2)\text{]} \\
h &= \text{IFFT2D}(\hat{h}) \quad \text{[Complexity: } O(N^2 \log N)\text{]}
\end{align}
Total complexity: $O(N^2 \log N)$

\textbf{Break-even point}: The spectral method is more efficient when $K > O(\log N)$.
\end{example}

\begin{example}[Practical Implementation: Fast Convolution]\label{ex:fast_convolution_implementation}
Fast convolution through the \gls{FFT} follows these algorithmic steps:

\textbf{Spectral convolution algorithm}:
\begin{enumerate}
    \item \textbf{Sizing}: Compute the optimal size as a power of 2 for \gls{FFT} efficiency
    \item \textbf{Transformation}: Apply FFT2D to both the image and the kernel with appropriate padding
    \item \textbf{Multiplication}: Element-wise product in the frequency domain
    \item \textbf{Inversion}: IFFT2D to recover the result in the spatial domain
    \item \textbf{Cropping}: Extract the valid region of the result
\end{enumerate}

\textbf{Comparative performance analysis}:
For different image and kernel sizes, typical results show:
\begin{itemize}
    \item \textbf{Small images (64$\times$64)}: Direct method more efficient for kernels $< 7 \times 7$
    \item \textbf{Medium images (256$\times$256)}: \gls{FFT} advantageous for kernels $> 5 \times 5$
    \item \textbf{Large images (512$\times$512)}: \gls{FFT} consistently faster with speedups of 2-10x
    \item \textbf{Numerical error}: Typically $< 10^{-12}$ with double precision
\end{itemize}
\end{example}

\begin{example}[Numerical Error Analysis]\label{ex:numerical_error_analysis}
The analysis of errors in \gls{FFT} convolutions reveals important patterns of numerical stability:

\textbf{Analysis methodology}:
\begin{enumerate}
    \item \textbf{Test signals}: Use random matrices with normal distribution
    \item \textbf{Reference kernel}: Discrete Laplacian as a standard test case
    \item \textbf{Comparison}: Direct convolution vs. \gls{FFT} method
    \item \textbf{Metrics}: Relative error in L2 norm
\end{enumerate}

\textbf{Observed empirical results}:
\begin{itemize}
    \item \textbf{Typical relative error}: $10^{-12}$ to $10^{-14}$ with double precision
    \item \textbf{Scalability}: The error grows logarithmically with image size
    \item \textbf{Kernel stability}: Well-conditioned kernels minimize propagation
    \item \textbf{Accumulation pattern}: Errors concentrate in high-frequency components
\end{itemize}

\textbf{Practical implications}:
The high numerical precision (error $< 10^{-12}$) makes the \gls{FFT} method practically equivalent to the direct method for most deep learning applications, where single precision ($\sim 10^{-7}$) is standard.
\end{example}

\textbf{Advanced Optimizations.} In practical implementations of spectral CNNs:

\textbf{1. Use of FFTW (Fastest Fourier Transform in the West)}:
\begin{itemize}
    \item Hardware-specific optimization (SIMD, multi-core)
    \item Automatic planning of the optimal \gls{FFT} algorithm
    \item Speedups of 2-5x over basic implementations
\end{itemize}

\textbf{2. Overlap-Add/Overlap-Save}:
For convolutions with large kernels in CNNs, block-processing techniques are used:
\begin{itemize}
    \item Divide the image into overlapping blocks
    \item FFT of small blocks (more cache-efficient)
    \item Reconstruction with overlap handling
\end{itemize}

\textbf{3. GPU Acceleration}:
\begin{itemize}
    \item cuFFT for CUDA implementations
    \item Use of shared memory to minimize transfers
    \item Batching of multiple simultaneous convolutions
\end{itemize}

\subsection{Harmonic Analysis on Groups}\label{subsec:analisis_armonico_grupos}

Harmonic analysis on groups generalizes the classical Fourier transform to more general spaces with group structure, providing the fundamental mathematical tools for geometric convolutions in neural networks.

We recall that the \textit{Haar measure} of a locally compact group $G$ is the unique Borel measure (up to a constant) invariant under left translations: $\mu(gE) = \mu(E)$ for all $g \in G$ (\autoref{def:medida_haar_gdl} of \autoref{ch:gdl}).\label{def:medida_haar}

\textbf{Existence and Uniqueness.} Haar's theorem guarantees that every Haar measure exists and is unique up to a multiplicative constant. For finite groups, the Haar measure is the counting measure.

The fundamental concepts of group representations (including the definition of representation (\autoref{def:representacion_grupo}), irreducible representations (\autoref{def:representacion_irreducible}), and the Peter-Weyl theorem (\autoref{thm:peter_weyl})) are developed in \autoref{ch:gdl}, Section~\ref{subsec:teoria_grupos_representaciones}. Here we briefly recall the specific tools of harmonic analysis that we will use for CNNs.

The \textit{character} of a representation $\rho: G \to GL(\mathbb{C}^d)$ is the function $\chi(g) = \operatorname{tr}(\rho(g))$, and the irreducible characters satisfy the orthogonality relation $\frac{1}{|G|} \sum_{g \in G} \chi_i(g) \overline{\chi_j(g)} = \delta_{ij}$. From these tools, the Fourier transform on finite groups is defined: for $f: G \to \mathbb{C}$ and an irreducible representation $\rho_i$ of dimension $d_i$,
\begin{equation}\label{eq:fourier_grupo_finito}
\hat{f}(\rho_i) = \sum_{g\in G} f(g)\rho_i(g) \in \mathbb{C}^{d_i \times d_i},
\end{equation}
with inversion $f(g) = \frac{1}{|G|} \sum_i d_i \operatorname{tr}(\hat{f}(\rho_i) \rho_i(g^{-1}))$. The fundamental result for CNNs is the \textbf{convolution theorem}: the Fourier transform converts group convolution into a matrix product,
\begin{equation}\label{eq:convolucion_fourier_grupos}
\widehat{f *_G h}(\rho_i) = \hat{f}(\rho_i) \hat{h}(\rho_i),
\end{equation}
generalizing the classical DFT result.

\begin{example}[Cyclic Group $\mathbb{Z}_n$]\label{ex:grupo_ciclico_fourier}
For the cyclic group of order $n$, the irreducible representations are one-dimensional:
\begin{equation}
\rho_k(j) = e^{-2\pi i jk/n}, \quad k = 0, 1, \ldots, n-1
\end{equation}
The Fourier transform reduces to the classical DFT:
\begin{equation}
\hat{f}(k) = \sum_{j=0}^{n-1} f(j) e^{-2\pi i jk/n}
\end{equation}
\end{example}

\begin{example}[Dihedral Group $D_n$]\label{ex:grupo_diedrico_fourier}
For the dihedral group of order $2n$ (symmetries of an $n$-gon), the representations include:
\begin{itemize}
    \item \textbf{Trivial representations}: $\rho(g) = 1$ for all $g$
    \item \textbf{Sign representations}: $\rho(r^j) = 1$, $\rho(sr^j) = -1$
    \item \textbf{Two-dimensional representations}: For rotations and reflections
\end{itemize}
\end{example}

\subsubsection{Detailed Examples of Harmonic Analysis on Finite Groups}

The following examples provide explicit computations that illustrate the practical construction of representations and Fourier transforms for groups relevant in GDL.

\begin{example}[Cyclic Group $\mathbb{Z}_4$: Explicit Computations]\label{ex:grupo_z4_detallado}
The cyclic group $\mathbb{Z}_4 = \{0, 1, 2, 3\}$ with addition modulo 4 represents $90^\circ$ rotations.

\textbf{Irreducible representations}:
The 1D representations are homomorphisms $\rho_k: \mathbb{Z}_4 \to \mathbb{C}^*$:
\begin{align}
\rho_0(j) &= 1 \quad \text{(trivial representation)} \\
\rho_1(j) &= i^j \quad \text{where } i = e^{2\pi i/4} \\
\rho_2(j) &= (-1)^j \\
\rho_3(j) &= (-i)^j
\end{align}

\textbf{Character matrix}:
\begin{equation}
\chi = \begin{bmatrix}
1 & 1 & 1 & 1 \\
1 & i & -1 & -i \\
1 & -1 & 1 & -1 \\
1 & -i & -1 & i
\end{bmatrix}
\end{equation}

\textbf{Discrete Fourier transform}:
For a function $f: \mathbb{Z}_4 \to \mathbb{C}$, $f = (f_0, f_1, f_2, f_3)^T$:
\begin{equation}
\hat{f} = \chi f = \begin{bmatrix}
f_0 + f_1 + f_2 + f_3 \\
f_0 + if_1 - f_2 - if_3 \\
f_0 - f_1 + f_2 - f_3 \\
f_0 - if_1 - f_2 + if_3
\end{bmatrix}
\end{equation}

\textbf{Application in G-CNNs}:
\begin{itemize}
    \item Each representation $\rho_k$ corresponds to a channel in convolutions equivariant to $90^\circ$ rotations
    \item The filters are expressed as linear combinations of the representations
    \item Equivariance is automatically preserved by the group structure
\end{itemize}
\end{example}

\begin{example}[Dihedral Group $D_3$: Symmetries of the Triangle]\label{ex:grupo_d3_detallado}
$D_3$ contains the 6 symmetries of the equilateral triangle: 3 rotations and 3 reflections.

\textbf{Group elements}:
\begin{itemize}
    \item $e$: identity
    \item $r, r^2$: rotations of $120^\circ, 240^\circ$
    \item $s, sr, sr^2$: reflections about the perpendicular bisectors
\end{itemize}

\textbf{Irreducible representations}:
\begin{enumerate}
    \item \textbf{Trivial} $\rho^{(1)}$: $\rho^{(1)}(g) = 1$ for all $g$
    \item \textbf{Alternating} $\rho^{(2)}$:
    \begin{align}
    \rho^{(2)}(r^k) &= 1, \quad k = 0,1,2 \\
    \rho^{(2)}(sr^k) &= -1, \quad k = 0,1,2
    \end{align}
    \item \textbf{Two-dimensional standard} $\rho^{(3)}$:
    \begin{align}
    \rho^{(3)}(e) &= \begin{bmatrix} 1 & 0 \\ 0 & 1 \end{bmatrix} \\
    \rho^{(3)}(r) &= \begin{bmatrix} -1/2 & -\sqrt{3}/2 \\ \sqrt{3}/2 & -1/2 \end{bmatrix} \\
    \rho^{(3)}(r^2) &= \begin{bmatrix} -1/2 & \sqrt{3}/2 \\ -\sqrt{3}/2 & -1/2 \end{bmatrix} \\
    \rho^{(3)}(s) &= \begin{bmatrix} 1 & 0 \\ 0 & -1 \end{bmatrix}
    \end{align}
\end{enumerate}

\textbf{Character table}:
\begin{equation}
\begin{array}{c|ccc}
 & e & r, r^2 & s, sr, sr^2 \\
\hline
\rho^{(1)} & 1 & 1 & 1 \\
\rho^{(2)} & 1 & 1 & -1 \\
\rho^{(3)} & 2 & -1 & 0
\end{array}
\end{equation}

\textbf{Computational implementation}:
The practical implementation of the dihedral group $D_3$ requires:

\textbf{Matrix representations}:
\begin{itemize}
    \item \textbf{Identity}: $2 \times 2$ identity matrix
    \item \textbf{Rotations}: Rotation matrices for $120^\circ$ and $240^\circ$
    \item \textbf{Reflections}: Combination of a basic reflection with rotations
\end{itemize}

\textbf{Fourier transform on $D_3$}:
For a function $f: D_3 \to \mathbb{C}$, the computational process involves:
\begin{enumerate}
    \item \textbf{Trivial projection}: Average of all function values
    \item \textbf{Alternating projection}: Difference between rotations and reflections
    \item \textbf{Two-dimensional projection}: A more complex linear combination
\end{enumerate}

\textbf{Analysis example}:
For the characteristic function of rotations $f(g) = 1$ if $g$ is a rotation, $f(g) = 0$ if it is a reflection:
\begin{itemize}
    \item Trivial component: $\frac{3}{6} = 0.5$
    \item Alternating component: $\frac{3-0}{6} = 0.5$
    \item 2D component: Vector with specific structure
\end{itemize}
\end{example}

\textbf{Connection with Practical G-CNNs.} The previous computations translate directly into neural network architectures:

\textbf{For $\mathbb{Z}_4$-CNNs}:
\begin{itemize}
    \item 4 filter orientations $(0^\circ, 90^\circ, 180^\circ, 270^\circ)$
    \item Each output channel corresponds to an irreducible representation
    \item The convolution is implemented as a sum over the group orbits
\end{itemize}

\textbf{For $D_3$-CNNs}:
\begin{itemize}
    \item 6 filter transformations (3 rotations $\times$ 2 reflections)
    \item Special handling of the 2D representation (coupled channels)
    \item Application in the analysis of triangular symmetries (crystallography, materials)
\end{itemize}

\textbf{Efficient implementation}:
The fast Fourier transform on finite groups is implemented through:
\begin{equation}
\text{Complexity}: O(|G|^{3/2}) \text{ vs } O(|G|^2) \text{ for the direct DFT}
\end{equation}

\begin{example}[Connection with Wavelets on Groups]\label{ex:wavelets_grupos_finitos}
Finite groups also admit wavelet analysis through appropriate mother functions.

\textbf{Mother wavelet on $\mathbb{Z}_4$}:
\begin{equation}
\psi(j) = \delta_{j,1} - \delta_{j,3} \quad \text{(difference function)}
\end{equation}

\textbf{Wavelet family}:
\begin{equation}
\psi_{k,j}(m) = \psi((m - j) \bmod 4) \cdot \omega_4^{km}
\end{equation}
where $\omega_4 = e^{2\pi i/4}$.

This construction enables multiresolution analysis in discrete spaces with rotational symmetry.
\end{example}

\begin{definition}[Fourier Transform on Compact Groups]\label{def:fourier_grupos_compactos}
For a compact group $G$ and a function $f \in L^2(G)$, the Fourier transform is defined through the Peter-Weyl theorem:
\begin{equation}
f = \sum_{\lambda \in \widehat{G}} d_\lambda \operatorname{tr}(\hat{f}(\lambda) U^\lambda)
\end{equation}
where $\widehat{G}$ indexes the irreducible representations.
\end{definition}

\subsubsection{Harmonic Analysis on SO(2) and SO(3)}

\begin{definition}[Circular Harmonics on SO(2)]\label{def:armonicos_circulares}
For the group of planar rotations $SO(2) \cong S^1$, the irreducible representations are:
\begin{equation}
U^n(\theta) = e^{in\theta}, \quad n \in \mathbb{Z}
\end{equation}
A function $f: S^1 \to \mathbb{C}$ is expanded as:
\begin{equation}
f(\theta) = \sum_{n=-\infty}^{\infty} c_n e^{in\theta}
\end{equation}
\end{definition}

\begin{definition}[Spherical Harmonics on SO(3)]\label{def:armonicos_esfericos}
For the group of three-dimensional rotations $SO(3)$, functions on the sphere $S^2$ are expanded in spherical harmonics:
\begin{equation}
f(\theta, \phi) = \sum_{l=0}^{\infty} \sum_{m=-l}^{l} c_{lm} Y_l^m(\theta, \phi)
\end{equation}
where $Y_l^m$ are the spherical harmonics of degree $l$ and order $m$.
\end{definition}

\textbf{Applications of the Peter-Weyl Theorem.} The Peter-Weyl theorem (\autoref{thm:peter_weyl}, \autoref{ch:gdl}) provides the theoretical basis for:
\begin{itemize}
    \item The channel decomposition in G-CNNs
    \item The design of equivariant filters on spheres
    \item The efficient parametrization of convolutions on Lie groups
\end{itemize}

\subsection{Spectral Analysis on Graphs}\label{subsec:analisis_espectral_grafos}

Spectral analysis on graphs extends the ideas of Fourier analysis to irregular domains, providing the mathematical foundations for spectral convolutions in neural networks on graphs. The theoretical foundations of the graph Laplacian, including definitions and spectral properties, are developed in \autoref{ch:gdl}, Section~\ref{subsec:geometria_topologica_grafos}. For the Laplacian matrix $L = D - A$ and its normalized variants, see \autoref{def:laplaciano_grafo} and \autoref{prop:propiedades_espectrales_laplaciano_grafos}.

In this section we focus on the applications of spectral analysis to the design of convolutions on graphs, assuming familiarity with the Laplacian $L$ and its normalized version $\mathcal{L} = D^{-1/2} L D^{-1/2}$.

\begin{theorem}[Spectral Theorem for the Laplacian]\label{thm:espectral_laplaciano}
Let $L$ be the Laplacian matrix of a connected graph. Then:
\begin{enumerate}
    \item The eigenvalues satisfy $0 = \lambda_0 < \lambda_1 \leq \lambda_2 \leq \cdots \leq \lambda_{n-1}$
    \item The eigenvector associated with $\lambda_0 = 0$ is $\phi_0 = \frac{1}{\sqrt{n}} \mathbf{1}$
    \item The eigenvectors $\{\phi_k\}_{k=0}^{n-1}$ form an orthonormal basis of $\mathbb{R}^n$
    \item The multiplicity of $\lambda_0$ equals the number of connected components
\end{enumerate}
For the proof, see \citep[Chapters~1--2]{Chung1997} and \autoref{prop:propiedades_espectrales_laplaciano_grafos}.
\end{theorem}

The Fourier transform on graphs (\autoref{def:transformada_fourier_grafos}, \autoref{ch:gdl}) projects signals onto the eigenvectors of the Laplacian: $\hat{f}(\lambda_k) = \langle \mathbf{f}, \phi_k \rangle$, with matrix form $\hat{f} = U^T f$. From this transform, spectral convolutions are defined:

\begin{definition}[Spectral Convolution on Graphs]\label{def:convolucion_espectral_grafos}
For a function $f: V \to \mathbb{R}$ and a spectral filter $g: \mathbb{R}_+ \to \mathbb{R}$:
\begin{equation}
(g \star f)(v_i) = \sum_{k=0}^{n-1} g(\lambda_k) \hat{f}(\lambda_k) \phi_k(v_i)
\end{equation}
In matrix form: $g \star f = \Phi g(\Lambda) \Phi^T f$, where $\Phi = [\phi_0, \ldots, \phi_{n-1}]$ and $\Lambda = \operatorname{diag}(\lambda_0, \ldots, \lambda_{n-1})$.
\end{definition}

\begin{theorem}[Frequency Interpretation]\label{thm:interpretacion_frecuencial_grafos}
The eigenvalues of the Laplacian have a frequency interpretation:
\begin{itemize}
    \item \textbf{Low frequencies} ($\lambda_k$ small): Eigenvectors vary smoothly over the graph
    \item \textbf{High frequencies} ($\lambda_k$ large): Eigenvectors oscillate rapidly between adjacent vertices
    \item \textbf{Low-pass filters}: $g(\lambda) \approx 1$ for small $\lambda$, $g(\lambda) \approx 0$ for large $\lambda$
\end{itemize}
\end{theorem}

\begin{proof}
The interpretation is justified through the Rayleigh quotient. For a signal $\mathbf{f} \in \mathbb{R}^n$ on the graph, the total variation is measured as:
\[
\mathbf{f}^T L \mathbf{f} = \sum_{(i,j) \in E} w_{ij}(f_i - f_j)^2
\]
The $k$-th eigenvector $\phi_k$ minimizes this total variation subject to orthogonality with the previous eigenvectors: $\lambda_k = \min_{\mathbf{f} \perp \phi_0, \ldots, \phi_{k-1}} \frac{\mathbf{f}^T L \mathbf{f}}{\mathbf{f}^T \mathbf{f}}$. Hence eigenvectors with small eigenvalues minimize the variation between adjacent nodes (smooth variation, low frequency), while large eigenvalues correspond to eigenvectors with high variation (rapid oscillations, high frequency).
\end{proof}

\begin{figure}[htbp]
    \centering
    \includegraphics[width=0.85\textwidth]{figures/cnn_graph_kernel_spectral}
    \caption[From CNN kernel to spectral representation on graphs]{Transformation of a CNN kernel to a spectral representation on graphs. The spatial kernel of a traditional CNN (left) can be interpreted as a spectral filter over the graph of the regular grid (right), where the eigenvalues of the Laplacian determine the frequencies and the filter modulates each frequency component.}
    \label{fig:cnn_graph_kernel_spectral}
\end{figure}

\begin{theorem}[Spectral Localization]\label{thm:localizacion_espectral}
If a spectral filter $g(\lambda)$ satisfies the decay conditions:
\begin{equation}
|g^{(k)}(\lambda)| \leq C_k (1 + \lambda)^{-k-\epsilon}
\end{equation}
for some $\epsilon > 0$ and constants $C_k$, where $g^{(k)}$ denotes the $k$-th derivative, then $g$ produces operations with local support on the graph.
\end{theorem}

\begin{proof}[Proof sketch]
The filter kernel in the graph domain is $K(i,j) = \sum_k g(\lambda_k) \phi_k(i) \phi_k(j)$, where $\{\lambda_k, \phi_k\}$ are the eigenpairs of the Laplacian. The decay condition on $g^{(k)}$ implies, through generalized integration by parts, that the kernel $K(i,j)$ decays exponentially with the geodesic distance $d(i,j)$ in the graph. Formally, for nodes $i,j$ at distance $d(i,j) > K$ one has $|K(i,j)| \leq C \cdot \alpha^{d(i,j)}$ for some $\alpha < 1$, which guarantees that the filtering operation is essentially local. See \citep[Theorem~5.5]{hammond2011wavelets} for the complete proof in the context of wavelets on graphs.
\end{proof}

\subsubsection{Polynomial Approximation of Filters}

To make spectral convolutions computationally efficient, filters are approximated by polynomials, avoiding the explicit diagonalization of the Laplacian.

\begin{definition}[Polynomial Filters]\label{def:filtros_polinomiales}
A spectral filter $g(\lambda)$ is approximated as:
\begin{equation}
g(\lambda) \approx \sum_{k=0}^K \theta_k p_k(\lambda)
\end{equation}
where $\{p_k\}$ is a polynomial basis and $\{\theta_k\}$ are trainable parameters.
\end{definition}

\begin{example}[Chebyshev Approximation]\label{ex:chebyshev_aproximacion}
Using Chebyshev polynomials $T_k$ on the interval $[-1, 1]$:
\begin{equation}
g(\tilde{\lambda}) = \sum_{k=0}^K \theta_k T_k(\tilde{\lambda})
\end{equation}
where $\tilde{\lambda} = 2\lambda/\lambda_{\max} - 1$ rescales the eigenvalues.
\end{example}

\begin{theorem}[Computational Efficiency]\label{thm:eficiencia_chebyshev}
The Chebyshev approximation allows computing $g(L)f$ in time $O(K|E|)$ without diagonalizing $L$, using the recurrence:
\begin{align}
\bar{x}_0 &= f \\
\bar{x}_1 &= \tilde{L} f \\
\bar{x}_k &= 2\tilde{L}\bar{x}_{k-1} - \bar{x}_{k-2}, \quad k \geq 2
\end{align}
where $\tilde{L} = 2L/\lambda_{\max} - I$.
\end{theorem}

\begin{proof}
By the Chebyshev recurrence, $T_k(\tilde{L})f$ is computed inductively: $\bar{x}_0 = f$ requires $O(n)$, $\bar{x}_1 = \tilde{L}f$ requires a sparse matrix-vector multiplication with cost $O(|E|)$ (since $\tilde{L}$ has at most $2|E| + n$ nonzero entries), and each step $\bar{x}_k = 2\tilde{L}\bar{x}_{k-1} - \bar{x}_{k-2}$ likewise requires $O(|E|)$. Since $K$ recurrence steps are performed, the total cost is $O(K|E|)$. Note that it is not necessary to compute the eigenvalues or eigenvectors of $L$, since the recurrence operates directly with the sparse matrix $\tilde{L}$.
\end{proof}

\subsubsection{Complete Example: Spectral Analysis of a Small Graph}

The following example illustrates all the previous concepts through explicit computations on a 4-node graph.

\begin{example}[Laplacian of a Triangular Graph with an Isolated Vertex]\label{ex:laplaciano_triangulo_4_nodos}
Consider the graph $G$ with 4 vertices and edges $E = \{(1,2), (1,3), (2,3), (3,4)\}$:

\begin{center}
\begin{tikzpicture}[scale=1.2]
    \node[circle, draw, fill=gdlblue!20] (1) at (0, 1) {1};
    \node[circle, draw, fill=gdlblue!20] (2) at (-1, -0.5) {2};
    \node[circle, draw, fill=gdlblue!20] (3) at (1, -0.5) {3};
    \node[circle, draw, fill=gdlgreen!20] (4) at (2.5, 0) {4};

    \draw[thick] (1) -- (2);
    \draw[thick] (1) -- (3);
    \draw[thick] (2) -- (3);
    \draw[thick] (3) -- (4);
\end{tikzpicture}
\end{center}

\textbf{1. Adjacency matrix}:
\begin{equation}
A = \begin{bmatrix}
0 & 1 & 1 & 0 \\
1 & 0 & 1 & 0 \\
1 & 1 & 0 & 1 \\
0 & 0 & 1 & 0
\end{bmatrix}
\end{equation}

\textbf{2. Degree matrix}:
\begin{equation}
D = \begin{bmatrix}
2 & 0 & 0 & 0 \\
0 & 2 & 0 & 0 \\
0 & 0 & 3 & 0 \\
0 & 0 & 0 & 1
\end{bmatrix}
\end{equation}

\textbf{3. Unnormalized Laplacian}:
\begin{equation}
L = D - A = \begin{bmatrix}
2 & -1 & -1 & 0 \\
-1 & 2 & -1 & 0 \\
-1 & -1 & 3 & -1 \\
0 & 0 & -1 & 1
\end{bmatrix}
\end{equation}

\textbf{4. Computation of eigenvalues and eigenvectors}:
The characteristic polynomial is:
\begin{equation}
\det(L - \lambda I) = \lambda(\lambda - 1)(\lambda - 3)(\lambda - 4)
\end{equation}

The eigenvalues are: $\lambda_0 = 0, \lambda_1 = 1, \lambda_2 = 3, \lambda_3 = 4$.

\textbf{Eigenvectors (computed and normalized)}:
\begin{align}
\phi_0 &= \tfrac{1}{2}(1, 1, 1, 1)^T \quad \text{(constant vector, $\lambda_0 = 0$)} \\
\phi_1 &= \tfrac{1}{\sqrt{6}}(-1, -1, 0, 2)^T \quad \text{(smooth variation, $\lambda_1 = 1$)} \\
\phi_2 &= \tfrac{1}{\sqrt{2}}(1, -1, 0, 0)^T \quad \text{(1-2 contrast, $\lambda_2 = 3$)} \\
\phi_3 &= \tfrac{1}{2\sqrt{3}}(1, 1, -3, 1)^T \quad \text{(high frequency, $\lambda_3 = 4$)}
\end{align}

It is verified that $L\phi_k = \lambda_k \phi_k$, $\|\phi_k\| = 1$ and $\langle \phi_i, \phi_j \rangle = \delta_{ij}$ for all $i, j$.

\textbf{5. Frequency interpretation}:
\begin{itemize}
    \item $\phi_0$ (DC): constant function over the entire graph
    \item $\phi_1$ (low frequency): smooth variation, vertex 4 contrasts with vertices 1-2
    \item $\phi_2$ (medium frequency): contrast between vertices 1 and 2
    \item $\phi_3$ (high frequency): rapid oscillation, vertex 3 contrasts strongly with its neighbors
\end{itemize}

\textbf{6. Spectral convolution - Practical example}:
Let $f = (1, 0, -1, 2)^T$ be a signal on the graph and $g(\lambda) = e^{-\lambda/2}$ a low-pass filter.

\textbf{Fourier transform of $f$}:
\begin{align}
\hat{f}_0 &= \phi_0^T f = \tfrac{1}{2}(1 + 0 - 1 + 2) = 1 \\
\hat{f}_1 &= \phi_1^T f = \tfrac{1}{\sqrt{6}}(-1 + 0 + 0 + 4) = \tfrac{3}{\sqrt{6}} = \tfrac{\sqrt{6}}{2} \\
\hat{f}_2 &= \phi_2^T f = \tfrac{1}{\sqrt{2}}(1 - 0 + 0 + 0) = \tfrac{1}{\sqrt{2}} = \tfrac{\sqrt{2}}{2} \\
\hat{f}_3 &= \phi_3^T f = \tfrac{1}{2\sqrt{3}}(1 + 0 + 3 + 2) = \tfrac{6}{2\sqrt{3}} = \sqrt{3}
\end{align}

\textbf{Filter application}:
\begin{align}
g(\lambda_0) \hat{f}_0 &= e^{0} \cdot 1 = 1 \\
g(\lambda_1) \hat{f}_1 &= e^{-1/2} \cdot \tfrac{\sqrt{6}}{2} \approx 0.74 \\
g(\lambda_2) \hat{f}_2 &= e^{-3/2} \cdot \tfrac{\sqrt{2}}{2} \approx 0.16 \\
g(\lambda_3) \hat{f}_3 &= e^{-2} \cdot \sqrt{3} \approx 0.23
\end{align}

\textbf{Filtered signal}:
\begin{equation}
f_{\text{filtered}} = \sum_{k=0}^{3} g(\lambda_k)\hat{f}_k \phi_k \approx (0.38, 0.15, 0.30, 1.17)^T
\end{equation}

\textbf{Interpretation}: The low-pass filter has smoothed the signal, especially reducing the high-frequency components ($\lambda_2 = 3$ and $\lambda_3 = 4$) while preserving the constant component.
\end{example}

\textbf{Computational Implementation of the Example.} The computational verification of the example follows this methodology:

\textbf{Graph construction}:
\begin{enumerate}
    \item Define 4 nodes with specific connectivity: $(1,2)$, $(1,3)$, $(2,3)$, $(3,4)$
    \item Compute adjacency and Laplacian matrices using standard libraries
    \item Verify basic properties: symmetry, positive semi-definiteness
\end{enumerate}

\textbf{Spectral analysis}:
\begin{itemize}
    \item \textbf{Eigenvalues}: $\{0, 1, 3, 4\}$ with clear frequency interpretation
    \item \textbf{Eigenvectors}: Orthogonal basis for representing signals on the graph
    \item \textbf{Verification}: Check orthogonality and normalization properties
\end{itemize}

\textbf{Spectral convolution}:
For the test signal $f = [1, 0, -1, 2]^T$ and low-pass filter $g(\lambda) = e^{-\lambda/2}$:
\begin{enumerate}
    \item \textbf{Forward transform}: $\hat{f} = U^T f$ using eigenvectors
    \item \textbf{Filtering}: Multiplication by $g(\lambda_k)$ for each component
    \item \textbf{Inverse transform}: $f_{filtered} = U \hat{f}_{filtered}$
\end{enumerate}

\textbf{Numerical results}:
The low-pass filter produces a smoothed signal $\approx [0.38, 0.15, 0.30, 1.17]^T$, demonstrating the attenuation of high-frequency components while preserving the global trends.

\textbf{Observations from the Example.} This example illustrates key concepts of spectral analysis on graphs:

\textbf{1. Connectivity and multiplicity}:
The graph is connected, so $\lambda_0 = 0$ has multiplicity 1.

\textbf{2. Geometric interpretation}:
- The low-frequency eigenvectors capture the global structure
- The high-frequency ones reveal local variations between adjacent nodes

\textbf{3. Effective filtering}:
The low-pass filter $e^{-\lambda/2}$ preserves the DC component and progressively attenuates the higher frequencies.

\textbf{4. Computational complexity}:
- Diagonalization: $O(n^3)$ for $n = 4$ (trivial)
- For large graphs, the Chebyshev polynomial approximation is preferred

\begin{example}[Other Types of Spectral Filters]\label{ex:otros_filtros_espectrales}
\leavevmode
\begin{enumerate}
    \item \textbf{Rational Filters}: $g(\lambda) = \frac{P(\lambda)}{Q(\lambda)}$ where $P, Q$ are polynomials
    \item \textbf{Lanczos Filters}: $g(\lambda) = \theta \cdot \operatorname{sinc}(\pi\lambda/\lambda_{\max})$
    \item \textbf{Splines}: Piecewise spline functions over the spectrum
    \item \textbf{Wavelet Filters}: Based on spectral wavelets on graphs
\end{enumerate}
\end{example}

\textbf{Localization vs Efficiency Trade-off.} There is a fundamental trade-off:
\begin{itemize}
    \item \textbf{Global spectral filters}: More expressive but $O(n^3)$ for diagonalization
    \item \textbf{Polynomial approximations}: Local and efficient $O(K|E|)$ but less expressive
    \item \textbf{Polynomial degree $K$}: Controls the radius of spatial localization
\end{itemize}

\subsection{Wavelet Bases and Multiresolution Analysis}\label{subsec:wavelets_multirresolucion}

Wavelets provide an alternative to the Fourier transform that allows simultaneous time-frequency analysis, overcoming the limitations of the uncertainty principle. Their application in neural networks offers natural hierarchical representations.

\begin{definition}[Mother Wavelet]\label{def:wavelet_madre}
A function $\psi \in L^2(\mathbb{R})$ is a \textbf{mother wavelet} if it satisfies the admissibility condition:
\begin{equation}
C_\psi = \int_{-\infty}^{\infty} \frac{|\hat{\psi}(\omega)|^2}{|\omega|} d\omega < \infty
\end{equation}
where $\hat{\psi}$ is the Fourier transform of $\psi$.
\end{definition}

\textbf{Implications of the Admissibility Condition.} The condition $C_\psi < \infty$ implies that $\hat{\psi}(0) = 0$, which means $\int_{-\infty}^{\infty} \psi(t) dt = 0$. This guarantees that wavelets have zero mean and provide local analysis.

\begin{definition}[Continuous Wavelet Transform]\label{def:transformada_wavelet_continua}
For a function $f \in L^2(\mathbb{R})$ and mother wavelet $\psi$, the continuous wavelet transform is defined as:
\begin{equation}
W_f(a, b) = \frac{1}{\sqrt{|a|}} \int_{-\infty}^{\infty} f(t) \psi^*\left(\frac{t-b}{a}\right) dt
\end{equation}
where $a \in \mathbb{R}\setminus\{0\}$ is the scale parameter and $b \in \mathbb{R}$ is the translation parameter.
\end{definition}

\begin{theorem}[Wavelet Reconstruction Formula]\label{thm:reconstruccion_wavelet}
If $\psi$ satisfies the admissibility condition, then:
\begin{equation}
f(t) = \frac{1}{C_\psi} \int_{-\infty}^{\infty} \int_{-\infty}^{\infty} W_f(a, b) \psi_{a,b}(t) \frac{da \, db}{a^2}
\end{equation}
where $\psi_{a,b}(t) = \frac{1}{\sqrt{|a|}} \psi\left(\frac{t-b}{a}\right)$.

See Section~\ref{sec:reconstruccion_wavelet} of \autoref{ch:teoremas_complementarios} for the complete proof.
\end{theorem}

\begin{definition}[Multiresolution Analysis (MRA)]\label{def:mra}
A \textbf{multiresolution analysis} of $L^2(\mathbb{R})$ is a sequence of closed subspaces $\{V_j\}_{j \in \mathbb{Z}}$ that satisfies:
\begin{enumerate}
    \item \textbf{Inclusion}: $\cdots \subset V_{-1} \subset V_0 \subset V_1 \subset \cdots$
    \item \textbf{Density}: $\overline{\bigcup_{j \in \mathbb{Z}} V_j} = L^2(\mathbb{R})$
    \item \textbf{Separation}: $\bigcap_{j \in \mathbb{Z}} V_j = \{0\}$
    \item \textbf{Scaling}: $f(x) \in V_j \Leftrightarrow f(2x) \in V_{j+1}$
    \item \textbf{Translation}: $f(x) \in V_0 \Rightarrow f(x-k) \in V_0$ for $k \in \mathbb{Z}$
    \item \textbf{Riesz basis}: There exists $\phi \in V_0$ such that $\{\phi(x-k)\}_{k \in \mathbb{Z}}$ is a Riesz basis of $V_0$
\end{enumerate}
\end{definition}

\begin{definition}[Scaling Function]\label{def:funcion_escala}
The function $\phi$ in condition 6 of the MRA is called the \textbf{scaling function} or \textbf{father wavelet}. It satisfies the scaling equation:
\begin{equation}
\phi(x) = \sqrt{2} \sum_{k \in \mathbb{Z}} h_k \phi(2x - k)
\end{equation}
where $\{h_k\}$ are the scaling coefficients.
\end{definition}

\begin{definition}[Orthogonal Wavelet]\label{def:wavelet_ortogonal}
Given an MRA with scaling function $\phi$, the associated \textbf{orthogonal wavelet} is defined as:
\begin{equation}
\psi(x) = \sqrt{2} \sum_{k \in \mathbb{Z}} g_k \phi(2x - k)
\end{equation}
where $g_k = (-1)^k h_{1-k}$ are related to the scaling coefficients.
\end{definition}

\begin{theorem}[Wavelet Decomposition]\label{thm:descomposicion_wavelet}
Every function $f \in L^2(\mathbb{R})$ can be decomposed as:
\begin{equation}
f(x) = \sum_{k \in \mathbb{Z}} c_{J,k} \phi_{J,k}(x) + \sum_{j=J}^{\infty} \sum_{k \in \mathbb{Z}} d_{j,k} \psi_{j,k}(x)
\end{equation}
where $\phi_{j,k}(x) = 2^{j/2} \phi(2^j x - k)$ and $\psi_{j,k}(x) = 2^{j/2} \psi(2^j x - k)$.

See Section~\ref{sec:descomposicion_wavelet} of \autoref{ch:teoremas_complementarios} for the proof.
\end{theorem}

\begin{example}[Haar Wavelets]\label{ex:wavelets_haar}
The Haar wavelets are the simplest:
\begin{align}
\phi(x) &= \begin{cases} 1 & \text{if } 0 \leq x < 1 \\ 0 & \text{otherwise} \end{cases} \\
\psi(x) &= \begin{cases} 1 & \text{if } 0 \leq x < 1/2 \\ -1 & \text{if } 1/2 \leq x < 1 \\ 0 & \text{otherwise} \end{cases}
\end{align}
\end{example}

\begin{example}[Daubechies Wavelets]\label{ex:wavelets_daubechies}
The Daubechies wavelets $D_N$ have compact support $[0, 2N-1]$ and $N$ vanishing moments:
\begin{equation}
\int_{-\infty}^{\infty} x^k \psi(x) dx = 0, \quad k = 0, 1, \ldots, N-1
\end{equation}
\end{example}

\subsubsection{Wavelets on Graphs}

\begin{definition}[Spectral Wavelets on Graphs]\label{def:wavelets_grafos}
For a graph with Laplacian $L$ and eigenvectors $\{\phi_k\}$, the spectral wavelets are defined as:
\begin{equation}
\psi_{s,i} = \sum_{k=0}^{n-1} g(s\lambda_k) \phi_k(i) \phi_k
\end{equation}
where $g$ is a mother wavelet function and $s$ is the scale parameter.
\end{definition}

\begin{example}[GraphWavelets]\label{ex:graph_wavelets}
A common choice for the mother function is:
\begin{equation}
g(x) = x e^{-x}
\end{equation}
which provides both spectral and spatial localization on the graph.
\end{example}

\textbf{Advantages of Wavelets in CNNs.} Wavelets offer specific advantages for neural networks:
\begin{itemize}
    \item \textbf{Multiresolution analysis}: Natural for hierarchical architectures
    \item \textbf{Localization}: Better than Fourier for local features
    \item \textbf{Parsimony}: Sparse representations for natural signals
    \item \textbf{Translation invariance}: Approximate depending on the type of wavelet
\end{itemize}

\subsubsection{Morlet Wavelet and Time-Frequency Analysis}

\begin{definition}[Morlet Wavelet]\label{def:wavelet_morlet}
The \textit{Morlet wavelet} is a wave modulated by a Gaussian:
\begin{equation}
\psi_{\omega_0}(t) = c_\sigma \pi^{-1/4} e^{i\omega_0 t} e^{-t^2/2}
\end{equation}
where $\omega_0 \geq 5$ is the central frequency and $c_\sigma$ is a normalization constant. Its Fourier transform is $\hat{\psi}_{\omega_0}(\omega) = c_\sigma \pi^{-1/4} e^{-(\omega - \omega_0)^2/2}$, a Gaussian centered at $\omega_0$.
\end{definition}

\begin{remark}[Optimal Time-Frequency Resolution]
The Morlet wavelet achieves equality in the Heisenberg-Gabor uncertainty principle:
\[
\Delta t \cdot \Delta \omega = \frac{1}{2}
\]
which makes it the wavelet with the best simultaneous resolution in time and frequency. At scale $a$, the wavelet $\psi_{a,b}(t) = a^{-1/2}\psi((t-b)/a)$ has temporal resolution $\Delta t \propto a$ and frequency resolution $\Delta \omega \propto 1/a$, exactly implementing the adaptive multiresolution property.
\end{remark}

\subsubsection{Scattering Transform: CNNs with Predefined Wavelets}

The deepest connection between wavelets and CNNs is established through Mallat's \textit{Scattering Transform}~\citep{Mallat2012scattering}, which shows that the early layers of a CNN can be understood as a bank of wavelet filters with a modulus nonlinearity.

\begin{definition}[Scattering Transform]\label{def:scattering_transform}
Let $\psi$ be a mother wavelet and $\phi$ a scaling function (low-pass). The \textit{scattering operator} of order $m$ is defined recursively:
\begin{align}
S_0 f &= f \star \phi_J \\
S_1 f &= \{|f \star \psi_{j_1}| \star \phi_J\}_{j_1} \\
S_m f &= \{|\cdots|f \star \psi_{j_1}| \star \psi_{j_2}| \cdots \star \psi_{j_m}| \star \phi_J\}_{j_1 < j_2 < \cdots < j_m}
\end{align}
where $\psi_j(t) = 2^{-j}\psi(2^{-j}t)$ are the dilated wavelets and $\phi_J$ is the scaling function at the coarsest resolution.
\end{definition}

\begin{proposition}[Properties of the Scattering Transform]\label{prop:propiedades_scattering}
The scattering operator $Sf = (S_0 f, S_1 f, S_2 f, \ldots)$ satisfies:
\begin{enumerate}
    \item \textbf{Translation invariance}: $S[\tau_c f] = S[f]$ for every translation $\tau_c f(t) = f(t - c)$.
    \item \textbf{Lipschitz stability to deformations}: For every diffeomorphism $\tau$ close to the identity:
    \begin{equation}
    \|Sf - S(f \circ \tau)\| \leq C\|\nabla \tau\|_\infty \|f\|
    \end{equation}
    \item \textbf{Energy conservation}: $\sum_{m=0}^{\infty} \|S_m f\|^2 = \|f\|^2$.
\end{enumerate}
\end{proposition}

\begin{proof}
(1) Translation commutes with convolution: $\tau_c f \star \psi_j = \tau_c(f \star \psi_j)$. The modulus preserves this translation: $|\tau_c(f \star \psi_j)| = \tau_c|f \star \psi_j|$. Finally, convolution with $\phi_J$ (low-pass at scale $2^J$) averages over translations smaller than $2^J$, eliminating the dependence on $c$.

(2) Since $\psi_j$ has frequency support concentrated at $|\omega| \sim 2^j$, the convolution $f \star \psi_j$ is sensitive to variations of $f$ at scale $2^{-j}$. A deformation $\tau$ with $\|\nabla \tau\|_\infty \leq \epsilon$ shifts the frequencies by at most $\epsilon \cdot 2^j$, which is small compared to the bandwidth $2^j$ of the filter, producing an error proportional to $\epsilon$.

(3) It follows from Parseval's identity for wavelets: $\|f\|^2 = \|f \star \phi_J\|^2 + \sum_j \|f \star \psi_j\|^2$. Applying this iteratively to the coefficients $|f \star \psi_j|$ and using that the modulus preserves the $L^2$ norm, total energy conservation is obtained.
\end{proof}

\begin{remark}[Scattering as a CNN with Fixed Weights]
The architecture of the scattering transform is isomorphic to a CNN with the following correspondences:
\begin{itemize}
    \item \textbf{Filters}: Predefined (not learned) wavelets $\psi_j$ $\leftrightarrow$ convolutional kernels
    \item \textbf{Nonlinearity}: Modulus $|{\cdot}|$ $\leftrightarrow$ ReLU
    \item \textbf{Pooling}: Convolution with $\phi_J$ $\leftrightarrow$ average pooling
    \item \textbf{Depth}: Scattering order $m$ $\leftrightarrow$ number of layers
\end{itemize}
The key difference is that the scattering transform has theoretical guarantees of stability and invariance, while trained CNNs learn them implicitly.
\end{remark}

\begin{figure}[htbp]
	\centering
	\includegraphics[width=0.85\textwidth]{figures/wavelet_multiresolution}
	\caption[Multiresolution wavelet decomposition]{Multiresolution wavelet decomposition: a signal is iteratively decomposed through low-pass ($h$) and high-pass ($g$) filters followed by downsampling ($\downarrow 2$), separating the approximation ($c_{j,k}$) and detail ($d_{j,k}$) components at multiple scales.}
	\label{fig:wavelet_multiresolution}
\end{figure}

\subsection{Spectral Convolutions in CNNs}\label{subsec:convoluciones_espectrales_cnns}

Spectral convolutions in CNNs leverage the properties of the frequency domain to implement convolutional operations efficiently, especially for large kernels or in applications with specific geometric structure.

\begin{definition}[Parametrized Spectral Filter]\label{def:filtro_espectral_parametrizado}
A \textbf{parametrized spectral filter} is a function $h_\theta: \mathbb{R}_+ \to \mathbb{R}$ where $\theta \in \mathbb{R}^p$ are trainable parameters. The filter acts on the frequency spectrum through:
\begin{equation}
h_\theta(\lambda) = \sum_{k=0}^{K-1} \theta_k \phi_k(\lambda)
\end{equation}
where $\{\phi_k\}$ is a basis of functions (polynomials, wavelets, etc.).
\end{definition}

\begin{definition}[Spectral Convolutional Layer]\label{def:capa_convolucional_espectral}
A spectral convolutional layer transforms an input signal $x \in \mathbb{R}^n$ through:
\begin{equation}
y = \sigma(h_\theta \star x + b)
\end{equation}
where:
\begin{itemize}
    \item $h_\theta \star x$ is the spectral convolution: $\mathcal{F}^{-1}[h_\theta(\cdot) \cdot \mathcal{F}[x]]$
    \item $\sigma$ is the activation function
    \item $b$ is the bias term
\end{itemize}
\end{definition}

\begin{theorem}[Spatial-Spectral Equivalence]\label{thm:equivalencia_espacial_espectral}
For domains with regular structure (grids, tori), spatial and spectral convolutions are mathematically equivalent:
\begin{equation}
(f * g)(x) = \mathcal{F}^{-1}[\mathcal{F}[f] \cdot \mathcal{F}[g]](x)
\end{equation}
However, they differ in:
\begin{itemize}
    \item \textbf{Computational complexity}: Spatial $O(n \cdot k^2)$ vs Spectral $O(n \log n)$
    \item \textbf{Localization}: Spatial (local) vs Spectral (global)
    \item \textbf{Interpretability}: Spatial (geometric) vs Spectral (frequency)
\end{itemize}
\end{theorem}

\begin{proof}
The equivalence follows from the convolution theorem for the Fourier transform. For functions $f, g \in L^1(\mathbb{R}^d)$ (or their discrete analogs on $\mathbb{Z}^d$), the Fourier transform satisfies $\mathcal{F}[f * g] = \mathcal{F}[f] \cdot \mathcal{F}[g]$. Applying $\mathcal{F}^{-1}$ to both sides yields the identity of the statement. The differences in complexity are: the direct spatial convolution performs $O(n \cdot k^2)$ operations (traversing $n$ positions applying a kernel of size $k \times k$), while the spectral convolution via FFT requires $O(n \log n)$ for the forward transform, $O(n)$ for the pointwise multiplication, and $O(n \log n)$ for the inverse transform, totaling $O(n \log n)$. The localization differs because a finite spatial kernel of size $k$ corresponds to a spectral filter with full support, and vice versa.
\end{proof}

\subsubsection{Practical Implementation of Spectral CNNs}

\textbf{Algorithm: Spectral Forward Pass.}\label{alg:forward_espectral} For an input $x \in \mathbb{R}^{H \times W}$ and parametrized filter $h_\theta$:
\begin{enumerate}
    \item \textbf{Forward transform}: $\hat{x} = \text{FFT2D}(x)$
    \item \textbf{Filter application}: $\hat{y} = h_\theta(\omega) \odot \hat{x}$ (element-wise product)
    \item \textbf{Inverse transform}: $y = \text{IFFT2D}(\hat{y})$
    \item \textbf{Activation}: $z = \sigma(y + b)$
\end{enumerate}

\textbf{Implementation Considerations.}
\begin{itemize}
    \item \textbf{Padding}: To avoid edge effects, circular padding or zero-padding is used
    \item \textbf{Numerical precision}: \gls{FFT}/IFFT can accumulate floating-point errors
    \item \textbf{Memory}: Store both the spatial and spectral representations
\end{itemize}

\begin{definition}[Transfer Function of a Linear CNN]\label{def:funcion_transferencia_cnn}
For a linear CNN (without nonlinear activation functions) with $L$ layers, the global transfer function in the frequency domain is:
\begin{equation}
H(\omega) = \prod_{l=1}^L H_l(\omega)
\end{equation}
where $H_l(\omega)$ is the frequency response of layer $l$. This factorization is only valid in the absence of nonlinearities between layers.
\end{definition}

\begin{example}[Simple Spectral CNN]\label{ex:cnn_espectral_simple}
Consider a two-layer CNN:
\begin{align}
\text{Layer 1}: \quad y_1 &= \sigma(\mathcal{F}^{-1}[h_1(\omega) \cdot \mathcal{F}[x]] + b_1) \\
\text{Layer 2}: \quad y_2 &= \mathcal{F}^{-1}[h_2(\omega) \cdot \mathcal{F}[y_1]] + b_2
\end{align}
The total frequency response is $H(\omega) = h_1(\omega) \cdot h_2(\omega)$ (without considering nonlinearities).
\end{example}

\subsubsection{Complexity Analysis}

\begin{theorem}[Computational Complexity]\label{thm:complejidad_convoluciones_espectrales}
For a 2D convolution of size $n \times n$ with a $k \times k$ kernel:
\begin{itemize}
    \item \textbf{Direct convolution}: $O(n^2 k^2)$
    \item \textbf{Spectral convolution}: $O(n^2 \log n)$
    \item \textbf{Break-even point}: Spectral is more efficient when $k > O(\log n)$
\end{itemize}
\end{theorem}

\begin{proof}
The direct 2D convolution computes, for each of the $n^2$ positions of the image, a sum of $k^2$ products, giving $O(n^2 k^2)$ operations. The spectral convolution requires: (i) 2D FFT of the image: $O(n^2 \log n)$; (ii) 2D FFT of the kernel (with zero-padding): $O(n^2 \log n)$; (iii) pointwise multiplication in frequency: $O(n^2)$; (iv) 2D IFFT: $O(n^2 \log n)$. The total is $O(n^2 \log n)$. The break-even point is obtained by equating $n^2 k^2 = n^2 \log n$, that is, $k^2 = \log n$, so the spectral convolution is preferable when $k > O(\sqrt{\log n})$.
\end{proof}

\begin{example}[Practical Efficiency Analysis]\label{ex:eficiencia_practica}
For a $512 \times 512$ image and a $k \times k$ kernel:
\begin{itemize}
    \item $k = 3$: Direct: $\sim 2.3 \times 10^6$ ops, Spectral: $\sim 2.6 \times 10^6$ ops
    \item $k = 15$: Direct: $\sim 5.9 \times 10^7$ ops, Spectral: $\sim 2.6 \times 10^6$ ops
    \item $k = 51$: Direct: $\sim 6.8 \times 10^8$ ops, Spectral: $\sim 2.6 \times 10^6$ ops
\end{itemize}
\end{example}

\begin{theorem}[Spectral Universal Approximation]\label{thm:aproximacion_universal_espectral}
Let $\mathcal{G} = (V, E)$ be a finite graph with Laplacian $L$ and spectrum contained in $[0, \lambda_{\max}]$. For any continuous spectral filter $h: [0, \lambda_{\max}] \to \mathbb{R}$ and $\epsilon > 0$, there exists a polynomial filter $p(\lambda) = \sum_{k=0}^{K} \theta_k T_k(\tilde{\lambda})$ of sufficiently large degree $K$ such that:
\begin{equation}
\|h(\lambda) - p(\lambda)\|_{\infty, [0, \lambda_{\max}]} < \epsilon
\end{equation}
Consequently, the filtered signal satisfies $\|h(L)f - p(L)f\| \leq \epsilon \|f\|$ for every signal $f$ on the graph.
\end{theorem}

\begin{proof}[Sketch]
The spectral interval $[0, \lambda_{\max}]$ is compact, so the result is a direct consequence of the Weierstrass approximation theorem: every continuous function on a compact interval can be uniformly approximated by polynomials. The bound on the filtered signal follows from $\|h(L)f - p(L)f\| = \|(h-p)(L)f\| \leq \|h - p\|_\infty \|f\|$, since the eigenvalues of $(h-p)(L)$ are bounded by $\|h-p\|_\infty$.
\end{proof}

\textbf{Limitations of Spectral Convolutions.}
\begin{itemize}
    \item \textbf{Non-locality}: Spectral filters affect the signal globally
    \item \textbf{Regularity}: They require domains with \gls{FFT} structure (regular grids)
    \item \textbf{Boundaries}: Difficulties with non-periodic boundary conditions
    \item \textbf{Interpretability}: Less intuitive than spatial convolution
\end{itemize}

\subsection{Frequency Interpretation of Classical CNNs}\label{subsec:interpretacion_frecuencial_cnns}

The frequency analysis of traditional CNNs reveals deep insights into what features they learn and how they process visual information.

\begin{definition}[Frequency Response of a Kernel]\label{def:respuesta_frecuencial_kernel}
For a convolutional kernel $K \in \mathbb{R}^{k \times k}$, its frequency response is:
\begin{equation}
H(\omega_x, \omega_y) = \sum_{i=0}^{k-1} \sum_{j=0}^{k-1} K[i,j] e^{-2\pi i(\omega_x i + \omega_y j)}
\end{equation}
\end{definition}

\begin{example}[Analysis of Common Kernels]\label{ex:analisis_kernels_comunes}
\textbf{1. Sobel Edge Detector}:
\begin{equation}
K_{\text{Sobel}} = \begin{bmatrix} -1 & 0 & 1 \\ -2 & 0 & 2 \\ -1 & 0 & 1 \end{bmatrix}
\end{equation}
Frequency response: high-pass with maximum at medium horizontal frequencies.

\textbf{2. Gaussian Filter}:
\begin{equation}
K_{\text{Gauss}}[i,j] = \frac{1}{2\pi\sigma^2} e^{-\frac{(i-c)^2 + (j-c)^2}{2\sigma^2}}
\end{equation}
Frequency response: low-pass with a smooth cutoff determined by $\sigma$.

\textbf{3. Laplacian}:
\begin{equation}
K_{\text{Lap}} = \begin{bmatrix} 0 & -1 & 0 \\ -1 & 4 & -1 \\ 0 & -1 & 0 \end{bmatrix}
\end{equation}
Frequency response: isotropic high-pass.
\end{example}

\begin{remark}[Spectral Evolution during Training]\label{rem:evolucion_espectral}
Empirical observations show that, during the training of a CNN, the spectral distribution of the filters evolves following a characteristic pattern:
\begin{enumerate}
    \item \textbf{Early layers}: Low-pass filters and edge detectors (high frequencies)
    \item \textbf{Intermediate layers}: Band-pass filters for specific textures
    \item \textbf{Deep layers}: Complex filters, less spectrally interpretable
\end{enumerate}
This behavior has been documented experimentally in multiple architectures, although it lacks a general formal proof.
\end{remark}

\textbf{Experiment: Spectral Bias of CNNs.}\label{exp:sesgo_espectral} Empirical studies show that CNNs exhibit a bias toward low frequencies:
\begin{itemize}
    \item They first learn low-frequency components
    \item They require more iterations to capture high-frequency details
    \item This bias can be improved with spectral initialization techniques
\end{itemize}

The spectral methods developed in this section find direct applications in domains where the frequency structure of the data is fundamental: medical image reconstruction (MRI, computed tomography), scientific signal processing (astronomy, seismology), and neural network compression through spectral pruning. The focus of this thesis, however, is centered on the mathematical foundation of these methods.

\subsection{Limitations and Challenges}\label{subsec:aplicaciones_limitaciones}

\textbf{Limitation: Dependence on the Domain Geometry.}\label{lim:geometria_dominio} Classical spectral convolutions require:
\begin{itemize}
    \item \textbf{Regular structure}: Uniform grids for efficient \gls{FFT}
    \item \textbf{Boundary conditions}: Periodic or with adequate padding
    \item \textbf{Fixed size}: Difficulties with variable-size data
\end{itemize}

\textbf{Consequence}: Restriction to regular Euclidean domains.

\textbf{Limitation: Localization-Efficiency Trade-off.}\label{lim:tradeoff_localizacion}
\begin{table}[h]
\centering
\begin{tabular}{|l|c|c|c|}
\hline
\textbf{Method} & \textbf{Localization} & \textbf{Efficiency} & \textbf{Interpretability} \\
\hline
Spatial convolution & Excellent & Medium & High \\
Spectral convolution & Poor & Excellent & Medium \\
Polynomial approximation & Good & Good & Medium \\
Wavelets & Excellent & Medium & High \\
\hline
\end{tabular}
\caption{Comparison of convolutional methods}
\label{tab:comparacion_metodos}
\end{table}

\textbf{Limitation: Numerical Stability.}\label{lim:estabilidad_numerica} Spectral methods can suffer from:
\begin{itemize}
    \item \textbf{Error accumulation}: In multiple \gls{FFT}/IFFT
    \item \textbf{Aliasing}: In downsampling without anti-aliasing filters
    \item \textbf{Floating-point precision}: Especially on GPU
\end{itemize}

\textbf{Limitation: Reduced Interpretability.}\label{lim:interpretabilidad} Compared to spatial convolutions:
\begin{itemize}
    \item Spectral filters are less visually intuitive
    \item The relationship with geometric features is indirect
    \item Debugging and error analysis are more complex
\end{itemize}

\textbf{Limitation: No Spectral Transferability between Domains.}\label{lim:transferibilidad} Spectral filters defined over a domain $\Omega_1$ cannot be directly transferred to another domain $\Omega_2$ with a different spectrum. Formally, if $\{\lambda_k^{(1)}, \phi_k^{(1)}\}$ and $\{\lambda_k^{(2)}, \phi_k^{(2)}\}$ are the spectral decompositions of the Laplacian operators on $\Omega_1$ and $\Omega_2$ respectively, a filter $h(\lambda)$ learned over $\{\lambda_k^{(1)}\}$ has no canonical interpretation over $\{\lambda_k^{(2)}\}$ unless there exists a natural correspondence between the eigenvalues. This limitation motivates the use of polynomial approximations such as Chebyshev filters (\autoref{thm:localizacion_espectral}), which define the filter as a function of the operator $h(L)$ instead of acting on individual eigenvalues.

\begin{proposition}[Fundamental Localization-Globality Trade-off]\label{prop:tradeoff_localizacion}
Let $h(\lambda)$ be a spectral filter with support in $[\lambda_{\min}, \lambda_{\max}]$ and let $K(x,y) = \sum_k h(\lambda_k) \phi_k(x)\phi_k(y)$ be its associated spatial kernel. Then:
\begin{equation}
\text{supp}_{\text{freq}}(h) \cdot \text{supp}_{\text{spatial}}(K) \geq C > 0
\end{equation}
where $C$ depends only on the dimension of the domain. In particular, a filter perfectly localized in frequency ($h = \delta_{\lambda_k}$) is completely delocalized in space ($K(x,y) = \phi_k(x)\phi_k(y)$, global support).
\end{proposition}

\begin{proof}
It is a direct consequence of the uncertainty principle for the spectral transform: the concentration of a function and its spectral representation cannot both be arbitrarily small simultaneously. For the case $h = \delta_{\lambda_k}$, the kernel reduces to $K(x,y) = \phi_k(x)\phi_k(y)$, which is nonzero over almost the entire domain (since the eigenvectors of the Laplacian vanish only on a set of measure zero). The general case follows by a duality argument analogous to the classical Heisenberg principle.
\end{proof}

\textbf{Synthesis: The Spectral Perspective as a Bridge.} Harmonic analysis provides a unifying framework that connects:
\begin{enumerate}
    \item \textbf{Classical CNNs} (Fourier analysis on regular grids)
    \item \textbf{Spectral GCNs} (spectral analysis on graphs)
    \item \textbf{Geometric CNNs} (harmonic analysis on groups and manifolds)
    \item \textbf{Transformers} (attention as adaptive convolution in the frequency domain)
\end{enumerate}

This perspective reveals that many seemingly different architectures share deep mathematical foundations, facilitating the development of new architectures that combine the best features of each approach.

The \gls{analisis_armonico} developed in the preceding section not only enriches our understanding of classical CNNs, but also provides the mathematical tools needed to establish the fundamental equivalence between convolutions on regular grids and spectral operations on graphs. This connection, which we will explore next, constitutes one of the pillars of the unifying framework of Geometric Deep Learning.

\section[From Convolutions to Graphs: The CNN-GNN Connection]{From Convolutions to Graphs: The \glspl{CNN}-\glspl{GNN} Connection}\label{sec:cnn_as_gnn}

\subsection{Preliminaries: Graph Theory}\label{subsec:preliminares_grafos_cnn}

For the basic definition of a graph, see \autoref{def:grafo_neuronal} in \autoref{ch:neural_networks}. In this section we focus on its specific relationship with convolutions.

\begin{definition}[Representation of Data as Graphs]\label{def:representacion_datos_grafos}
Data is represented in a graph by means of:
\begin{itemize}
	\item \textbf{Node features}: \(X \in \mathbb{R}^{n \times d}\) where \(x_i \in \mathbb{R}^d\) is the feature of node \(i\)
	\item \textbf{Adjacency matrix}: \(A \in \{0,1\}^{n \times n}\) where \(A_{ij} = 1\) if \((v_i, v_j) \in E\)
	\item \textbf{Edge features} (optional): \(E \in \mathbb{R}^{|E| \times d_e}\)
\end{itemize}
\end{definition}

\begin{example}[Different Types of Data as Graphs]
\leavevmode
\begin{enumerate}
	\item \textbf{Molecules}: atoms = nodes, bonds = edges
	\item \textbf{Social networks}: people = nodes, relationships = edges
	\item \textbf{Images}: pixels = nodes, neighborhood = edges (special case!)
	\item \textbf{Point clouds}: points = nodes, k-NN = edges
\end{enumerate}
\end{example}

\begin{figure}[htbp]
	\centering
	\includegraphics[width=0.85\textwidth]{figures/image_as_graph}
	\caption[Image as a regular grid graph]{An image can be represented as a regular grid graph with 4-connectivity: each pixel is a node and the edges connect horizontally and vertically neighboring pixels. The neighborhood $\mathcal{N}(v_i)$ corresponds exactly to the support of the convolutional kernel.}
	\label{fig:image_as_graph}
\end{figure}

\begin{example}[Explicit Construction: 3×3 Image as a Graph]\label{ex:imagen_3x3_grafo}
Consider an RGB image of size $3 \times 3$ with 3 color channels. The transformation into a graph proceeds in three steps:

\textbf{Step 1: Node identification}. Each pixel $(i,j)$ is mapped to a node $v_{i,j}$ of the graph:
\[
V = \{v_{0,0}, v_{0,1}, v_{0,2}, v_{1,0}, v_{1,1}, v_{1,2}, v_{2,0}, v_{2,1}, v_{2,2}\}
\]
with $|V| = 9$ nodes.

\textbf{Step 2: Edge construction}. For 4-connectivity, we connect horizontally and vertically adjacent pixels:
\begin{align}
E = \{&(v_{0,0}, v_{0,1}), (v_{0,1}, v_{0,2}), \quad \text{(row 0)} \\
    &(v_{1,0}, v_{1,1}), (v_{1,1}, v_{1,2}), \quad \text{(row 1)} \\
    &(v_{2,0}, v_{2,1}), (v_{2,1}, v_{2,2}), \quad \text{(row 2)} \\
    &(v_{0,0}, v_{1,0}), (v_{1,0}, v_{2,0}), \quad \text{(column 0)} \\
    &(v_{0,1}, v_{1,1}), (v_{1,1}, v_{2,1}), \quad \text{(column 1)} \\
    &(v_{0,2}, v_{1,2}), (v_{1,2}, v_{2,2}) \} \quad \text{(column 2)}
\end{align}
We obtain $|E| = 12$ undirected edges.

\textbf{Step 3: Feature assignment}. Each node $v_{i,j}$ receives a feature vector $\mathbf{x}_{i,j} \in \mathbb{R}^3$ corresponding to the RGB values of pixel $(i,j)$:
\[
X = \begin{bmatrix} \mathbf{x}_{0,0}^T \\ \mathbf{x}_{0,1}^T \\ \vdots \\ \mathbf{x}_{2,2}^T \end{bmatrix} \in \mathbb{R}^{9 \times 3}
\]

The adjacency matrix $A \in \{0,1\}^{9 \times 9}$ has the block-tridiagonal form characteristic of regular grids.
\end{example}

\textbf{Special Properties of Grid Graphs.} Grid graphs derived from images possess unique structural properties:
\begin{enumerate}
    \item \textbf{Regularity}: Nearly constant degree ($\deg(v) = 4$ except at borders)
    \item \textbf{Spatial structure}: Positions $(i,j)$ define a natural embedding in $\mathbb{R}^2$
    \item \textbf{Symmetry}: Invariance to discrete translations
    \item \textbf{Predictable spectrum}: The eigenvalues of the Laplacian are known analytically
\end{enumerate}

These properties make CNNs especially effective on grids, but they also limit their applicability to general graphs.

\subsection[CNNs as GNNs on Regular Grids]{\glspl{CNN} as \glspl{GNN} on Regular Grids}\label{subsec:cnns_como_gnns_grillas}

\begin{theorem}[A \gls{CNN} is a Special Case of a \gls{GNN}]\label{thm:cnn_caso_especial_gnn}
A 2D \gls{CNN} is exactly equivalent to a \gls{GNN} where:
\begin{itemize}
	\item \textbf{Graph}: Regular grid \(G = (\mathbb{Z}^2 \cap \Omega, E)\)
	\item \textbf{Vertices}: Pixels at positions \((i,j)\)
	\item \textbf{Edges}: 4- or 8-neighbor connectivity
	\item \textbf{Weights}: Shared according to relative position (kernel)
\end{itemize}
\end{theorem}

\begin{proof}
The equivalence is established through the natural correspondence:
\begin{enumerate}
	\item Each pixel \((i,j)\) corresponds to a node \(v_{i,j}\) in the graph
	\item The spatial connections define the edges: \((v_{i,j}, v_{i',j'}) \in E\) if and only if \(|i-i'| + |j-j'| \leq 1\)
	\item The convolutional kernel \(K\) determines the edge weights according to the relative distance
	\item The aggregation operation in \glspl{GNN} reduces exactly to 2D convolution
\end{enumerate}
\end{proof}

\begin{definition}[Grid as a Graph]\label{def:grilla_como_grafo}
For an image of size \(H \times W\):
\begin{itemize}
	\item \(V = \{(i,j) : 0 \leq i < H, 0 \leq j < W\}\)
	\item \(E = \{((i,j), (i',j')) : |i-i'| + |j-j'| \leq 1\}\) (4-connectivity)
	\item Constant degree: \(\deg(v) = 4\) (except at borders)
\end{itemize}
\end{definition}

\begin{theorem}[Constructive CNN-GNN Equivalence]\label{thm:equivalencia_constructiva_cnn_gnn}
Given a CNN with an architecture of convolutional layers $\{L_\ell\}_{\ell=1}^L$ operating on $H \times W$ grids, there exists an explicit construction of an equivalent GNN that produces identical outputs for every input.

\textbf{Construction}:
\begin{enumerate}
    \item \textbf{Graph}: $G = (V, E)$ where $V = \{(i,j) : 0 \leq i < H, 0 \leq j < W\}$ and
    \[
    E = \{((i,j), (i',j')) : |i-i'| + |j-j'| = 1\}
    \]

    \item \textbf{Initial features}: For each node $v_{i,j}$, assign $h_{v_{i,j}}^{(0)} = x_{i,j}$ where $x_{i,j}$ is the pixel value at $(i,j)$

    \item \textbf{Message function}: For a CNN kernel $K \in \mathbb{R}^{k \times k \times c_{\text{in}} \times c_{\text{out}}}$, define
    \[
    M_\ell(h_v^{(\ell)}, h_u^{(\ell)}, e_{vu}) = W_{v-u}^{(\ell)} h_u^{(\ell)}
    \]
    where $W_{v-u}^{(\ell)} \in \mathbb{R}^{c_{\text{out}} \times c_{\text{in}}}$ is the submatrix of the kernel corresponding to the displacement $(v-u)$

    \item \textbf{Aggregation}:
    \[
    m_v^{(\ell+1)} = \sum_{u \in \mathcal{N}(v)} M_\ell(h_v^{(\ell)}, h_u^{(\ell)}, e_{vu})
    \]

    \item \textbf{Update}:
    \[
    h_v^{(\ell+1)} = \sigma(m_v^{(\ell+1)} + b^{(\ell)})
    \]
    where $\sigma$ is the activation function and $b^{(\ell)}$ the bias
\end{enumerate}

This construction guarantees that $h_v^{(\ell)} = y_{i,j}^{(\ell)}$ for every node $v = (i,j)$ and layer $\ell$, where $y_{i,j}^{(\ell)}$ is the CNN activation at position $(i,j)$ of layer $\ell$.
\end{theorem}

\begin{proof}
We proceed by induction on the depth of the network.

\textbf{Base case} ($\ell = 0$): By construction, $h_v^{(0)} = x_{i,j} = y_{i,j}^{(0)}$ for every node $v = (i,j)$.

\textbf{Inductive step}: Suppose $h_v^{(\ell)} = y_{i,j}^{(\ell)}$ for every $v = (i,j)$. We must show that $h_v^{(\ell+1)} = y_{i,j}^{(\ell+1)}$.

For the CNN, the output at position $(i,j)$ of layer $\ell+1$ is:
\begin{align}
y_{i,j}^{(\ell+1)} &= \sigma\left(\sum_{(u,v) \in \mathcal{K}} K_{u,v} \cdot x_{i+u,j+v}^{(\ell)} + b^{(\ell)}\right) \\
&= \sigma\left(\sum_{(i',j') \in \mathcal{N}(i,j)} K_{i'-i,j'-j} \cdot y_{i',j'}^{(\ell)} + b^{(\ell)}\right)
\end{align}

where $\mathcal{K}$ is the support of the kernel and $\mathcal{N}(i,j)$ the neighborhood of pixel $(i,j)$.

For the GNN, using the proposed construction:
\begin{align}
h_v^{(\ell+1)} &= \sigma\left(\sum_{u \in \mathcal{N}(v)} W_{v-u}^{(\ell)} h_u^{(\ell)} + b^{(\ell)}\right) \\
&= \sigma\left(\sum_{(i',j') \in \mathcal{N}(i,j)} K_{i'-i,j'-j} \cdot h_{(i',j')}^{(\ell)} + b^{(\ell)}\right) \\
&\overset{\text{I.H.}}{=} \sigma\left(\sum_{(i',j') \in \mathcal{N}(i,j)} K_{i'-i,j'-j} \cdot y_{i',j'}^{(\ell)} + b^{(\ell)}\right) \\
&= y_{i,j}^{(\ell+1)}
\end{align}

where we used the inductive hypothesis (I.H.) in the third equality. $\square$
\end{proof}

\begin{example}[Numerical Verification: 3×3 Sobel Kernel]\label{ex:verificacion_numerica_sobel}
Consider the horizontal Sobel kernel:
\[
K = \begin{bmatrix} -1 & 0 & 1 \\ -2 & 0 & 2 \\ -1 & 0 & 1 \end{bmatrix}
\]

Applied to the $3 \times 3$ image:
\[
I = \begin{bmatrix} 0 & 1 & 2 \\ 0 & 1 & 2 \\ 0 & 1 & 2 \end{bmatrix}
\]

\textbf{CNN method}: The convolution at the central pixel $(1,1)$ produces:
\begin{align}
y_{1,1} &= -1 \cdot 0 + 0 \cdot 1 + 1 \cdot 2 \\
        &\quad + (-2) \cdot 0 + 0 \cdot 1 + 2 \cdot 2 \\
        &\quad + (-1) \cdot 0 + 0 \cdot 1 + 1 \cdot 2 \\
        &= 0 + 2 + 0 + 4 + 0 + 2 = 8
\end{align}

\textbf{GNN method}: We construct the graph with $v_{1,1}$ connected to its 8 neighbors. The messages are:
\begin{align}
m_{1,1} &= K_{-1,-1} \cdot h_{0,0} + K_{-1,0} \cdot h_{0,1} + K_{-1,1} \cdot h_{0,2} \\
        &\quad + K_{0,-1} \cdot h_{1,0} + K_{0,1} \cdot h_{1,2} \\
        &\quad + K_{1,-1} \cdot h_{2,0} + K_{1,0} \cdot h_{2,1} + K_{1,1} \cdot h_{2,2} \\
        &= (-1)(0) + (0)(1) + (1)(2) + (-2)(0) + (2)(2) \\
        &\quad + (-1)(0) + (0)(1) + (1)(2) = 8
\end{align}

Both methods produce the identical result $y_{1,1} = 8$, verifying the equivalence.
\end{example}

\subsection{Convolution-Message Passing Equivalence}\label{subsec:equivalencia_conv_mp}

\begin{theorem}[Fundamental Equivalence]\label{thm:equivalencia_fundamental_conv_mp}
The 2D convolution operation:
\[
y_{i,j} = \sum_{(u,v) \in \mathcal{N}} K_{u,v} \cdot x_{i+u,j+v}
\]

is exactly equivalent to message passing on the grid:
\[
m_{(i,j)} = \sum_{(i',j') \in N(i,j)} K_{i'-i,j'-j} \cdot x_{i',j'}
\]

where \(N(i,j)\) are the neighbors of pixel \((i,j)\).
\end{theorem}

\begin{proof}
The convolution sums over a fixed neighborhood relative to the central pixel.
On the regular grid, this neighborhood corresponds exactly to the
neighbors in the graph with weights determined by the relative position.
The kernel weights \(K_{u,v}\) correspond to the weighted messages. \(\square\)
\end{proof}

\begin{figure}[htbp]
	\centering
	\includegraphics[width=0.90\textwidth]{figures/conv_message_passing}
	\caption[Convolution-message passing equivalence]{Equivalence between convolution and message passing: the three fundamental steps (message, aggregation, update) correspond exactly to the convolutional operation when the graph is a regular grid.}
	\label{fig:conv_message_passing}
\end{figure}

\begin{example}[Step-by-Step Demonstration: Convolution = Message Passing]\label{ex:paso_a_paso_conv_mp}
Consider a 1D signal $f = [1, 3, 2, 4]$ and a kernel $K = [0.5, 1, 0.5]$ (smoothing filter).

\textbf{Interpretation as convolution}:
\begin{align}
(f * K)[1] &= 0.5 \cdot f[0] + 1 \cdot f[1] + 0.5 \cdot f[2] \\
           &= 0.5 \cdot 1 + 1 \cdot 3 + 0.5 \cdot 2 = 4.5
\end{align}

\textbf{Interpretation as message passing}:
We construct a linear graph $v_0 \leftrightarrow v_1 \leftrightarrow v_2 \leftrightarrow v_3$ with initial features $h_{v_i}^{(0)} = f[i]$.

The messages received by node $v_1$ are:
\begin{itemize}
    \item From $v_0$: $m_{v_1 \leftarrow v_0} = W_{-1} \cdot h_{v_0}^{(0)} = 0.5 \cdot 1 = 0.5$
    \item Self-message: $m_{v_1 \leftarrow v_1} = W_0 \cdot h_{v_1}^{(0)} = 1 \cdot 3 = 3$
    \item From $v_2$: $m_{v_1 \leftarrow v_2} = W_{+1} \cdot h_{v_2}^{(0)} = 0.5 \cdot 2 = 1$
\end{itemize}

Aggregation produces:
\[
m_{v_1}^{(1)} = \sum_{u \in \mathcal{N}(v_1)} m_{v_1 \leftarrow u} = 0.5 + 3 + 1 = 4.5
\]

Once again, both approaches produce the same result: $4.5$.
\end{example}

\textbf{Connection with the Unified MPNN Framework.} The convolution-message passing equivalence established in this section is a special case of the \textit{Message Passing Neural Networks} (\glspl{MPNN}) framework introduced in \autoref{def:mpnn_unificador} of \autoref{ch:gdl}.

Specifically, CNNs correspond to MPNNs with:
\begin{itemize}
    \item \textbf{Fixed graph}: Regular grid with local connectivity
    \item \textbf{Message function}: $M_\ell(h_v, h_u, e_{vu}) = W_{v-u}^{(\ell)} h_u$ (dependence only on the relative displacement)
    \item \textbf{Aggregation}: Unweighted sum $m_v = \sum_{u \in \mathcal{N}(v)} M_\ell(\cdot)$
    \item \textbf{Update}: $U_\ell(h_v, m_v) = \sigma(m_v + b^{(\ell)})$
\end{itemize}

This formulation reveals that the fundamental restriction of CNNs is \textbf{weight sharing based on relative position} ($W_{v-u}$), which only makes sense in domains with a translation structure. For general graphs lacking such structure, we need more flexible message functions, as explored in \autoref{subsec:message_passing_unificador} of \autoref{ch:gdl}.

\begin{proposition}[Computational Complexity of CNN vs GNN on Grids]\label{prop:complejidad_cnn_vs_gnn_grilla}
For an image of size $H \times W$ with $C$ channels and a $k \times k$ kernel:
\begin{itemize}
    \item \textbf{Standard CNN}: $O(H \cdot W \cdot k^2 \cdot C)$ operations
    \item \textbf{GNN on a grid}: $O(|E| \cdot C) = O(4HW \cdot C)$ operations (for 4-connectivity)
\end{itemize}

When $k$ is small (typically $k \in \{3, 5, 7\}$), both complexities are equivalent: $O(HWC)$. However, the GNN formulation reveals that the complexity fundamentally depends on the number of edges $|E|$, not on the kernel size, providing a unified perspective for analyzing efficiency across different graph topologies.
\end{proposition}

\begin{proof}
\textbf{Standard CNN}: The image has $H \cdot W$ spatial positions. At each position, the $k \times k$ kernel performs $k^2$ multiplications over each of the $C$ channels, giving a total of $O(H \cdot W \cdot k^2 \cdot C)$ operations.

\textbf{GNN on a grid}: Each node (pixel) has at most 4 neighbors in 4-connectivity (up, down, left, right), so $|E| = O(4HW)$. The message passing step processes $C$ channels per edge, resulting in $O(4HW \cdot C) = O(|E| \cdot C)$ operations.
\end{proof}

\textbf{Dual Perspective: Spatial vs. Spectral.} The equivalence established in this section admits two complementary interpretations:

\textbf{Spatial perspective}: The convolution aggregates information from local neighbors with shared weights determined by the relative position. This is the natural formulation of CNNs.

\textbf{Spectral perspective}: As developed in \autoref{sec:analisis_armonico}, the convolution corresponds to multiplication in the Fourier domain. For regular grids, the eigenvectors of the Laplacian are exactly the discrete Fourier bases, establishing the connection with \autoref{def:convolucion_espectral_grafos}.

This spatial-spectral duality is fundamental: classical CNNs operate in the spatial domain for computational efficiency, while spectral formulations provide the theoretical tools to generalize to arbitrary graphs through the Chebyshev approximation (\autoref{ex:chebyshev_aproximacion}).

\subsection{Grid Laplacian and Filters}\label{subsec:laplaciano_grilla_filtros}

\begin{definition}[Discrete Laplacian on a 2D Grid]\label{def:laplaciano_discreto_grilla}
The Laplacian operator on a grid:
\[
\mathcal{L} = D - A
\]
where \(D\) is the degree matrix and \(A\) the adjacency matrix.

For a 2D grid with 4-connectivity:
\begin{align}
\mathcal{L}_{(i,j),(i,j)} &= 4 \\
\mathcal{L}_{(i,j),(i \pm 1,j)} = \mathcal{L}_{(i,j),(i,j \pm 1)} &= -1
\end{align}
\end{definition}

\begin{example}[Laplacian Kernel as a CNN Filter]
The discrete Laplacian corresponds to the kernel:
\[
K = \begin{bmatrix}
 0 & -1 &  0 \\
-1 &  4 & -1 \\
 0 & -1 &  0
\end{bmatrix}
\]
This is a classical edge detector in image processing, and its spectral interpretation as a high-pass filter is formalized in \autoref{thm:cnns_filtros_espectrales_grillas}.
\end{example}

\begin{proposition}[Eigenvalues of the Laplacian on a Rectangular Grid]\label{prop:eigenvalores_laplaciano_grilla}
For a 2D grid of size $n_1 \times n_2$ with periodic boundary conditions, the eigenvalues of the Laplacian are:
\begin{equation}
\lambda_{k_1,k_2} = 4 - 2\cos\!\left(\frac{2\pi k_1}{n_1}\right) - 2\cos\!\left(\frac{2\pi k_2}{n_2}\right), \quad 0 \leq k_i < n_i
\end{equation}
with eigenvectors $\phi_{k_1,k_2}(i,j) = \frac{1}{\sqrt{n_1 n_2}} e^{2\pi \imath(k_1 i/n_1 + k_2 j/n_2)}$, which are exactly the discrete Fourier bases.
\end{proposition}

\begin{proof}
For periodic boundary conditions, the functions $\phi_{k_1,k_2}$ satisfy $\phi_{k_1,k_2}(i + n_1, j) = \phi_{k_1,k_2}(i,j)$. Applying the Laplacian:
\begin{align}
(\mathcal{L}\phi_{k_1,k_2})(i,j) &= 4\phi_{k_1,k_2}(i,j) - \phi_{k_1,k_2}(i+1,j) - \phi_{k_1,k_2}(i-1,j) - \phi_{k_1,k_2}(i,j+1) - \phi_{k_1,k_2}(i,j-1) \\
&= \left(4 - e^{2\pi \imath k_1/n_1} - e^{-2\pi \imath k_1/n_1} - e^{2\pi \imath k_2/n_2} - e^{-2\pi \imath k_2/n_2}\right)\phi_{k_1,k_2}(i,j) \\
&= \left(4 - 2\cos\!\left(\frac{2\pi k_1}{n_1}\right) - 2\cos\!\left(\frac{2\pi k_2}{n_2}\right)\right)\phi_{k_1,k_2}(i,j)
\end{align}
confirming that $\phi_{k_1,k_2}$ is an eigenvector with eigenvalue $\lambda_{k_1,k_2}$.
\end{proof}

\begin{remark}[Frequency Interpretation of the Eigenvalues]\label{rem:interpretacion_frecuencial_eigenvalores}
The eigenvalues of the Laplacian on the grid admit a direct frequency interpretation:
\begin{itemize}
    \item \textbf{Zero frequency}: $\lambda_{0,0} = 0$ corresponds to the DC component (constant signal); the eigenvector $\phi_{0,0} = 1/\sqrt{n_1 n_2}$ is uniform.
    \item \textbf{Low frequencies}: small eigenvalues ($k_1, k_2$ close to 0) correspond to smooth variations on the grid. A low-pass filter selects these modes.
    \item \textbf{High frequencies}: large eigenvalues (close to $\lambda_{\max} = 8$ for the square grid) correspond to rapid variations. A high-pass filter (such as the Laplacian) selects these modes.
    \item \textbf{Maximum frequency}: $\lambda_{n/2,n/2} = 8$ corresponds to the maximum possible oscillation (checkerboard pattern $(-1)^{i+j}$).
\end{itemize}
\end{remark}

\begin{example}[Low-Pass and High-Pass Filters on the Grid]\label{ex:filtros_paso_bajo_alto}
Classical image-processing filters have precise spectral interpretations:

\textbf{Smoothing filter ($3 \times 3$ mean)}: The kernel $K_{\text{media}} = \frac{1}{9}\begin{bsmallmatrix} 1 & 1 & 1 \\ 1 & 1 & 1 \\ 1 & 1 & 1 \end{bsmallmatrix}$ acts as a low-pass filter, attenuating the high-frequency components (large eigenvalues of the Laplacian) while preserving the low-frequency components.

\textbf{Sharpening filter}: The kernel $K_{\text{sharp}} = \begin{bsmallmatrix} 0 & -1 & 0 \\ -1 & 5 & -1 \\ 0 & -1 & 0 \end{bsmallmatrix} = I + \mathcal{L}$ amplifies the high-frequency components, where $I$ preserves the original signal and $\mathcal{L}$ adds the local differences.

In the spectral domain, the response of these filters is:
\begin{align}
\hat{K}_{\text{media}}(\lambda) &\approx 1 - c\lambda + O(\lambda^2) \quad \text{(low-pass)} \\
\hat{K}_{\text{sharp}}(\lambda) &= 1 + \lambda \quad \text{(amplification of high frequencies)}
\end{align}
\end{example}

\begin{proposition}[Connection between the Discrete and Continuous Laplacian]\label{prop:laplaciano_discreto_continuo}
The discrete grid Laplacian is a finite-difference approximation of the continuous Laplacian $\Delta f = \frac{\partial^2 f}{\partial x^2} + \frac{\partial^2 f}{\partial y^2}$. For a sufficiently smooth function $f$ defined on the grid with step $h$:
\[
(\mathcal{L}f)_{i,j} = h^2 \Delta f(ih, jh) + O(h^4)
\]
\end{proposition}

\begin{proof}
By the second-order Taylor expansion:
\begin{align}
f(i+1,j) &= f(i,j) + h f_x + \frac{h^2}{2}f_{xx} + \frac{h^3}{6}f_{xxx} + O(h^4) \\
f(i-1,j) &= f(i,j) - h f_x + \frac{h^2}{2}f_{xx} - \frac{h^3}{6}f_{xxx} + O(h^4)
\end{align}
Adding: $f(i+1,j) + f(i-1,j) = 2f(i,j) + h^2 f_{xx} + O(h^4)$. Analogously for the $j$ direction. Therefore:
\[
4f(i,j) - f(i+1,j) - f(i-1,j) - f(i,j+1) - f(i,j-1) = -h^2(f_{xx} + f_{yy}) + O(h^4) = -h^2 \Delta f + O(h^4)
\]
which shows that $(\mathcal{L}f)_{i,j} = -h^2 \Delta f + O(h^4)$, with negative sign consistent with the convention $\mathcal{L} = D - A$ (positive semidefinite Laplacian).
\end{proof}

\subsection{Spectral Generalization}\label{subsec:generalizacion_espectral}

\begin{theorem}[CNNs as Spectral Filters on Grids]\label{thm:cnns_filtros_espectrales_grillas}
Convolutional filters can be interpreted as spectral
filters over the grid Laplacian:
\[
h_\theta * f = \sum_k h_\theta(\lambda_k)\langle f, \phi_k\rangle\phi_k
\]
where \(\{\lambda_k, \phi_k\}\) are eigenvalues/eigenvectors of the grid Laplacian.
\end{theorem}

\begin{proof}
Let $\mathcal{L} = U \Lambda U^T$ be the spectral decomposition of the grid Laplacian, with $\Lambda = \operatorname{diag}(\lambda_0, \ldots, \lambda_{n-1})$ and $U = [\phi_0 | \cdots | \phi_{n-1}]$. A convolutional filter with kernel $h_\theta$ acts on the signal $f$ in the spatial domain as $(h_\theta * f)(i) = \sum_j h_\theta(i-j) f(j)$. In the spectral domain, by the convolution theorem, this is equivalent to $\widehat{h_\theta * f}(k) = \hat{h}_\theta(k) \cdot \hat{f}(k)$, where $\hat{f}(k) = \langle f, \phi_k \rangle$. Defining $h_\theta(\lambda_k) := \hat{h}_\theta(k)$, the filtered signal is reconstructed as:
\[
(h_\theta * f)(i) = \sum_k h_\theta(\lambda_k) \langle f, \phi_k \rangle \phi_k(i)
\]
which is exactly the spectral convolution of the statement.
\end{proof}

\begin{proposition}[From Regular Grids to Arbitrary Graphs]\label{prop:grillas_a_grafos}
The natural generalization of CNNs to arbitrary graphs is:
\begin{enumerate}
	\item Replace the regular grid by a general graph
	\item Replace the spatial kernel by functions of the Laplacian
	\item Maintain the idea of local aggregation of information
\end{enumerate}
\end{proposition}

\begin{proof}
The generalization is justified by the spectral equivalence established in \autoref{thm:cnns_filtros_espectrales_grillas}: since convolutional filters on grids are functions of the Laplacian $h_\theta(\mathcal{L})$, and the Laplacian $\mathcal{L} = D - A$ is defined for any graph (not only grids), the formulation $h_\theta(\mathcal{L})f = \sum_k h_\theta(\lambda_k)\langle f, \phi_k \rangle \phi_k$ extends directly to arbitrary graphs. The three steps of the statement correspond to: (1) replacing the grid by a graph with its own adjacency matrix, (2) using the spectral decomposition of the new Laplacian as a functional basis, and (3) preserving spatial locality through polynomial filters of degree $K$ that operate on neighborhoods of radius $K$.
\end{proof}

Spectral convolution on graphs (\autoref{def:convolucion_espectral_grafos}) allows the convolutional operation to be carried over to the spectral domain: $\text{conv}_\theta \star x = U g_\theta(\Lambda) U^T x$, where $\mathcal{L} = U\Lambda U^T$ is the spectral decomposition of the Laplacian and $g_\theta$ is the trainable filter.

\subsubsection{Detailed Spectral Analysis}

Spectral analysis provides the fundamental mathematical tools for understanding and designing convolutions over graphs as a generalization of classical \glspl{CNN}.

The spectral decomposition of the normalized Laplacian $\mathcal{L} = U \Lambda U^T$ and the graph Fourier transform $\hat{f} = U^T f$ (\autoref{def:transformada_fourier_grafos}, \autoref{ch:gdl}) allow spectral filters to be defined as functions $g: [0,2] \to \mathbb{R}$ that act on signals via $(g \star f) = U g(\Lambda) U^T f$. A $K$-localized filter concentrates its support on the first $K$ eigenvalues, inducing local operations with complexity $O(K|V|)$ that generalize the local kernels of \glspl{CNN}.

The Chebyshev polynomial approximation (\autoref{ex:chebyshev_aproximacion}) makes these spectral convolutions computationally tractable. The recurrence of \autoref{thm:eficiencia_chebyshev} reduces the complexity to $O(K|E|)$ without requiring the diagonalization of the Laplacian. See \citep{defferrard2017convolutionalneuralnetworksgraphs} for the original formulation in the context of \glspl{GNN}.

\textbf{Conceptual Bridge: Unifying Harmonic Analysis.} The harmonic analysis developed in \autoref{sec:analisis_armonico} provides the mathematical language to unify CNNs and GNNs:
\begin{enumerate}
    \item \textbf{Euclidean domain}: Classical Fourier transform (\autoref{def:transformada_fourier_continua})
    \item \textbf{Discrete grids}: DFT and FFT (Definitions~\ref{def:dft}, \ref{def:fft})
    \item \textbf{Finite/compact groups}: Fourier transform on groups (Equation~\ref{eq:fourier_grupo_finito})
    \item \textbf{General graphs}: Graph Fourier transform (\autoref{def:transformada_fourier_grafos}, \autoref{ch:gdl})
\end{enumerate}

Each level generalizes the previous one while preserving the fundamental intuition: low frequencies capture smooth/global structure, high frequencies capture local details/variations.

\subsubsection{From Spectral CNNs to Graph Neural Networks}

The preceding spectral theory establishes the formal bridge between traditional \glspl{CNN} and graph networks, demonstrating that both are special cases of the same fundamental mathematical principle.

\begin{theorem}[Fundamental CNN-GNN Connection]\label{thm:conexion_fundamental_cnn_gnn}
Every \gls{CNN} can be represented exactly as a spectral \gls{GNN} over a regular grid, where:
\begin{enumerate}
    \item The graph $G = (V, E)$ represents the pixel structure with $V = \{(i,j) : 1 \leq i,j \leq n\}$
    \item The connectivity $E$ is determined by the spatial neighborhood
    \item The spectral filters correspond to convolutional kernels via the discrete Fourier transform
\end{enumerate}
\end{theorem}

\begin{proof}[Proof sketch]
For a 2D grid of size $n \times n$, consider the graph $G$ where each pixel $(i,j)$ is a node and the edges connect adjacent pixels. The Laplacian of this graph has the form:
\[
\mathcal{L}_{(i,j),(i',j')} = \begin{cases}
4 & \text{if } (i,j) = (i',j') \\
-1 & \text{if } |i-i'| + |j-j'| = 1 \\
0 & \text{otherwise}
\end{cases}
\]

Its eigenvectors are exactly the two-dimensional discrete Fourier functions:
\[
\phi_{k,\ell}(i,j) = \frac{1}{n} e^{2\pi \imath (ki + \ell j)/n}
\]

Therefore, the spatial convolution $f * g$ corresponds exactly to the spectral multiplication $\hat{f} \odot \hat{g}$ in the graph's frequency domain. $\square$
\end{proof}

\subsubsection{Computationally Efficient Implementations}

The pioneering work of \citet{defferrard2017convolutionalneuralnetworksgraphs} solved the computational problem of spectral convolutions through Chebyshev polynomial approximations.

\begin{definition}[ChebNet: Chebyshev Convolution]\label{def:chebnet}
A spectral filter approximated by Chebyshev polynomials of order $K$ is defined as:
\[
g_\theta(\Lambda) = \sum_{k=0}^{K} \theta_k T_k(\tilde{\Lambda})
\]
where $T_k$ are the Chebyshev polynomials of the first kind, $\tilde{\Lambda} = 2\Lambda/\lambda_{\max} - I$ is the scaled Laplacian, and $\theta_k$ are trainable parameters.
\end{definition}

\begin{theorem}[Localization of Chebyshev Filters]\label{thm:localizacion_chebyshev}
A Chebyshev filter of order $K$ is strictly $K$-localized: the response at node $i$ depends only on nodes at maximum distance $K$ from $i$ in the graph.
\end{theorem}

\begin{proof}
A Chebyshev polynomial $T_k$ of degree $k$ satisfies $T_k(\tilde{\mathcal{L}}) = 2\tilde{\mathcal{L}} T_{k-1}(\tilde{\mathcal{L}}) - T_{k-2}(\tilde{\mathcal{L}})$. Since $\tilde{\mathcal{L}}_{ij} \neq 0$ only if nodes $i$ and $j$ are adjacent (distance 1), multiplication by $\tilde{\mathcal{L}}$ extends the support of a signal by exactly one hop. Thus, $T_1(\tilde{\mathcal{L}})f$ depends on neighbors at distance $\leq 1$, and by induction, $T_k(\tilde{\mathcal{L}})f$ at node $i$ depends only on nodes at distance $\leq k$. Since the filter is $g_\theta(\tilde{\mathcal{L}}) = \sum_{k=0}^K \theta_k T_k(\tilde{\mathcal{L}})$, the response at each node depends on nodes at distance $\leq K$.
\end{proof}

The efficient evaluation of these filters is obtained through the Chebyshev recurrence established in \autoref{thm:eficiencia_chebyshev}, which allows $g_\theta(\tilde{\mathcal{L}}) \mathbf{x}$ to be computed in time $O(K|\mathcal{E}|)$ without diagonalizing the Laplacian, with only $O(K)$ trainable parameters independent of the graph size.

This computational improvement eliminates the need to compute the full spectral decomposition, making spectral \glspl{GNN} practical for large-scale graphs.

\subsubsection{Concrete Examples of CNN $\rightarrow$ GNN Translation}

\begin{example}[Sobel Kernel as a Spectral Filter]
The Sobel edge-detection kernel:
\[
K_{\text{Sobel}} = \begin{bmatrix} -1 & 0 & 1 \\ -2 & 0 & 2 \\ -1 & 0 & 1 \end{bmatrix}
\]

In the spectral domain of the grid graph, this filter corresponds to:
\[
\hat{g}(\lambda_{k,\ell}) = |\sin(2\pi k/n)| \cdot \text{vertical frequency function}
\]

This example demonstrates how feature-detection filters in \glspl{CNN} have natural spectral interpretations on graphs.
\end{example}

\begin{example}[Pooling as Graph Coarsening]
The $2 \times 2$ max-pooling of \glspl{CNN} corresponds to:
\begin{enumerate}
    \item \textbf{Grouping}: Partition each $2 \times 2$ block of pixels
    \item \textbf{Aggregation}: Apply the max function over each partition
    \item \textbf{New graph}: Build a coarser graph with $n/4$ nodes
\end{enumerate}

This generalizes to arbitrary graphs through algebraic coarsening algorithms that preserve the local spectral structure.
\end{example}

\subsection{Implications and Transition}\label{subsec:implicaciones_transicion_cnn_gnn}

This \gls{CNN}-\gls{GNN} equivalence has profound implications that transcend mere mathematical analogy and provide a conceptual framework for understanding the limitations and extensions of convolutional architectures.

\subsubsection{Conceptual Unification}

\begin{proposition}[Shared CNN-GNN Principles]\label{prop:principios_compartidos}
\glspl{CNN} and \glspl{GNN} share three fundamental architectural principles:
\begin{enumerate}
    \item \textbf{Locality}: Aggregation of information from local neighborhoods
    \item \textbf{Weight sharing}: Reuse of parameters across the domain
    \item \textbf{Hierarchical composition}: Construction of multi-scale representations
\end{enumerate}
\end{proposition}

\begin{proof}
(1) \textit{Locality}: In a CNN, each convolutional layer aggregates information from a neighborhood of size $k$ (the kernel). In a GNN, the message passing step aggregates information from the 1-hop neighbors $\mathcal{N}(i)$. In both cases, the operation is local and is determined by the adjacency structure of the domain.

(2) \textit{Weight sharing}: In a CNN, the same kernel is applied at all spatial positions. In a GNN, the same message function $\phi$ and update function $\gamma$ are applied at all nodes. In both cases, the parameters are independent of position.

(3) \textit{Hierarchical composition}: In a CNN, pooling operations reduce the spatial resolution, creating multi-scale representations. In a GNN, coarsening algorithms (such as Graclus or DiffPool) collapse nodes, producing progressively smaller graphs with higher-level representations.
\end{proof}

\subsubsection{Limitations of CNNs and Their Overcoming through GNNs}

\begin{remark}[Structural Limitations of CNNs]\label{rem:limitaciones_estructurales_cnns}
Traditional \glspl{CNN} present the following fundamental restrictions, a direct consequence of their exclusive equivariance to translations (\autoref{theorem:convolution_translation_equivariance}):
\begin{enumerate}
    \item \textbf{Regular-grid dependence}: They require a fixed Euclidean structure
    \item \textbf{Topological rigidity}: They cannot directly process irregular domains
    \item \textbf{Loss of geometric information}: The Euclidean representation can destroy intrinsic relationships
    \item \textbf{Limited scalability}: Quadratic complexity with resolution for non-local structures
\end{enumerate}
\end{remark}

\begin{example}[Cases where CNNs Fail and GNNs Succeed]
\leavevmode
\begin{enumerate}
    \item \textbf{Molecular analysis}: Molecules are naturally graphs, not grids
    \item \textbf{Social networks}: Social connectivity does not respect Euclidean geometry
    \item \textbf{3D meshes}: Complex surfaces with variable topology
    \item \textbf{Transportation systems}: Networks with non-uniform connectivity
\end{enumerate}
\end{example}

\subsubsection{Bidirectional Transfer of Techniques}

\begin{proposition}[Transferability of Innovations]
Techniques developed for \glspl{CNN} adapt naturally to \glspl{GNN} and vice versa:
\begin{itemize}
    \item \textbf{CNN $\rightarrow$ GNN}: Pooling, dropout, batch normalization, skip connections
    \item \textbf{GNN $\rightarrow$ CNN}: Attention mechanisms, adaptive pooling, graph-inspired regularization
\end{itemize}
\end{proposition}

\begin{proof}
Transferability is grounded in the structural equivalence established in \autoref{thm:conexion_fundamental_cnn_gnn}. CNN techniques that depend only on local operations and weight sharing transfer to GNNs because both architectures share these principles (see \autoref{prop:principios_compartidos}): pooling corresponds to graph coarsening, dropout is applied to nodes/edges, batch normalization is adapted per node, and skip connections preserve the structure. In the reverse direction, the attention mechanisms of GNNs (which weight neighbors adaptively) transfer to CNNs as spatial attention, and adaptive pooling techniques on graphs inspire non-uniform pooling methods on images.
\end{proof}

\subsubsection{Framework of Progressive Generalization}

\begin{theorem}[Generalization Hierarchy]\label{thm:jerarquia_generalizacion_cnn}
Neural architectures admit the following generalization hierarchy:
\begin{align}
\text{Perceptron} &\subset \text{MLP} \subset \text{CNN (grid)} \subset \text{spectral CNN} \nonumber \\
&\subset \text{spectral GNN} \subset \text{spatial GNN} \subset \text{general MPNN}
\end{align}
where each inclusion represents the relaxation of a specific architectural restriction.
\end{theorem}

\begin{corollary}[Increasing Expressive Power]\label{cor:poder_expresivo_creciente}
Each level of the above hierarchy can express strictly more functions than the previous levels, at the cost of greater computational complexity and risk of overfitting.
\end{corollary}

\begin{proof}
Each inclusion in the hierarchy relaxes a specific architectural restriction: CNN $\subset$ spectral CNN relaxes the spatial locality of the kernel, spectral CNN $\subset$ spectral GNN relaxes the regular-grid structure, etc. The inclusion is strict in each case because the relaxation allows the expression of functions that violate the previous restriction. For example, there exist functions equivariant with respect to the symmetry group of the graph that cannot be implemented as convolutions over a grid, since the grid imposes translation symmetries that the general graph does not possess. The increase in expressive power entails greater computational complexity (more parameters and operations per layer) and greater risk of overfitting (a larger hypothesis space with the same amount of data).
\end{proof}

\subsubsection{Interpretability from a Graph Perspective}

\textbf{New Interpretation of CNNs.} The graph perspective provides new tools for interpreting \glspl{CNN}:
\begin{itemize}
    \item \textbf{Spectral analysis}: Understanding filters as spectral functions
    \item \textbf{Localization}: Quantifying the receptive field using graph distances
    \item \textbf{Robustness}: Analyzing stability using graph perturbations
    \item \textbf{Transferability}: Evaluating generalization through spectral properties
\end{itemize}

\subsubsection{Bridge towards Geometric Extensions}

This fundamental \gls{CNN}-\gls{GNN} equivalence establishes the conceptual foundation for the geometric extensions we will explore:

\begin{definition}[Taxonomy of Geometric Extensions]\label{def:taxonomia_extensiones}
Starting from the CNN-GNN basis, we identify three main directions of generalization:
\begin{enumerate}
    \item \textbf{Group Convolutions}: Generalize translations to richer transformation groups
    \item \textbf{Convolutions on Manifolds}: Extend to domains with non-trivial curvature
    \item \textbf{Gauge Convolutions}: Incorporate gauge invariances for fiber-bundle structures
\end{enumerate}
\end{definition}

Each of these directions maintains the fundamental principles established in this section while extending the expressive power of convolutional architectures to increasingly sophisticated geometric domains.

\subsubsection{Historical Context and Development}

\textbf{Historical CNN-GNN Evolution.} The connection between \glspl{CNN} and \glspl{GNN} developed through several foundational works:

\begin{enumerate}
    \item \textbf{2013}: \citet{bruna2014spectralnetworkslocallyconnected} establish the first spectral convolutions on graphs, demonstrating the theoretical equivalence with traditional \glspl{CNN}

    \item \textbf{2016}: \citet{defferrard2017convolutionalneuralnetworksgraphs} solve the computational problems through Chebyshev approximations, making spectral \glspl{GNN} practical

    \item \textbf{2017}: The works on GraphSAGE and GAT introduce spatial approaches, complementing the spectral techniques

    \item \textbf{2017-present}: Synthesis within the \gls{GDL} framework, unifying both approaches under common geometric principles
\end{enumerate}

\begin{remark}[Impact of the Unified Framework]\label{rem:impacto_marco_unificado}
The CNN-GNN unification has had three fundamental impacts:
\begin{enumerate}
    \item \textbf{Theoretical}: It provided rigorous mathematical foundations for the design of architectures
    \item \textbf{Practical}: It enabled the systematic transfer of techniques between domains
    \item \textbf{Applied}: It enabled the efficient processing of non-Euclidean data in real applications
\end{enumerate}
\end{remark}

\subsubsection{Mathematical Construction of the CNN-GNN Correspondence}

The mathematical translation between CNNs and GNNs is grounded in two main structural transformations:

\begin{definition}[Embedding of an Image into a Graph]\label{def:inmersion_imagen_grafo}
An image $I \in \mathbb{R}^{H \times W}$ is canonically embedded into a graph $G = (V, E, W)$ by means of:
\begin{align}
V &= \{(i,j) : 1 \leq i \leq H, 1 \leq j \leq W\} \\
E &= \{((i,j), (i',j')) : |i-i'| + |j-j'| \leq \delta\}
\end{align}
where $\delta \in \{1, 2\}$ determines the connectivity (4 or 8 neighbors) and $W$ may incorporate weights based on intensity similarity to capture non-uniform local relationships.
\end{definition}

\begin{theorem}[Spectral CNN-GNN Correspondence]\label{thm:correspondencia_espectral_cnn_gnn}
Given a CNN with a layer architecture $\{L_\ell\}$ and kernels $\{K_\ell\}$, there exists a canonical transformation to an equivalent GNN by means of:
\begin{enumerate}
    \item \textbf{Spectral transform}: Each kernel $K_\ell$ induces a spectral filter $g_\ell(\lambda) = \mathcal{F}\{K_\ell\}(\lambda)$ where $\mathcal{F}$ denotes the discrete Fourier transform
    \item \textbf{Polynomial approximation}: The filter $g_\ell$ is approximated by $\hat{g}_\ell(\lambda) = \sum_{k=0}^{K_\ell} \theta_{\ell,k} T_k(\lambda)$ with coefficients $\theta_{\ell,k}$ determined by least squares
    \item \textbf{Architectural preservation}: The composition of layers remains identical, guaranteeing exact functional equivalence
\end{enumerate}
\end{theorem}

\begin{figure}[htbp]
	\centering
	\includegraphics[width=0.95\textwidth]{figures/cnn_to_spectral_graph}
	\caption[Transformation of a CNN kernel to its spectral representation on a grid graph]{Transformation of a CNN kernel to its spectral representation on a grid graph. Left: $3 \times 3$ kernel applied over an image grid with values proportional to the intensity. Center: representation of the same domain as a graph with 4-connectivity. Right: spectral filter $g(\lambda)$ and its approximation through Chebyshev polynomials $\hat{g}(\lambda) = \sum_k \theta_k T_k(\lambda)$.}
	\label{fig:cnn_to_spectral_graph}
\end{figure}

The next section will begin with group convolutions, where the regular grid structure is generalized to more complex homogeneous spaces, extending the unifying CNN-GNN framework towards richer geometric transformations.
\section{Group Convolutions}\label{sec:group_convolutions}

The first significant extension of classical \glspl{CNN} emerges upon recognizing that translations represent only a special case of a broader group of geometric transformations. Group convolutions provide a rigorous theoretical framework for incorporating additional symmetries such as rotations, reflections, and more complex transformations, while preserving the equivariance properties that make convolutional architectures powerful.

This section explores how group theory and harmonic analysis provide the mathematical tools needed to design neural architectures that automatically respect the symmetries inherent to different data domains. The foundations of group representation theory are found in the classical works of \citet{Peter1927} and \citet{Serre1977}.

\subsection{Theoretical Foundations}\label{subsec:fundamentos_teoricos}

The group-theoretic foundations needed for group convolutions (groups, group actions, equivariance, homogeneous spaces) are fully developed in \autoref{ch:gdl}, Sections~\ref{subsec:teoria_grupos_representaciones} and~\ref{subsec:equivarianza_invarianza}. This section focuses on the specific application of these concepts to geometric convolutions.

Recall that a function $\Phi$ between function spaces is \textit{$G$-equivariant} (\autoref{def:equivariance_gdl}, \autoref{ch:gdl}) if $\Phi(L_g f) = L_g' (\Phi(f))$ for all $g \in G$, and \textit{$G$-invariant} if $\Phi(L_g f) = \Phi(f)$. Equivariance generalizes the translation invariance of classical \glspl{CNN}: if the data possess symmetries under $G$, $G$-equivariant functions preserve these symmetries throughout processing, resulting in more robust and parameter-efficient representations.

\subsubsection{Convolution on Groups}

The generalization of the convolution operation to arbitrary groups constitutes the theoretical core of the geometric extensions of \glspl{CNN}.

\begin{definition}[Group Convolution]
For functions $f, \psi: G \to \mathbb{R}$ defined over a group $G$, the \textbf{group convolution} is defined as:
\begin{equation}
(f *_G \psi)(u) = \int_G f(uv^{-1}) \psi(v) \, d\mu(v)
\label{eq:group_convolution}
\end{equation}
where $\mu$ is the Haar measure on $G$.
\end{definition}

This definition generalizes the classical convolution by replacing the subtraction operation $u - v$ with the group operation $uv^{-1}$. For discrete groups, the integral becomes a sum.

The following theorem establishes the fundamental connection between equivariance and convolution in the context of neural networks:

\begin{theorem}[Kondor-Trivedi, \citep{kondor2018generalizationconvolutionalneural}]
A feedforward neural network $N$ is equivariant to the action of a compact group $G$ on its inputs if and only if each layer of $N$ implements a generalized form of convolution derived from equation \ref{eq:group_convolution}.
\end{theorem}

\begin{proof}[Sketch of the proof]
The proof is based on two directions:

\textbf{Forward direction:} If each layer implements group convolution, then by construction it preserves the group structure, which implies equivariance.

\textbf{Converse direction:} If the network is equivariant, then via representation-theoretic techniques one can show that the linear transformations between layers must have the form of generalized convolutions.

The complete proof uses harmonic analysis and representation theory to characterize all possible equivariant linear transformations.
\end{proof}

\subsubsection{Representation Theory}

Representation theory provides the mathematical tools to analyze how groups act on vector spaces, which is essential for the design of equivariant architectures. The fundamental concepts of group representations, including irreducible representations and the decomposition theorem, are developed in \autoref{ch:gdl}, Section~\ref{subsec:teoria_grupos_representaciones}.

In summary, a \textit{representation} is a homomorphism $\rho: G \to GL(V)$ that allows one to ``linearize'' the action of the group, and the \textit{irreducible representations} (\autoref{def:representacion_irreducible}) are the fundamental building blocks from which all other representations are constructed.

\subsubsection{Harmonic Analysis on Groups}

Harmonic analysis generalizes concepts such as the Fourier transform to non-abelian groups, providing powerful tools for frequency analysis in more general domains.

The foundations of this harmonic analysis on groups are based on representation theory. The basic concepts of groups, homomorphisms, and Lie groups are established in \autoref{ch:gdl}, Section~\ref{subsec:teoria_grupos_representaciones}, while here we develop the specific aspects of representations in the context of geometric convolutions.

For the complete theory of group representations, including definitions, decomposition theorems, and the Peter-Weyl theorem, see \autoref{ch:gdl}, Section~\ref{subsec:teoria_grupos_representaciones}. Here we focus on its specific application to the harmonic analysis of convolutions.

In the context of equivariant \glspl{CNN}, each layer processes features that transform under specific representations of the symmetry group (\autoref{def:representacion_grupo}, \autoref{ch:gdl}). The choice of these representations (that is, of the homomorphism $\rho: G \to GL(V)$ where $V$ is the feature space) determines what types of patterns the network can detect equivariantly.

\begin{example}[Regular Representation in \glspl{GCNN}]\label{ex:representacion_regular_cnn}
For a finite group $G = \{g_1, g_2, \ldots, g_n\}$, the \textbf{left regular representation} acts on functions $f: G \rightarrow \mathbb{R}^d$ via:
\begin{equation}
(\rho(g)f)(h) = f(g^{-1}h)
\end{equation}
This representation is fundamental in G-\glspl{CNN}, where the channels correspond to evaluations of the function at different group elements.
\end{example}

\textbf{Decomposition of Representations.} By the decomposition theorem (\autoref{thm:descomposicion_representaciones}, \autoref{ch:gdl}), every representation of a compact group decomposes into a direct sum of irreducible representations. In the context of G-\glspl{CNN}, this has direct consequences for architecture design.

\textbf{Implications for CNN Design.} The Peter-Weyl theorem (\autoref{thm:peter_weyl}, \autoref{ch:gdl}) guarantees that $L^2(G) = \bigoplus_{\lambda \in \widehat{G}} n_\lambda \mathcal{H}_\lambda$, where each $\mathcal{H}_\lambda$ is the space of an irreducible representation. This has direct consequences for G-\glspl{CNN}: the decomposition into irreducible representations $\mathcal{H}_\lambda$ determines the optimal channel structure. Each $\mathcal{H}_\lambda$ corresponds to a channel type, and $n_\lambda$ indicates how many channels of that type to include in order to fully capture the symmetries of the group.

\subsubsection{Fourier Analysis on Groups}

The Fourier transform on groups (Equation~\ref{eq:fourier_grupo_finito}) generalizes to the continuous case by replacing the sum with an integral over the Haar measure: $\hat{f}(\rho_i) = \int_G f(g) \rho_i(g) \, d\mu(g)$. The convolution theorem (Equation~\ref{eq:convolucion_fourier_grupos}) extends identically, turning the group convolution into a matrix product of the transforms.

\subsubsection{Homogeneous Spaces}

In practical applications, functions are defined over homogeneous spaces $X = G/H$ where $H$ is a subgroup of $G$, rather than over the full group.

\begin{definition}[Homogeneous Space]
A space $X$ is \textbf{homogeneous} under the action of $G$ if for any $x, y \in X$, there exists $g \in G$ such that $g \cdot x = y$.
\end{definition}

The sphere $S^2$ is homogeneous under $SO(3)$, and regular images are homogeneous under translations. The generalization of convolutions to homogeneous spaces requires additional lifting and projection techniques.

\subsubsection{Connection with Classical CNNs}

Classical \glspl{CNN} are a special case of group convolutions where $G = (\mathbb{Z}^2, +)$ acts by translations on discrete images. In this case:

\begin{itemize}
    \item The irreducible representations are one-dimensional (characters)
    \item The group convolution reduces to the classical convolution
    \item The filters are translation-invariant
\end{itemize}

The geometric extensions incorporate richer groups such as $\gls{SO2}$ for rotations, $SE(2)$ for planar rigid motions, or $\gls{SO}$ for three-dimensional rotations.

\subsubsection{Unified Theoretical Framework}

The theoretical framework presented provides the foundations for developing specific architectures:

\begin{enumerate}
    \item \textbf{Group identification}: Determine which symmetries are present in the data
    \item \textbf{Representation analysis}: Characterize the relevant irreducible representations
    \item \textbf{Convolution design}: Implement equivariant operations efficiently
    \item \textbf{Optimization}: Exploit the group structure to reduce parameters
\end{enumerate}

This theoretical framework establishes the rigorous mathematical foundations on which the various extensions of group convolutions we will explore in the following subsections are built, from finite groups to complex manifolds, always preserving the fundamental principles of equivariance and computational efficiency.

A deep understanding of these theoretical foundations is essential to appreciate both the capabilities and the limitations of each specific extension, as well as for the design of new architectures that exploit particular symmetries present in specific application domains.

\subsection{Equivariant Universality Theorem}\label{subsec:teorema_universalidad_equivariante}

A fundamental result in the theory of group convolutions establishes that equivariant architectures can approximate any continuous equivariant function. This theorem provides theoretical justification for the use of group convolutions and guarantees their expressive power.

\begin{theorem}[Equivariant Universality for CNNs]\label{theorem:equivariant_universality_cnns}
    Let $G$ be a compact group acting on a compact domain $\mathcal{D}$, and let $\mathcal{F}_G$ be the space of continuous $G$-equivariant functions from $L^2(\mathcal{D})$ to $L^2(\mathcal{D})$. For any $f \in \mathcal{F}_G$ and $\epsilon > 0$, there exists a $G$-equivariant convolutional network $\Phi_{\text{G-CNN}}$ such that:
    \begin{equation}
        \|f - \Phi_{\text{G-CNN}}\|_{L^2} < \epsilon
    \end{equation}
    That is, \glspl{GCNN} are universal approximators in the space of equivariant functions.
\end{theorem}

\subsubsection{Constructive Proof for the Discrete Group $\mathbb{Z}^2$}

We present a simplified proof for the case of the discrete translation group $\mathbb{Z}^2$, which illustrates the main ideas of the general theorem.

\textbf{Step 1: Spectral Decomposition}

For a function $f: \mathbb{Z}^2 \to \mathbb{R}$ of compact support, the discrete Fourier transform allows us to write:
\begin{equation}
    f(x) = \sum_{k \in \mathbb{Z}^2} \hat{f}(k) e^{2\pi i \langle k, x \rangle / N}
\end{equation}

\textbf{Step 2: Approximation by Convolutions}

Each frequency component can be approximated by a convolution with an appropriate filter. For frequency $k$, we define:
\begin{equation}
    g_k(x) = \frac{1}{N^2} e^{2\pi i \langle k, x \rangle / N}
\end{equation}

\textbf{Step 3: Construction of the Network}

A convolutional network with $M$ filters in the first layer can be expressed as:
\begin{equation}
    \Phi_{\text{CNN}}(f) = \sigma\left(\sum_{j=1}^M w_j (f * h_j) + b\right)
\end{equation}

where $h_j$ are the learnable filters, $w_j$ are weights, $b$ is the bias, and $\sigma$ is the activation function.

\textbf{Step 4: Equivariance by Construction}

Equivariance is automatically preserved:
\begin{align}
    \Phi_{\text{CNN}}(L_g f) &= \sigma\left(\sum_{j=1}^M w_j (L_g f * h_j) + b\right) \\
    &= \sigma\left(\sum_{j=1}^M w_j L_g(f * h_j) + b\right) \\
    &= L_g \Phi_{\text{CNN}}(f)
\end{align}

where we use the property $L_g(f * h) = (L_g f) * h = L_g(f * h)$ for convolutions.

\subsubsection{Concrete Example: Equivariant Edge Detection}

Consider the function $f: \mathbb{Z}^2 \to \mathbb{R}$ that detects vertical edges. A classical CNN with a vertical Sobel filter:
\begin{equation}
    h = \begin{pmatrix} -1 & 0 & 1 \\ -2 & 0 & 2 \\ -1 & 0 & 1 \end{pmatrix}
\end{equation}

is equivariant to translations: if we translate the input image, the feature map also translates in the same way.

\subsubsection{Extension to More General Groups}

For non-abelian groups such as $SO(2)$ or finite groups such as $D_4$, the proof extends using:

\begin{enumerate}
    \item \textbf{Harmonic analysis}: Decomposition into irreducible representations
    \item \textbf{Group convolutions}: Replacement of the classical convolution by the corresponding group operation
    \item \textbf{Representation theory}: Guarantee that the transformations preserve the equivariant structure
\end{enumerate}

\begin{corollary}[Parametric Efficiency]\label{corollary:parametric_efficiency}
    $G$-equivariant networks require $|G|$ times fewer parameters than fully connected networks to approximate $G$-equivariant functions with the same accuracy, where $|G|$ is the order of the group.
\end{corollary}

\begin{proof}
Let $f: X \to Y$ be a $G$-equivariant function. A fully connected network (without symmetry constraints) must independently learn the value of $f$ at each point $x$ and at all its transformations $\{g \cdot x : g \in G\}$. This requires $N$ parameters to represent $f$ with accuracy $\epsilon$.

A $G$-equivariant network, by construction, satisfies $\Phi(g \cdot x) = g \cdot \Phi(x)$ for all $g \in G$. Therefore, it suffices to learn $\Phi$ on a fundamental domain of the action of $G$, which has size $|X|/|G|$ of the original domain. The values on the remaining $|G| - 1$ copies are obtained automatically by equivariance, which reduces the number of parameters needed to $N/|G|$ for the same approximation accuracy $\epsilon$.
\end{proof}

This theoretical result mathematically justifies why equivariant architectures are more efficient: they automatically exploit the symmetries present in the data, reducing the hypothesis space without loss of expressive power.

\subsection{Practical Implementation of Group Convolutions}\label{subsec:implementacion_practica_gcnn}

Before exploring the specific applications to finite and continuous groups, it is essential to understand how the abstract concepts of group theory translate into concrete computational implementations. This section presents practical implementations that demonstrate the direct connection between the mathematical theory developed in the previous subsections and functional deep learning architectures.

\subsubsection{Modular Implementation Framework}\label{subsubsec:marco_implementacion_modular}

The implementation of \glspl{GCNN} requires a modular architecture that separates the abstract definition of the group from its specific application in convolution operations. The design presented below facilitates extension to different groups while preserving the mathematical correctness of the equivariant operations.

\paragraph{Group Abstraction}

The mathematical implementation of group convolutions requires an \textbf{algebraic abstraction} that captures the essential structure of any group $G$. This abstraction must provide:

\begin{enumerate}
    \item \textbf{Set of elements}: A mathematical representation of the set $G$ with its group operation
    \item \textbf{Group action}: A function $\cdot: G \times X \to X$ that defines how the group elements transform the elements of the domain
    \item \textbf{Group operation}: The composition law $*: G \times G \to G$ that defines the multiplication in the group
    \item \textbf{Inverse element}: A function $\text{inv}: G \to G$ that assigns to each element its inverse
\end{enumerate}

This mathematical abstraction makes it possible to develop the theory of group convolutions uniformly, guaranteeing that the operations respect the specific algebraic structure of each group and preserve the required equivariance properties.

\paragraph{Mathematical Structure of the Dihedral Group $D_4$}

The dihedral group $D_4$ represents the symmetries of the square and serves as a paradigmatic example for understanding the theory of group convolutions for finite groups. Its algebraic structure is characterized by:

\begin{definition}[Elements of $D_4$]
The dihedral group $D_4$ contains 8 elements that can be represented as:
\begin{equation}
D_4 = \{e, r, r^2, r^3, s, sr, sr^2, sr^3\}
\end{equation}
where:
\begin{itemize}
    \item $e$ is the identity element
    \item $r$ is a $90^\circ$ rotation (with $r^4 = e$)
    \item $s$ is a reflection (with $s^2 = e$)
    \item The fundamental relation is $srs = r^{-1}$
\end{itemize}
\end{definition}

\begin{proposition}[Action of $D_4$ on the Plane]
The natural action of $D_4$ on $\mathbb{R}^2$ is defined by:
\begin{align}
r \cdot (x,y) &= (-y, x) \quad \text{($90^\circ$ rotation)} \\
s \cdot (x,y) &= (x, -y) \quad \text{(horizontal reflection)}
\end{align}
All other elements are obtained by composition of these basic transformations.
\end{proposition}

\begin{proof}
The rotation $r$ by $90°$ acts on $(x,y)$ via the matrix $R_{90} = \begin{psmallmatrix} 0 & -1 \\ 1 & 0 \end{psmallmatrix}$, giving $R_{90}(x,y)^T = (-y,x)^T$. The reflection $s$ about the $x$-axis acts via $S = \begin{psmallmatrix} 1 & 0 \\ 0 & -1 \end{psmallmatrix}$, giving $S(x,y)^T = (x,-y)^T$. The 8 elements of $D_4$ are generated by composition: $r^2 \cdot (x,y) = (-x,-y)$, $r^3 \cdot (x,y) = (y,-x)$, $sr \cdot (x,y) = (-y,-x)$, $sr^2 \cdot (x,y) = (-x,y)$, $sr^3 \cdot (x,y) = (y,x)$. These 8 transformations are the only linear isometries of $\mathbb{R}^2$ that preserve the unit square.
\end{proof}

\paragraph{Mathematical Verification of the Group Axioms}

The group structure of $D_4$ is verified through rigorous checking of the fundamental axioms:

\begin{theorem}[Axioms of $D_4$]
The set $D_4 = \{e, r, r^2, r^3, s, sr, sr^2, sr^3\}$ with the composition operation satisfies:

\begin{enumerate}
    \item \textbf{Closure}: For any $g_1, g_2 \in D_4$, we have $g_1 \circ g_2 \in D_4$
    \item \textbf{Associativity}: $(g_1 \circ g_2) \circ g_3 = g_1 \circ (g_2 \circ g_3)$ for all $g_1, g_2, g_3 \in D_4$
    \item \textbf{Identity element}: There exists $e \in D_4$ such that $e \circ g = g \circ e = g$ for all $g \in D_4$
    \item \textbf{Inverse element}: For each $g \in D_4$, there exists $g^{-1} \in D_4$ such that $g \circ g^{-1} = g^{-1} \circ g = e$
\end{enumerate}
\end{theorem}

\begin{proof}[Sketch of proof]
The verification is carried out by explicitly constructing the multiplication table and directly checking each property. Closure and associativity are checked exhaustively for all $8^3 = 512$ possible triples, while the existence of the identity and inverses is established by direct construction.
\end{proof}

\subsubsection{Mathematical Theory of Equivariant Convolution}\label{subsubsec:convolucion_equivariante_matematica}

Equivariant convolution requires a rigorous mathematical formulation that captures the essence of the lifting and group convolution operations.

\paragraph{Mathematical Lifting Operation}

The first fundamental operation in a group convolution is the \textbf{lifting}, which transforms signals defined on the Euclidean space $\mathbb{R}^2$ into functions defined over the group $G$.

\begin{definition}[Lifting Convolution]\label{def:convolucion_elevacion}
Given a function $f: \mathbb{R}^2 \to \mathbb{R}$ and a filter $\psi: \mathbb{R}^2 \to \mathbb{R}$, the \textbf{lifting convolution} produces a function $F: G \to \mathbb{R}$ defined by:
\begin{equation}
(f \star \psi)(g) = \int_{\mathbb{R}^2} f(x) \psi(g^{-1} \cdot x) \, dx
\end{equation}
where $g \in G$ and $g^{-1} \cdot x$ denotes the action of the inverse element $g^{-1}$ on the point $x \in \mathbb{R}^2$.
\end{definition}

\begin{theorem}[Equivariance of Lifting]\label{thm:equivarianza_elevacion}
The lifting operation is equivariant in the sense that:
\begin{equation}
(L_h f \star \psi)(g) = (f \star \psi)(h^{-1}g)
\end{equation}
where $L_h f(x) = f(h^{-1} \cdot x)$ is the action of the group on the input function.
\end{theorem}

\begin{proof}
By definition:
\begin{align}
(L_h f \star \psi)(g) &= \int_{\mathbb{R}^2} (L_h f)(x) \psi(g^{-1} \cdot x) \, dx \\
&= \int_{\mathbb{R}^2} f(h^{-1} \cdot x) \psi(g^{-1} \cdot x) \, dx
\end{align}
Performing the change of variable $y = h^{-1} \cdot x$ (with constant Jacobian by the group action), we obtain:
\begin{align}
&= \int_{\mathbb{R}^2} f(y) \psi(g^{-1} \cdot h \cdot y) \, dy \\
&= \int_{\mathbb{R}^2} f(y) \psi((h^{-1}g)^{-1} \cdot y) \, dy \\
&= (f \star \psi)(h^{-1}g)
\end{align}
\end{proof}

\paragraph{Mathematical Group Convolution}

The operations after lifting implement the \textbf{full group convolution}, which operates between functions defined over the group.

\begin{definition}[Group Convolution]\label{def:convolucion_grupo}
For functions $f, \psi: G \to \mathbb{R}$, the \textbf{group convolution} is defined as:
\begin{equation}
(f \star_G \psi)(u) = \sum_{g \in G} f(g) \psi(g^{-1}u)
\end{equation}
where the sum is over all elements of the group $G$. This definition is consistent with the continuous convolution (\autoref{eq:group_convolution}), where the change of variable $v \mapsto g = uv^{-1}$ transforms $\int_G f(uv^{-1})\psi(v)\,d\mu(v)$ into $\sum_g f(g)\psi(g^{-1}u)$.
\end{definition}

\begin{theorem}[Equivariance of Group Convolution]\label{thm:equivarianza_convolucion_grupo}
The group convolution preserves equivariance:
\begin{equation}
(L_h f) \star_G \psi = L_h (f \star_G \psi)
\end{equation}
where $L_h$ denotes the action of the group by left translation.
\end{theorem}

\begin{proof}
By the definition of the action by left translation, $(L_h f)(g) = f(h^{-1}g)$. Then:
\begin{align}
((L_h f) \star_G \psi)(u) &= \sum_{g \in G} (L_h f)(g) \psi(g^{-1}u) \\
&= \sum_{g \in G} f(h^{-1}g) \psi(g^{-1}u)
\end{align}
Performing the change of variable $g' = h^{-1}g$ (equivalently $g = hg'$):
\begin{align}
&= \sum_{g' \in G} f(g') \psi((hg')^{-1}u) \\
&= \sum_{g' \in G} f(g') \psi(g'^{-1}h^{-1}u) \\
&= (f \star_G \psi)(h^{-1}u) \\
&= (L_h (f \star_G \psi))(u)
\end{align}
\end{proof}

\textbf{Algebraic Structure.} The group convolution inherits the structure of the group algebra $\mathbb{R}[G]$, where multiplication is given by convolution. This algebraic structure is fundamental to guaranteeing that the composition of layers preserves equivariance.

\subsubsection{Mathematical D4-Equivariant Architecture}\label{subsubsec:arquitectura_d4_matematica}

The combination of the previous theoretical components makes it possible to construct an architecture that is fully equivariant to the group $D_4$ through the mathematical composition of equivariant transformations.

\begin{definition}[D4-Equivariant Network]\label{def:red_d4_equivariante}
A neural network $\Phi: \mathcal{F}(\mathbb{R}^2) \to \mathbb{R}^C$ is \textit{$D_4$-equivariant} if it can be decomposed as:
\begin{equation}
\Phi = \Phi_{\text{inv}} \circ \Phi_{G} \circ \Phi_{G} \circ \Phi_{\text{lift}}
\end{equation}
where:
\begin{itemize}
    \item $\Phi_{\text{lift}}: \mathcal{F}(\mathbb{R}^2) \to \mathcal{F}(D_4 \times \mathbb{R}^2)$ is the lifting layer
    \item $\Phi_{G}: \mathcal{F}(D_4 \times \mathbb{R}^2) \to \mathcal{F}(D_4 \times \mathbb{R}^2)$ are equivariant group layers
    \item $\Phi_{\text{inv}}: \mathcal{F}(D_4 \times \mathbb{R}^2) \to \mathbb{R}^C$ is a final invariant layer
\end{itemize}
\end{definition}

\begin{theorem}[Composition toward Invariance]\label{thm:composicion_equivarianza_d4}
If each component $\Phi_i$ of the above decomposition is $D_4$-equivariant and the final layer $\Phi_{\text{inv}}$ is $D_4$-invariant, then the full composition $\Phi$ is $D_4$-invariant:
\begin{equation}
\Phi(L_g f) = \Phi(f) \quad \forall g \in D_4, f \in \mathcal{F}(\mathbb{R}^2)
\end{equation}
\end{theorem}

\begin{proof}
By the composition property of equivariant functions established in \autoref{theorem:equivariant_composition}, the composition of equivariant transformations preserves equivariance. The final invariant layer eliminates the dependence on the group, producing invariant outputs.
\end{proof}

\textbf{Feature Dimensionality.} In each equivariant layer, features are organized as functions $f: D_4 \times \mathbb{R}^2 \to \mathbb{R}^d$, where the dimension of the group ($|D_4| = 8$) introduces a multiplicative factor in the feature space. This structure preserves all the information about orientations and reflections.

\begin{figure}[htbp]
	\centering
	\includegraphics[width=0.85\textwidth]{figures/d4_group_actions}
	\caption[Elements of the group $D_4$ acting on a patch]{The 8 elements of the dihedral group $D_4$ acting on an asymmetric patch. Top row: the 4 rotations ($e$, $r$, $r^2$, $r^3$). Bottom row: the 4 reflections ($s$, $sr$, $sr^2$, $sr^3$). The $D_4$-equivariant network produces the same output regardless of the applied transformation.}
	\label{fig:d4_group_actions}
\end{figure}

\subsubsection{Mathematical Verification of Equivariance}\label{subsubsec:verificacion_matematica_equivarianza}

The theoretical correctness of equivariant architectures must be complemented with rigorous methods of mathematical verification that validate the fundamental property of equivariance.

\begin{definition}[Equivariance Error Measure]\label{def:error_equivarianza}
For a function $\Phi$ and a group $G$, we define the \textit{equivariance error measure} as:
\begin{equation}
\mathcal{E}_G(\Phi, x) = \frac{1}{|G|} \sum_{g \in G} \|\Phi(g \cdot x) - g \cdot \Phi(x)\|^2
\end{equation}
where $\|\cdot\|$ is an appropriate norm in the output space.
\end{definition}

\begin{theorem}[Equivariance Criterion]\label{thm:criterio_equivarianza}
A function $\Phi$ is $G$-equivariant if and only if $\mathcal{E}_G(\Phi, x) = 0$ for every input $x$ in the domain.
\end{theorem}

\begin{proof}
$(\Rightarrow)$ If $\Phi$ is $G$-equivariant, then $\Phi(g \cdot x) = g \cdot \Phi(x)$ for all $g \in G$, which implies $\|\Phi(g \cdot x) - g \cdot \Phi(x)\|^2 = 0$ for all $g$. Therefore, $\mathcal{E}_G(\Phi, x) = \frac{1}{|G|} \sum_{g \in G} 0 = 0$.

$(\Leftarrow)$ If $\mathcal{E}_G(\Phi, x) = 0$, then $\sum_{g \in G} \|\Phi(g \cdot x) - g \cdot \Phi(x)\|^2 = 0$. Since each summand is non-negative (squared norm), they must all be zero: $\|\Phi(g \cdot x) - g \cdot \Phi(x)\|^2 = 0$ for all $g \in G$. By the positive definiteness of the norm, $\Phi(g \cdot x) = g \cdot \Phi(x)$ for all $g \in G$, that is, $\Phi$ is $G$-equivariant.
\end{proof}

\begin{proposition}[Error Propagation]\label{prop:propagacion_error}
For the composition of functions $\Phi = \Phi_n \circ \cdots \circ \Phi_1$, where each $\Phi_i$ has Lipschitz constant $L_i$ and $x_i = (\Phi_{i-1} \circ \cdots \circ \Phi_1)(x)$ are the intermediate activations (with $x_0 = x$), the maximum equivariance deviation at each layer propagates as:
\begin{equation}
\sup_{g \in G}\|\Phi(g \cdot x) - g \cdot \Phi(x)\| \leq \sum_{i=1}^{n} \left(\prod_{j=i+1}^{n} L_j\right) \sup_{g \in G}\|\Phi_i(g \cdot x_i) - g \cdot \Phi_i(x_i)\|
\end{equation}
where the empty product (when $i = n$) equals $1$.
\end{proposition}

\begin{proof}
For $n = 1$ the inequality is trivial. For the inductive step, let $\Psi = \Phi_{n-1} \circ \cdots \circ \Phi_1$. Fixing $g \in G$, we write:
\begin{align}
\|\Phi(g \cdot x) - g \cdot \Phi(x)\| &= \|\Phi_n(\Psi(g \cdot x)) - g \cdot \Phi_n(\Psi(x))\| \\
&\leq \|\Phi_n(\Psi(g \cdot x)) - \Phi_n(g \cdot \Psi(x))\| + \|\Phi_n(g \cdot \Psi(x)) - g \cdot \Phi_n(\Psi(x))\|
\end{align}
The first term is bounded by the Lipschitz property: $L_n \|\Psi(g \cdot x) - g \cdot \Psi(x)\|$. The second is the equivariance deviation of $\Phi_n$ evaluated at $x_n = \Psi(x)$. Taking the supremum over $g$ and applying the inductive hypothesis to $\Psi$ yields the desired bound.
\end{proof}

\begin{definition}[Practical Equivariance]\label{cor:tolerancia_error}
In numerical implementations with finite precision, a function is considered \textit{practically equivariant} if:
\begin{equation}
\mathcal{E}_G(\Phi, x) < \epsilon
\end{equation}
where $\epsilon$ is an appropriate tolerance (typically $\epsilon \sim 10^{-6}$ for floating-point precision).
\end{definition}

\paragraph{Analysis of Mathematical Precision}

The theoretical analysis of equivariance precision reveals that architectures based on group theory maintain exact equivariance properties, subject only to numerical precision limitations.

\begin{remark}[Numerical Stability of Equivariance]\label{rem:estabilidad_numerica}
For a $D_4$-equivariant architecture implemented with finite-precision arithmetic, the equivariance error is of the order:
\begin{equation}
\sup_{g \in G}\|\Phi(g \cdot x) - g \cdot \Phi(x)\| = O\!\left(\epsilon_{\text{mach}} \cdot \text{depth}(\Phi) \cdot \|x\|\right)
\end{equation}
where $\epsilon_{\text{mach}}$ is the machine precision and $\text{depth}(\Phi)$ is the depth of the network. This heuristic estimate reflects that each layer introduces a rounding error proportional to $\epsilon_{\text{mach}}$, and these errors accumulate linearly across the layers of the network, analogously to the standard analysis of error propagation in floating-point arithmetic.
\end{remark}

\textbf{Mathematical Meaning of the Errors.} The observed errors of the order of $10^{-6}$ in numerical implementations are consistent with the expected precision of 32-bit floating-point arithmetic, confirming that the theoretical equivariance is preserved up to the limiting precision of the hardware.

In particular, when the numerical precision tends to zero ($\epsilon_{\text{mach}} \to 0$), the implementations converge toward exact equivariance, since the equivariance error tends to zero with $\epsilon_{\text{mach}}$.

This theoretical grounding demonstrates how the abstract concepts of group theory materialize in mathematically rigorous architectures, providing the conceptual basis for addressing more complex groups in the following subsections.

\subsection[GCNNs for Finite Groups]{\glspl{GCNN} for Finite Groups}\label{subsec:gcnns_grupos_finitos}

The application of the theory of group convolutions to finite groups represents one of the most direct and computationally tractable extensions of classical \glspl{CNN}. Finite groups appear naturally in many computer vision domains, where discrete symmetries such as rotations by multiples of 90 degrees or reflections are relevant to the interpretation of the data.

\subsubsection{Finite Groups in Computer Vision}

In the context of image processing, the most relevant finite groups are those that preserve the discrete grid structure while incorporating significant geometric symmetries.

\begin{definition}[Cyclic Group of Rotations]
The cyclic group $C_4 = \{e, r, r^2, r^3\}$ represents the rotations by multiples of $90^\circ$ in the plane, where $r$ denotes a $90^\circ$ counterclockwise rotation and $e$ is the identity.
\end{definition}

\begin{definition}[Dihedral Group of the Square]
The dihedral group $D_4 = \{e, r, r^2, r^3, s, sr, sr^2, sr^3\}$ includes the four rotations of $C_4$ plus four reflections $s, sr, sr^2, sr^3$, where $s$ represents a reflection about an axis.
\end{definition}

These groups are particularly important because:
\begin{itemize}
    \item \textbf{Grid preservation}: The transformations maintain the discrete pixel structure
    \item \textbf{Practical relevance}: Many textures and patterns exhibit these symmetries
    \item \textbf{Computational efficiency}: The finite number of elements allows exact implementations
\end{itemize}

The most common symmetry groups in practical applications are:

\begin{table}[h]
\centering
\begin{tabular}{|l|c|l|l|}
\hline
\textbf{Group} & \textbf{Order} & \textbf{Elements} & \textbf{Applications} \\
\hline
$C_4$ (p4) & 4 & Rotations $0^\circ, 90^\circ, 180^\circ, 270^\circ$ & Textures, regular patterns \\
$D_4$ (p4m) & 8 & Rotations + reflections & Symmetric objects, crystals \\
$C_2$ & 2 & Identity, $180^\circ$ rotation & Bilateral symmetry \\
$D_2$ & 4 & Identity + 3 reflections & Shape analysis \\
\hline
\end{tabular}
\caption{Common finite groups in computer vision}
\label{tab:grupos_finitos}
\end{table}

\subsubsection[Construction of GCNNs for Discrete Groups]{Construction of \glspl{GCNN} for Discrete Groups}

The construction of a \gls{GCNN} for a finite group $G$ follows the theoretical framework established in \autoref{subsec:fundamentos_teoricos}, but exploits the discrete structure for significant computational simplifications.

\paragraph{Lifting from the Input Space}

The first layer of a \gls{GCNN} must transform functions defined on the Euclidean space $\mathbb{R}^2$ (or its discretization) into functions defined over the group $G$. This operation is called \textbf{lifting}.

\begin{definition}[Lifting Convolution]
For an input image $f: \mathbb{Z}^2 \to \mathbb{R}$ and a filter $\psi: \mathbb{Z}^2 \to \mathbb{R}$, the lifting convolution produces a function on $G$ defined by:
\begin{equation}
(f \star \psi)(g) = \sum_{x \in \mathbb{Z}^2} f(x) \psi(g^{-1} \cdot x)
\label{eq:lifting_convolution}
\end{equation}
where $g \in G$ and the sum is over a local window.
\end{definition}

For the group $C_4$, this means that for each spatial position and each rotation $g \in \{e, r, r^2, r^3\}$, we compute the response of the filter rotated by $g^{-1}$.

\paragraph{Group Convolutions in Intermediate Layers}

The subsequent layers operate on functions $f: G \to \mathbb{R}^C$ (where $C$ is the number of channels) using full group convolutions:

\begin{equation}
(f \star_G \psi)(u) = \sum_{g \in G} f(g) \psi(u^{-1}g)
\label{eq:group_conv_discrete}
\end{equation}

\begin{algorithm}[h]
\caption{G-CNN Convolution for a Finite Group}
\label{alg:gcnn_finite}
\begin{algorithmic}[1]
\Require Input function $f$, filters $\{\psi_i\}$, group $G = \{g_1, \ldots, g_n\}$
\Ensure Feature map $\text{output}$
\State Initialize $\text{output}$ with dimensions $(|G|, H, W, C_{\text{out}})$
\For{each filter $i$ in $\{1, \ldots, C_{\text{out}}\}$}
    \For{each group element $u \in G$}
        \For{each position $(h, w)$}
            \State $\text{sum} \leftarrow 0$
            \For{each group element $g \in G$}
                \For{each channel $c$ in input}
                    \State $\text{sum} \leftarrow \text{sum} + f[g][h][w][c] \cdot \psi_i[u^{-1}g][c]$
                \EndFor
            \EndFor
            \State $\text{output}[u][h][w][i] \leftarrow \text{sum}$
        \EndFor
    \EndFor
\EndFor
\Return $\text{output}$
\end{algorithmic}
\end{algorithm}

\subsubsection{Mathematical Theory for the Rotation Group $C_4$}

The cyclic group $C_4 = \{e, r, r^2, r^3\}$ provides a fundamental case for understanding the mathematical structure of group convolutions.

\paragraph{Structure of the Feature Space}

The feature maps in a $C_4$-CNN form a vector space $\mathcal{V} = \mathbb{R}^{C_4 \times \Omega \times \mathbb{R}^C}$, where:
\begin{itemize}
    \item $C_4$: Group space with 4 elements
    \item $\Omega \subset \mathbb{Z}^2$: Discrete spatial domain
    \item $\mathbb{R}^C$: Feature channel space
\end{itemize}

\begin{definition}[Action of $C_4$ on Filters]
The action of the group $C_4$ on filters $\psi: \mathbb{Z}^2 \to \mathbb{R}$ is defined via the matrix representation of rotations:
\begin{equation}
(r^k \cdot \psi)(i,j) = \psi(R_k^{-1} \cdot (i,j))
\end{equation}
where $R_k$ is the rotation matrix by $k \cdot 90^\circ$:
\begin{equation}
R_k = \begin{pmatrix}
\cos(k\pi/2) & -\sin(k\pi/2) \\
\sin(k\pi/2) & \cos(k\pi/2)
\end{pmatrix}
\end{equation}
\end{definition}

\begin{theorem}[Parametric Efficiency of $C_4$-CNNs]
A $C_4$-CNN requires exactly $1/4$ of the parameters that a standard CNN requires to represent the same space of $C_4$-equivariant functions, due to the parameter sharing induced by the group structure.
\end{theorem}

\begin{proof}
In a standard CNN, each filter $\psi: \mathbb{Z}^2 \to \mathbb{R}$ has independent parameters for each kernel position. In a $C_4$-CNN, equivariance requires that the 4 rotated versions of the filter $\{\psi, r \cdot \psi, r^2 \cdot \psi, r^3 \cdot \psi\}$ be generated from the same set of base parameters $\psi$. To represent a $C_4$-equivariant function, a standard CNN would need 4 independent filters (one for each orientation), with a total of $4P$ parameters (where $P$ is the kernel size). The $C_4$-CNN uses a single base filter with $P$ parameters and generates the other 3 via the group action, reaching the same function space with $P = (4P)/4$ parameters.
\end{proof}

\subsubsection{Theory of the Dihedral Group $D_4$}

The extension to the dihedral group $D_4$ incorporates reflections in addition to rotations, significantly enriching the algebraic structure.

\paragraph{Mathematical Structure of Reflections}

The reflections in $D_4$ are characterized by orthogonal linear transformations with determinant $-1$:

\begin{definition}[Matrix Representation of Reflections in $D_4$]
The fundamental reflections are expressed as:
\begin{align}
s_h: \begin{pmatrix} x \\ y \end{pmatrix} &\mapsto \begin{pmatrix} 1 & 0 \\ 0 & -1 \end{pmatrix} \begin{pmatrix} x \\ y \end{pmatrix} = \begin{pmatrix} x \\ -y \end{pmatrix} \\
s_v: \begin{pmatrix} x \\ y \end{pmatrix} &\mapsto \begin{pmatrix} -1 & 0 \\ 0 & 1 \end{pmatrix} \begin{pmatrix} x \\ y \end{pmatrix} = \begin{pmatrix} -x \\ y \end{pmatrix} \\
s_d: \begin{pmatrix} x \\ y \end{pmatrix} &\mapsto \begin{pmatrix} 0 & 1 \\ 1 & 0 \end{pmatrix} \begin{pmatrix} x \\ y \end{pmatrix} = \begin{pmatrix} y \\ x \end{pmatrix}
\end{align}
\end{definition}

\paragraph{Complete Algebraic Structure}

The multiplication table of $D_4$ encodes the complete algebraic structure of the group:

\begin{table}[h]
\centering
\small
\begin{tabular}{|c|c|c|c|c|c|c|c|c|}
\hline
$\cdot$ & $e$ & $r$ & $r^2$ & $r^3$ & $s$ & $sr$ & $sr^2$ & $sr^3$ \\
\hline
$e$ & $e$ & $r$ & $r^2$ & $r^3$ & $s$ & $sr$ & $sr^2$ & $sr^3$ \\
$r$ & $r$ & $r^2$ & $r^3$ & $e$ & $sr$ & $sr^2$ & $sr^3$ & $s$ \\
$r^2$ & $r^2$ & $r^3$ & $e$ & $r$ & $sr^2$ & $sr^3$ & $s$ & $sr$ \\
$r^3$ & $r^3$ & $e$ & $r$ & $r^2$ & $sr^3$ & $s$ & $sr$ & $sr^2$ \\
$s$ & $s$ & $sr^3$ & $sr^2$ & $sr$ & $e$ & $r^3$ & $r^2$ & $r$ \\
$sr$ & $sr$ & $s$ & $sr^3$ & $sr^2$ & $r$ & $e$ & $r^3$ & $r^2$ \\
$sr^2$ & $sr^2$ & $sr$ & $s$ & $sr^3$ & $r^2$ & $r$ & $e$ & $r^3$ \\
$sr^3$ & $sr^3$ & $sr^2$ & $sr$ & $s$ & $r^3$ & $r^2$ & $r$ & $e$ \\
\hline
\end{tabular}
\caption{Multiplication table of the dihedral group $D_4$}
\label{tab:d4_multiplication}
\end{table}

\subsubsection{Fundamental Theoretical Considerations}

\paragraph{Dimensionality Analysis}

The dimensional scaling by the factor $|G|$ has profound theoretical implications for expressive capacity:

\begin{remark}[Dimensionality-Expressivity Trade-off]\label{thm:tradeoff_dimensional}
For a G-CNN that processes functions over $G \times \Omega$, the dimension of the feature space scales as $|G| \cdot \text{dim}(\Omega) \cdot C$. However, the weight sharing induced by equivariance reduces the number of free parameters. Intuitively, if the data possess additional symmetries (that is, the stabilizer $\text{Stab}(x) = \{g \in G : g \cdot x = x\}$ is non-trivial for typical inputs), the effective dimension of the parameter space is further reduced. By the orbit-stabilizer theorem, $|G| = |\text{Orb}(x)| \cdot |\text{Stab}(x)|$, which heuristically suggests that the effective dimension is of the order of $|\text{Orb}(x)| \cdot \text{dim}(\Omega) \cdot C$ rather than $|G| \cdot \text{dim}(\Omega) \cdot C$.
\end{remark}

\paragraph{Implicit Geometric Regularization}

Equivariance induces a natural form of regularization:

\begin{remark}[Regularization by Equivariance]\label{prop:regularizacion_equivarianza}
In architectures where equivariance is not built in by construction, it can be imposed as a soft constraint via the penalty:
\begin{equation}
\mathcal{R}_{\text{equiv}} = \frac{1}{|G|} \sum_{g \in G} \|\theta - g \cdot \theta\|^2
\end{equation}
where $g \cdot \theta$ denotes the transformation of parameters under the action of $g$. Note that $\mathcal{R}_{\text{equiv}} = 0$ if and only if $\theta = g \cdot \theta$ for all $g \in G$, which corresponds exactly to the weight sharing that defines an equivariant layer. The penalty $\mathcal{R}_{\text{equiv}}$ acts as a continuous relaxation of this constraint: minimizing it encourages the parameters to approach the subspace of shared weights, thereby reducing the effective capacity of the model analogously to an $L_2$ regularization.
\end{remark}

\subsubsection{Theoretical Foundations of the Empirical Advantages}

\paragraph{Experiments in Texture Recognition}

The experiments of Cohen and Welling~\citep{cohen2016groupequivariantconvolutional} on the \textit{rotated CIFAR-10} dataset demonstrate significant improvements:

\begin{table}[h]
\centering
\begin{tabular}{|l|c|c|}
\hline
\textbf{Model} & \textbf{CIFAR-10} & \textbf{Rotated CIFAR-10} \\
\hline
Standard CNN & 92.3\% & 47.2\% \\
\gls{GCNN} ($C_4$) & 92.8\% & 71.8\% \\
\gls{GCNN} ($D_4$) & 93.1\% & 73.4\% \\
\hline
\end{tabular}
\caption{Recognition accuracy with arbitrary rotations}
\label{tab:resultados_rotaciones}
\end{table}

\paragraph{Analysis of Crystallographic Patterns}

In materials science applications, \glspl{GCNN} have shown advantages in the analysis of crystal structures that exhibit symmetries specific to the corresponding space group~\citep{weiler2018learningsteerable}.

\paragraph{Medical Segmentation}

In medical imaging, where the orientation of the patient can vary, \glspl{GCNN} provide automatic robustness without the need for extensive data augmentation.

\subsubsection{Limitations and Considerations}

\paragraph{Computational Scalability}

The computational cost scales linearly with $|G|$:
\begin{equation}
\text{FLOPs}_{\text{G-CNN}} = |G| \times \text{FLOPs}_{\text{CNN}}
\end{equation}

For large groups, this can be prohibitive in applications with real-time constraints.

\paragraph{Approximate Symmetries}

In real data, symmetries are often approximate. \glspl{GCNN} can be sensitive to deviations from exact symmetry, requiring care in the selection of the appropriate group.

\paragraph{Compatibility with Existing Architectures}

Integration with techniques such as attention, skip connections, or batch normalization requires careful adaptations to preserve equivariance.

\subsubsection{Extensions and Future Directions}

\paragraph{More General Finite Groups}

The extension to symmetry groups of regular polyhedra (tetrahedral, octahedral, icosahedral) opens up applications in 3D data analysis and computational chemistry.

\paragraph{Combination with Other Symmetries}

The composition of discrete symmetries with continuous symmetries (such as scaling) represents an active area of research.

\paragraph{Symmetry Learning}

Emerging methods seek to automatically discover the relevant symmetries instead of specifying them a priori, combining \glspl{GCNN} with representation learning techniques.

\glspl{GCNN} for finite groups represent a mature and practical extension of classical convolutions, providing theoretically guaranteed robustness to discrete symmetries common in computer vision applications. Their relatively straightforward implementation and their demonstrated benefits make them a valuable tool for domains where discrete symmetries are relevant. Understanding these principles establishes the basis for extensions to more complex groups, as we will see in the following subsections on spherical convolutions and convolutions on manifolds.


Steerable \glspl{CNN} represent a natural extension of discrete \glspl{GCNN} toward the continuous regime, preserving the theoretical guarantees of equivariance while providing greater representational flexibility. Their formulation based on representation theory provides an elegant mathematical framework that connects classical harmonic analysis with modern deep learning architectures.

\subsection{Compact Continuous Groups (SO(2), SO(3))}\label{subsec:grupos_continuos_compactos}

The compact continuous groups $SO(2)$ and $SO(3)$ represent fundamental extensions that allow exact equivariance to continuous rotations in 2D and 3D respectively. This subsection examines the specialized architectures that exploit these rotational symmetries, from Steerable CNNs for SO(2) to spherical CNNs and sophisticated extensions that incorporate scales and irregular structures.

\subsubsection{The Group SO(2): Continuous 2D Rotations and Steerable CNNs}\label{subsubsec:steerable_cnns_so2}

The transition from finite groups of discrete rotations (such as $C_4$ or $D_4$ studied in \autoref{subsec:gcnns_grupos_finitos}) to the continuous rotation group $SO(2)$ represents one of the most significant advances in the development of geometrically aware architectures. While \glspl{GCNN} for finite groups provide equivariance to a discrete set of transformations, Steerable CNNs allow equivariance to continuous rotations, offering a richer and more flexible representation of rotational symmetries.

\paragraph{Motivation: Limitations of Discretizations}

\glspl{GCNN} for finite groups, although powerful, present inherent limitations when the natural symmetries of the data do not align perfectly with the discrete transformations of the group.

\begin{example}[Limitation of $C_4$ on Natural Data]
Consider an image that contains an object rotated by $45^{\circ}$. A \gls{GCNN} based on $C_4$ ($90^{\circ}$ rotations) cannot capture this exact orientation, producing suboptimal responses that are interpolations between the available discrete orientations.
\end{example}

Discretization introduces rotational \textit{aliasing} effects, where intermediate rotations are poorly approximated by the nearest discrete orientations:

\begin{equation}
\text{Error}(\theta) = \min_{k} |\theta - k \cdot \frac{2\pi}{|G|}|
\end{equation}

For $C_4$, this error can reach up to $45^{\circ}$, which represents a significant approximation for many practical applications.

\paragraph{The Continuous Rotation Group SO(2)}

The special orthogonal group $SO(2)$ consists of all rotations in the plane about the origin:

\begin{definition}[Group SO(2)]
$SO(2) = \{R_\theta : \theta \in [0, 2\pi)\}$ where each element $R_\theta$ is the rotation matrix:
\begin{equation}
R_\theta = \begin{pmatrix}
\cos \theta & -\sin \theta \\
\sin \theta & \cos \theta
\end{pmatrix}
\end{equation}
\end{definition}

\begin{proposition}[Properties of SO(2)]
The group $SO(2)$ is:
\begin{enumerate}
    \item \textbf{Compact}: Topologically equivalent to $S^1$
    \item \textbf{Abelian}: $R_\alpha R_\beta = R_{\alpha + \beta} = R_\beta R_\alpha$
    \item \textbf{One-parameter}: Parametrized by a single angle $\theta$
\end{enumerate}
\end{proposition}

\begin{proof}
(1) \textit{Compact}: $SO(2) = \{R_\theta : \theta \in [0, 2\pi)\}$ with $R_\theta = \begin{psmallmatrix} \cos\theta & -\sin\theta \\ \sin\theta & \cos\theta \end{psmallmatrix}$. The map $\theta \mapsto R_\theta$ is a homeomorphism between $[0, 2\pi) / (0 \sim 2\pi)$ and $SO(2)$, identifying it with $S^1$. Since $S^1$ is closed and bounded in $\mathbb{R}^2$, it is compact by Heine-Borel.

(2) \textit{Abelian}: $R_\alpha R_\beta$ corresponds to the rotation by angle $\alpha + \beta$ (composition of rotations in the plane). Since addition in $\mathbb{R}/2\pi\mathbb{Z}$ is commutative, $R_\alpha R_\beta = R_{\alpha+\beta} = R_{\beta+\alpha} = R_\beta R_\alpha$.

(3) \textit{One-parameter}: $\dim(SO(2)) = 1$, since the condition $R^TR = I$ with $\det R = 1$ imposes 3 independent constraints on the 4 parameters of a $2 \times 2$ matrix, leaving $4 - 3 = 1$ degree of freedom.
\end{proof}

This structure allows the use of \textbf{circular harmonic analysis} to construct efficient representations of equivariant functions.

\paragraph{Theory of Steerable Filters}

Steerable filters constitute the theoretical foundation of Steerable CNNs, allowing responses to arbitrary rotations to be generated from a base set of filters.

\begin{figure}[htbp]
	\centering
	\includegraphics[width=0.85\textwidth]{figures/steerable_filters}
	\caption[Steerable filters and steerability]{Steerable filters: any rotation of the base filter $\psi_0$ can be expressed as a linear combination of filters from a finite basis. This property allows equivariance to continuous rotations with a finite number of parameters.}
	\label{fig:steerable_filters}
\end{figure}

\begin{definition}[Steerable Filter]
A filter $h: \mathbb{R}^2 \to \mathbb{R}$ is \textbf{steerable} of order $N$ if there exist base filters $\{h_0, h_1, \ldots, h_{N-1}\}$ and interpolation functions $\{k_n(\theta)\}_{n=0}^{N-1}$ such that:
\begin{equation}
h_\theta(x, y) = \sum_{n=0}^{N-1} k_n(\theta) h_n(x, y)
\end{equation}
where $h_\theta$ is the version of the base filter rotated by angle $\theta$.
\end{definition}

The theory of circular harmonics provides a natural basis for constructing steerable filters:

\begin{theorem}[Representation by Circular Harmonics]
Every function $f: S^1 \to \mathbb{C}$ (where $S^1$ represents orientations) admits a Fourier series expansion:
\begin{equation}
f(\theta) = \sum_{k=-\infty}^{\infty} c_k e^{ik\theta}
\end{equation}
The coefficients $c_k$ transform simply under rotations: if $g_\phi(\theta) = f(\theta + \phi)$, then the coefficients of $g_\phi$ are $c_k e^{ik\phi}$.
\end{theorem}

\begin{proof}
The Fourier series expansion of $f$ is a consequence of the completeness of the orthonormal system $\{e^{ik\theta}\}_{k \in \mathbb{Z}}$ in $L^2(S^1)$, with $\langle e^{ik\theta}, e^{il\theta} \rangle = \frac{1}{2\pi}\int_0^{2\pi} e^{i(k-l)\theta} d\theta = \delta_{kl}$. The coefficients are computed as $c_k = \frac{1}{2\pi}\int_0^{2\pi} f(\theta) e^{-ik\theta} d\theta$. For the transformation under rotations, if $g_\phi(\theta) = f(\theta + \phi)$, its coefficients are:
\[
\hat{g}_k = \frac{1}{2\pi}\int_0^{2\pi} f(\theta + \phi) e^{-ik\theta} d\theta = \frac{1}{2\pi}\int_0^{2\pi} f(\alpha) e^{-ik(\alpha - \phi)} d\alpha = e^{ik\phi} c_k
\]
where the substitution $\alpha = \theta + \phi$ and the periodicity of the integral were used.
\end{proof}

\begin{definition}[Circular Harmonic Filter]
A circular harmonic filter of frequency $k$ in polar coordinates $(r, \theta)$ is defined as:
\begin{equation}
\psi_k(r, \theta) = R_k(r) e^{ik\theta}
\end{equation}
where $R_k(r)$ is the radial part and $e^{ik\theta}$ is the angular part.
\end{definition}

\paragraph{Architecture of Steerable CNNs}

The construction of Steerable CNNs follows the framework developed by Weiler and Welling~\citep{weiler2019general}, generalizing the principles of group equivariance to continuous representations.

Instead of discretizing orientations, Steerable CNNs operate on \textbf{geometric fields} that transform according to irreducible representations of $SO(2)$.

\begin{definition}[Geometric Field]
A geometric field of type $\rho$ is a function $f: \mathbb{R}^2 \to V$ where $V$ is a vector space equipped with a representation $\rho: SO(2) \to GL(V)$ such that:
\begin{equation}
(g \cdot f)(x) = \rho(g) f(g^{-1} x)
\end{equation}
for all $g \in SO(2)$.
\end{definition}

\begin{definition}[Irreducible Representations of SO(2)]\label{def:irreps_so2}
The \textit{irreducible representations} of $SO(2)$ are classified by integers $k \in \mathbb{Z}$:

\begin{equation}
\rho_k(R_\theta) = e^{ik\theta} \quad \text{for } k \neq 0
\end{equation}

and the trivial representation for $k = 0$: $\rho_0(R_\theta) = 1$.

For real representations, frequencies $\pm k$ are combined:
\begin{equation}
\rho_k^{\text{real}}(R_\theta) = \begin{pmatrix}
\cos(k\theta) & -\sin(k\theta) \\
\sin(k\theta) & \cos(k\theta)
\end{pmatrix}
\end{equation}
\end{definition}

\paragraph{Applications in Natural Image Analysis}

Steerable CNNs have demonstrated superiority in tasks where exact orientation is critical:

\begin{itemize}
    \item \textbf{Fiber analysis}: In medical images of connective tissue
    \item \textbf{Fingerprint recognition}: Where ridges have precise orientations
    \item \textbf{Materials analysis}: Detection of cracks and structural defects
\end{itemize}

Experiments on datasets augmented with continuous rotations show consistent advantages over discrete approaches, providing a natural transition toward more complex groups such as $SO(3)$ for three-dimensional data.

\subsubsection{Practical Implementation of Steerable CNNs for SO(2)}\label{subsubsec:implementacion_steerable_cnns}

The transition from the theory of steerable filters to efficient computational implementations requires addressing several technical challenges specific to the continuous group $SO(2)$. Unlike finite groups, where transformations are discrete and exact, Steerable CNNs must handle continuous rotations using finite-dimensional representations.

\paragraph{Irreducible Representations of SO(2)}

The practical implementation begins with the explicit construction of the irreducible representations:

The irreducible representations of $SO(2)$ are implemented mathematically as already defined in \autoref{def:irreps_so2}.

\paragraph{Generation of Circular Harmonic Filters}

The core of Steerable CNNs lies in the generation of base filters using circular harmonics:

\begin{listing}[h]
\begin{lstlisting}[language=Python]
class CircularHarmonics:
    """Circular harmonics for SO(2) steerable filters."""

    @staticmethod
    def generate_filter_basis(size: int, max_frequency: int,
                             sigma: float = 0.5) -> torch.Tensor:
        """Generate basis of circular harmonic filters."""
        # Create coordinate grid
        coords = torch.linspace(-1, 1, size)
        y, x = torch.meshgrid(coords, coords, indexing='ij')

        # Convert to polar coordinates
        r = torch.sqrt(x**2 + y**2)
        theta = torch.atan2(y, x)

        # Gaussian radial envelope
        radial_envelope = torch.exp(-0.5 * (r / sigma) ** 2)

        # Generate basis functions
        basis_filters = []

        # k=0 (trivial representation)
        basis_0 = radial_envelope
        basis_filters.append(basis_0)

        # k>0 (non-trivial representations)
        for k in range(1, max_frequency + 1):
            # Cosine component
            cos_k = radial_envelope * torch.cos(k * theta)
            basis_filters.append(cos_k)

            # Sine component
            sin_k = radial_envelope * torch.sin(k * theta)
            basis_filters.append(sin_k)

        return torch.stack(basis_filters)

    @staticmethod
    def rotate_harmonics(harmonics: torch.Tensor, angle: float,
                        max_frequency: int) -> torch.Tensor:
        """Rotate circular harmonics by a given angle."""
        result = harmonics.clone()

        # k=0 unchanged
        # k>0 components
        idx = 1
        for k in range(1, max_frequency + 1):
            # Extract cosine and sine components
            cos_comp = harmonics[idx]
            sin_comp = harmonics[idx + 1]

            # Apply rotation matrix
            c = math.cos(k * angle)
            s = math.sin(k * angle)

            result[idx] = c * cos_comp - s * sin_comp
            result[idx + 1] = s * cos_comp + c * sin_comp

            idx += 2

        return result
\end{lstlisting}
\caption{Generation and rotation of circular harmonics}
\label{code:circular_harmonics}
\end{listing}

\paragraph{Steerable Convolutional Layer}

The implementation of the steerable convolution combines the base filters with learnable coefficients:

\begin{listing}[h]
\begin{lstlisting}[language=Python]
class SteerableConv2d(nn.Module):
    """2D steerable convolution layer for SO(2) equivariance."""

    def __init__(self, in_irreps: SO2_Irreps, out_irreps: SO2_Irreps,
                 kernel_size: int, max_frequency: int = 4):
        super().__init__()

        self.in_irreps = in_irreps
        self.out_irreps = out_irreps
        self.kernel_size = kernel_size
        self.max_frequency = max_frequency

        # Generate harmonic basis
        self.register_buffer(
            'basis_filters',
            CircularHarmonics.generate_filter_basis(
                kernel_size, max_frequency
            )
        )

        num_basis = self.basis_filters.shape[0]

        # Learnable coefficients to combine basis functions
        self.coefficients = nn.Parameter(
            torch.randn(out_irreps.dim(), in_irreps.dim(), num_basis)
        )

        self.reset_parameters()

    def reset_parameters(self):
        """Initialize parameters."""
        std = math.sqrt(2.0 / (self.in_irreps.dim() + self.out_irreps.dim()))
        nn.init.normal_(self.coefficients, 0, std)

    def forward(self, x: torch.Tensor) -> torch.Tensor:
        """Forward pass of steerable convolution."""
        # Generate steerable kernels
        kernels = self._generate_kernels()  # (out_dim, in_dim, K, K)

        # Apply convolution
        output = F.conv2d(x, kernels, padding=self.kernel_size//2)

        return output

    def _generate_kernels(self) -> torch.Tensor:
        """Generate steerable kernels from basis functions."""
        # Combine basis functions with learned coefficients
        kernels = torch.einsum(
            'oib,bkl->oikl',
            self.coefficients,
            self.basis_filters
        )

        return kernels
\end{lstlisting}
\caption{Implementation of steerable convolution}
\label{code:steerable_conv}
\end{listing}

\paragraph{Equivariant Activation Functions}

A critical aspect of Steerable CNNs is the design of activation functions that preserve the structure of irreducible representations:

\begin{listing}[h]
\begin{lstlisting}[language=Python]
def steerable_relu(x: torch.Tensor, irreps: SO2_Irreps) -> torch.Tensor:
    """ReLU that preserves irrep structure."""
    result = []
    channel_idx = 0

    for freq, mult in irreps.irreps:
        if freq == 0:
            # Trivial representation - standard ReLU
            channels = x[:, channel_idx:channel_idx+mult]
            result.append(F.relu(channels))
        else:
            # Non-trivial representation - norm-ReLU
            channels = x[:, channel_idx:channel_idx+mult]  # (N, 2, H, W)

            # Compute norm
            norm = torch.norm(channels, dim=1, keepdim=True)  # (N, 1, H, W)

            # Apply gating
            gated_norm = F.relu(norm)

            # Normalize and scale
            eps = 1e-8
            unit_channels = channels / (norm + eps)
            result.append(unit_channels * gated_norm)

        channel_idx += mult

    return torch.cat(result, dim=1)
\end{lstlisting}
\caption{Activation function that preserves equivariance}
\label{code:steerable_relu}
\end{listing}

\paragraph{Complete SO(2)-CNN Architecture}

The combination of all the components makes it possible to construct a CNN that is fully equivariant to continuous rotations:

\begin{listing}[h]
\begin{lstlisting}[language=Python]
class SO2SteerableCNN(nn.Module):
    """Fully steerable CNN for SO(2)."""

    def __init__(self, num_classes: int = 10, max_frequency: int = 3):
        super().__init__()

        self.max_frequency = max_frequency

        # Define irreps progression
        self.irreps_hidden = SO2_Irreps(max_frequency)

        # Lifting layer: R^2 -> SO(2)
        self.lift = nn.Conv2d(3, self.irreps_hidden.dim() * 16,
                              kernel_size=3, padding=1)

        # Steerable layers
        self.steerable1 = SteerableLayer(
            self.irreps_hidden, self.irreps_hidden,
            kernel_size=3, padding=1
        )

        self.steerable2 = SteerableLayer(
            self.irreps_hidden, self.irreps_hidden,
            kernel_size=3, padding=1
        )

        # Pooling and classification
        self.pool = nn.MaxPool2d(2, 2)
        self.adaptive_pool = nn.AdaptiveAvgPool2d(1)

        hidden_dim = self.irreps_hidden.dim() * 16
        self.classifier = nn.Sequential(
            nn.Linear(hidden_dim, 128),
            nn.ReLU(),
            nn.Dropout(0.5),
            nn.Linear(128, num_classes)
        )

    def forward(self, x: torch.Tensor) -> torch.Tensor:
        # Lifting
        x = F.relu(self.lift(x))
        x = self.pool(x)

        # Steerable convolutions
        x = self.steerable1(x)
        x = self.pool(x)

        x = self.steerable2(x)

        # Global pooling
        x = self.adaptive_pool(x)
        x = x.view(x.size(0), -1)

        # Classification
        x = self.classifier(x)

        return x
\end{lstlisting}
\caption{Complete SO(2) Steerable CNN architecture}
\label{code:so2_steerable_cnn}
\end{listing}

\paragraph{Equivariance Verification for Continuous Rotations}

Empirical verification for continuous groups requires sampling rotations from the continuum:

\begin{listing}[h]
\begin{lstlisting}[language=Python]
def test_continuous_equivariance(model: nn.Module, x: torch.Tensor,
                               angles: List[float] = None) -> Dict[str, float]:
    """Test equivariance for continuous rotations."""
    if angles is None:
        angles = [0, np.pi/4, np.pi/2, np.pi, 3*np.pi/2]

    model.eval()
    results = {}

    with torch.no_grad():
        # Original output
        original_output = model(x)

        for angle in angles:
            # Rotate input using bilinear interpolation
            rotated_x = rotate_tensor_continuous(x, angle)

            # Get output of rotated input
            rotated_output = model(rotated_x)

            # For classification, we expect similar outputs (invariance)
            diff = torch.mean((original_output - rotated_output) ** 2).item()
            results[f'rotation_{angle:.2f}'] = diff

    return results

def rotate_tensor_continuous(x: torch.Tensor, angle: float) -> torch.Tensor:
    """Rotate tensor by a given angle using interpolation."""
    # Implementation using grid_sample for continuous rotation
    N, C, H, W = x.shape

    # Create rotation matrix
    cos_a, sin_a = math.cos(angle), math.sin(angle)
    rotation_matrix = torch.tensor([
        [cos_a, -sin_a],
        [sin_a, cos_a]
    ], dtype=x.dtype, device=x.device)

    # Create coordinate grid and apply rotation
    coords = torch.linspace(-1, 1, H, device=x.device)
    y, x_coords = torch.meshgrid(coords, coords, indexing='ij')
    grid = torch.stack([x_coords, y], dim=-1)

    # Apply rotation
    grid_flat = grid.view(-1, 2)
    rotated_grid = grid_flat @ rotation_matrix.T
    rotated_grid = rotated_grid.view(H, W, 2)
    rotated_grid = rotated_grid.unsqueeze(0).repeat(N, 1, 1, 1)

    # Sample
    return F.grid_sample(x, rotated_grid, mode='bilinear',
                        padding_mode='border', align_corners=True)
\end{lstlisting}
\caption{Equivariance verification for continuous rotations}
\label{code:continuous_equivariance_test}
\end{listing}

\paragraph{Empirical Analysis and Comparisons}

The empirical experiments demonstrate the advantages of Steerable CNNs over finite discretizations:

\begin{table}[h]
\centering
\begin{tabular}{|l|c|c|c|}
\hline
\textbf{Method} & \textbf{Rotated MNIST} & \textbf{Texture Dataset} & \textbf{Parameters} \\
\hline
Standard CNN & 67.3\% & 52.1\% & 125K \\
G-CNN (C4) & 91.2\% & 76.8\% & 95K \\
G-CNN (D4) & 93.5\% & 79.2\% & 87K \\
Steerable CNN & \textbf{96.1\%} & \textbf{84.7\%} & \textbf{73K} \\
\hline
\end{tabular}
\caption{Comparison of accuracy and parametric efficiency}
\label{tab:steerable_comparison}
\end{table}

The results show that Steerable CNNs not only outperform finite discretizations in accuracy, but also require fewer parameters due to their more efficient representation of continuous symmetries.

\paragraph{Visualization of Steerable Filters}

A notable feature of Steerable CNNs is the ability to visualize how the filters adapt continuously and smoothly under rotations:

\begin{listing}[h]
\begin{lstlisting}[language=Python]
def visualize_steerable_filters(model: SO2SteerableCNN,
                              layer_name: str = 'steerable1',
                              num_rotations: int = 8):
    """Visualize steerable filters under rotations."""
    layer = getattr(model, layer_name)
    angles = np.linspace(0, 2*np.pi, num_rotations, endpoint=False)

    fig, axes = plt.subplots(2, num_rotations//2, figsize=(16, 8))
    axes = axes.flatten()

    # Get base filters
    basis_filters = layer.conv.basis_filters

    for i, angle in enumerate(angles):
        # Rotate base harmonics
        rotated_basis = CircularHarmonics.rotate_harmonics(
            basis_filters, angle, layer.conv.max_frequency
        )

        # Combine with learned coefficients
        coeffs = layer.conv.coefficients[0, 0, :]  # First channel
        kernel = torch.sum(coeffs.unsqueeze(-1).unsqueeze(-1) * rotated_basis, dim=0)

        # Visualize
        axes[i].imshow(kernel.detach().cpu().numpy(), cmap='RdBu')
        axes[i].set_title(f'\theta = {angle:.2f}')
        axes[i].axis('off')

    plt.suptitle('Steerable Filter under Continuous Rotations')
    plt.tight_layout()
    plt.show()
\end{lstlisting}
\caption{Visualization of steerable filters}
\label{code:visualize_steerable}
\end{listing}

This practical implementation demonstrates how the representation theory of $SO(2)$ translates into efficient computational architectures that handle exact equivariance to continuous rotations, providing the basis for extensions to more complex groups such as $SO(3)$.
\subsection{Spherical CNNs and SO(3) Equivariance}\label{subsec:spherical_cnns_so3}

The natural extension of Steerable CNNs from the planar SO(2) domain toward spherical and three-dimensional volumetric data requires the development of equivariant architectures under the rotation group SO(3). This transition represents one of the most significant advances in Geometric Deep Learning, enabling the direct processing of data that exhibit complete rotational symmetries in three-dimensional space.

\subsubsection{Motivation: From SO(2) to SO(3)}

While Steerable CNNs solve the problem of rotational equivariance in the plane, many real-world applications involve data with intrinsic spherical structure or complete three-dimensional rotational symmetries.

\begin{example}[Omnidirectional Data]
$360^\circ$ images captured by omnidirectional cameras are naturally defined on the sphere $S^2$. A conventional CNN applied to the equirectangular projection of such an image suffers from severe geometric distortions, especially near the poles, where physical rotations are mapped non-uniformly.
\end{example}

\begin{example}[Global Climate Data]
Climate models generate data over the Earth's surface, which is topologically spherical. Weather predictions require equivariance to rotations of the planet, a property that standard CNNs on planar projections cannot guarantee.
\end{example}

\paragraph{Limitations of Planar Projections}

Traditional projections of spherical data to rectangular domains introduce fundamental distortions:

\begin{enumerate}
    \item \textbf{Area distortion}: The equirectangular projection assigns disproportionate area to polar regions
    \item \textbf{Artificial discontinuities}: The left and right edges of the projection are artificially discontinuous
    \item \textbf{Loss of equivariance}: Spherical rotations become complex transformations in the projected domain
\end{enumerate}

\begin{figure}[htbp]
	\centering
	\includegraphics[width=0.80\textwidth]{figures/spherical_harmonics}
	\caption[Spherical harmonics $Y_l^m$]{Spherical harmonics $Y_l^m(\theta, \varphi)$ for degrees $l = 0, 1, 2$. The degree $l$ controls the frequency (number of oscillations) and the order $m$ determines the orientation. Blue and red regions represent positive and negative values respectively.}
	\label{fig:spherical_harmonics}
\end{figure}

\subsubsection{The Rotation Group SO(3)}

The special orthogonal group SO(3) consists of all rotations in three-dimensional space that preserve the origin:

\begin{definition}[Group SO(3)]
$SO(3) = \{R \in \mathbb{R}^{3 \times 3} : R^T R = I, \det(R) = 1\}$

Each element can be parameterized by Euler angles $(\alpha, \beta, \gamma)$ or by the rotation axis $\mathbf{n} \in S^2$ and angle $\theta \in [0, 2\pi)$.
\end{definition}

\begin{proposition}[Properties of SO(3)]
The group SO(3) is:
\begin{enumerate}
    \item \textbf{Non-abelian}: In general, $R_1 R_2 \neq R_2 R_1$
    \item \textbf{Compact}: Topologically equivalent to $\mathbb{R}P^3$ (real projective space)
    \item \textbf{Three-parameter}: Requires three independent parameters to completely specify a rotation
\end{enumerate}
\end{proposition}

\begin{proof}
(1) \textit{Non-abelian}: It suffices to exhibit a counterexample. Let $R_x$ be the $90°$ rotation about the $x$ axis and $R_z$ the $90°$ rotation about the $z$ axis. Then $R_x R_z (1,0,0)^T = R_x(0,1,0)^T = (0,0,1)^T$ but $R_z R_x (1,0,0)^T = R_z(1,0,0)^T = (0,1,0)^T$, hence $R_x R_z \neq R_z R_x$.

(2) \textit{Compact}: $SO(3) = \{R \in \mathbb{R}^{3\times 3} : R^TR = I, \det R = 1\}$ is closed in $\mathbb{R}^{9}$ (preimage of regular values under continuous functions) and bounded ($\|R\|_F^2 = \operatorname{tr}(R^TR) = 3$ for all $R \in SO(3)$), hence compact by Heine-Borel. Topological equivalence with $\mathbb{R}P^3$ is established via the axis-angle parameterization: each rotation is identified with a vector $\theta \mathbf{n}$ with $\|\mathbf{n}\| = 1$ and $\theta \in [0, \pi]$, where $(\pi, \mathbf{n})$ and $(\pi, -\mathbf{n})$ represent the same rotation.

(3) \textit{Three-parameter}: $SO(3) \subset \mathbb{R}^{3 \times 3}$ has 9 entries. The conditions $R^TR = I$ impose 6 independent equations (the matrix $R^TR - I$ is symmetric, with 6 independent entries), and $\det R = 1$ is redundant (it follows from $R^TR = I$ and connectedness). Thus, $\dim(SO(3)) = 9 - 6 = 3$.
\end{proof}

This non-commutative structure introduces significant complexities compared to SO(2), requiring more sophisticated tools from harmonic analysis.

\subsubsection{Spherical Harmonic Analysis}

Harmonic analysis on the sphere $S^2$ provides the theoretical basis for constructing SO(3)-equivariant convolutions.

\paragraph{Spherical Harmonics}

\begin{definition}[Spherical Harmonics]
The spherical harmonics $Y_l^m(\theta, \phi)$ form a complete orthonormal basis for $L^2(S^2)$:
\begin{equation}
Y_l^m(\theta, \phi) = \sqrt{\frac{2l+1}{4\pi}\frac{(l-|m|)!}{(l+|m|)!}} P_l^{|m|}(\cos \theta) e^{im\phi}
\end{equation}
where $P_l^m$ are the associated Legendre polynomials, $l \geq 0$ is the degree and $|m| \leq l$ is the order.
\end{definition}

\begin{theorem}[Spherical Harmonic Expansion]
Every function $f \in L^2(S^2)$ admits the expansion:
\begin{equation}
f(\theta, \phi) = \sum_{l=0}^{\infty} \sum_{m=-l}^{l} f_l^m Y_l^m(\theta, \phi)
\end{equation}
with coefficients $f_l^m = \int_{S^2} f(\mathbf{x}) \overline{Y_l^m(\mathbf{x})} d\mathbf{x}$.

See Section~\ref{sec:expansion_armonicos_esfericos} of \autoref{ch:teoremas_complementarios} for the proof of completeness.
\end{theorem}

\paragraph{Transformation under Rotations}

The fundamental property that makes spherical harmonics useful is their behavior under rotations:

\begin{theorem}[Transformation of Spherical Harmonics]
Under a rotation $R \in SO(3)$, spherical harmonics transform as:
\begin{equation}
Y_l^m(R^{-1}\mathbf{x}) = \sum_{m'=-l}^{l} D_{m'm}^{(l)}(R) Y_l^{m'}(\mathbf{x})
\end{equation}
where $D^{(l)}(R)$ is the irreducible representation matrix of dimension $(2l+1) \times (2l+1)$ of SO(3).
\end{theorem}

\begin{proof}
The spherical harmonics $\{Y_l^m\}_{m=-l}^{l}$ form a basis of the subspace $\mathcal{H}_l$ of homogeneous harmonic polynomials of degree $l$ restricted to $S^2$. This subspace is invariant under the action of $SO(3)$: if $R \in SO(3)$ and $Y_l^m \in \mathcal{H}_l$, then $Y_l^m \circ R^{-1} \in \mathcal{H}_l$ (since composition with a rotation preserves both homogeneity and harmonicity, because $\Delta$ commutes with rotations). Since $\mathcal{H}_l$ is of dimension $2l+1$ and invariant, the function $Y_l^m(R^{-1}\mathbf{x})$ is expressed as a linear combination of the basis: $Y_l^m(R^{-1}\mathbf{x}) = \sum_{m'} D_{m'm}^{(l)}(R) Y_l^{m'}(\mathbf{x})$, where the coefficients $D_{m'm}^{(l)}(R) = \int_{S^2} Y_l^m(R^{-1}\mathbf{x}) \overline{Y_l^{m'}(\mathbf{x})} \, d\mathbf{x}$ define the representation matrix. The irreducibility of $D^{(l)}$ follows from the fact that $\mathcal{H}_l$ admits no proper invariant subspaces under $SO(3)$.
\end{proof}

This property is the key to constructing equivariant filters: the coefficients of different degrees $l$ mix only among harmonics of the same degree.

\subsubsection{Architecture of Spherical CNNs}

Spherical CNNs, introduced by Cohen et al.~\citep{cohen2018sphericalcnns}, extend the group equivariance framework to the sphere through the direct use of spherical harmonic analysis.

\begin{example}[Construction of an Equivariant Layer on $S^2$]\label{ex:spherical_equivariant_layer_cnn}
Consider the sphere $S^2 = SO(3)/SO(2)$ for three-dimensional directional data. A function $f: S^2 \rightarrow \mathbb{R}^{d'}$ is $SO(3)$-equivariant if $f(R \cdot \mathbf{x}) = \rho(R) f(\mathbf{x})$ for every rotation $R \in SO(3)$.

To construct a spherical convolutional layer:
\begin{enumerate}
    \item \textbf{Representation}: We choose the representation $\rho_\ell: SO(3) \rightarrow GL(\mathbb{R}^{d_\ell})$ for layer $\ell$
    \item \textbf{Spherical kernel}: The kernel $\psi: S^2 \rightarrow \mathbb{R}^{d_{\ell+1} \times d_\ell}$ must satisfy:
    \begin{equation}
    \psi(R \cdot \omega) = \rho_{\ell+1}(R) \psi(\omega) \rho_\ell(R)^{-1}
    \end{equation}
    \item \textbf{Spherical convolution}: The output is computed as:
    \begin{equation}
    (f * \psi)(\mathbf{x}) = \int_{S^2} f(\omega) \psi(R_{\mathbf{x} \leftarrow \omega}) \, d\omega
    \end{equation}
\end{enumerate}

This construction underpins Spherical \glspl{CNN}, where the spherical harmonics $Y_\ell^m$ provide the irreducible representations of $SO(3)$.
\end{example}

\paragraph{Representation of Spherical Signals}

Instead of operating on pixels in a rectangular grid, Spherical CNNs directly process functions defined on $S^2$:

\begin{definition}[Spherical Signal]
A spherical signal with $C$ channels is a function $f: S^2 \to \mathbb{R}^C$ that can be represented as:
\begin{equation}
f(\mathbf{x}) = \sum_{l=0}^{L} \sum_{m=-l}^{l} \mathbf{f}_l^m Y_l^m(\mathbf{x})
\end{equation}
where $\mathbf{f}_l^m \in \mathbb{R}^C$ are per-channel coefficient vectors.
\end{definition}

\paragraph{Equivariant Spherical Convolution}

Convolution on the sphere is defined using the group structure of SO(3):

\begin{definition}[SO(3) Convolution]
For functions $f: SO(3) \to \mathbb{R}^{C_{\text{in}}}$ and kernel $\psi: SO(3) \to \mathbb{R}^{C_{\text{out}} \times C_{\text{in}}}$:
\begin{equation}
(f \star \psi)(g) = \int_{SO(3)} f(h) \psi(g^{-1}h) dh
\end{equation}
where $dh$ is the Haar measure on SO(3).
\end{definition}

\begin{theorem}[Equivariance of Spherical Convolution]
The SO(3) convolution is equivariant: if $g \cdot f(x) = f(g^{-1}x)$, then:
\begin{equation}
g \cdot (f \star \psi) = (g \cdot f) \star \psi
\end{equation}
\end{theorem}

\begin{proof}
Let $g \in SO(3)$. We compute $(g \cdot f) \star \psi$ using the definition of convolution:
\begin{align}
[(g \cdot f) \star \psi](x) &= \int_{SO(3)} (g \cdot f)(h) \, \psi(x^{-1}h) \, dh = \int_{SO(3)} f(g^{-1}h) \, \psi(x^{-1}h) \, dh
\end{align}
With the change of variable $h' = g^{-1}h$ (i.e., $h = gh'$) and using the left-invariance of the Haar measure ($dh = dh'$):
\begin{align}
&= \int_{SO(3)} f(h') \, \psi(x^{-1}gh') \, dh' = \int_{SO(3)} f(h') \, \psi((g^{-1}x)^{-1}h') \, dh' = (f \star \psi)(g^{-1}x)
\end{align}
This is exactly $[g \cdot (f \star \psi)](x) = (f \star \psi)(g^{-1}x)$, which proves equivariance.
\end{proof}

\subsubsection{Implementation via Spherical FFT}

The efficient implementation of Spherical CNNs requires the use of spherical fast Fourier transforms (S-FFT) to accelerate convolutions in the frequency domain.

\paragraph{Spherical Fourier Transform}

\begin{definition}[Forward S-FFT]
The forward transform computes the spherical harmonic coefficients:
\begin{equation}
\hat{f}_l^m = \int_{S^2} f(\theta, \phi) \overline{Y_l^m(\theta, \phi)} \sin \theta \, d\theta \, d\phi
\end{equation}
\end{definition}

\begin{definition}[Inverse S-FFT]
The inverse transform reconstructs the spatial function:
\begin{equation}
f(\theta, \phi) = \sum_{l=0}^{L} \sum_{m=-l}^{l} \hat{f}_l^m Y_l^m(\theta, \phi)
\end{equation}
\end{definition}

\paragraph{Convolution in the Spectral Domain}

The convolution theorem extends naturally to the spherical case:

\begin{theorem}[Spherical Convolution Theorem]
In the spectral domain, convolution reduces to pointwise multiplication within each degree-$l$ subspace:
\begin{equation}
\widehat{(f \star \psi)}_l^m = \sum_{m'=-l}^{l} \hat{f}_l^{m'} \hat{\psi}_l^{m'm}
\end{equation}
where $\hat{\psi}_l^{m'm}$ are the spectral components of the kernel.
\end{theorem}

\begin{proof}
We expand $f$ and $\psi$ in spherical harmonics. The convolution $(f \star \psi)(g) = \int_{SO(3)} f(h) \psi(g^{-1}h) \, dh$ is expressed in the spectral domain using the Fourier transform over $SO(3)$: $\hat{F}^{(l)}_{mm'} = \int_{SO(3)} F(g) D^{(l)}_{mm'}(g) \, dg$. By the convolution theorem for compact groups, the transform of the convolution is the matrix product of the transforms:
\[
\widehat{(f \star \psi)}^{(l)} = \hat{f}^{(l)} \cdot \hat{\psi}^{(l)}
\]
where the product is matrix multiplication in the indices $m, m'$ within each degree-$l$ block. Componentwise, this gives $\widehat{(f \star \psi)}_l^m = \sum_{m'=-l}^{l} \hat{f}_l^{m'} \hat{\psi}_l^{m'm}$.
\end{proof}

\begin{algorithm}[h]
\caption{Efficient Spherical CNN Convolution}
\label{alg:spherical_conv}
\begin{algorithmic}[1]
\Require Spherical signal $f$, spectral kernel $\hat{\psi}$, maximum degree $L$
\Ensure Convolved signal in the spectral domain
\State Compute S-FFT: $\hat{f}_l^m = \text{S-FFT}(f)$
\For{$l = 0$ to $L$}
    \For{$m = -l$ to $l$}
        \State $\hat{g}_l^m = \sum_{m'} \hat{f}_l^{m'} \hat{\psi}_l^{m'm}$
    \EndFor
\EndFor
\State \textbf{return} $\{\hat{g}_l^m\}$ (keep spectral representation for subsequent layers)
\end{algorithmic}
\end{algorithm}

\subsubsection{Extension to 3D Volumetric Data}

Spherical CNNs extend naturally to three-dimensional volumetric data by combining spherical harmonic analysis with radial processing.

\paragraph{3D Spherical Coordinates}

For volumetric data $f: \mathbb{R}^3 \to \mathbb{R}$, the decomposition in spherical coordinates $(r, \theta, \phi)$ is used:

\begin{equation}
f(r, \theta, \phi) = \sum_{l=0}^{L} \sum_{m=-l}^{l} f_l^m(r) Y_l^m(\theta, \phi)
\end{equation}

where $f_l^m(r)$ are radial functions that capture the variation with respect to the distance from the origin.

\paragraph{3D Steerable CNNs}

The complete formulation for 3D data, developed by Weiler et al.~\citep{weiler2018steerable}, combines:

\begin{enumerate}
    \item \textbf{Radial processing}: 1D convolutions along radial rays
    \item \textbf{Angular mixing}: Linear transformations between spherical harmonic coefficients
    \item \textbf{SO(3) equivariance}: Preservation of the irreducible representation structure
\end{enumerate}

\begin{definition}[3D Steerable Kernel]
A 3D steerable kernel is parameterized as:
\begin{equation}
K(r, \theta, \phi) = \sum_{l=0}^{L} \sum_{i=1}^{I} R_i(r) \sum_{m=-l}^{l} W_{l,i}^m Y_l^m(\theta, \phi)
\end{equation}
where $R_i(r)$ are radial basis functions and $W_{l,i}^m$ are learnable parameters.
\end{definition}

\subsubsection{Applications in Spherical Domains}

\paragraph{Analysis of Omnidirectional Images}

Spherical CNNs have demonstrated significant advantages in the processing of $360^\circ$ images:

\begin{itemize}
    \item \textbf{Object recognition}: Robust detection independent of camera orientation
    \item \textbf{Virtual reality}: Efficient processing of immersive content
    \item \textbf{Robotics}: Navigation based on omnidirectional vision
\end{itemize}

\paragraph{Climate and Geophysical Modeling}

In Earth sciences, Spherical CNNs provide natural tools for:

\begin{itemize}
    \item \textbf{Weather prediction}: Models that respect the spherical geometry of the planet
    \item \textbf{Ocean current analysis}: Processing of vector data on the sphere
    \item \textbf{Satellite monitoring}: Analysis of Earth observation data
\end{itemize}

\paragraph{Computational Astrophysics}

The natural spherical structure of astronomical data makes Spherical CNNs particularly appropriate:

\begin{example}[Cosmic Microwave Background]
The analysis of anisotropies in the cosmic microwave background (CMB) requires processing of signals defined over the complete celestial sphere. Spherical CNNs can detect cosmological patterns while preserving all rotational symmetries.
\end{example}

\subsubsection{Performance Comparison}

Comparative experiments demonstrate the advantages of Spherical CNNs over conventional methods:

\begin{table}[h]
\centering
\begin{tabular}{|l|c|c|c|}
\hline
\textbf{Method} & \textbf{Omnidirectional MNIST} & \textbf{Climate Pattern} & \textbf{Spherical CIFAR} \\
\hline
CNN + Equirectangular & 89.2\% & 76.4\% & 72.8\% \\
CNN + Cubemap & 91.7\% & 79.1\% & 76.2\% \\
Spherical CNN & \textbf{96.4\%} & \textbf{87.3\%} & \textbf{83.1\%} \\
\hline
\end{tabular}
\caption{Accuracy on synthetic and real spherical datasets}
\label{tab:spherical_performance}
\end{table}

\paragraph{Computational Efficiency}

The complexity of Spherical CNNs depends on the maximum degree $L$ of the spherical harmonics:

\begin{itemize}
    \item \textbf{Space complexity}: $O(L^2)$ coefficients per spherical location
    \item \textbf{Time complexity}: $O(L^3)$ for S-FFT, $O(L^2)$ for spectral convolutions
    \item \textbf{Parallelization}: Spectral transforms parallelize efficiently
\end{itemize}

\subsubsection{Extensions and Recent Developments}

\paragraph{Spherical CNNs on Unstructured Grids}

Recent developments have extended Spherical CNNs to non-uniform grids~\citep{cohen2021spherical}:

\begin{itemize}
    \item \textbf{Adaptive meshes}: Support for variable resolutions depending on the region
    \item \textbf{Icosahedral grids}: Spherical discretizations with improved symmetry properties
    \item \textbf{Spectral interpolation}: Techniques for handling irregular samplings
\end{itemize}

\paragraph{Hybrid Architectures}

Current research explores combinations with other approaches:

\begin{enumerate}
    \item \textbf{Gauge Equivariant Networks}: Combination with local gauge symmetries~\citep{cohen2019gaugeequivariantconvolutional}
    \item \textbf{Spherical transformers}: Integration of attention mechanisms with SO(3) equivariance
    \item \textbf{Graph Neural Networks}: Treatment of spherical data as graphs on the sphere
\end{enumerate}

\subsubsection{Limitations and Challenges}

\paragraph{Computational Challenges}

\begin{enumerate}
    \item \textbf{Spectral complexity}: The $O(L^3)$ cost of S-FFTs limits the angular resolution
    \item \textbf{Numerical precision}: High-degree spherical harmonics are sensitive to floating-point errors
    \item \textbf{Memory}: Storing spectral coefficients can be prohibitive for large $L$
\end{enumerate}

\paragraph{Theoretical Limitations}

\begin{itemize}
    \item \textbf{Bandlimit assumption}: Real data may not be spectrally band-limited
    \item \textbf{Sampling artifacts}: Discretizations introduce spectral aliasing
    \item \textbf{Boundary effects}: Handling partial regions of the sphere requires specialized techniques
\end{itemize}

\subsubsection{Unifying Perspective}

Spherical CNNs represent a natural and powerful extension of the group equivariance framework from SO(2) toward SO(3), demonstrating how the principles of harmonic analysis can be integrated directly into deep learning architectures. Their success illustrates a fundamental principle: the alignment between the intrinsic symmetries of the data and the architectural inductive biases produces significant improvements in both efficiency and performance.

The progression from standard \glspl{CNN} → discrete \glspl{GCNN} → SO(2) Steerable \glspl{CNN} → SO(3) Spherical \glspl{CNN} illustrates an evolution toward increasingly geometrically aware architectures. This trend continues toward even more general formulations, including networks equivariant to arbitrary Lie groups and convolutions over general Riemannian manifolds.

The theoretical framework developed for Spherical CNNs provides the conceptual tools needed to address data with even more complex geometries, establishing the foundations for the development of Geometric Deep Learning in arbitrary non-Euclidean domains.

\subsection{Non-Compact Groups (E(3), SE(3))}\label{subsec:grupos_no_compactos}

The Euclidean groups $E(n)$ and rigid-motion groups $SE(3)$ represent fundamental extensions toward non-compact groups that incorporate translations together with rotations. These architectures are especially relevant for the processing of 3D point clouds, molecular systems, and particle analysis, where the symmetries include both rotations and translations in three-dimensional space.

The transition from compact groups such as $SO(2)$ and $SO(3)$ toward non-compact groups introduces significant mathematical and computational challenges. While compact groups admit finite Haar measures that allow direct integration over the group, non-compact groups require more sophisticated techniques to construct efficient equivariant representations.

\subsubsection{Mathematical Structure of Non-Compact Groups}

\paragraph{The Euclidean Group E(n)}

The \textit{Euclidean group} $E(n) = \mathbb{R}^n \rtimes O(n)$ is the group of all isometries of $\mathbb{R}^n$ (\autoref{def:grupo_e_n} of \autoref{ch:gdl}).\label{def:grupo_euclidiano}

Each element $g \in E(n)$ is represented as a pair $(R, \mathbf{t})$ with $R \in O(n)$ and $\mathbf{t} \in \mathbb{R}^n$, acting on $\mathbf{x} \in \mathbb{R}^n$ via:
\begin{equation}
g \cdot \mathbf{x} = R\mathbf{x} + \mathbf{t}
\end{equation}

The group operation is defined as:
\begin{equation}
(R_1, \mathbf{t}_1) \circ (R_2, \mathbf{t}_2) = (R_1 R_2, R_1 \mathbf{t}_2 + \mathbf{t}_1)
\end{equation}

\begin{proposition}[Properties of $E(n)$]
The Euclidean group $E(n)$ satisfies:
\begin{enumerate}
    \item \textbf{Non-compact}: The subgroup of translations $\mathbb{R}^n$ is non-compact
    \item \textbf{Lie group}: Differentiable manifold of dimension $n + \frac{n(n-1)}{2}$
    \item \textbf{Semidirect product}: $\mathbb{R}^n$ is a normal subgroup
\end{enumerate}
\end{proposition}

\begin{proof}
(1) \textit{Non-compact}: The subgroup of pure translations $\{(I, \mathbf{t}) : \mathbf{t} \in \mathbb{R}^n\} \cong \mathbb{R}^n$ is unbounded, hence $E(n)$ is not compact.

(2) \textit{Lie group}: $E(n)$ is realized as a closed subgroup of $GL(n+1, \mathbb{R})$ via the matrix representation $(R, \mathbf{t}) \mapsto \begin{psmallmatrix} R & \mathbf{t} \\ 0 & 1 \end{psmallmatrix}$. The dimension is $n$ (translations) $+ \frac{n(n-1)}{2}$ (dimension of $O(n)$, by the count $n^2 - \frac{n(n+1)}{2}$ constraints from orthogonality $R^TR = I$).

(3) \textit{Semidirect product}: For $(R_1, \mathbf{t}_1), (R_2, \mathbf{t}_2) \in E(n)$, conjugation of a pure translation $(I, \mathbf{t})$ gives $(R, \mathbf{s})(I, \mathbf{t})(R, \mathbf{s})^{-1} = (I, R\mathbf{t})$, which remains a pure translation. Thus, the translations form a normal subgroup, and the group law $(R_1, \mathbf{t}_1) \cdot (R_2, \mathbf{t}_2) = (R_1 R_2, R_1 \mathbf{t}_2 + \mathbf{t}_1)$ is exactly that of the semidirect product $\mathbb{R}^n \rtimes O(n)$.
\end{proof}

\paragraph{The Group of Rigid Motions SE(n)}

The \textit{special Euclidean group} $SE(n) = \mathbb{R}^n \rtimes SO(n)$ is the subgroup of $E(n)$ that preserves orientation (\autoref{def:grupo_se_n} of \autoref{ch:gdl}).\label{def:grupo_se_n_cnn}

For $n = 3$, the group $SE(3)$ is fundamental in robotics, computer vision and computational chemistry, since it captures all rigid transformations that preserve distances and orientation in three-dimensional space.

\paragraph{Representations of Non-Compact Groups}

Unlike compact groups, where all unitary irreducible representations are finite-dimensional, non-compact groups admit infinite-dimensional irreducible representations. However, for practical applications in deep learning, we restrict ourselves to finite-dimensional representations that capture the rotational components.

\begin{proposition}[Decomposition of Representations]\label{prop:descomp_repr_en}
Let $E(n) = \mathbb{R}^n \rtimes O(n)$ be the Euclidean group (semidirect product). In every finite-dimensional representation of $E(n)$ where the translations act trivially, the representation factors through the quotient $E(n)/\mathbb{R}^n \cong O(n)$:
\begin{equation}
\rho(R, \mathbf{t}) = \rho_{O(n)}(R)
\end{equation}
where $\rho_{O(n)}$ is a representation of $O(n)$. Analogously, for $SE(n) = \mathbb{R}^n \rtimes SO(n)$, the representation factors through $SO(n)$. See \citep{DBLP:journals/corr/abs-2104-13478} for a discussion in the context of GDL.
\end{proposition}

\begin{proof}
Let $\rho: E(n) \to GL(V)$ be a finite-dimensional representation where the translations act trivially: $\rho(I, \mathbf{t}) = \operatorname{Id}_V$ for all $\mathbf{t} \in \mathbb{R}^n$. For $(R, \mathbf{t}) \in E(n)$, since $(R, \mathbf{t}) = (R, \mathbf{0}) \cdot (I, R^{-1}\mathbf{t})$, we have $\rho(R, \mathbf{t}) = \rho(R, \mathbf{0}) \cdot \rho(I, R^{-1}\mathbf{t}) = \rho(R, \mathbf{0})$. Defining $\rho_{O(n)}(R) := \rho(R, \mathbf{0})$, one verifies that $\rho_{O(n)}$ is a group homomorphism: $\rho_{O(n)}(R_1 R_2) = \rho(R_1 R_2, \mathbf{0}) = \rho(R_1, \mathbf{0}) \rho(R_2, \mathbf{0}) = \rho_{O(n)}(R_1) \rho_{O(n)}(R_2)$. This proves the factorization $\rho = \rho_{O(n)} \circ \pi$, where $\pi: E(n) \to E(n)/\mathbb{R}^n \cong O(n)$ is the canonical projection.
\end{proof}

This factorization motivates the design of architectures that explicitly separate invariant information (scalar features) from equivariant information (vector coordinates).

\subsubsection{E(n) Equivariant Graph Neural Networks}

The E(n) Equivariant Graph Neural Networks of Satorras et al.~\citep{satorras2021enequivariant} introduce a general framework for Euclidean group equivariance via an architecture that explicitly processes both invariant features and equivariant coordinates.

\paragraph{Foundations of the EGNN Design}

The central idea is to construct messages that depend only on quantities invariant under $E(n)$, while coordinate updates are performed via vector differences that preserve equivariance.

\paragraph{EGNN Architecture}

EGNNs maintain equivariance through the explicit separation of scalar information and vector information. In each layer, the following are processed simultaneously:
\begin{itemize}
    \item \textbf{Scalar features} $h_i \in \mathbb{R}^d$ (invariant under $E(n)$)
    \item \textbf{Coordinates} $\mathbf{x}_i \in \mathbb{R}^n$ (equivariant under $E(n)$)
\end{itemize}

The update in an EGNN layer proceeds in three stages:

\textbf{Stage 1: Construction of invariant messages}

Messages between nodes are constructed using only $E(n)$-invariant quantities:
\begin{equation}\label{eq:egnn_message}
\mathbf{m}_{ij} = \phi_e(h_i, h_j, \|\mathbf{x}_i - \mathbf{x}_j\|^2, a_{ij})
\end{equation}

where:
\begin{itemize}
    \item $h_i, h_j \in \mathbb{R}^d$ are scalar features (inherently invariant)
    \item $\|\mathbf{x}_i - \mathbf{x}_j\|^2$ is the squared Euclidean distance (invariant under isometries)
    \item $a_{ij}$ are optional edge attributes
    \item $\phi_e: \mathbb{R}^{2d+2} \to \mathbb{R}^{d'}$ is an MLP
\end{itemize}

\textbf{Invariance of Distances.} The squared distance is invariant under $E(n)$: for any $(R, \mathbf{t}) \in E(n)$,
\begin{equation}
\|(R\mathbf{x}_i + \mathbf{t}) - (R\mathbf{x}_j + \mathbf{t})\|^2 = \|R(\mathbf{x}_i - \mathbf{x}_j)\|^2 = \|\mathbf{x}_i - \mathbf{x}_j\|^2
\end{equation}
since $R$ is orthogonal and the translations cancel.

\textbf{Stage 2: Equivariant coordinate update}

The coordinates are updated using weighted vector differences:
\begin{equation}\label{eq:egnn_coord_update}
\mathbf{x}_i' = \mathbf{x}_i + C \sum_{j \in \mathcal{N}(i)} (\mathbf{x}_i - \mathbf{x}_j) \phi_x(\mathbf{m}_{ij})
\end{equation}

where:
\begin{itemize}
    \item $\phi_x: \mathbb{R}^{d'} \to \mathbb{R}$ is an MLP that produces scalar coefficients
    \item $C$ is a normalization constant (typically $C = 1/|\mathcal{N}(i)|$)
    \item $\mathcal{N}(i)$ is the neighborhood of node $i$
\end{itemize}

\begin{proposition}[Equivariance of the Coordinate Update]\label{prop:egnn_coord_equivariance}
The update~\eqref{eq:egnn_coord_update} is equivariant under $E(n)$. If we apply $(R, \mathbf{t}) \in E(n)$ to all coordinates, then:
\begin{equation}
\mathbf{x}_i'^{(g)} = R\mathbf{x}_i' + \mathbf{t}
\end{equation}
where $\mathbf{x}_i'^{(g)}$ denotes the updated coordinate after transforming the inputs.
\end{proposition}

\begin{proof}
Consider the transformed coordinates $\tilde{\mathbf{x}}_i = R\mathbf{x}_i + \mathbf{t}$. The messages $\mathbf{m}_{ij}$ are invariant under this transformation (by construction using distances). The update of the transformed coordinates is:
\begin{align}
\tilde{\mathbf{x}}_i' &= \tilde{\mathbf{x}}_i + C \sum_{j \in \mathcal{N}(i)} (\tilde{\mathbf{x}}_i - \tilde{\mathbf{x}}_j) \phi_x(\mathbf{m}_{ij}) \\
&= R\mathbf{x}_i + \mathbf{t} + C \sum_{j \in \mathcal{N}(i)} ((R\mathbf{x}_i + \mathbf{t}) - (R\mathbf{x}_j + \mathbf{t})) \phi_x(\mathbf{m}_{ij}) \\
&= R\mathbf{x}_i + \mathbf{t} + C \sum_{j \in \mathcal{N}(i)} R(\mathbf{x}_i - \mathbf{x}_j) \phi_x(\mathbf{m}_{ij}) \\
&= R\left(\mathbf{x}_i + C \sum_{j \in \mathcal{N}(i)} (\mathbf{x}_i - \mathbf{x}_j) \phi_x(\mathbf{m}_{ij})\right) + \mathbf{t} \\
&= R\mathbf{x}_i' + \mathbf{t}
\end{align}
demonstrating equivariance.
\end{proof}

\textbf{Stage 3: Invariant feature update}

The scalar features are updated using aggregated messages:
\begin{align}
\mathbf{m}_i &= \sum_{j \in \mathcal{N}(i)} \mathbf{m}_{ij} \label{eq:egnn_aggregate}\\
h_i' &= \phi_h(h_i, \mathbf{m}_i) \label{eq:egnn_feature_update}
\end{align}

where $\phi_h: \mathbb{R}^{d + d'} \to \mathbb{R}^d$ is an MLP. This update is trivially $E(n)$-invariant since it operates only on invariant scalar quantities.

\paragraph{Equivariance Properties}

The architecture guarantees that for any Euclidean transformation $g = (R, \mathbf{t}) \in E(n)$:
\begin{align}
h_i'(g \cdot \mathbf{X}) &= h_i'(\mathbf{X}) \quad \text{(invariance)} \\
\mathbf{x}_i'(g \cdot \mathbf{X}) &= g \cdot \mathbf{x}_i'(\mathbf{X}) \quad \text{(equivariance)}
\end{align}

where $g \cdot \mathbf{X} = \{R\mathbf{x}_i + \mathbf{t}\}_{i=1}^N$ denotes the transformation of all coordinates.

\begin{example}[EGNN Message Passing for a Molecular System]\label{ex:egnn_molecular}
Consider a small molecule with three atoms located at:
\begin{equation}
\mathbf{x}_1 = \begin{pmatrix} 0 \\ 0 \\ 0 \end{pmatrix}, \quad
\mathbf{x}_2 = \begin{pmatrix} 1.5 \\ 0 \\ 0 \end{pmatrix}, \quad
\mathbf{x}_3 = \begin{pmatrix} 0.75 \\ 1.3 \\ 0 \end{pmatrix}
\end{equation}

With initial scalar features $h_1 = [6.0]$ (carbon), $h_2 = [1.0]$ (hydrogen), $h_3 = [8.0]$ (oxygen).

\textbf{Step 1: Distance computation}
\begin{align}
d_{12}^2 &= \|\mathbf{x}_1 - \mathbf{x}_2\|^2 = 2.25 \\
d_{13}^2 &= \|\mathbf{x}_1 - \mathbf{x}_3\|^2 = 0.75^2 + 1.3^2 = 2.2525 \\
d_{23}^2 &= \|\mathbf{x}_2 - \mathbf{x}_3\|^2 = 0.75^2 + 1.3^2 = 2.2525
\end{align}

\textbf{Step 2: Messages} (using a simplified $\phi_e$)
\begin{equation}
\mathbf{m}_{12} = \phi_e(h_1, h_2, d_{12}^2) = \text{MLP}([6.0, 1.0, 2.25])
\end{equation}

\textbf{Step 3: Coordinate update}
\begin{equation}
\mathbf{x}_1' = \mathbf{x}_1 + \frac{1}{2}\left[(\mathbf{x}_1 - \mathbf{x}_2)\phi_x(\mathbf{m}_{12}) + (\mathbf{x}_1 - \mathbf{x}_3)\phi_x(\mathbf{m}_{13})\right]
\end{equation}

This update automatically preserves equivariance: if we rotate the entire molecule, the updated coordinates rotate in the same way.
\end{example}

\subsubsection{SE(3)-Transformers}

Fuchs et al.~\citep{fuchs2020se3transformers} develop the SE(3)-Transformers, which combine equivariance to the group of rigid motions $SE(3)$ with attention mechanisms, providing a more expressive architecture than EGNNs by using complete tensor representations of $SO(3)$.

\paragraph{The SE(3) Group and its Properties}

\begin{definition}[Group SE(3)]
The group of rigid motions $SE(3) = \mathbb{R}^3 \rtimes SO(3)$ consists of all transformations that preserve distances and orientation in $\mathbb{R}^3$.
\end{definition}

Each element of $SE(3)$ is written as $(R, \mathbf{t})$ where $R \in SO(3)$ and $\mathbf{t} \in \mathbb{R}^3$, acting on points via:
\begin{equation}
(R, \mathbf{t}) \cdot \mathbf{x} = R\mathbf{x} + \mathbf{t}
\end{equation}

\begin{proposition}[Differences with E(3)]
$SE(3) \subset E(3)$ is the subgroup that preserves orientation (it does not include reflections). This restriction is crucial in applications where chirality matters, such as in molecular chemistry.
\end{proposition}

\begin{proof}
$SE(3) = \mathbb{R}^3 \rtimes SO(3)$ consists of pairs $(R, \mathbf{t})$ with $R \in SO(3)$ (i.e., $\det R = +1$), while $E(3) = \mathbb{R}^3 \rtimes O(3)$ also includes elements with $\det R = -1$ (reflections). The inclusion $SE(3) \subset E(3)$ is strict: for example, the reflection $(x,y,z) \mapsto (-x,y,z)$ belongs to $E(3)$ but not to $SE(3)$. In molecular chemistry, a molecule and its mirror image (enantiomer) may have different properties (chirality), so it is necessary to distinguish between transformations that preserve ($SE(3)$) and those that invert ($O(3) \setminus SO(3)$) orientation.
\end{proof}

\paragraph{Irreducible Representations of SO(3)}

Unlike EGNNs, which use only scalar features and position vectors, SE(3)-Transformers work with complete tensor representations of $SO(3)$.

\begin{definition}[Irreducible Representations of SO(3)]\label{def:irreps_so3}
The irreducible representations of $SO(3)$ are classified by non-negative integers $l = 0, 1, 2, \ldots$ (degree or angular momentum). The degree-$l$ representation has dimension $2l + 1$ and is denoted $\rho^{(l)}: SO(3) \to \mathbb{R}^{(2l+1) \times (2l+1)}$.

For a rotation $R \in SO(3)$ parameterized by Euler angles $(\alpha, \beta, \gamma)$, the matrix elements are:
\begin{equation}
D_{mm'}^{(l)}(\alpha, \beta, \gamma) = e^{-im\alpha} d_{mm'}^{(l)}(\beta) e^{-im'\gamma}
\end{equation}
where $d_{mm'}^{(l)}$ are the Wigner small functions.
\end{definition}

\begin{example}[Low Degrees of SO(3)]
\leavevmode
\begin{itemize}
    \item $l = 0$: Scalars (dimension 1), invariant under rotations
    \item $l = 1$: Vectors (dimension 3), transform as $\mathbf{v}' = R\mathbf{v}$
    \item $l = 2$: Symmetric traceless second-order tensors (dimension 5)
\end{itemize}
\end{example}

\paragraph{Feature Tensors}

SE(3)-Transformers use tensor representations that transform according to irreducible representations of $SO(3)$:

\begin{definition}[Feature Tensor]\label{def:feature_tensor_se3}
A \textit{feature tensor} of type $l$ at node $i$ is a vector $\mathbf{f}_i^{(l)} \in \mathbb{R}^{2l+1}$ that transforms under rotations $R \in SO(3)$ as:
\begin{equation}
(\mathbf{f}_i^{(l)})' = D^{(l)}(R) \mathbf{f}_i^{(l)}
\end{equation}
where $D^{(l)}(R)$ is the degree-$l$ irreducible representation matrix.
\end{definition}

In practice, each node maintains multiple copies (multiplicities) of each irreducible type:
\begin{equation}
\mathbf{F}_i = \bigoplus_{l=0}^{L_{\max}} \bigoplus_{k=1}^{n_l} \mathbf{f}_{i,k}^{(l)}
\end{equation}
where $n_l$ is the multiplicity of type $l$ and $L_{\max}$ is the maximum degree.

\paragraph{Tensor Products and Channel Mixing}

The combination of feature tensors of different types requires the use of Clebsch-Gordan tensor products, which guarantee equivariance under $SO(3)$.

\begin{theorem}[Clebsch-Gordan Decomposition]\label{thm:clebsch_gordan}
The tensor product of two irreducible representations of $SO(3)$ decomposes as a direct sum of irreducibles:
\begin{equation}
\rho^{(l_1)} \otimes \rho^{(l_2)} = \bigoplus_{l=|l_1-l_2|}^{l_1+l_2} \rho^{(l)}
\end{equation}

The Clebsch-Gordan coefficients $C_{l_1 m_1, l_2 m_2}^{l m}$ implement this decomposition:
\begin{equation}
(f^{(l_1)} \otimes g^{(l_2)})^{(l)}_m = \sum_{m_1=-l_1}^{l_1} \sum_{m_2=-l_2}^{l_2} C_{l_1 m_1, l_2 m_2}^{l m} f^{(l_1)}_{m_1} g^{(l_2)}_{m_2}
\end{equation}

See Section~\ref{sec:clebsch_gordan} of \autoref{ch:teoremas_complementarios} for the complete proof.
\end{theorem}

\textbf{Selection Rule.} The tensor product produces degrees in the range $|l_1 - l_2| \leq l \leq l_1 + l_2$. For example:
\begin{itemize}
    \item Vector $\times$ Vector ($l=1 \times l=1$): generates $l \in \{0, 1, 2\}$ (scalar + vector + tensor)
    \item Scalar $\times$ Vector ($l=0 \times l=1$): generates only $l=1$ (vector)
\end{itemize}

\paragraph{Equivariant Attention Mechanism}

The key challenge is to design attention weights that are invariant under $SE(3)$, allowing the weighted combination of tensor features to preserve equivariance.

The attention weights are computed using only the scalar component ($l=0$) of the feature tensors:
\begin{equation}\label{eq:se3_attention_weights}
\alpha_{ij} = \text{softmax}_j\left(\text{MLP}\left(\mathbf{f}_i^{(0)} \oplus \mathbf{f}_j^{(0)} \oplus h(r_{ij})\right)\right)
\end{equation}

where:
\begin{itemize}
    \item $\mathbf{f}_i^{(0)}, \mathbf{f}_j^{(0)}$ are the scalar components (invariant)
    \item $r_{ij} = \|\mathbf{x}_i - \mathbf{x}_j\|$ is the distance (invariant under $SE(3)$)
    \item $h(r_{ij})$ is a radial embedding of the distance
    \item $\oplus$ denotes concatenation
\end{itemize}

\begin{proposition}[Invariance of the Attention Weights]
The weights $\alpha_{ij}$ defined in~\eqref{eq:se3_attention_weights} are invariant under $SE(3)$: for any $(R, \mathbf{t}) \in SE(3)$,
\begin{equation}
\alpha_{ij}((R, \mathbf{t}) \cdot \{\mathbf{x}_k\}, \{\mathbf{F}_k\}) = \alpha_{ij}(\{\mathbf{x}_k\}, \{\mathbf{F}_k\})
\end{equation}
\end{proposition}

\begin{proof}
The attention weights depend on the positions $\mathbf{x}_k$ only through the differences $\mathbf{x}_i - \mathbf{x}_j$ and the norms $\|\mathbf{x}_i - \mathbf{x}_j\|$. Under the action $(R, \mathbf{t}) \in SE(3)$: the transformed difference is $R\mathbf{x}_i + \mathbf{t} - (R\mathbf{x}_j + \mathbf{t}) = R(\mathbf{x}_i - \mathbf{x}_j)$, and the norm $\|R(\mathbf{x}_i - \mathbf{x}_j)\| = \|\mathbf{x}_i - \mathbf{x}_j\|$ by the orthogonality of $R$. Since $\alpha_{ij}$ is computed as a function of $\|\mathbf{x}_i - \mathbf{x}_j\|$ and of the invariant features $\mathbf{F}_k$ (which by hypothesis are invariant under $SE(3)$), the attention weights are invariant.
\end{proof}

\paragraph{Equivariant Linear Combination}

Once the invariant attention weights are computed, the update of the tensor features is performed via a weighted linear combination within each irreducible type:

\begin{equation}\label{eq:se3_feature_update}
\mathbf{f}_i^{(l), \text{new}} = \sum_{j \in \mathcal{N}(i)} \alpha_{ij} \cdot W_{ij}^{(l)} \mathbf{f}_j^{(l)}
\end{equation}

where $W_{ij}^{(l)}$ are learnable mixing matrices that act within the type-$l$ subspace.

\begin{theorem}[Equivariance of SE(3)-Transformers]\label{thm:se3_transformer_equivariance}
The update~\eqref{eq:se3_feature_update} with weights~\eqref{eq:se3_attention_weights} is equivariant: if we transform all positions and features by $(R, \mathbf{t}) \in SE(3)$, then
\begin{equation}
(\mathbf{f}_i^{(l), \text{new}})' = D^{(l)}(R) \mathbf{f}_i^{(l), \text{new}}
\end{equation}
\end{theorem}

\begin{proof}
The attention weights $\alpha_{ij}$ are invariant under $SE(3)$ (by the previous Proposition). The key point is that the matrices $W_{ij}^{(l)}$ in equation~\eqref{eq:se3_feature_update} act as scalar mixing coefficients between channels of the same type $l$, that is, they are multiples of the identity in the $(2l+1)$-dimensional representation subspace. In the implementation of \citet{fuchs2020se3transformers}, the value messages are constructed via TFN-type equivariant kernels~\citep{thomas2018tensorfieldnetworks}, which use Clebsch-Gordan coefficients and learnable radial functions. Therefore, each value message $\mathbf{v}_j^{(l)} = W_{ij}^{(l)} \mathbf{f}_j^{(l)}$ transforms as $(\mathbf{v}_j^{(l)})' = D^{(l)}(R) \mathbf{v}_j^{(l)}$. Then:
\begin{align}
(\mathbf{f}_i^{(l), \text{new}})' &= \sum_j \alpha_{ij} \, (\mathbf{v}_j^{(l)})' = \sum_j \alpha_{ij} \, D^{(l)}(R) \mathbf{v}_j^{(l)} \\
&= D^{(l)}(R) \sum_j \alpha_{ij} \, \mathbf{v}_j^{(l)} = D^{(l)}(R) \mathbf{f}_i^{(l), \text{new}}
\end{align}
where in the second step we use the invariance of $\alpha_{ij}$ and the linearity of $D^{(l)}(R)$. Translations do not affect the tensor features (only the positions), completing the equivariance under $SE(3)$.
\end{proof}

\paragraph{Radial Filters and Spatial Dependence}

To incorporate spatial information, radial filters are used that modulate the weights according to distance:

\begin{equation}
W_{ij}^{(l)} = \sum_{k=1}^{K} w_k^{(l)} R_k(r_{ij})
\end{equation}

where $R_k(r)$ are radial basis functions (typically Gaussians or Bessels) and $w_k^{(l)}$ are learnable parameters.

\subsubsection{Tensor Field Networks}

Thomas et al.~\citep{thomas2018tensorfieldnetworks} introduce Tensor Field Networks (TFNs), providing the fundamental theoretical framework for $SE(3)$-equivariant networks via an elegant formulation based on tensor fields and spherical harmonics.

\paragraph{Tensor Fields as Functions over Points}

\begin{definition}[Tensor Field of Order $l$]\label{def:tensor_field}
A \textit{tensor field} of order $l$ associated with a set of points $\{\mathbf{x}_i\}_{i=1}^N$ is a collection of tensors $\{T_i^{(l)}\}_{i=1}^N$ where each $T_i^{(l)} \in \mathbb{R}^{2l+1}$ transforms according to the degree-$l$ irreducible representation of $SO(3)$:
\begin{equation}
(T_i^{(l)})' = D^{(l)}(R) T_i^{(l)} \quad \text{for } R \in SO(3)
\end{equation}
\end{definition}

In TFNs, each point in a 3D point cloud carries multiple tensor fields of different orders, capturing rich geometric information in an equivariant manner.

\paragraph{Convolution of Tensor Fields}

The fundamental operation in TFNs is the convolution between tensor fields, which must preserve equivariance under $SE(3)$.

\begin{definition}[TFN Convolution]\label{def:tfn_convolution}
The convolution between a tensor field $T^{(l_{\text{in}})}$ and a kernel $W^{(l_{\text{out}}, l_{\text{in}})}$ at point $i$ is defined as:
\begin{equation}\label{eq:tfn_conv}
(T * W)_i^{(l_{\text{out}})} = \sum_{j \in \mathcal{N}(i)} \sum_{l_{\text{in}}} K_{ij}^{(l_{\text{out}}, l_{\text{in}})} T_j^{(l_{\text{in}})}
\end{equation}
where $K_{ij}^{(l_{\text{out}}, l_{\text{in}})}$ is an SE(3)-equivariant kernel that depends on the relative position $\mathbf{r}_{ij} = \mathbf{x}_i - \mathbf{x}_j$.
\end{definition}

\paragraph{Radial-Angular Factorization of the Kernel}

The key design of TFNs is to factorize the kernel into radial (learnable) and angular (fixed, given by spherical harmonics) components:

\begin{equation}\label{eq:tfn_kernel_factorization}
K_{ij}^{(l_{\text{out}}, l_{\text{in}})} = \sum_{l_{\text{filter}}} R^{l_{\text{filter}}}_{l_{\text{out}}, l_{\text{in}}}(r_{ij}) \mathcal{Y}^{(l_{\text{filter}})}(\hat{\mathbf{r}}_{ij})
\end{equation}

where:
\begin{itemize}
    \item $r_{ij} = \|\mathbf{r}_{ij}\|$ is the distance (invariant scalar)
    \item $\hat{\mathbf{r}}_{ij} = \mathbf{r}_{ij} / r_{ij}$ is the unit direction
    \item $\mathcal{Y}^{(l_{\text{filter}})}$ are spherical harmonics of order $l_{\text{filter}}$
    \item $R^{l_{\text{filter}}}_{l_{\text{out}}, l_{\text{in}}}(r)$ are learnable radial functions
\end{itemize}

\begin{proposition}[Equivariance of the Factorization]\label{prop:tfn_kernel_equivariance}
The kernel~\eqref{eq:tfn_kernel_factorization} is SE(3)-equivariant: under a rotation $R \in SO(3)$,
\begin{equation}
K_{ij}^{(l_{\text{out}}, l_{\text{in}})}(R\mathbf{r}_{ij}) = D^{(l_{\text{out}})}(R) K_{ij}^{(l_{\text{out}}, l_{\text{in}})}(\mathbf{r}_{ij}) (D^{(l_{\text{in}})}(R))^{-1}
\end{equation}
\end{proposition}

\begin{proof}
The kernel factorizes as $K_{ij}^{(l_{\text{out}}, l_{\text{in}})}(\mathbf{r}) = \sum_{J} \varphi_J(\|\mathbf{r}\|) \, C^{l_{\text{out}}}_{l_{\text{in}}, J} \, Y_J(\hat{\mathbf{r}})$, where $\varphi_J$ are radial functions (invariant under rotation), $C$ are Clebsch-Gordan coefficients and $Y_J$ are spherical harmonics. Under a rotation $R$: the radial part is invariant ($\|R\mathbf{r}\| = \|\mathbf{r}\|$), and the spherical harmonics transform as $Y_J^m(R^{-1}\hat{\mathbf{r}}) = \sum_{m'} D_{m'm}^{(J)}(R) Y_J^{m'}(\hat{\mathbf{r}})$. The properties of the Clebsch-Gordan coefficients guarantee that this transformation of the harmonics translates into the tensor transformation law $D^{(l_{\text{out}})}(R) K (D^{(l_{\text{in}})}(R))^{-1}$, as equivariance requires.
\end{proof}

\paragraph{Coupling via Clebsch-Gordan Products}

The combination of input tensors with the filter kernel is performed via Clebsch-Gordan tensor products:

\begin{equation}\label{eq:tfn_cg_coupling}
(T^{(l_{\text{in}})} \otimes \mathcal{Y}^{(l_{\text{filter}})})^{(l_{\text{out}})}_m = \sum_{m_{\text{in}}=-l_{\text{in}}}^{l_{\text{in}}} \sum_{m_f=-l_{\text{filter}}}^{l_{\text{filter}}} C_{l_{\text{in}} m_{\text{in}}, l_{\text{filter}} m_f}^{l_{\text{out}} m} T^{(l_{\text{in}})}_{m_{\text{in}}} \mathcal{Y}^{(l_{\text{filter}})}_{m_f}(\hat{\mathbf{r}}_{ij})
\end{equation}

The angular-momentum selection rule determines which combinations are permitted:
\begin{equation}
|l_{\text{in}} - l_{\text{filter}}| \leq l_{\text{out}} \leq l_{\text{in}} + l_{\text{filter}}
\end{equation}

\begin{example}[Specific Tensor Product]\label{ex:tfn_tensor_product}
Consider the coupling of a vector field ($l_{\text{in}} = 1$) with a quadrupolar spherical-harmonic filter ($l_{\text{filter}} = 2$):

Possible products: $|1 - 2| \leq l_{\text{out}} \leq 1 + 2 \Rightarrow l_{\text{out}} \in \{1, 2, 3\}$

\begin{itemize}
    \item $l_{\text{out}} = 1$: Generates a new vector field (3 components)
    \item $l_{\text{out}} = 2$: Generates a second-order tensor field (5 components)
    \item $l_{\text{out}} = 3$: Generates a third-order tensor field (7 components)
\end{itemize}

Each component is computed explicitly using:
\begin{equation}
T^{(l_{\text{out}})}_m = \sum_{m_1=-1}^{1} \sum_{m_2=-2}^{2} C_{1 m_1, 2 m_2}^{l_{\text{out}} m} T^{(1)}_{m_1} \mathcal{Y}^{(2)}_{m_2}(\hat{\mathbf{r}}_{ij})
\end{equation}
\end{example}

\paragraph{Radial Basis Functions}

The radial functions $R^{l_{\text{filter}}}_{l_{\text{out}}, l_{\text{in}}}(r)$ capture distance dependencies and are the only learnable components of the kernel. Common options include:

\textbf{Gaussian Functions:}
\begin{equation}
R_k(r) = \exp\left(-\frac{(r - \mu_k)^2}{2\sigma^2}\right)
\end{equation}
where $\{\mu_k\}$ are fixed centers over a range of distances.

\textbf{Bessel Functions:}
\begin{equation}
R_k(r) = \sqrt{\frac{2}{r_{\text{cut}}}} \frac{\sin(k\pi r / r_{\text{cut}})}{r}
\end{equation}
which vanish at the cutoff radius $r_{\text{cut}}$.

\textbf{Polynomials with Envelope:}
\begin{equation}
R_k(r) = r^k e^{-\alpha r} \quad \text{for } k = 0, 1, 2, \ldots
\end{equation}

\paragraph{Complete TFN Layer}

A TFN layer combines convolutions over multiple tensor types:

\begin{equation}\label{eq:tfn_layer}
T_i^{(l), \text{out}} = \sum_{l' \in \mathcal{L}_{\text{in}}} \sum_{j \in \mathcal{N}(i)} \sum_{l_f} W^{l_f}_{l,l'}(r_{ij}) \left(T_j^{(l')} \otimes \mathcal{Y}^{(l_f)}(\hat{\mathbf{r}}_{ij})\right)^{(l)}
\end{equation}

where:
\begin{itemize}
    \item $\mathcal{L}_{\text{in}}$ is the set of input tensor types
    \item $W^{l_f}_{l,l'}(r)$ are learnable distance-dependent weights
    \item The tensor product $\otimes$ is implemented via Clebsch-Gordan coefficients
\end{itemize}

\begin{algorithm}[h]
\caption{Forward Pass of a TFN Layer}
\label{alg:tfn_forward}
\begin{algorithmic}[1]
\Require Input tensor fields $\{T_i^{(l)}\}$, positions $\{\mathbf{x}_i\}$, cutoff radius $r_{\text{cut}}$
\Ensure Output tensor fields $\{T_i^{(l), \text{out}}\}$
\For{each point $i$}
    \State Initialize $T_i^{(l), \text{out}} \leftarrow \mathbf{0}$ for all $l$
    \For{each neighbor $j$ with $r_{ij} < r_{\text{cut}}$}
        \State Compute $r_{ij} \leftarrow \|\mathbf{x}_i - \mathbf{x}_j\|$
        \State Compute $\hat{\mathbf{r}}_{ij} \leftarrow (\mathbf{x}_i - \mathbf{x}_j) / r_{ij}$
        \State Evaluate spherical harmonics $\mathcal{Y}^{(l_f)}(\hat{\mathbf{r}}_{ij})$ for $l_f \in \{0, 1, \ldots, L_{\max}\}$
        \For{each output type $l_{\text{out}}$}
            \For{each input type $l_{\text{in}}$}
                \For{each filter order $l_f$ with $|l_{\text{in}} - l_f| \leq l_{\text{out}} \leq l_{\text{in}} + l_f$}
                    \State Evaluate radial weight $w \leftarrow W^{l_f}_{l_{\text{out}}, l_{\text{in}}}(r_{ij})$
                    \State Compute CG product: $P \leftarrow w \cdot (T_j^{(l_{\text{in}})} \otimes \mathcal{Y}^{(l_f)}(\hat{\mathbf{r}}_{ij}))^{(l_{\text{out}})}$
                    \State $T_i^{(l_{\text{out}}), \text{out}} \leftarrow T_i^{(l_{\text{out}}), \text{out}} + P$
                \EndFor
            \EndFor
        \EndFor
    \EndFor
\EndFor
\State \textbf{return} $\{T_i^{(l), \text{out}}\}$
\end{algorithmic}
\end{algorithm}

\subsubsection{Comparison: Compact SO(3) vs. Non-Compact SE(3)}\label{subsubsec:so3_vs_se3_comparison}

It is instructive to contrast the fundamental differences between architectures for compact groups (SO(3)) and non-compact groups (SE(3), E(3)):

\begin{table}[h]
\centering
\begin{tabular}{|l|p{5cm}|p{5cm}|}
\hline
\textbf{Property} & \textbf{SO(3) (Compact)} & \textbf{SE(3) (Non-Compact)} \\
\hline
Group structure & Rotations only & Rotations + translations \\
\hline
Compactness & Compact ($\sim \mathbb{RP}^3$) & Non-compact (unbounded $\mathbb{R}^3$ subgroup) \\
\hline
Haar measure & Finite and normalizable & Infinite (non-normalizable on $\mathbb{R}^3$) \\
\hline
Irreducible representations & All finite-dimensional: $\dim(\rho^{(l)}) = 2l+1$ & Mixed: finite ($SO(3)$ part) and infinite ($\mathbb{R}^3$ part) \\
\hline
Integration over the group & Directly feasible & Requires restriction to a compact subgroup \\
\hline
Basic invariants & None (transitive action on $S^2$) & Distances $\|\mathbf{x}_i - \mathbf{x}_j\|$ \\
\hline
Architecture design & Pure tensor fields & Scalar/vector separation essential \\
\hline
Computational complexity & $O(L^3)$ per CG product & $O(L^3)$ + coordinate-handling overhead \\
\hline
Main applications & Spherical data, CMB, climate & Molecules, point clouds, robotics \\
\hline
\end{tabular}
\caption{Comparison between equivariant SO(3) and SE(3) architectures}
\label{tab:so3_vs_se3}
\end{table}

\paragraph{Implications for Design}

The mathematical differences have important practical consequences:

\begin{enumerate}
    \item \textbf{Processing of invariants}: In SE(3)/E(3), relative distances provide natural invariant quantities that can be processed with standard MLPs. In SO(3), no such intrinsic invariants exist.

    \item \textbf{Information separation}: SE(3) architectures (such as EGNNs) explicitly separate invariant scalar features from equivariant coordinates. Pure SO(3) architectures (such as Spherical CNNs) work with homogeneous tensor fields.

    \item \textbf{Parametric efficiency}: EGNNs can be more efficient than full TFNs by leveraging distance invariants, sacrificing expressivity for simplicity.
\end{enumerate}

\subsubsection{Computational Complexity Analysis}

\paragraph{Time Complexity}

For an equivariant layer processing $N$ points with neighborhoods of average size $k$ and maximum degree $L_{\max}$:

\textbf{EGNN:}
\begin{itemize}
    \item Message computation: $O(N \cdot k \cdot d^2)$ where $d$ is the feature dimension
    \item Coordinate update: $O(N \cdot k \cdot n)$ where $n$ is the spatial dimension
    \item \textbf{Total}: $O(N \cdot k \cdot \max(d^2, n))$
\end{itemize}

\textbf{SE(3)-Transformers:}
\begin{itemize}
    \item Attention computation: $O(N \cdot k \cdot L_{\max}^2)$
    \item Tensor products: $O(N \cdot k \cdot L_{\max}^3 \cdot n_{\text{mult}})$ per CG product
    \item \textbf{Total}: $O(N \cdot k \cdot L_{\max}^3 \cdot n_{\text{mult}})$
\end{itemize}

\textbf{Tensor Field Networks:}
\begin{itemize}
    \item Spherical harmonic evaluation: $O(N \cdot k \cdot L_{\max}^2)$
    \item Clebsch-Gordan products: $O(N \cdot k \cdot L_{\max}^3 \cdot n_{l_{\text{in}}} \cdot n_{l_{\text{out}}})$
    \item Radial functions: $O(N \cdot k \cdot K_{\text{radial}})$ negligible
    \item \textbf{Total}: $O(N \cdot k \cdot L_{\max}^3 \cdot n_{\text{types}})$
\end{itemize}

\paragraph{Space Complexity}

\textbf{EGNN:}
\begin{equation}
\text{Memory} = O(N \cdot d + N \cdot n) = O(N \cdot (d + n))
\end{equation}

\textbf{TFN/SE(3)-Transformers:}
\begin{equation}
\text{Memory} = O\left(N \cdot \sum_{l=0}^{L_{\max}} n_l (2l+1)\right) = O(N \cdot L_{\max}^2 \cdot n_{\text{avg}})
\end{equation}

where $n_{\text{avg}}$ is the average multiplicity per irreducible type.

\paragraph{Scalability and Trade-offs}

\begin{itemize}
    \item \textbf{EGNNs}: More scalable, suitable for large systems (thousands of atoms). Less expressive since they do not use high-order tensor types.

    \item \textbf{TFNs/SE(3)-Transformers}: More expressive, capturing complex directional interactions. Limited to medium-sized systems (hundreds of atoms) with $L_{\max} \leq 4$.

    \item \textbf{Practical compromise}: Use EGNNs for pre-processing and TFNs for refinement in critical regions.
\end{itemize}

\subsubsection{Enhancements with Geometric and Physical Quantities}

The work of Tholke and De Fabritiis~\citep{tholke2022geometric} demonstrates that incorporating domain-specific geometric and physical quantities significantly improves the performance of $E(3)$-equivariant networks without breaking the symmetry.

\paragraph{Additional Geometric Features}

Beyond Euclidean distances, other geometric invariants can be incorporated:

\textbf{Angles between triplets} (invariant under $SE(3)$):
\begin{equation}
\theta_{ijk} = \arccos\left(\frac{(\mathbf{x}_i - \mathbf{x}_j) \cdot (\mathbf{x}_k - \mathbf{x}_j)}{\|\mathbf{x}_i - \mathbf{x}_j\| \|\mathbf{x}_k - \mathbf{x}_j\|}\right)
\end{equation}

\textbf{Dihedral angles} (quadruple of points, captures chirality):
\begin{equation}
\phi_{ijkl} = \arctan2\left(\frac{(\mathbf{n}_1 \times \mathbf{n}_2) \cdot \mathbf{b}}{\|\mathbf{b}\|}, \mathbf{n}_1 \cdot \mathbf{n}_2\right)
\end{equation}
where $\mathbf{n}_1, \mathbf{n}_2$ are normals to planes and $\mathbf{b} = \mathbf{x}_k - \mathbf{x}_j$ is the axis.

\textbf{Oriented volumes} (sensitive to chirality):
\begin{equation}
V_{ijkl} = \frac{1}{6} \left|(\mathbf{x}_j - \mathbf{x}_i) \cdot \left((\mathbf{x}_k - \mathbf{x}_i) \times (\mathbf{x}_l - \mathbf{x}_i)\right)\right|
\end{equation}

\textbf{Chirality and E(3) vs. SE(3).} Dihedral angles are invariant under $SE(3)$ (which preserves orientation) but change sign under reflections. If complete equivariance under $E(3)$ (including reflections) is required, $|\phi_{ijkl}|$ or pairs $(\cos \phi, |\sin \phi|)$ must be used.

\paragraph{Physical Invariants for Molecular Systems}

In computational chemistry, physically meaningful quantities improve the inductive bias:

\textbf{Coulomb energy} (invariant):
\begin{equation}
E_{\text{Coulomb}} = \sum_{i < j} \frac{q_i q_j}{4\pi \epsilon_0 \|\mathbf{x}_i - \mathbf{x}_j\|}
\end{equation}

\textbf{Lennard-Jones potential} (invariant):
\begin{equation}
V_{LJ}(r) = 4\epsilon \left[\left(\frac{\sigma}{r}\right)^{12} - \left(\frac{\sigma}{r}\right)^6\right]
\end{equation}

\textbf{Inertia tensor} (covariant under rotations):
\begin{equation}
I_{ab} = \sum_i m_i \left(\|\mathbf{x}_i\|^2 \delta_{ab} - x_i^a x_i^b\right)
\end{equation}

\paragraph{Hybrid Architecture}

Incorporation is performed by concatenating invariant features with the learned features:
\begin{equation}
\mathbf{h}_i^{(l+1)} = \phi_{\text{MLP}}\left(\mathbf{h}_i^{(l)} \oplus \mathbf{f}_{\text{geo}}(\{\mathbf{x}_k\}_{k \in \mathcal{N}(i)}) \oplus \mathbf{f}_{\text{phys}}(\{q_k, m_k\})\right)
\end{equation}
where:
\begin{itemize}
    \item $\mathbf{f}_{\text{geo}}$ extracts geometric invariants (angles, distances)
    \item $\mathbf{f}_{\text{phys}}$ computes physical invariants (energies, potentials)
    \item $\oplus$ denotes concatenation
\end{itemize}

\begin{example}[Molecular System with Enriched Features]
For a water molecule (H$_2$O), the additional geometric invariants include:
\begin{itemize}
    \item O-H distances: $d_1 \approx 0.96$ Å, $d_2 \approx 0.96$ Å
    \item H-O-H angle: $\theta \approx 104.5^\circ$
    \item Dipole moment (norm, invariant): $|\boldsymbol{\mu}| \approx 1.85$ Debye
    \item Principal moment of inertia (spectral invariant)
\end{itemize}

These features complement the basic distances, improving the prediction of properties such as bond energies and vibrational frequencies.
\end{example}

$E(n)$- and $\text{SE}(3)$-equivariant networks have demonstrated consistent improvements in data efficiency (a typical reduction of one order of magnitude in training requirements), out-of-distribution generalization and transferability to unseen configurations, in domains where Euclidean or rigid symmetries are relevant: molecular systems, three-dimensional point clouds and particle analysis. The main computational challenge lies in computing Clebsch-Gordan products, whose cost grows with $L_{\max}$, and in the numerical stability of high-order spherical harmonic sums.

\subsection{Permutation Groups}\label{subsec:grupos_permutaciones}

Permutation groups $S_n$ represent one of the most fundamental cases in the design of equivariant neural architectures, where invariance to the ordering of elements is crucial for processing unordered sets and graph structures. Unlike continuous symmetries such as rotations or translations, permutations are purely combinatorial, transforming the arrangement of elements discretely without altering their intrinsic values.

\subsubsection{The Symmetric Group: Algebraic Foundations}

The symmetric group $S_n$ constitutes one of the most studied mathematical objects in group theory, with a rich algebraic structure that underpins its application in deep learning.

\begin{definition}[Symmetric Group]\label{def:symmetric_group}
The \textit{symmetric group} $S_n$ is the group of all bijections (permutations) of the set $\{1, 2, \ldots, n\}$ under the operation of composition. Formally:
\begin{equation}
S_n = \{\pi: \{1,\ldots,n\} \to \{1,\ldots,n\} \mid \pi \text{ is bijective}\}
\end{equation}
with group operation given by composition $(\pi_1 \cdot \pi_2)(i) = \pi_1(\pi_2(i))$.
\end{definition}

\begin{proposition}[Fundamental Properties of $S_n$]\label{prop:symmetric_group_properties}
The symmetric group $S_n$ satisfies:
\begin{enumerate}
    \item \textbf{Order}: $|S_n| = n!$ (factorial of $n$)
    \item \textbf{Non-commutativity}: $S_n$ is non-abelian for $n \geq 3$
    \item \textbf{Generation}: $S_n$ is generated by the transpositions $(i, i+1)$ for $i=1,\ldots,n-1$
    \item \textbf{Presentation}: $S_n = \langle \sigma_1, \ldots, \sigma_{n-1} \mid \sigma_i^2 = e, (\sigma_i\sigma_{i+1})^3 = e, (\sigma_i\sigma_j)^2 = e \text{ for } |i-j| > 1 \rangle$
\end{enumerate}
\end{proposition}

\begin{proof}[Partial proof]
(1) There are $n$ choices for $\pi(1)$, then $n-1$ for $\pi(2)$, etc., giving $n!$ total permutations. (2) For $n=3$, the permutations $(12)$ and $(23)$ do not commute: $(12)(23) = (123) \neq (132) = (23)(12)$. (3) and (4) are proven by explicit construction; see \citep{dummit2004abstract} for full details.
\end{proof}

\paragraph{Cycle Notation and Decomposition}

Every permutation $\pi \in S_n$ decomposes uniquely (up to ordering) as a product of disjoint cycles. This decomposition is fundamental to understanding the algebraic and representational structure of $S_n$.

\begin{example}[Cycle Notation]
The permutation $\pi \in S_5$ defined by $\pi(1)=3, \pi(2)=2, \pi(3)=5, \pi(4)=1, \pi(5)=4$ is written in cycle notation as:
\begin{equation}
\pi = (1\,3\,5\,4) = (1\,3\,5\,4)(2)
\end{equation}
indicating that $1 \to 3 \to 5 \to 4 \to 1$ forms a 4-cycle and $2$ is a fixed point.
\end{example}

\paragraph{Matrix Representation}

For computational implementation, permutations are represented via permutation matrices.

\begin{definition}[Permutation Matrix]\label{def:permutation_matrix}
For $\pi \in S_n$, the \textit{permutation matrix} $P_\pi \in \mathbb{R}^{n \times n}$ is defined by:
\begin{equation}
(P_\pi)_{ij} = \begin{cases}
1 & \text{if } j = \pi(i) \\
0 & \text{otherwise}
\end{cases}
\end{equation}
Equivalently, $P_\pi$ is the matrix whose $i$-th row is the vector $\mathbf{e}_{\pi(i)}$.
\end{definition}

\begin{proposition}[Properties of Permutation Matrices]\label{prop:permutation_matrix_properties}
Permutation matrices satisfy:
\begin{enumerate}
    \item $P_{\pi_1 \cdot \pi_2} = P_{\pi_1} P_{\pi_2}$ (homomorphism)
    \item $P_\pi^{-1} = P_\pi^T$ (orthogonality)
    \item $P_\pi \mathbf{x}$ permutes the components of the vector $\mathbf{x}$
\end{enumerate}
\end{proposition}

\begin{proof}
(1) Let $(P_\pi)_{ij} = \delta_{i, \pi(j)}$. Then $(P_{\pi_1} P_{\pi_2})_{ij} = \sum_k (P_{\pi_1})_{ik}(P_{\pi_2})_{kj} = \sum_k \delta_{i, \pi_1(k)} \delta_{k, \pi_2(j)} = \delta_{i, \pi_1(\pi_2(j))} = (P_{\pi_1 \cdot \pi_2})_{ij}$.

(2) $(P_\pi^T)_{ij} = (P_\pi)_{ji} = \delta_{j, \pi(i)} = \delta_{\pi^{-1}(j), i} = (P_{\pi^{-1}})_{ij}$. Since $P_\pi P_{\pi^{-1}} = P_{\pi \cdot \pi^{-1}} = P_e = I$, we have $P_\pi^{-1} = P_{\pi^{-1}} = P_\pi^T$.

(3) $(P_\pi \mathbf{x})_i = \sum_j (P_\pi)_{ij} x_j = \sum_j \delta_{i,\pi(j)} x_j = x_{\pi^{-1}(i)}$, which is the $\pi^{-1}(i)$-th component of $\mathbf{x}$, that is, $P_\pi$ permutes the components according to $\pi$.
\end{proof}

\begin{theorem}[Cayley's Theorem for $S_n$]\label{thm:cayley_symmetric}
Every finite group $G$ of order $n$ is isomorphic to a subgroup of $S_n$. In particular, $S_n$ contains copies of all finite groups of order $\leq n!$.

See Section~\ref{sec:cayley} of \autoref{ch:teoremas_complementarios} for the proof.
\end{theorem}

\textbf{Implications for neural architectures.} Cayley's theorem implies that any finite discrete symmetry can be encoded as a subgroup of permutations, justifying the general approach of studying architectures invariant/equivariant to $S_n$.

\subsubsection{Foundations of Learning over Sets}

Machine learning over sets requires functions that do not depend on the arbitrary order in which the elements are presented.

\begin{definition}[Permutation-Invariant Function]\label{def:permutation_invariant_function}
A function $f: \mathbb{R}^{n \times d} \to \mathbb{R}^k$ is \textit{permutation-invariant} if for every permutation $\pi \in S_n$:
\begin{equation}
f(\pi \cdot X) = f(X)
\end{equation}
where $\pi \cdot X$ permutes the rows of the matrix $X \in \mathbb{R}^{n \times d}$.
\end{definition}

\begin{definition}[Permutation-Equivariant Function]\label{def:permutation_equivariant_function}
A function $f: \mathbb{R}^{n \times d} \to \mathbb{R}^{n \times k}$ is \textit{permutation-equivariant} if:
\begin{equation}
f(\pi \cdot X) = \pi \cdot f(X) \quad \forall \pi \in S_n
\end{equation}
\end{definition}

\begin{example}[Basic Invariant Functions]
The following functions are permutation-invariant:
\begin{align}
f_{\text{sum}}(X) &= \sum_{i=1}^n \mathbf{x}_i \\
f_{\text{max}}(X) &= \max_{i=1,\ldots,n} \mathbf{x}_i \quad \text{(elementwise)} \\
f_{\text{mean}}(X) &= \frac{1}{n}\sum_{i=1}^n \mathbf{x}_i
\end{align}
\end{example}

\subsubsection{Representation Theory of $S_n$ and Young Tableaux}

The representation theory of the symmetric group provides a deep understanding of how to construct equivariant layers for permutational data. Although modern architectures such as DeepSets use more elementary constructions, the underlying theory is fundamental for universal approximation results and advanced architecture design.

\paragraph{Partitions and Conjugacy Classes}

The structure of $S_n$ is intimately related to the partitions of the integer $n$.

\begin{definition}[Partition]\label{def:partition_integer}
A \textit{partition} of $n \in \mathbb{N}$ is a decreasing sequence $\lambda = (\lambda_1, \lambda_2, \ldots, \lambda_k)$ with $\lambda_1 \geq \lambda_2 \geq \cdots \geq \lambda_k > 0$ and $\sum_{i=1}^k \lambda_i = n$. We write $\lambda \vdash n$.
\end{definition}

\begin{example}[Partitions of $n=4$]
The partitions of 4 are: $(4)$, $(3,1)$, $(2,2)$, $(2,1,1)$, $(1,1,1,1)$, corresponding to the 5 types of conjugacy classes in $S_4$.
\end{example}

\begin{theorem}[Conjugacy Classes of $S_n$]\label{thm:conjugacy_classes_sn}
The conjugacy classes of $S_n$ are in bijection with the partitions of $n$. Two permutations are conjugate if and only if they have the same cycle structure.
\end{theorem}

\begin{proof}[Sketch]
Conjugation preserves cycle length: if $\pi = (i_1 \cdots i_k)$ is a $k$-cycle and $\sigma \in S_n$, then $\sigma \pi \sigma^{-1} = (\sigma(i_1) \cdots \sigma(i_k))$ is also a $k$-cycle. Therefore, the cycle type (partition formed by the cycle lengths) is invariant under conjugation. See \citep{sagan2001symmetric} for the complete proof.
\end{proof}

\paragraph{Young Tableaux and Irreducible Representations}

Young tableaux provide an explicit combinatorial construction of the irreducible representations of $S_n$.

\begin{definition}[Young Diagram]\label{def:young_diagram}
A \textit{Young diagram} of shape $\lambda \vdash n$ is an arrangement of $n$ boxes organized in left-justified rows, with $\lambda_i$ boxes in the $i$-th row. A \textit{Young tableau} is a Young diagram whose cells are filled with the numbers $\{1, 2, \ldots, n\}$.
\end{definition}

\begin{example}[Young Tableau for $\lambda = (3,2,1) \vdash 6$]
An example of a standard Young tableau:
\begin{equation*}
\ytableaushort{135,24,6}
\end{equation*}
where the first row contains 1, 3, 5; the second row 2, 4; and the third row 6.
\end{example}

\begin{definition}[Standard Young Tableau]\label{def:standard_young_tableau}
A \textit{standard Young tableau} (SYT) is a tableau where the numbers increase strictly from left to right in each row and from top to bottom in each column.
\end{definition}

\begin{theorem}[Construction of Irreducible Representations]\label{thm:irrep_construction_sn}
For each partition $\lambda \vdash n$, there exists an irreducible representation $V_\lambda$ of $S_n$ of dimension equal to the number of standard Young tableaux of shape $\lambda$. Every irreducible representation of $S_n$ is obtained in this way, and:
\begin{equation}
\sum_{\lambda \vdash n} (\dim V_\lambda)^2 = n!
\end{equation}
\end{theorem}

\begin{proof}[Idea of the construction]
For each Young tableau $T$ of shape $\lambda$, the \textit{polytabloid} $e_T = a_T b_T$ is defined, where $a_T$ antisymmetrizes over columns and $b_T$ symmetrizes over rows. The space $V_\lambda$ is the span of all polytabloids. The complete proof uses the group algebra $\mathbb{C}[S_n]$ and Schur's lemma; see \citep{fulton1997young} for details.
\end{proof}

\begin{example}[Dimensions for $n=3$]
For $S_3$, the partitions and dimensions are:
\begin{itemize}
    \item $\lambda = (3)$: Trivial representation, $\dim V_{(3)} = 1$
    \item $\lambda = (2,1)$: Standard representation, $\dim V_{(2,1)} = 2$
    \item $\lambda = (1,1,1)$: Sign representation, $\dim V_{(1,1,1)} = 1$
\end{itemize}
Verification: $1^2 + 2^2 + 1^2 = 6 = 3!$
\end{example}

\paragraph{Connection with Group Convolutions}

The representation theory of $S_n$ connects with neural architectures via the harmonic analysis on finite groups developed in Section~\ref{subsec:gcnns_grupos_finitos}.

\textbf{Convolution Theorem for $S_n$.}\label{rem:convolution_sn} By the Peter-Weyl theorem for finite groups, every function $f: S_n \to \mathbb{R}^d$ decomposes as:
\begin{equation}
f = \bigoplus_{\lambda \vdash n} f_\lambda
\end{equation}
where $f_\lambda$ lives in the isotypic space corresponding to the irreducible representation $V_\lambda$. A convolutional layer over $S_n$ acts independently in each isotypic space.

\textbf{Practical Limitations.} Although mathematically elegant, the explicit construction of convolutions over $S_n$ via irreducible representations is computationally intractable for large $n$ ($|S_n| = n!$ grows rapidly). Practical architectures such as DeepSets avoid working directly with $S_n$, operating instead over the space of sets via invariant/equivariant functions constructed elementarily.

\subsubsection{Universal Approximation Theorem: Deep Sets}

The fundamental result for permutation-invariant architectures is the DeepSets theorem of Zaheer et al.~\citep{zaheer2017deepsets}, which completely characterizes the form of continuous permutation-invariant functions.

\begin{theorem}[DeepSets Theorem - Universal Approximation]\label{thm:deepsets_universal}
Let $f: \mathcal{X}^n \to \mathcal{Y}$ be a continuous permutation-invariant function, where $\mathcal{X} \subset \mathbb{R}^d$ and $\mathcal{Y} \subset \mathbb{R}^k$ are compact. Then, for any $\epsilon > 0$, there exist continuous functions $\phi: \mathcal{X} \to \mathbb{R}^m$ and $\rho: \mathbb{R}^m \to \mathcal{Y}$ (with $m$ sufficiently large) such that:
\begin{equation}
\sup_{X \in \mathcal{X}^n} \left\|f(X) - \rho\left(\sum_{i=1}^n \phi(x_i)\right)\right\| < \epsilon
\label{eq:deepsets_representation}
\end{equation}

Moreover, both $\phi$ and $\rho$ can be implemented via neural networks.
\end{theorem}

\begin{remark}
\citet{zaheer2017deepsets} also prove an exact representation $f(X) = \rho(\sum_i \phi(x_i))$ for sets from a countable universe, but in that case the functions $\phi$ and $\rho$ may necessarily be discontinuous. \citet{wagstaff2019limitations} showed that, for the exact representation with continuous functions, the latent dimension $m$ must satisfy $m \geq n$ (the size of the input set).
\end{remark}

\begin{proof}[Complete proof]
The proof proceeds in three fundamental steps.

\textbf{Step 1: Reduction to functions over measures.}
We identify an unordered set $\{x_1, \ldots, x_n\} \subset \mathcal{X}$ with the empirical measure:
\begin{equation}
\mu_X = \frac{1}{n}\sum_{i=1}^n \delta_{x_i}
\end{equation}
where $\delta_x$ is the Dirac delta at $x$. Then $f$ induces a function $\tilde{f}$ over discrete probability measures with finite support in $\mathcal{X}$.

\textbf{Step 2: Characterization via moments.}
By the Weierstrass approximation theorem on spaces of measures, every continuous function over probability measures can be approximated arbitrarily well via functionals of the form:
\begin{equation}
\tilde{f}(\mu) \approx F\left(\int_\mathcal{X} g_1(x)\, d\mu(x), \ldots, \int_\mathcal{X} g_M(x)\, d\mu(x)\right)
\end{equation}
for basis functions $g_1, \ldots, g_M: \mathcal{X} \to \mathbb{R}$ and continuous function $F: \mathbb{R}^M \to \mathcal{Y}$.

For our empirical measure:
\begin{equation}
\int_\mathcal{X} g_j(x)\, d\mu_X(x) = \frac{1}{n}\sum_{i=1}^n g_j(x_i)
\end{equation}

\textbf{Step 3: Construction of $\phi$ and $\rho$.}
We define $\phi: \mathcal{X} \to \mathbb{R}^M$ as:
\begin{equation}
\phi(x) = (g_1(x), g_2(x), \ldots, g_M(x))
\end{equation}

Then:
\begin{equation}
\sum_{i=1}^n \phi(x_i) = \sum_{i=1}^n (g_1(x_i), \ldots, g_M(x_i)) = \left(\sum_{i=1}^n g_1(x_i), \ldots, \sum_{i=1}^n g_M(x_i)\right)
\end{equation}

We define $\rho: \mathbb{R}^M \to \mathcal{Y}$ as:
\begin{equation}
\rho(z_1, \ldots, z_M) = F\left(\frac{z_1}{n}, \ldots, \frac{z_M}{n}\right)
\end{equation}

Then:
\begin{align}
\rho\left(\sum_{i=1}^n \phi(x_i)\right) &= F\left(\frac{1}{n}\sum_{i=1}^n g_1(x_i), \ldots, \frac{1}{n}\sum_{i=1}^n g_M(x_i)\right) \\
&= F\left(\int_\mathcal{X} g_1\, d\mu_X, \ldots, \int_\mathcal{X} g_M\, d\mu_X\right) \\
&\approx \tilde{f}(\mu_X) = f(\{x_1, \ldots, x_n\})
\end{align}

By the universal approximation theorem for neural networks, both $\phi$ and $\rho$ can be approximated arbitrarily well via neural networks, completing the proof.
\end{proof}

\textbf{Uniqueness of the sum.}\label{rem:uniqueness_sum} The sum is the only aggregation operation that guarantees universal approximation for all continuous permutation-invariant functions. Other aggregators such as $\max$ or $\text{mean}$ have limited expressive power. For example, $\max$ cannot distinguish between $\{(1,0), (0,1)\}$ and $\{(1,1)\}$ after applying $\phi(x) = x$.

\textbf{Extension to variable-size sets.} For variable-size sets, the representation extends naturally by dividing by $n$ in the aggregation, or by using architectures with weighted attention as we will see in Set Transformers.

\subsubsection{PointNet: Architecture and Mathematical Analysis}

PointNet~\citep{qi2017pointnet} revolutionized the processing of 3D point clouds via an architecture explicitly designed for exact permutation invariance. While the DeepSets theorem provides the theoretical basis, PointNet demonstrates its practical viability for three-dimensional geometric data.

\paragraph{PointNet Architecture}

PointNet processes a point cloud $\mathcal{P} = \{\mathbf{p}_1, \ldots, \mathbf{p}_n\} \subset \mathbb{R}^3$ via three main components:

\begin{definition}[PointNet Architecture]\label{def:pointnet_architecture}
The PointNet architecture consists of:
\begin{enumerate}
    \item \textbf{Input transformation}: $T$-Net that predicts an affine transformation $T \in \mathbb{R}^{3 \times 3}$ to normalize the input
    \item \textbf{Per-point feature encoder}: Shared MLP $\phi_\theta: \mathbb{R}^3 \to \mathbb{R}^{1024}$
    \item \textbf{Global symmetric pooling}: Aggregation via max-pooling
    \item \textbf{Global decoder}: MLP $\rho_\psi: \mathbb{R}^{1024} \to \mathbb{R}^k$ for the final task
\end{enumerate}
\end{definition}

Mathematically, the implemented function is:
\begin{equation}
f_{\text{PointNet}}(\mathcal{P}) = \rho_\psi\left(\max_{i=1,\ldots,n} \phi_\theta(T \mathbf{p}_i)\right)
\label{eq:pointnet_formula}
\end{equation}

\paragraph{Analysis of Equivariance and Invariance}

\begin{proposition}[Invariance of PointNet]\label{prop:pointnet_invariance}
The function $f_{\text{PointNet}}$ defined in~\eqref{eq:pointnet_formula} is exactly invariant to permutations of the input points.
\end{proposition}

\begin{proof}
Let $\pi \in S_n$ be a permutation. Then:
\begin{align}
f_{\text{PointNet}}(\pi \cdot \mathcal{P}) &= \rho_\psi\left(\max_{i=1,\ldots,n} \phi_\theta(T \mathbf{p}_{\pi(i)})\right) \\
&= \rho_\psi\left(\max_{j=1,\ldots,n} \phi_\theta(T \mathbf{p}_j)\right) \quad \text{(reindexing)} \\
&= f_{\text{PointNet}}(\mathcal{P})
\end{align}
where we use that $\max$ is a symmetric operation: $\max\{a_{\pi(1)}, \ldots, a_{\pi(n)}\} = \max\{a_1, \ldots, a_n\}$.
\end{proof}

\textbf{Max vs Sum.} Unlike the DeepSets theorem, which uses sum, PointNet uses max-pooling. Although max has lower theoretical expressive power (as noted in \autoref{rem:uniqueness_sum}), in practice it works well for point clouds where extremal geometry is informative. PointNet++ extends this with local hierarchies.

\paragraph{Spatial Transformation and Approximate Equivariance}

PointNet's T-Net predicts a transformation that normalizes the point cloud, providing approximate invariance to Euclidean transformations.

\begin{definition}[Spatial Transformer Network for Points]\label{def:stn_points}
The T-Net is a sub-network that predicts a transformation matrix:
\begin{equation}
T = \text{T-Net}(\mathcal{P}) = \rho_{\text{T}}\left(\max_{i=1,\ldots,n} \phi_{\text{T}}(\mathbf{p}_i)\right) \in \mathbb{R}^{3 \times 3}
\end{equation}
where $\rho_{\text{T}}$ is regularized to produce approximately orthogonal transformations.
\end{definition}

\begin{proposition}[Approximation to SE(3) Invariance]
If $T$ learns the inverse transformation $g^{-1}$ for each transformed input $g \cdot \mathcal{P}$, then PointNet is approximately $SE(3)$-invariant. However, this is not theoretically guaranteed and depends on training.
\end{proposition}

\begin{proof}[Proof sketch]
Suppose the T-Net learns a transformation $T$ such that $T(g \cdot \mathcal{P}) = g^{-1}$ for all $g \in SE(3)$. Then the composition $T(g \cdot \mathcal{P}) \cdot (g \cdot \mathcal{P}) = g^{-1} g \cdot \mathcal{P} = \mathcal{P}$, that is, the input is ``canonicalized'' independently of the transformation $g$ applied. Since $\phi \circ T$ processes the canonicalized input, the output is invariant: $\phi(T(g \cdot \mathcal{P}) \cdot g \cdot \mathcal{P}) = \phi(\mathcal{P})$. However, the condition $T(g \cdot \mathcal{P}) = g^{-1}$ is a hypothesis about the learning capacity of the T-Net that is not guaranteed: the network must learn to invert arbitrary $SE(3)$ transformations, which is a non-trivial regression problem and is only achieved approximately.
\end{proof}

\paragraph{Critical Function and Theoretical Analysis}

A key result of the analysis of PointNet is the concept of the critical set.

\begin{definition}[Critical Set]\label{def:critical_set_pointnet}
The \textit{critical set} $\mathcal{C} \subset \mathcal{P}$ is the minimal subset of points that completely determines the max-pooling:
\begin{equation}
\mathcal{C} = \{\mathbf{p}_i \in \mathcal{P} \mid \exists k: (\phi_\theta(\mathbf{p}_i))_k = \max_{j}(\phi_\theta(\mathbf{p}_j))_k\}
\end{equation}
\end{definition}

\begin{proposition}[Robustness via the Critical Set]
PointNet is robust to perturbations and irregular sampling because the critical set typically contains points representative of the global structure of the object.
\end{proposition}

\begin{proof}
By the definition of the critical set (\autoref{def:critical_set_pointnet}), the global max-pooling produces $\mathbf{g} = \max_j \phi_\theta(\mathbf{p}_j)$ (componentwise). Adding or perturbing a point $\mathbf{p}_i \notin \mathcal{C}$ does not alter the result of the max-pooling, since $(\phi_\theta(\mathbf{p}_i))_k < \max_j (\phi_\theta(\mathbf{p}_j))_k$ for all $k$. Only the points in $\mathcal{C}$ contribute to the global vector, and $|\mathcal{C}| \leq K$ (dimension of the pooling layer). Thus, PointNet is robust to perturbations of up to $|\mathcal{P}| - |\mathcal{C}|$ non-critical points. The set $\mathcal{C}$ captures the ``skeletal structure'' of the object, providing natural robustness to subsampling.
\end{proof}

\paragraph{Limitations of PointNet}

Despite its success, PointNet has fundamental limitations:

\begin{itemize}
    \item \textbf{Does not capture local structure}: It treats each point independently until the global pooling
    \item \textbf{Loss of geometric information}: Max-pooling discards local spatial relations
    \item \textbf{Sensitivity to density}: Performance depends on the sampling density
\end{itemize}

These limitations motivated PointNet++~\citep{qi2017pointnetplus}, which introduces hierarchical local processing while preserving permutation invariance.

\subsubsection{Set Transformers and Permutation-Equivariant Attention}

A fundamental limitation of PointNet and DeepSets is that the global aggregation (sum or max) discards information about relations between elements of the set. Set Transformers introduce attention mechanisms that allow element-element interactions while preserving permutation invariance.

\paragraph{Foundations of Set Transformers}

Attention mechanisms, which will be studied in depth in \autoref{ch:attention}, can be adapted to operate over sets in a permutation-equivariant manner.

\begin{definition}[Multihead Attention over Sets]\label{def:set_multihead_attention}
For a set of embeddings $X = \{\mathbf{x}_1, \ldots, \mathbf{x}_n\} \in \mathbb{R}^{n \times d}$, multi-head attention is:
\begin{equation}
\text{MAB}(X) = \text{Concat}(\text{head}_1, \ldots, \text{head}_h)W^O
\end{equation}
where each head is:
\begin{equation}
\text{head}_i = \text{Attention}(XW_i^Q, XW_i^K, XW_i^V)
\end{equation}
and standard attention is:
\begin{equation}
\text{Attention}(Q, K, V) = \text{softmax}\left(\frac{QK^T}{\sqrt{d_k}}\right)V
\end{equation}
\end{definition}

\begin{proposition}[Equivariance of Attention]
The MAB operation is permutation-equivariant: for every $\pi \in S_n$,
\begin{equation}
\text{MAB}(\pi \cdot X) = \pi \cdot \text{MAB}(X)
\end{equation}
\end{proposition}

\begin{proof}
The matrix multiplication $QK^T$ and the subsequent multiplication by $V$ permute rows consistently with the permutation of $X$. Softmax operates independently on each row, preserving the permutation. Therefore, the composition is equivariant.
\end{proof}

\paragraph{Induced Set Attention Block (ISAB)}

To improve computational efficiency, Set Transformers introduce induction points that reduce the complexity from $O(n^2)$ to $O(nm)$ where $m \ll n$.

\begin{definition}[ISAB]\label{def:isab}
An Induced Set Attention Block uses $m$ learnable induction points $I \in \mathbb{R}^{m \times d}$:
\begin{equation}
\text{ISAB}_m(X) = \text{MAB}(X, \text{MAB}(I, X))
\end{equation}
where $\text{MAB}(A, B)$ denotes attention with queries from $A$ and keys/values from $B$.
\end{definition}

\paragraph{Pooling by Multihead Attention (PMA)}

To obtain fixed-size outputs invariant to permutations, PMA uses learnable seeds.

\begin{definition}[PMA]\label{def:pma}
Pooling by Multihead Attention with $k$ seeds $S \in \mathbb{R}^{k \times d}$:
\begin{equation}
\text{PMA}_k(X) = \text{MAB}(S, \text{rFF}(X))
\end{equation}
where rFF is a row-wise feedforward. The result is invariant to permutations of $X$.
\end{definition}

\textbf{Connection with \autoref{ch:attention}.} Set Transformers are a special case of Transformers applied to sets. In \autoref{ch:attention} we will delve into the theory of attention, including geometric and equivariant variants for non-Euclidean data. The mechanisms introduced here are the permutation-equivariant version of standard attention.

\subsubsection{Numerical Example: Processing a Set with DeepSets}

To make the theory concrete, we develop a complete numerical example showing how a DeepSets layer processes a small set.

\begin{example}[DeepSets in Action]\label{ex:deepsets_numerical}
Consider a set of 3 two-dimensional vectors:
\begin{equation}
\mathcal{X} = \{(2, 1), (1, 3), (3, 2)\}
\end{equation}

\textbf{Step 1: Per-element encoder $\phi$.} We use a simple linear transformation $\phi(\mathbf{x}) = W\mathbf{x} + \mathbf{b}$ with:
\begin{equation}
W = \begin{pmatrix} 1 & -1 \\ 0 & 2 \\ 1 & 1 \end{pmatrix}, \quad \mathbf{b} = \begin{pmatrix} 0 \\ 1 \\ 0 \end{pmatrix}
\end{equation}

We apply $\phi$ to each element:
\begin{align}
\phi((2,1)) &= \begin{pmatrix} 1 & -1 \\ 0 & 2 \\ 1 & 1 \end{pmatrix} \begin{pmatrix} 2 \\ 1 \end{pmatrix} + \begin{pmatrix} 0 \\ 1 \\ 0 \end{pmatrix} = \begin{pmatrix} 1 \\ 3 \\ 3 \end{pmatrix} \\
\phi((1,3)) &= \begin{pmatrix} 1 & -1 \\ 0 & 2 \\ 1 & 1 \end{pmatrix} \begin{pmatrix} 1 \\ 3 \end{pmatrix} + \begin{pmatrix} 0 \\ 1 \\ 0 \end{pmatrix} = \begin{pmatrix} -2 \\ 7 \\ 4 \end{pmatrix} \\
\phi((3,2)) &= \begin{pmatrix} 1 & -1 \\ 0 & 2 \\ 1 & 1 \end{pmatrix} \begin{pmatrix} 3 \\ 2 \end{pmatrix} + \begin{pmatrix} 0 \\ 1 \\ 0 \end{pmatrix} = \begin{pmatrix} 1 \\ 5 \\ 5 \end{pmatrix}
\end{align}

\textbf{Step 2: Symmetric aggregation.} We sum the embeddings:
\begin{equation}
\sum_{i=1}^3 \phi(\mathbf{x}_i) = \begin{pmatrix} 1 \\ 3 \\ 3 \end{pmatrix} + \begin{pmatrix} -2 \\ 7 \\ 4 \end{pmatrix} + \begin{pmatrix} 1 \\ 5 \\ 5 \end{pmatrix} = \begin{pmatrix} 0 \\ 15 \\ 12 \end{pmatrix}
\end{equation}

\textbf{Step 3: Global decoder $\rho$.} We apply a final transformation $\rho(\mathbf{z}) = V\mathbf{z}$ with:
\begin{equation}
V = \begin{pmatrix} 1 & 0 & 0 \\ 0 & 0.1 & 0 \end{pmatrix}
\end{equation}

Final result:
\begin{equation}
f(\mathcal{X}) = \rho\left(\sum_{i=1}^3 \phi(\mathbf{x}_i)\right) = \begin{pmatrix} 1 & 0 & 0 \\ 0 & 0.1 & 0 \end{pmatrix} \begin{pmatrix} 0 \\ 15 \\ 12 \end{pmatrix} = \begin{pmatrix} 0 \\ 1.5 \end{pmatrix}
\end{equation}

\textbf{Invariance verification.} Let us permute the set as $\mathcal{X}' = \{(3,2), (2,1), (1,3)\}$ (permutation $(1\,3)$):

The individual representations are the same (in a different order), and the sum is commutative:
\begin{equation}
\sum \phi(\mathcal{X}') = \begin{pmatrix} 1 \\ 5 \\ 5 \end{pmatrix} + \begin{pmatrix} 1 \\ 3 \\ 3 \end{pmatrix} + \begin{pmatrix} -2 \\ 7 \\ 4 \end{pmatrix} = \begin{pmatrix} 0 \\ 15 \\ 12 \end{pmatrix}
\end{equation}

Therefore, $f(\mathcal{X}') = f(\mathcal{X}) = (0, 1.5)^T$, verifying permutation invariance.
\end{example}

\subsubsection{Applications of Permutation-Invariant Architectures}

Permutation-invariant and equivariant architectures have transformed multiple domains where the data lack a natural intrinsic order.

\paragraph{3D Point Cloud Processing}

3D point clouds obtained via LiDAR, laser scanning or photogrammetric reconstruction are unordered sets of spatial positions. PointNet and its variants have revolutionized:

\begin{itemize}
    \item \textbf{3D object classification}: ModelNet40 (98.4\% accuracy with PointNet++)
    \item \textbf{Semantic segmentation}: Segmentation of urban point clouds (ScanNet, S3DIS)
    \item \textbf{Object detection}: 3D detection for autonomous vehicles (KITTI)
    \item \textbf{Scene reconstruction}: Understanding of complete 3D geometry
\end{itemize}

Permutation invariance allows these architectures to process point clouds of variable density and irregular sampling without preprocessing.

\paragraph{Molecular Modeling and Drug Discovery}

Molecules are naturally graphs with invariance to atom permutations. Permutation-equivariant architectures enable:

\begin{itemize}
    \item \textbf{Molecular property prediction}: Solubility, toxicity, biological activity
    \item \textbf{Molecule generation}: Design of new pharmacological compounds
    \item \textbf{Conformation prediction}: Energetic states of flexible molecules
    \item \textbf{Retrosynthetic synthesis}: Planning of chemical synthesis routes
\end{itemize}

Invariant representations capture chemical properties independently of the arbitrary labeling of atoms.

\paragraph{High-Energy Particle Physics}

Particle collision events in accelerators such as the LHC produce variable sets of detected particles. Permutation-invariant networks process:

\begin{itemize}
    \item \textbf{Jet identification}: Classification of hadronic jets
    \item \textbf{Anomaly detection}: Search for new physics beyond the Standard Model
    \item \textbf{Event reconstruction}: Identification of initial particles
\end{itemize}

Permutation invariance is physical: the order of detection of particles is arbitrary.

\paragraph{Reasoning over Sets and Logic}

Problems of abstract reasoning over sets benefit from DeepSets architectures:

\begin{itemize}
    \item \textbf{Question answering over sets}: ``How many blue elements are there?''
    \item \textbf{Anomaly detection}: Identification of outliers in populations
    \item \textbf{Aggregation tasks}: Statistics over empirical distributions
    \item \textbf{Combinatorial optimization}: NP-hard problems over sets (TSP, bin packing)
\end{itemize}

\paragraph{Processing of Graphs as Permutations}

Although graph networks will be studied in depth in later chapters, it is important to note that invariance to node permutations is fundamental for \glspl{GNN}. Permutation-equivariant architectures provide the theoretical basis for:

\begin{itemize}
    \item \textbf{Message-passing neural networks}: Message passing invariant to relabeling
    \item \textbf{Graph classification}: Prediction over complete graphs
    \item \textbf{Node embeddings}: Representations independent of node order
\end{itemize}

\subsubsection{Synthesis: The Landscape of Permutation-Symmetric Architectures}

The theory of permutation groups provides the unifying framework for understanding architectures over unordered data. The theoretical progression presented illustrates:

\begin{enumerate}
    \item \textbf{Algebraic foundations}: $S_n$, irreducible representations, Young theory
    \item \textbf{Approximation theory}: The DeepSets theorem characterizes universal functions
    \item \textbf{Practical implementations}: PointNet, Set Transformers, hierarchical variants
    \item \textbf{Extensions with attention}: Capturing relations via permutation-equivariant attention
    \item \textbf{Transformative applications}: From point clouds to particle physics
\end{enumerate}

The deep understanding of permutation symmetries not only underpins existing architectures, but also provides design principles for new architectures that respect the intrinsic combinatorial structure of the data. This perspective extends naturally to irregular graphs (later chapters) and to multi-modal data where different components have different symmetries.

The transition from discrete finite groups ($C_4$, $D_4$) to continuous groups ($SO(2)$, $SO(3)$) and finally permutation groups illustrates the versatility of the group convolution framework, demonstrating how unified mathematical principles underlie the diversity of architectures in geometric deep learning.

\subsection{Comparative Synthesis and Unifying Framework}\label{subsec:sintesis_comparativa_gcnn}

We now synthesize the equivariant convolution approaches presented, from discrete finite groups to complex continuous groups, into a unifying comparative framework.

\subsubsection{Comparative Framework of Equivariant Architectures}\label{subsubsec:marco_comparativo_arquitecturas}

The choice of symmetry group determines both the representational capabilities and the computational complexity:

\begin{table}[h]
\centering
\small
\begin{tabular}{|l|c|c|c|c|c|}
\hline
\textbf{Group} & \textbf{Order} & \textbf{Type} & \textbf{Complexity} & \textbf{Memory} & \textbf{Accuracy} \\
\hline
Standard CNN & - & Translational & $O(N^2)$ & 1x & Baseline \\
$C_4$ & 4 & Finite & $O(4N^2)$ & 4x & +15-25\% \\
$D_4$ & 8 & Finite & $O(8N^2)$ & 8x & +20-30\% \\
$SO(2)$ Discrete & $L$ & Discretized & $O(LN^2)$ & $L$x & +25-35\% \\
$SO(2)$ Steerable & $\infty$ & Continuous & $O(F^2N^2)$ & $F^2$x & +30-40\% \\
$SO(3)$ Spherical & $\infty$ & Continuous & $O(L^3N^2)$ & $L^3$x & +35-50\% \\
$SE(3)$ & $\infty$ & Non-compact & $O(L^4N^2)$ & $L^4$x & +40-60\% \\
\hline
\end{tabular}
\caption{Comparison of complexity and performance for different groups. $N$: spatial resolution, $L$: maximum degree of spherical representations, $F$: maximum frequency in harmonic representations.}
\label{tab:group_complexity_comparison}
\end{table}

\subsubsection{Empirical Evaluation}\label{subsubsec:framework_evaluacion_empirica}

The evaluation of equivariant architectures requires measuring three complementary dimensions: the equivariance error (how faithfully the symmetry is preserved), the accuracy on the target task, and the computational efficiency. The following code illustrates a benchmarking framework that integrates these metrics:

\begin{listing}[h]
\begin{lstlisting}[language=Python]
class EquivarianceBenchmark:
    """Unified framework for equivariance evaluation."""

    def __init__(self, model, group, n_transforms=20):
        self.model = model
        self.group = group
        self.transforms = group.sample_elements(n_transforms)

    def measure_equivariance_error(self, x: torch.Tensor) -> float:
        """Measures mean squared equivariance error."""
        self.model.eval()
        with torch.no_grad():
            y = self.model(x)
            errors = []
            for g in self.transforms:
                gx = self.group.apply(g, x)
                gy_expected = self.group.apply_output(g, y)
                errors.append(torch.mean((self.model(gx) - gy_expected)**2).item())
        return np.mean(errors)

    def benchmark(self, test_loader, input_shape):
        """Complete benchmark: equivariance + efficiency + accuracy."""
        sample = next(iter(test_loader))[0][:1]
        eq_error = self.measure_equivariance_error(sample)

        # Efficiency
        x = torch.randn(*input_shape)
        t0 = time.time()
        for _ in range(100):
            self.model(x)
        fwd_ms = (time.time() - t0) / 100 * 1000
        n_params = sum(p.numel() for p in self.model.parameters())

        return {'eq_error': eq_error, 'fwd_ms': fwd_ms, 'params': n_params}
\end{lstlisting}
\caption{Benchmarking framework for equivariance}
\label{code:equivariance_benchmark}
\end{listing}

Applying this framework to standard implementations:

\begin{table}[h]
\centering
\footnotesize
\begin{tabular}{|l|c|c|c|c|c|c|}
\hline
\textbf{Architecture} & \textbf{CIFAR-10} & \textbf{Rot-MNIST} & \textbf{Params} & \textbf{Time} & \textbf{Memory} & \textbf{Equiv. Error} \\
 & \textbf{(\%)} & \textbf{(\%)} & \textbf{(M)} & \textbf{(ms)} & \textbf{(MB)} & \textbf{($10^{-6}$)} \\
\hline
Standard CNN & 92.3 & 67.1 & 1.2 & 2.1 & 45 & - \\
G-CNN (C4) & 93.8 & 91.5 & 1.1 & 3.2 & 52 & 3.2 \\
G-CNN (D4) & 94.2 & 93.1 & 1.0 & 4.1 & 58 & 4.7 \\
SO(2) Steerable & 95.1 & 96.8 & 0.9 & 5.8 & 67 & 2.1 \\
Spherical CNN & 94.8 & 97.2 & 1.4 & 12.3 & 128 & 1.8 \\
SE(3) Transformer & 95.6 & 98.1 & 2.1 & 18.7 & 256 & 1.2 \\
\hline
\end{tabular}
\caption{Comparative results on standard benchmarks}
\label{tab:comprehensive_benchmark_results}
\end{table}

\subsubsection{Practical Recommendations}\label{subsubsec:tradeoffs_recomendaciones}

The empirical results reveal a fundamental trade-off between equivariance accuracy and computational efficiency. The architectures on the Pareto frontier (those not simultaneously dominated in accuracy, efficiency and equivariance error) suggest the following guidelines:

\begin{enumerate}
    \item \textbf{Real-time production}: Discrete G-CNNs ($C_4$, $D_4$) offer the best balance between equivariance improvements and computational overhead.
    \item \textbf{High accuracy}: Steerable CNNs and SE(3) architectures justify their higher computational cost with superior accuracy.
    \item \textbf{Unknown symmetries}: Start with simple G-CNNs to validate the presence of symmetries before adopting more complex architectures.
    \item \textbf{Iterative development}: Standard CNN as baseline $\to$ discrete G-CNN $\to$ continuous extensions according to benchmarking results.
\end{enumerate}

\subsubsection{Future Perspectives}\label{subsubsec:perspectivas_desarrollo_futuro}

The open challenges include: approximate algorithms that preserve equivariance guarantees with lower complexity, adaptive representations that dynamically adjust the group according to the data, extension to temporal and stochastic symmetries, and integration of symmetries at multiple scales simultaneously.

The \glspl{convolucion_grupo} developed so far represent a significant advance over classical CNNs, but they still operate under the premise of exact algebraic symmetries. The next natural step in the hierarchy of Geometric Deep Learning is to address data residing in general geometric spaces where the symmetries are local or approximate, leading naturally to the study of convolutions on Riemannian manifolds.
\section{Convolutions on Manifolds}\label{sec:manifold_cnns}

While group convolutions exploit specific algebraic symmetries, convolutions on manifolds address the fundamental challenge of extending convolutional operations to general geometric spaces. This section examines how the concepts of differential geometry provide the tools necessary to generalize \glspl{CNN} beyond regular grids, preserving the intrinsic geometric structure of the data.

Riemannian manifolds appear naturally in numerous applications: 3D surfaces in computer graphics, configuration spaces in robotics, and high-dimensional data structures in machine learning. The ability to process data directly on these spaces, without resorting to Euclidean approximations that can introduce distortions, represents a significant conceptual advance.

\subsection{Geodesic CNNs}\label{subsec:geodesic_cnns}

Group convolutions provide an elegant framework for extending classical \glspl{CNN} while maintaining specific symmetries such as rotations and translations. However, much real-world data resides in spaces that do not admit a natural group structure, but rather possess a more general geometry characterized by local curvatures and varying metrics. This observation motivates the extension towards \textbf{Riemannian manifolds}, where geodesic CNNs emerge as the first systematic generalization of convolutional operations to non-Euclidean spaces.

\subsubsection{Transition: From Groups to Manifolds}

\paragraph{Limitations of Group Convolutions}

Although group convolutions elegantly capture discrete and continuous symmetries, they face fundamental limitations when the data does not exhibit clear group symmetries:

\begin{enumerate}
    \item \textbf{Variable local geometry}: Data on curved surfaces does not admit globally consistent translations
    \item \textbf{Intrinsic curvature}: The local geometry varies from point to point, invalidating uniform convolution kernels
    \item \textbf{Natural coordinates}: Many domains (surfaces, manifolds) lack canonical global coordinate systems
\end{enumerate}

\begin{example}[Curved Surfaces]
Consider the analysis of 3D shapes represented as triangulated surfaces. While a sphere admits rotations as global symmetries, an arbitrary surface such as a facial model lacks evident symmetries. The local curvature varies dramatically between different regions (convex nose vs. eye sockets), requiring local geometric adaptation.
\end{example}

\begin{example}[Construction of Geodesic Coordinates on Facial Surfaces]
To illustrate the theoretical construction of geodesic coordinates, let us consider the mathematical analysis of a 3D facial model represented as a Riemannian manifold $(M, g)$:

\textbf{Geometric Configuration}:
The facial surface is characterized by:
\begin{itemize}
    \item Genus-0 topology with approximately $10^4$ discrete points
    \item Variable curvature: nose ($K > 0$), eye sockets ($K < 0$), forehead ($K \approx 0$)
    \item Riemannian metric induced by the embedding in $\mathbb{R}^3$
\end{itemize}

\textbf{Mathematical Construction of the Coordinate System}:

\begin{algorithm}[H]
\caption{Rigorous Construction of Geodesic Charts}
\begin{algorithmic}[1]
    \State Given a center point $p \in M$, determine the injectivity radius $r_{\text{inj}}(p)$
    \State Construct an orthonormal basis $\{e_1, e_2\}$ in $T_p M$ using principal curvature directions
    \State Define the exponential map $\exp_p: B_r(0) \subset T_p M \to M$ with $r < r_{\text{inj}}(p)$
    \State For polar coordinates $(\rho, \theta) \in [0,r] \times [0,2\pi)$:
    \State \quad Map via $q = \exp_p(\rho(\cos\theta \, e_1 + \sin\theta \, e_2))$
    \State Correct by the Jacobian factor $J_p(\rho,\theta) = \sqrt{\det g_{\exp_p(\rho e_\theta)}}$
\end{algorithmic}
\end{algorithm}

\textbf{Fundamental Geometric Properties}:
\begin{itemize}
    \item \textbf{Positive curvature (nose)}: $\det(\text{d}\exp_p) < 1$, geodesic concentration
    \item \textbf{Negative curvature (sockets)}: $\det(\text{d}\exp_p) > 1$, geodesic divergence
    \item \textbf{Zero curvature (forehead)}: $\det(\text{d}\exp_p) \approx 1$, local Euclidean behavior
    \item \textbf{Automatic adaptation}: The metric $g$ determines the structure of the geodesic chart
\end{itemize}
\end{example}

\paragraph{Riemannian Manifolds as a Generalization}

A \textbf{Riemannian manifold} $(M, g)$ consists of a smooth manifold $M$ equipped with a Riemannian metric $g$ that defines inner products on each tangent space $T_p M$.

\begin{definition}[Riemannian Manifold]
A Riemannian manifold $(M, g)$ is a differentiable manifold $M$ together with a metric tensor $g$ that assigns to each point $p \in M$ an inner product $g_p: T_p M \times T_p M \to \mathbb{R}$ that varies smoothly with $p$.
\end{definition}

This structure allows defining fundamental geometric concepts:
\begin{itemize}
    \item \textbf{Geodesic distances}: $d_g(p,q) = \inf_{\gamma} \int_0^1 \sqrt{g_{\gamma(t)}(\dot{\gamma}(t), \dot{\gamma}(t))} dt$
    \item \textbf{Geodesic neighborhoods}: Regions locally isometric to Euclidean spaces
    \item \textbf{Parallel transport}: Mechanism for comparing vectors at different points
\end{itemize}

\subsubsection{Fundamental Architecture: Geodesic CNN}

The seminal work of \citet{masci2018geodesicconvolutionalneural} introduces \textbf{Geodesic CNNs} as the first rigorous extension of convolutional operations to general Riemannian manifolds.

\paragraph{Central Principle: Geodesic Coordinate Systems}

The fundamental idea consists of establishing local coordinate systems around each point of the manifold using \textbf{geodesic coordinates}.

\begin{definition}[Geodesic Coordinates]
Given a point $p \in M$ on a Riemannian manifold, the geodesic coordinates in a neighborhood $U_p$ are defined via the exponential map:
\begin{equation}
\exp_p: T_p M \supset B_\epsilon(0) \to U_p \subset M
\end{equation}
where $B_\epsilon(0)$ is a ball of radius $\epsilon$ in the tangent space $T_p M$.
\end{definition}

\paragraph{Construction of the Geodesic Chart}

For each point $p \in M$, a \textbf{geodesic chart} is constructed that allows defining local convolutional operations:

\begin{algorithm}[H]
\caption{Construction of a Geodesic Chart}
\begin{algorithmic}[1]
    \State \textbf{Input:} Point $p \in M$, geodesic radius $r > 0$
    \State Compute orthonormal basis $\{e_1, e_2\}$ in $T_p M$
    \State Define polar grid: $\{(r_i \cos \theta_j, r_i \sin \theta_j)\}$ in $T_p M$
    \State Map via $\exp_p$ to points $\{q_{ij}\}$ in $M$
    \State Interpolate feature values at the mapped points
    \State \textbf{Output:} Regular geodesic patch around $p$
\end{algorithmic}
\end{algorithm}

\begin{figure}[htbp]
	\centering
	\includegraphics[width=0.80\textwidth]{figures/geodesic_convolution}
	\caption[Geodesic convolution on a surface]{Geodesic convolution: the kernel is applied on the geodesic ball $B_r(p)$ around a point $p$ on the surface $\mathcal{M}$, parallel-transporting the features of the neighbors $q_j$ to the tangent space $T_p\mathcal{M}$.}
	\label{fig:geodesic_convolution}
\end{figure}

\paragraph{Geodesic Convolution Operation}

The convolution at point $p$ is defined as:

\begin{equation}
(f * k)_g(p) = \sum_{i,j} k_{ij} \cdot f(\exp_p(r_i e_{\theta_j})) \cdot J_p(r_i, \theta_j)
\label{eq:geodesic_convolution}
\end{equation}

where:
\begin{itemize}
    \item $k_{ij}$ are the learnable kernel weights
    \item $f$ is the feature function on the manifold
    \item $J_p(r_i, \theta_j)$ is the Jacobian correction factor
    \item $e_{\theta_j} = \cos\theta_j e_1 + \sin\theta_j e_2$ are angular directions
\end{itemize}

\subsubsection{Critical Technical Aspects}

\paragraph{Jacobian Correction}

The Riemannian metric introduces distortions that must be compensated. The Jacobian factor corrects the non-uniformity of the exponential map:

\begin{equation}
J_p(r, \theta) = \sqrt{\det g_{\exp_p(r e_\theta)}(\partial_r \exp_p, \partial_\theta \exp_p)}
\end{equation}

\paragraph{Choice of the Tangent Basis}

The construction requires selecting an orthonormal basis in each tangent space $T_p M$. The options include:
\begin{enumerate}
    \item \textbf{Principal directions}: Use eigenvectors of the curvature operator
    \item \textbf{Local consistency}: Minimize variation with respect to neighboring points
    \item \textbf{Parallel transport}: Propagate the basis from a reference point
\end{enumerate}

\subsubsection{Anisotropic Diffusion Descriptors}

Complementing geodesic CNNs, \citet{Boscaini2016AnisotropicDD} introduce \textbf{Anisotropic Diffusion Descriptors} (ADs) that overcome the limitations of isotropic spectral methods.

\paragraph{Motivation: Isotropic Limitations}

Classical spectral descriptors such as Heat Kernel Signatures (HKS) assume isotropic diffusion, losing crucial directional information in shapes with elongated structures or prominent directional features.

\paragraph{Mathematical Formulation}

ADs use an anisotropic diffusion operator parameterized by a conductivity tensor $C_p$:

\begin{equation}
\frac{\partial u}{\partial t} = -\text{div}(C \cdot \nabla u)
\end{equation}

where $C_p$ adapts the diffusion to the local geometry. The solution $u_t(p)$ for different times $t$ and directions generates the descriptor:

\begin{equation}
\text{AD}_p = \{u_t(p, \theta_k) : t \in T, \theta_k \in \Theta\}
\end{equation}

\paragraph{Computational Advantages}

\begin{enumerate}
    \item \textbf{Directional sensitivity}: Captures local orientations of geometric features
    \item \textbf{Robustness}: Stability under isometric deformations
    \item \textbf{Efficiency}: Computation via optimized spectral formulation
\end{enumerate}

\subsubsection{Implementation and Practical Considerations}

\paragraph{Discretization on Triangular Meshes}

In practice, manifolds are represented as triangular meshes. The implementation requires:

\begin{algorithm}[H]
\caption{Geodesic CNN on a Triangular Mesh}
\begin{algorithmic}[1]
    \State Compute the discrete Laplace-Beltrami operator $L$
    \State For each vertex $v_i$:
    \State \quad Estimate the local metric via triangle areas
    \State \quad Construct an orthonormal basis in the discrete tangent space
    \State \quad Compute approximate geodesic distances
    \State \quad Apply bilinear interpolation for regular sampling
    \State Propagate gradients respecting the discrete geometry
\end{algorithmic}
\end{algorithm}

\paragraph{Numerical Stability}

Critical aspects for robust implementations:
\begin{itemize}
    \item \textbf{Adaptive radius}: Adjust $r$ according to local curvature to avoid singularities
    \item \textbf{Geodesic regularization}: Smooth distortions in regions of high curvature
    \item \textbf{Boundary handling}: Special treatment of boundary conditions
\end{itemize}

\subsubsection{Limitations and Motivation for Extensions}

Geodesic CNNs, although revolutionary, present limitations that motivate later developments:

\paragraph{Dependence on Local Coordinates}
The construction of geodesic charts introduces dependencies on the choice of tangent bases, creating potential inconsistency between adjacent regions.

\paragraph{Computational Scalability}
The computation of geodesic coordinates and Jacobian corrections is computationally intensive, especially for high-resolution manifolds.

\paragraph{Limitation to Smooth Geometries}
The exponential map requires sufficient differentiability, limiting applicability to manifolds with singularities or discontinuities.

\paragraph{Lack of Globality}
By processing local patches independently, important global geometric correlations are potentially lost.

\textbf{Transition towards General Methods.} These limitations motivate the development of more general approaches that:
\begin{enumerate}
    \item Avoid explicit dependence on coordinate systems
    \item Incorporate global information efficiently
    \item Handle manifolds with complex or discontinuous geometries
    \item Provide theoretical guarantees of intrinsic equivariance
\end{enumerate}
This motivation sets the context for the methods that will be developed in the following subsections.

\subsection{General Methods: MoNet}\label{subsec:monet}

MoNet~\citep{DBLP:journals/corr/MontiBMRSB16} introduces a unifying framework that abstracts geometric relationships via \textbf{pseudo-local coordinates}, providing greater flexibility than approaches specific to each type of manifold.

\subsubsection{Theoretical Framework}

\begin{definition}[Pseudo-local Coordinates]\label{def:coordenadas_pseudo_locales}
Pseudo-local coordinates are functions $u: M \times M \to \mathbb{R}^d$ that encode local geometric relationships between pairs of points $(x, y) \in M \times M$, satisfying:
\begin{enumerate}
    \item \textbf{Locality}: $u(x, y) = 0$ if $x$ and $y$ are not in the same local neighborhood
    \item \textbf{Invariance}: $u$ is invariant under isometries of the manifold
    \item \textbf{Differentiability}: $u$ is smooth in the interior of its domain
\end{enumerate}
\end{definition}

\textbf{Fundamental examples of pseudo-local coordinates}:

\begin{enumerate}
    \item \textbf{Geodesic distance and angle}:
    \begin{equation}
    u_{\text{geod}}(x, y) = (d_g(x, y), \angle(x, y))
    \end{equation}
    where $\angle(x, y)$ is the polar angle in geodesic coordinates centered at $x$.

    \item \textbf{Local spectral coordinates}:
    \begin{equation}
    u_{\text{spec}}(x, y) = (\phi_1(x) - \phi_1(y), \phi_2(x) - \phi_2(y), \ldots, \phi_k(x) - \phi_k(y))
    \end{equation}
    using eigenfunctions of the Laplace-Beltrami operator.

    \item \textbf{Heat kernel encoding}:
    \begin{equation}
    u_{\text{heat}}(x, y) = (H_t(x, y))_{t \in T}
    \end{equation}
    where $H_t(x, y) = \sum_{i=0}^{\infty} e^{-\lambda_i t} \phi_i(x) \phi_i(y)$ is the heat kernel.
\end{enumerate}

\begin{definition}[Generalized MoNet Kernel]\label{def:kernel_monet}
MoNet kernels are defined as linear combinations of parametric basis functions:
\begin{equation}
w_\theta(u) = \sum_{j=1}^J w_j(\theta) \psi_j(u)
\end{equation}
where:
\begin{itemize}
    \item $w_j(\theta)$ are learnable weights dependent on parameters $\theta$
    \item $\psi_j(u)$ are basis functions (typically Gaussian)
    \item $u$ are pseudo-local coordinates
\end{itemize}

For the specific case of Gaussian mixtures:
\begin{equation}
w_\theta(u) = \sum_{j=1}^J \alpha_j(\theta) \exp\left(-\frac{1}{2}(u - \mu_j(\theta))^T \Sigma_j^{-1}(\theta) (u - \mu_j(\theta))\right)
\end{equation}
\end{definition}

\begin{theorem}[Universal Approximation of MoNet Kernels]\label{thm:universal_approximation_monet}
For any continuous function $k: \mathbb{R}^d \to \mathbb{R}$ with compact support, and for any $\epsilon > 0$, there exists a finite number $J$ and parameters $\{\alpha_j, \mu_j, \Sigma_j\}_{j=1}^J$ such that:
\begin{equation}
\sup_{u \in \text{supp}(k)} |k(u) - w_\theta(u)| < \epsilon
\end{equation}
\end{theorem}

\begin{proof}[Proof sketch]
The proof follows from the density of Gaussian mixtures in $L^2(\mathbb{R}^d)$ and the compactness of the support. The technical details involve:
\begin{enumerate}
    \item Approximation in $L^2$ norm by the Weierstrass approximation theorem
    \item Uniform convergence on compact sets
    \item Control of the approximation error via the number of Gaussian components
\end{enumerate}
\end{proof}

\textbf{MoNet convolution operation}: The convolution at a point $x \in M$ is defined as:
\begin{equation}
(f * w)(x) = \int_{\mathcal{N}(x)} w_\theta(u(x, y)) f(y) dV_g(y)
\end{equation}
where:
\begin{itemize}
    \item $\mathcal{N}(x)$ is a local neighborhood of $x$
    \item $dV_g$ is the Riemannian volume element
    \item $u(x, y)$ are the pseudo-local coordinates
\end{itemize}

\textbf{Discretization on meshes}: For discrete manifolds represented as meshes:
\begin{equation}
(f * w)(x_i) = \sum_{j \in \mathcal{N}(i)} w_\theta(u(x_i, x_j)) f(x_j) A_j
\end{equation}
where $A_j$ are local area weights (e.g., Voronoi areas).

\subsubsection{Unification of Architectures}

MoNet unifies multiple existing architectures via specific choices of coordinates:

\begin{itemize}
    \item \textbf{Geodesic CNNs}: $u(p,q) = (\rho(p,q), \theta(p,q))$
    \item \textbf{Graph CNNs}: $u(p,q) = (\frac{1}{\sqrt{d_p}}, \frac{1}{\sqrt{d_q}})$
    \item \textbf{SplineCNNs}: Coordinates based on B-splines
\end{itemize}

This unification allows transferring intuitions and techniques between different geometric domains.

\begin{example}[Unified Theoretical Framework of MoNet]

To illustrate the versatility of the MoNet framework, we analyze the mathematical formulation of adaptive kernels for three distinct geometric domains:

\textbf{General Mathematical Formulation}:

Let $u: M \times M \to \mathbb{R}^d$ be a pseudo-local coordinate function. MoNet kernels are defined as Gaussian mixtures in the coordinate space:

\begin{equation}
w_\theta(u) = \sum_{j=1}^J \alpha_j \exp\left(-\frac{1}{2}(u - \mu_j)^T \Sigma_j^{-1} (u - \mu_j)\right)
\end{equation}

where $\theta = \{\alpha_j, \mu_j, \Sigma_j\}_{j=1}^J$ are learnable parameters.

\textbf{Specializations by Geometric Domain}:

\begin{enumerate}
    \item \textbf{Riemannian Surfaces}:
    \begin{equation}
    u_{\text{geo}}(p,q) = (\rho_g(p,q), \theta_p(q))
    \end{equation}
    where $\rho_g$ is the geodesic distance and $\theta_p(q)$ is the angle in geodesic polar coordinates centered at $p$.

    \item \textbf{Discrete Graphs}:
    \begin{equation}
    u_{\text{graph}}(i,j) = \left(\frac{1}{\sqrt{d_i}}, \frac{1}{\sqrt{d_j}}\right)
    \end{equation}
    where $d_i, d_j$ are the degrees of vertices $i, j$.

    \item \textbf{Euclidean Grids}:
    \begin{equation}
    u_{\text{grid}}(p,q) = q - p
    \end{equation}
    representing relative Cartesian differences.
\end{enumerate}

\textbf{Unifying Properties of the Framework}:
\begin{enumerate}
    \item \textbf{Mathematical consistency}: All domains reduce to the same Gaussian kernel formulation
    \item \textbf{Theoretical transferability}: The parameters $\theta$ can be transferred between domains with compatible dimensions
    \item \textbf{Architectural flexibility}: New domains only require defining the function $u$
    \item \textbf{Approximation guarantees}: Under regularity conditions, Gaussian mixtures can approximate any continuous kernel
\end{enumerate}
\end{example}

\subsection{Parallel Transport}\label{subsec:transporte_paralelo}

The work of Schonsheck et al.~\citep{schonsheck2018parallel} introduces Parallel Transport Convolutions (PTC), which address the coordinatization limitations of geodesic CNNs through the use of parallel transport.

\subsubsection{Geometric Foundations}

Parallel transport provides a canonical way to compare vectors at different points of a manifold, preserving the intrinsic geometric structure. The foundations of connections, parallel transport, and curvature are developed in \autoref{ch:gdl}: the Levi-Civita connection and the Christoffel symbols in Section~\ref{subsec:transporte_paralelo_gdl}, and parallel transport in principal bundles in \autoref{def:transporte_paralelo_formal}. Let us recall the essential properties for CNNs: parallel transport $\Gamma_\gamma: T_{\gamma(0)}M \to T_{\gamma(1)}M$ is a linear isometry that preserves the metric and satisfies the composition property $\Gamma_{\gamma_1 * \gamma_2} = \Gamma_{\gamma_2} \circ \Gamma_{\gamma_1}$.

Holonomy (the composition of parallel transports along closed loops) is intimately related to curvature. The Ambrose-Singer theorem establishes that the Lie algebra of the holonomy group is generated by the curvature operators $R(X,Y)$, which implies that the complete geometric information of the manifold can be captured via parallel transport, theoretically justifying Parallel Transport Convolutions.

\subsubsection{Parallel Transport Convolution}

\begin{definition}[PTC]
For a function $f: M \to \mathbb{R}^d$ that assigns vectors to points of the manifold, the PTC is defined as:
\begin{equation}
(f *_{PTC} \psi)(p) = \int_{\mathcal{N}(p)} \psi(\rho(p,q)) \langle \Gamma_{\gamma_{pq}}(f(q)), \mathbf{e}_i \rangle
\end{equation}
where $\mathcal{N}(p)$ is a neighborhood of $p$, $\gamma_{pq}$ is the geodesic from $p$ to $q$, and $\{\mathbf{e}_i\}$ is a fixed basis in $T_p M$.
\end{definition}

\subsubsection{Advantages of the Approach}

PTCs offer several advantages over previous methods:

\begin{itemize}
    \item \textbf{Coordinate invariance}: They do not depend on arbitrary choices of bases
    \item \textbf{Orientation preservation}: Parallel transport maintains relative angles and orientations
    \item \textbf{Geometric flexibility}: Applicable to manifolds with variable curvature
\end{itemize}

\begin{example}[Exact Parallel Transport on the Sphere $S^2$]

To illustrate Parallel Transport Convolutions, we develop the mathematical theory of parallel transport on the unit sphere $S^2$, where the exact formulation is possible.

\textbf{Mathematical Formulation of Parallel Transport}:

Let $\gamma: [0,1] \to S^2$ be a geodesic (great circle) connecting points $p = \gamma(0)$ and $q = \gamma(1)$. For a tangent vector $v \in T_p S^2$, the parallel transport $\Gamma_\gamma(v) \in T_q S^2$ is obtained via:

\begin{equation}
\nabla_{\dot{\gamma}(t)} V(t) = 0
\end{equation}

where $V(t)$ is the parallel extension of $v$ along $\gamma$.

\textbf{Explicit Formula for $S^2$}:

For non-antipodal points $p, q \in S^2$, parallel transport admits the closed-form expression:

\begin{equation}
\Gamma_{\gamma_{pq}}(v) = v - \frac{\langle v, p \rangle}{1 + \langle p, q \rangle}(p + q) + \frac{\langle v, q \rangle}{1 + \langle p, q \rangle}(p - q)
\end{equation}

\textbf{Special Cases}:
\begin{enumerate}
    \item \textbf{Identical points} ($p = q$): $\Gamma(v) = v$ (identity)
    \item \textbf{Antipodal points} ($q = -p$): The transport is not uniquely determined; a choice of great circle is required
    \item \textbf{Nearby points} ($\|p - q\| \ll 1$): First-order approximation $\Gamma(v) \approx v + O(\|p-q\|^2)$
\end{enumerate}

\textbf{Fundamental Properties of the Transport}:
\begin{enumerate}
    \item \textbf{Metric preservation}: $\langle \Gamma(v), \Gamma(w) \rangle_q = \langle v, w \rangle_p$
    \item \textbf{Linearity}: $\Gamma(\alpha v + \beta w) = \alpha \Gamma(v) + \beta \Gamma(w)$
    \item \textbf{Rotational invariance}: For rotation $R \in SO(3)$: $\Gamma_{R\gamma}(Rv) = R\Gamma_\gamma(v)$
    \item \textbf{Group property}: $\Gamma_{\gamma_2} \circ \Gamma_{\gamma_1} = \Gamma_{\gamma_2 * \gamma_1}$
\end{enumerate}

\textbf{Application to Convolutions}:

The parallel transport convolution on $S^2$ is defined as:
\begin{equation}
(f *_{PTC} \psi)(p) = \int_{S^2} \psi(d(p,q)) \langle \Gamma_{\gamma_{qp}}(f(q)), e_i \rangle \, d\sigma(q)
\end{equation}

where $d\sigma$ is the standard area measure on $S^2$ and $\{e_i\}$ is a fixed orthonormal basis in $T_p S^2$.

\textbf{Theoretical Advantages}:
\begin{enumerate}
    \item \textbf{Intrinsic invariance}: Independent of choices of local coordinates
    \item \textbf{Orientation preservation}: Maintains relative angles between tangent vectors
    \item \textbf{Geometric foundation}: Based on the canonical Levi-Civita connection
    \item \textbf{Stability}: Numerical conditioning controlled by the curvature
\end{enumerate}

\textbf{Theoretical Limitations}:
\begin{itemize}
    \item Computational complexity $O(n^2)$ for exact transport
    \item Singularities in antipodal configurations
    \item Dependence on the specific geometry of $S^2$
\end{itemize}
\end{example}

\subsection{Specific Manifolds}\label{subsec:variedades_especificas}

This subsection examines specialized extensions for manifolds with particular structures, leveraging specific geometric properties to achieve greater efficiency and expressiveness.

\subsubsection{Riemannian Networks}

Riemannian Graph Neural Networks~\citep{gu2021riemanniangraph} extend GNNs by incorporating Riemannian structure into the feature space:

\begin{definition}[Riemannian Aggregation]\label{def:agregacion_riemanniana}
In the hyperbolic Poincaré space $\mathbb{D}^n = \{x \in \mathbb{R}^n : \|x\| < 1\}$ with metric $g_x = \left(\frac{2}{1 - \|x\|^2}\right)^2 g_E$, aggregation is performed via:
\begin{equation}
\mathbf{h}_i^{(l+1)} = \sigma\left(\frac{1}{|\mathcal{N}(i)|} \bigoplus_{j \in \mathcal{N}(i)} \mathbf{W}^{(l)} \mathbf{h}_j^{(l)}\right)
\end{equation}
where $\bigoplus$ denotes the Möbius addition in the Poincaré model, defined by:
\begin{equation}
x \oplus y = \frac{(1 + 2\langle x, y \rangle + \|y\|^2)x + (1 - \|x\|^2)y}{1 + 2\langle x, y \rangle + \|x\|^2\|y\|^2}
\end{equation}
\end{definition}

\begin{proposition}[Geometric Consistency of the Aggregation]\label{prop:consistencia_agregacion_riemanniana}
The Riemannian aggregation of \autoref{def:agregacion_riemanniana} satisfies:
\begin{enumerate}
    \item \textbf{Closure}: If $\mathbf{h}_j^{(l)} \in \mathbb{D}^n$ for all $j \in \mathcal{N}(i)$, then $\mathbf{h}_i^{(l+1)} \in \mathbb{D}^n$.
    \item \textbf{Euclidean reduction}: When the curvature tends to zero ($\kappa \to 0$), the Möbius addition reduces to ordinary vector addition, recovering the standard GNN aggregation.
\end{enumerate}
\end{proposition}

\begin{proof}
For (1), the Möbius addition $x \oplus y$ maps $\mathbb{D}^n \times \mathbb{D}^n \to \mathbb{D}^n$: we verify that $\|x \oplus y\| < 1$ using the identity $\|x \oplus y\|^2 = \frac{\|x\|^2 + 2\langle x,y\rangle + \|y\|^2}{1 + 2\langle x,y\rangle + \|x\|^2\|y\|^2}$, and the inequality $\|x\|^2\|y\|^2 \leq \|x\|^2 + \|y\|^2$ (which follows from $(\|x\|\|y\| - 1)^2 \geq 0$ and $\|x\|, \|y\| < 1$) implies that the numerator is strictly less than the denominator.

For (2), the Poincaré model with curvature $\kappa < 0$ uses radius $r = 1/\sqrt{-\kappa}$. When $\kappa \to 0^-$, the radius tends to infinity, the hyperbolic distances approach the Euclidean ones, and the Möbius formula linearizes: $x \oplus y \to x + y$.
\end{proof}

\subsubsection{Hyperbolic Spaces}

Hyperbolic spaces provide natural representations for hierarchical data. The fundamental reason is that the hyperbolic space $\mathbb{H}^n$ has a volume that grows exponentially with the geodesic radius, reflecting the exponential growth in the number of nodes with depth in a tree.

\begin{definition}[Hyperbolic Distance in the Poincaré Model]\label{def:distancia_hiperbolica_poincare}
In the Poincaré disk $\mathbb{D}^n$, the geodesic distance between $x, y \in \mathbb{D}^n$ is:
\begin{equation}
d_{\mathbb{H}}(x, y) = \operatorname{arcosh}\left(1 + 2\frac{\|x - y\|^2}{(1 - \|x\|^2)(1 - \|y\|^2)}\right)
\end{equation}
\end{definition}

Hyperbolic Attention Networks~\citep{gulcehre2019hyperbolic} incorporate attention mechanisms that respect hyperbolic geometry:

\begin{equation}
\alpha_{ij} = \frac{\exp(-d_{\mathbb{H}}(\mathbf{q}_i, \mathbf{k}_j)/\tau)}{\sum_{k} \exp(-d_{\mathbb{H}}(\mathbf{q}_i, \mathbf{k}_k)/\tau)}
\end{equation}

where $d_{\mathbb{H}}$ is the hyperbolic distance of \autoref{def:distancia_hiperbolica_poincare} and $\tau > 0$ is a temperature parameter.

\begin{proposition}[Representational Capacity of Hyperbolic Embeddings]\label{prop:capacidad_hiperbolica}
A tree with $n$ nodes and branching factor $b$ can be embedded in $\mathbb{H}^2$ with $O(1)$ distortion (independent of $n$), whereas any embedding in $\mathbb{R}^d$ incurs distortion $\Omega(\sqrt{\log n / d})$.
\end{proposition}

\begin{proof}
For the hyperbolic space: we place the root at the origin and distribute the children of the node at depth $k$ on a geodesic circle of radius proportional to $k$. Since the perimeter of a hyperbolic circle grows as $2\pi \sinh(r) \sim \pi e^r$, there is sufficient space for $b^k$ nodes at depth $k$ using uniform angular separation, with distortion bounded independently of $n$.

For $\mathbb{R}^d$: by the Johnson-Lindenstrauss lemma, preserving the $\binom{n}{2}$ distances with distortion $1 + \epsilon$ requires dimension $d = \Omega(\log n / \epsilon^2)$, and the metric structure of the tree imposes the lower bound $\Omega(\sqrt{\log n / d})$ on the distortion.
\end{proof}

\subsubsection{Riemannian Residual Networks}

The work of Katsman et al.~\citep{katsman2023riemannian} extends ResNet architectures to general Riemannian manifolds, replacing the Euclidean addition $x + F(x)$ with its intrinsic analog via the exponential map.

\begin{definition}[Riemannian Residual]\label{def:residual_riemanniano}
The residual connection on a manifold $(\mathcal{M}, g)$ is defined as:
\begin{equation}
\mathbf{h}_{l+1} = \operatorname{Exp}_{\mathbf{h}_l}(\mathcal{F}(\mathbf{h}_l))
\end{equation}
where $\operatorname{Exp}_p: T_p\mathcal{M} \to \mathcal{M}$ is the exponential map (which follows the geodesic from $p$ with initial velocity $v$ for unit time) and $\mathcal{F}: T_p\mathcal{M} \to T_p\mathcal{M}$ is a parameterized differentiable function.
\end{definition}

\begin{proposition}[Stability of the Riemannian Residual]\label{prop:estabilidad_residual_riemanniano}
Let $(\mathcal{M}, g)$ be a complete Riemannian manifold with bounded sectional curvature $|K| \leq \kappa$. If $\|\mathcal{F}(\mathbf{h}_l)\|_g < \operatorname{inj}(\mathcal{M})$ (the injectivity radius), then:
\begin{enumerate}
    \item The update $\mathbf{h}_{l+1} = \operatorname{Exp}_{\mathbf{h}_l}(\mathcal{F}(\mathbf{h}_l))$ is well defined and is a local diffeomorphism.
    \item The geodesic distance satisfies:
    \begin{equation}
    d(\mathbf{h}_l, \mathbf{h}_{l+1}) = \|\mathcal{F}(\mathbf{h}_l)\|_g
    \end{equation}
\end{enumerate}
\end{proposition}

\begin{proof}
(1) By the Hopf-Rinow theorem, the completeness of $\mathcal{M}$ guarantees that $\operatorname{Exp}_p$ is defined on all of $T_p\mathcal{M}$. The condition $\|v\| < \operatorname{inj}(\mathcal{M})$ ensures that $\operatorname{Exp}_p$ is a diffeomorphism on the ball $B(0, \operatorname{inj}(\mathcal{M})) \subset T_p\mathcal{M}$.

(2) The curve $\gamma(t) = \operatorname{Exp}_{\mathbf{h}_l}(t \cdot \mathcal{F}(\mathbf{h}_l))$ for $t \in [0,1]$ is a geodesic. Since $\|\mathcal{F}(\mathbf{h}_l)\|_g < \operatorname{inj}(\mathcal{M})$, this geodesic is minimizing, hence $d(\mathbf{h}_l, \mathbf{h}_{l+1}) = L(\gamma) = \|\dot{\gamma}(0)\| = \|\mathcal{F}(\mathbf{h}_l)\|_g$.
\end{proof}

\begin{remark}[Connection with ODEs on Manifolds]
The Riemannian residual iteration can be interpreted as the Euler discretization of an ordinary differential equation on $\mathcal{M}$:
\begin{equation}
\frac{D\mathbf{h}}{dt} = \mathcal{F}_t(\mathbf{h}(t))
\end{equation}
where $D/dt$ is the covariant derivative along the curve $\mathbf{h}(t)$. This perspective connects directly with the theory of Neural ODEs~\citep{chen2019neural}, generalizing it to non-Euclidean domains.
\end{remark}

\subsubsection{Integration and Perspectives}

The extensions of neural networks to specific manifolds presented in this subsection share a common pattern: each adapts the fundamental algebraic operations (addition, inner product, aggregation) to the underlying geometric space, replacing the Euclidean operations with their intrinsically defined analogs.

Riemannian aggregation via Möbius addition, hyperbolic attention mechanisms based on the distance $d_{\mathbb{H}}$, and residual connections formulated with the exponential map $\operatorname{Exp}$ constitute instances of the same principle: respecting the geometry of the data space when designing each component of the network. This principle is formalized in the general framework of \autoref{ch:gdl}, where equivariance with respect to the isometry group of the manifold guarantees that the learned representations are intrinsic.

From a practical point of view, the choice between these formulations depends on the natural geometry of the data: hyperbolic spaces for hierarchical data (trees, taxonomies), SPD manifolds for covariance matrices, and general Riemannian manifolds when the curvature is variable. The following subsection formalizes these principles with the rigor of Riemannian geometry.

\subsection{CNNs on Riemannian Manifolds: Rigorous Foundations}

\subsubsection{Foundations of Riemannian Geometry}

For the foundations of differential manifolds and Riemannian metrics, including definitions of topological space, topological manifold, differentiable manifold, and Riemannian metric, see \autoref{subsec:fundamentos_variedades} of \autoref{ch:gdl}.

\begin{definition}[Levi-Civita Connection]
The Levi-Civita connection \(\nabla\) is the unique affine connection on \(\mathcal{M}\) that:
\begin{enumerate}
	\item Is torsion-free: \(\nabla_X Y - \nabla_Y X = [X,Y]\)
	\item Is compatible with the metric: \(X \langle Y, Z \rangle = \langle \nabla_X Y, Z \rangle + \langle Y, \nabla_X Z \rangle\)
\end{enumerate}
\end{definition}

\begin{definition}[Parallel Transport]
The parallel transport \(\Gamma_\gamma: T_p \mathcal{M} \to T_q \mathcal{M}\) along a
curve \(\gamma: [0,1] \to \mathcal{M}\) with \(\gamma(0) = p, \gamma(1) = q\) preserves the inner product:
\[
\langle\Gamma_\gamma(v), \Gamma_\gamma(w)\rangle_q = \langle v, w \rangle_p
\]
\end{definition}

\subsubsection{Rigorous Parallel Transport Convolution}

\begin{figure}[htbp]
	\centering
	\includegraphics[width=0.75\textwidth]{figures/parallel_transport_surface}
	\caption[Parallel transport on a curved surface]{Parallel transport of a vector $v \in T_p\mathcal{M}$ along a curve $\gamma$ on a curved surface. The transported vector $\Gamma_\gamma(v) \in T_q\mathcal{M}$ preserves the norm and angles, but its orientation with respect to the surface changes according to the curvature.}
	\label{fig:parallel_transport_surface}
\end{figure}

\begin{definition}[Convolution on a Manifold via Parallel Transport]
For \(f: \mathcal{M} \to \mathbb{R}^d\) and kernel \(K: \mathbb{R}_+ \to \mathbb{R}\):
\[
(f *_\mathcal{M} K)(p) = \int_{B_r(p)} K(d(p,q)) \Gamma_{\gamma_{pq}}(f(q)) \, d\text{Vol}(q)
\]
where:
\begin{itemize}
	\item \(B_r(p)\) is the geodesic ball of radius \(r\) centered at \(p\)
	\item \(d(p,q)\) is the geodesic distance
	\item \(\gamma_{pq}\) is the geodesic from \(q\) to \(p\)
	\item \(d\text{Vol}\) is the Riemannian volume element
\end{itemize}
\end{definition}

\begin{theorem}[Preservation of Isometries]
The parallel transport convolution is equivariant to
the isometries of the Riemannian manifold:
\[
\phi_*(f *_\mathcal{M} K) = (\phi_* f) *_\mathcal{M} K
\]
for every isometry \(\phi: \mathcal{M} \to \mathcal{M}\).
\end{theorem}

\begin{proof}
Isometries preserve geodesic distances, parallel transport, and the volume element:
\begin{align}
[(f *_\mathcal{M} K) \circ \phi](p) &= \int_{B_r(\phi(p))} K(d(\phi(p),q)) \Gamma_{\gamma_{\phi(p)q}}(f(q)) \, d\text{Vol}(q) \\
&= \int_{B_r(p)} K(d(p,\phi^{-1}(q))) \Gamma_{\gamma_{p\phi^{-1}(q)}}(f(\phi^{-1}(q))) \, d\text{Vol}(\phi^{-1}(q)) \\
&= [(\phi_* f) *_\mathcal{M} K](p) \qquad \square
\end{align}
\end{proof}

\subsection{Riemannian Residual Architectures}

\subsubsection{Geometric Residual Connections}

\begin{definition}[Riemannian Residual Connection]
On a manifold \(\mathcal{M}\), the residual update is performed via the exponential map:
\[
x_{l+1} = \text{Exp}_{x_l}(F_l(x_l))
\]
where:
\begin{itemize}
	\item \(x_l \in \mathcal{M}\) is the representation at layer \(l\)
	\item \(F_l: \mathcal{M} \to T\mathcal{M}\) maps points to tangent vectors
	\item \(\text{Exp}: T\mathcal{M} \to \mathcal{M}\) is the exponential map
\end{itemize}
\end{definition}

\begin{theorem}[Geometric Stability of Riemannian ResNets]
Riemannian residual connections preserve:
\begin{enumerate}
	\item \textbf{The metric structure of the manifold}: \(d(x_{l+1}, y_{l+1}) \approx d(x_l, y_l)\) for small updates
	\item \textbf{Equivariance to isometries}: \(\phi(x_{l+1}) = \text{Exp}_{\phi(x_l)}(d\phi(F_l(x_l)))\)
	\item \textbf{Numerical stability}: Avoids gradient problems on curved manifolds
\end{enumerate}
\end{theorem}

\begin{proof}[Proof sketch]
(1) Let $x_{l+1} = \operatorname{Exp}_{x_l}(\mathcal{F}(x_l))$ be the residual update. For small $\|\mathcal{F}(x_l)\|$, the exponential map is a first-order approximation: $d(x_{l+1}, x_l) \approx \|\mathcal{F}(x_l)\|$. If $\mathcal{F}$ is Lipschitz with constant $L_\mathcal{F}$, then $d(x_{l+1}, y_{l+1}) \leq d(x_l, y_l)(1 + L_\mathcal{F} \|\cdot\|)$, approximately preserving the distances.

(2) Let $\phi: \mathcal{M} \to \mathcal{M}$ be an isometry. Since $\phi$ preserves the connection, it commutes with $\operatorname{Exp}$: $\phi(\operatorname{Exp}_p(v)) = \operatorname{Exp}_{\phi(p)}(d\phi_p(v))$. If $\mathcal{F}$ is equivariant ($\mathcal{F}(\phi(x)) = d\phi_x(\mathcal{F}(x))$), then $\phi(x_{l+1}) = \phi(\operatorname{Exp}_{x_l}(\mathcal{F}(x_l))) = \operatorname{Exp}_{\phi(x_l)}(d\phi(\mathcal{F}(x_l))) = \operatorname{Exp}_{\phi(x_l)}(\mathcal{F}(\phi(x_l)))$.

(3) The residual formulation $x_{l+1} = \operatorname{Exp}_{x_l}(\mathcal{F}(x_l))$ avoids the accumulation of errors from direct composition: instead of composing nonlinear functions over the manifold, each layer only adds a tangent perturbation. The Jacobian of the composition satisfies $\|J_{x_{L}} \circ \cdots \circ J_{x_1}\| \leq \prod_l (1 + \|d\mathcal{F}_l\|)$, which grows moderately compared to $\prod_l \|J_l\|$ in architectures without residual connections.
\end{proof}

\begin{definition}[Riemannian Batch Normalization]
For data on a Riemannian manifold, normalization is performed in the tangent space:
\[
\text{BN}_\mathcal{M}(x) = \text{Exp}_{\mu} \left( \frac{\text{Log}_\mu(x) - \mathbb{E}[\text{Log}_\mu(X)]}{\sqrt{\text{Var}[\text{Log}_\mu(X)] + \epsilon}} \right)
\]
where \(\mu\) is the Fréchet mean of the batch.
\end{definition}

\subsubsection{Implementations in Specific Spaces}

\begin{example}[Hyperbolic Spaces (\(\mathbb{H}^n\))]
In the Poincaré model, the operations are performed via:
\begin{align}
\text{Möbius addition}: \quad x \oplus y &= \frac{(1+2\langle x,y\rangle + \|y\|^2)x + (1-\|x\|^2)y}{1 + 2\langle x,y\rangle + \|x\|^2\|y\|^2} \\
\text{Multiplication}: \quad \lambda \otimes x &= \tanh(\lambda \tanh^{-1}(\|x\|)) \frac{x}{\|x\|}
\end{align}
Ideal for hierarchical data with exponential growth.
\end{example}

\begin{example}[Spheres (\(S^n\))]
For directional data on \(S^n\), the updates preserve the norm:
\[
x_{l+1} = \text{Exp}_{x_l}(P_{T_{x_l}S^n}(F_l(x_l)))
\]
where \(P_{T_{x_l}S^n}\) is the projection onto the tangent space.
\end{example}

\begin{example}[SPD Manifolds]
For positive definite matrices \(\text{SPD}(n)\):
\[
\text{Exp}_X(V) = X^{1/2} \exp(X^{-1/2}VX^{-1/2}) X^{1/2}
\]
Applicable to covariance data and affine transformations.
\end{example}

\subsection{Advantages and Applications}

Methods on manifolds have demonstrated advantages in multiple domains:

\begin{itemize}
    \item \textbf{3D shape analysis}: Classification and segmentation of surfaces
    \item \textbf{Robotics}: Navigation in configuration spaces
    \item \textbf{Computer vision}: Analysis of omnidirectional images
    \item \textbf{Bioinformatics}: Analysis of protein structures
    \item \textbf{Natural language processing}: Hyperbolic embeddings for hierarchies
    \item \textbf{Quantitative finance}: Modeling of correlations in SPD matrices
\end{itemize}

\begin{remark}[Computational Advantages of Riemannian Architectures]\label{rem:ventajas_computacionales_riemannianas}
CNNs on manifolds offer:
\begin{enumerate}
	\item \textbf{Representational efficiency}: Lower effective dimensionality for data with intrinsic symmetries
	\item \textbf{Better generalization}: Respect for the natural geometry of the data
	\item \textbf{Geometric interpretability}: Direct correspondence between operations and geometric properties
\end{enumerate}
\end{remark}

The progression from specific geodesic CNNs towards general frameworks such as MoNet, parallel transport techniques, and Riemannian residual architectures illustrates the maturation of the field, providing increasingly sophisticated tools for processing non-Euclidean data.

\subsection{Case Studies in Specific Applications}\label{subsec:casos_estudio_variedades}

This subsection examines concrete applications of CNNs on manifolds, providing detailed examples of practical implementations and experimental comparisons with traditional Euclidean methods.

\subsubsection{3D Shape Analysis: FAUST Dataset}\label{subsubsec:formas_3d_faust}

The FAUST dataset (FAUST: Dataset and Evaluation for 3D Mesh Registration) provides a standard benchmark for evaluating geometric methods on 3D human surfaces. This case study illustrates how geodesic CNNs overcome the limitations of Euclidean projections.

\paragraph{Experimental Setup}

The FAUST dataset contains 100 triangular meshes of 10 human subjects in 10 different poses each. Each mesh has approximately 6890 vertices, representing genus-0 surfaces with significant variable curvature.

\begin{example}[Geodesic CNN Implementation for FAUST]
The practical implementation follows these steps:

\begin{algorithm}[H]
\caption{Geodesic CNN for FAUST Pose Classification}
\begin{algorithmic}[1]
    \State \textbf{Input:} Triangular mesh $M$ with vertices $V$ and faces $F$
    \State \textbf{Preprocessing:}
    \State \quad Normalize total surface area to 1
    \State \quad Compute discrete Laplace-Beltrami operator $L$
    \State \quad Compute HKS descriptors as initial features
    \State \textbf{For each vertex} $v_i \in V$:
    \State \quad \quad Construct local geodesic coordinate system
    \State \quad \quad Extract geodesic patch of radius $r = 0.1 \times \sqrt{\text{area}}$
    \State \quad \quad Sample on polar grid: 8 directions $\times$ 5 radii
    \State \quad \quad Apply $3 \times 3$ convolutional kernel in geodesic coordinates
    \State \quad \quad Correct by local Jacobian factor
    \State \textbf{Aggregation:} Global pooling over all vertices
    \State \textbf{Output:} Pose classification (10 classes)
\end{algorithmic}
\end{algorithm}

\textbf{Critical Technical Features}:
\begin{itemize}
    \item \textbf{Construction of geodesic coordinates}: For each vertex $v$, an orthonormal basis is selected in the tangent space using eigenvectors of the local curvature operator
    \item \textbf{Jacobian correction}: The factor $J(r, \theta) = \sqrt{\det g_r}$ where $g_r$ is the metric induced in polar coordinates
    \item \textbf{Boundary handling}: Vertices near boundaries use adaptive patches with reduced radii
\end{itemize}
\end{example}

\paragraph{Comparative Results}

\begin{table}[h]
\centering
\caption{Comparison of Methods on the FAUST Dataset}
\label{tab:faust_comparison}
\begin{tabular}{lcccc}
\hline
\textbf{Method} & \textbf{Accuracy} & \textbf{Time/Mesh} & \textbf{Parameters} & \textbf{Robust to Deform.} \\
\hline
Euclidean CNN (projection) & 73.2\% & 15ms & 2.1M & No \\
SpectralCNN & 81.7\% & 45ms & 1.8M & Yes (partial) \\
Geodesic CNN & \textbf{89.4\%} & 120ms & 1.5M & Yes \\
MoNet & 87.1\% & 95ms & 1.7M & Yes \\
PTC (Parallel Transport) & 88.6\% & 180ms & 1.4M & Yes \\
\hline
\end{tabular}
\end{table}

\paragraph{Analysis of Results}

The results reveal significant advantages of the geometric methods:

\begin{enumerate}
    \item \textbf{Superior accuracy}: Geodesic CNN outperforms Euclidean methods by 16.2 percentage points
    \item \textbf{Invariance to deformations}: Geometric methods maintain accuracy under isometric deformations
    \item \textbf{Parameter efficiency}: Lower number of parameters due to the intrinsic geometric structure
    \item \textbf{Computational trade-off}: Higher computational cost due to the construction of geodesic coordinates
\end{enumerate}

\textbf{Insights from the FAUST Case.} The detailed analysis of the FAUST dataset reveals that:
\begin{itemize}
    \item High-curvature regions (joints) benefit most from geodesic coordinates
    \item The Jacobian correction is critical near geometric singularities
    \item The optimal geodesic radius depends on the local mesh density
    \item The choice of initial features (HKS vs. Cartesian coordinates) significantly affects performance
\end{itemize}

\subsubsection{Omnidirectional Image Processing}\label{subsubsec:imagenes_omnidireccionales}

Omnidirectional images $(360^\circ)$ present unique challenges due to the distortions inherent in any planar projection. This case study compares approaches on the sphere $S^2$ versus Euclidean projections.

\paragraph{Datasets and Setup}

\begin{itemize}
    \item \textbf{SUN360}: 67,448 panoramic images for scene classification
    \item \textbf{Stanford 2D-3D-S}: Multimodal dataset with semantic segmentation annotations
    \item \textbf{Resolution}: $2048 \times 1024$ pixels in equirectangular projection
\end{itemize}

\begin{example}[Comparison of Representations for $360^\circ$ Images]

\textbf{1. Equirectangular Projection + Standard CNN}:
\begin{itemize}
    \item \textbf{Advantages}: Direct implementation, compatibility with existing architectures
    \item \textbf{Disadvantages}: Severe polar distortions, loss of metric information
    \item \textbf{Accuracy on SUN360}: 76.3\%
\end{itemize}

\textbf{2. Icosahedral Projection + Geodesic CNN}:
\begin{itemize}
    \item \textbf{Method}: Project the sphere onto an icosahedron, apply geodesic CNNs on each face
    \item \textbf{Advantages}: Uniform distortion, preservation of metric properties
    \item \textbf{Disadvantages}: Complexity in handling edges between faces
    \item \textbf{Accuracy on SUN360}: 82.1\%
\end{itemize}

\textbf{3. Direct Spherical CNNs}:
\begin{itemize}
    \item \textbf{Method}: Convolutions defined directly on spherical harmonics
    \item \textbf{Advantages}: Mathematically elegant, exact rotational equivariance
    \item \textbf{Disadvantages}: High computational complexity, resolution limitations
    \item \textbf{Accuracy on SUN360}: 84.7\%
\end{itemize}
\end{example}

\paragraph{Detailed Analysis: Distortion Effects}

\begin{figure}[h]
\centering
% Esta sería una figura ilustrativa de las distorsiones
\begin{minipage}{0.45\textwidth}
\textbf{Equirectangular Projection}:
\begin{itemize}
    \item Maximum distortion at poles: $\times 10$ actual area
    \item Circular objects → elliptical towards poles
    \item Loss of rotational invariance
\end{itemize}
\end{minipage}
\hfill
\begin{minipage}{0.45\textwidth}
\textbf{Geodesic Method}:
\begin{itemize}
    \item Maximum distortion: $\times 1.15$ actual area
    \item Preservation of local shapes
    \item Approximate rotational invariance
\end{itemize}
\end{minipage}
\caption[Comparison of distortions in omnidirectional images]{Comparison of distortions in omnidirectional images}
\label{fig:omnidirectional_distortion}
\end{figure}

\subsubsection{Analysis of Medical Surfaces: Cerebral Cortex}\label{subsubsec:corteza_cerebral}

The analysis of the cerebral cortex represents one of the most challenging use cases for CNNs on manifolds, due to the complex genus-0 topology with intricate foldings and extreme variable curvature.

\paragraph{Dataset and Challenges}

\begin{itemize}
    \item \textbf{Human Connectome Project (HCP)}: 1200 subjects with high-resolution magnetic resonance images
    \item \textbf{Preprocessing}: Extraction of cortical surfaces using FreeSurfer
    \item \textbf{Technical challenges}:
        \begin{itemize}
            \item Complex topology with deep sulci and gyri
            \item Variable curvature: from flat regions to $180^\circ$ foldings
            \item Significant inter-subject anatomical variability
            \item Need for precise anatomical correspondence
        \end{itemize}
\end{itemize}

\begin{example}[Pipeline for Cortical Analysis]

\textbf{Task}: Automatic segmentation of cortical areas (parcellation) into 180 anatomical regions.

\begin{algorithm}[H]
\caption{Geodesic CNN for Cortical Parcellation}
\begin{algorithmic}[1]
    \State \textbf{Input:} Cortical surface $S$ with $\sim$150k vertices
    \State \textbf{Local features:}
    \State \quad Mean and Gaussian curvature at each vertex
    \State \quad Cortical thickness (distance to white matter)
    \State \quad Spectral descriptors (first 10 eigenvalues of the Laplacian)
    \State \textbf{Construction of geodesic patches}:
    \State \quad Adaptive radius: $r = \max(2mm, k \times \sqrt{\text{local area}})$
    \State \quad Orientation based on principal gradient direction
    \State \quad Sampling: polar grid 12 directions $\times$ 6 radii
    \State \textbf{CNN architecture}:
    \State \quad 3 geodesic convolutional layers (32, 64, 128 channels)
    \State \quad Batch normalization adapted to manifolds
    \State \quad Skip connections preserving geometric structure
    \State \textbf{Output:} Probabilities for 180 anatomical classes
\end{algorithmic}
\end{algorithm}
\end{example}

\paragraph{Clinical Results}

\begin{proposition}[Optimality of Geometric Segmentation in the Cerebral Cortex]\label{prop:optimalidad_segmentacion}
Segmentation based on geodesic coordinates achieves theoretical optimality under the following conditions:

\begin{enumerate}
    \item \textbf{Regularity hypothesis}: Anatomical boundaries follow curves of minimal curvature in the induced cortical metric
    \item \textbf{Variational principle}: The optimal segmentation minimizes the functional
    \begin{equation}
    E[\Omega] = \int_{\partial \Omega} \kappa_g ds + \lambda \int_{\Omega} \|\nabla f\|_g^2 dA
    \end{equation}
    where $\kappa_g$ is the geodesic curvature of the boundary and $f$ are anatomical features
    \item \textbf{Topological consistency}: The regions satisfy genus and connectivity constraints determined by the anatomy
\end{enumerate}
\end{proposition}

\begin{proof}[Proof sketch]
Under the regularity hypothesis, the anatomical boundaries are geodesics or curves of minimal geodesic curvature in the cortical metric. The functional $E[\Omega]$ combines a geometric regularization term $\int_{\partial \Omega} \kappa_g \, ds$ (which penalizes boundaries with high geodesic curvature) and a fidelity term $\lambda \int_{\Omega} \|\nabla f\|_g^2 \, dA$ (which penalizes regions with high feature variation). By the calculus of variations on manifolds, the minimizers of $E$ satisfy the Euler-Lagrange equation:
\[
\kappa_g = \lambda \left[\|\nabla f\|_g^2\right]_{\partial \Omega}
\]
where the right-hand side represents the jump of $\|\nabla f\|_g^2$ across the boundary. Geodesic coordinates provide the natural coordinate system for solving this variational problem, since in them the metric takes the form $ds^2 = d\rho^2 + J(\rho, \theta)^2 d\theta^2$ with $J(0, \theta) = 0$ and $\partial_\rho J(0,\theta) = 1$, simplifying the computation of $\kappa_g$ and $\|\nabla f\|_g$.
\end{proof}

\begin{remark}[Convergence of Geometric Methods]
Under regularity conditions, geodesic CNNs converge to the optimal segmentation with rate $O(h^{p})$ where $h$ is the mesh resolution and $p$ depends on the regularity of the anatomical boundaries. This observation is analogous to the classical convergence rates of finite element methods on manifolds.
\end{remark}

\paragraph{Theoretical Foundations of the Clinical Impact}

The clinical advantages of geometric methods are grounded in solid mathematical principles:
\begin{itemize}
    \item \textbf{Principle of Bayesian optimality}: The incorporation of a geometric prior reduces the variance of the estimator by the factor $\frac{\sigma^2}{\sigma^2 + \sigma_{\text{prior}}^2}$
    \item \textbf{Regularization theorem}: The $H^s$ norm on manifolds provides natural control over spurious oscillations
    \item \textbf{Invariance to transformations}: The intrinsic formulation eliminates dependence on arbitrary scanner orientations
    \item \textbf{Automatic adaptation}: The Riemannian metric adapts naturally to individual anatomical variability
\end{itemize}

\subsection{Practical Implementation and Optimizations}\label{subsec:implementacion_practica_variedades}

This subsection provides essential technical details for robust implementations of CNNs on manifolds, including optimized algorithms, computational efficiency considerations, and numerical stabilization techniques.

\subsubsection{Efficient Discretization of Differential Operators}\label{subsubsec:discretizacion_operadores}

The practical implementation of CNNs on manifolds requires careful discretizations of differential geometry concepts.

\paragraph{Computation of the Mesh Laplacian}

The discrete Laplace-Beltrami operator is fundamental to many operations. The cotangent discretization provides the most accurate approximation:

\begin{definition}[Discrete Cotangent Laplacian]
For a triangular mesh with vertices $V$ and faces $F$, the discrete Laplacian is defined as:
\begin{equation}
L_{ij} = \begin{cases}
\frac{1}{2A_i}(\cot \alpha_{ij} + \cot \beta_{ij}) & \text{if } (i,j) \text{ are adjacent} \\
-\sum_{k \neq i} L_{ik} & \text{if } i = j \\
0 & \text{otherwise}
\end{cases}
\end{equation}
where $A_i$ is the Voronoi area of vertex $i$, and $\alpha_{ij}, \beta_{ij}$ are the angles opposite to edge $(i,j)$.
\end{definition}

\begin{algorithm}[H]
\caption{Efficient Construction of the Mesh Laplacian}
\begin{algorithmic}[1]
    \State \textbf{Input:} Vertices $V \in \mathbb{R}^{n \times 3}$, faces $F \in \mathbb{Z}^{m \times 3}$
    \State Initialize sparse matrix $L \in \mathbb{R}^{n \times n}$
    \State Compute edge vectors for each face
    \State \textbf{For each face} $f = (i, j, k) \in F$:
    \State \quad Compute edge lengths: $a = \|v_j - v_k\|$, $b = \|v_k - v_i\|$, $c = \|v_i - v_j\|$
    \State \quad Compute cotangents using the law of cosines:
    \State \quad \quad $\cot \alpha = \frac{b^2 + c^2 - a^2}{4 \times \text{area}(f)}$
    \State \quad Update matrix entries: $L[i,j] += \cot \alpha$
    \State Normalize by Voronoi areas
    \State Impose the zero-sum condition: $L_{ii} = -\sum_{j \neq i} L_{ij}$
    \State \textbf{Output:} Sparse Laplacian matrix
\end{algorithmic}
\end{algorithm}

\paragraph{Optimized GPU Implementation}

\textbf{Optimization of the Mesh Laplacian for Parallel Computation.} The efficient discretization of the Laplace-Beltrami operator requires special theoretical considerations for implementation on parallel architectures:

\textbf{Computational Properties of the Discrete Laplacian}:
\begin{enumerate}
    \item \textbf{Sparsity}: The Laplacian matrix has $O(n)$ nonzero entries for $n$ vertices
    \item \textbf{Symmetry}: $L_{ij} = L_{ji}$ guaranteed by geometric construction
    \item \textbf{Negative definiteness}: All eigenvalues are $\leq 0$, with $\lambda_0 = 0$
    \item \textbf{Locality}: Each entry depends only on the local geometry of adjacent triangles
\end{enumerate}

\textbf{Considerations for Parallelization}:
\begin{itemize}
    \item \textbf{Vectorization of cotangents}: The computation of $\cot \alpha_{ij} = \frac{\langle e_{ik}, e_{jk} \rangle}{\|e_{ik} \times e_{jk}\|}$ can be parallelized over all faces simultaneously
    \item \textbf{Sparse assembly}: The construction of the matrix requires coherent access patterns to avoid race conditions
    \item \textbf{Preconditioning}: Adaptive multigrid methods for large linear systems
    \item \textbf{Numerical stability}: Adaptive regularization in regions of high curvature or degenerate triangles
\end{itemize}

\textbf{Algorithmic Complexity}:
\begin{itemize}
    \item Construction: $O(|F|)$ where $|F|$ is the number of faces
    \item Matrix-vector product: $O(n \cdot \text{average degree})$
    \item Memory: $O(n + |E|)$ for sparse storage
\end{itemize}

\subsubsection{Algorithms for Geodesic Computation}\label{subsubsec:calculo_geodesicas}

The efficient computation of geodesic distances is critical for the performance of geodesic CNNs.

\paragraph{Approximation via Dijkstra on Meshes}

For dense meshes, the modified Dijkstra algorithm provides efficient geodesic approximations:

\begin{algorithm}[H]
\caption{Geodesic Dijkstra on a Triangular Mesh}
\begin{algorithmic}[1]
    \State \textbf{Input:} Mesh $(V, E)$, source vertex $s$, maximum radius $r_{max}$
    \State Initialize distances: $d[s] = 0$, $d[v] = \infty$ for $v \neq s$
    \State Initialize priority queue $Q$ with all vertices
    \State \textbf{While} $Q$ is not empty \textbf{and} $\min(Q) < r_{max}$:
    \State \quad Extract vertex $u$ with minimum distance
    \State \quad \textbf{For each} neighbor $v$ of $u$:
    \State \quad \quad Compute geodesic weight: $w = \|p_u - p_v\|$ (Euclidean distance)
    \State \quad \quad \textbf{If} $d[u] + w < d[v]$:
    \State \quad \quad \quad $d[v] = d[u] + w$
    \State \quad \quad \quad Update predecessor: $\pi[v] = u$
    \State \textbf{Output:} Geodesic distances $d$ and predecessors $\pi$
\end{algorithmic}
\end{algorithm}

\paragraph{Fast Marching Method for Greater Accuracy}

For applications that require greater geodesic accuracy, the Fast Marching Method provides solutions of the Eikonal equation:

\begin{equation}
\|\nabla T(x)\| = f(x)
\end{equation}

where $T(x)$ is the arrival time and $f(x) = 1/v(x)$ is the local velocity function.

\textbf{Theoretical Formulation of the Fast Marching Method.} The Fast Marching Method provides an efficient solution of the Eikonal equation $\|\nabla T(x)\| = f(x)$ on triangular meshes, where $T(x)$ represents the arrival time from a source.

\textbf{Theoretical Foundations}:
\begin{enumerate}
    \item \textbf{Eikonal equation}: $\|\nabla T\|^2 = f^2(x)$ with boundary condition $T(\partial \Omega) = g$
    \item \textbf{Causality principle}: Information propagates from regions of lower time towards higher ones
    \item \textbf{Monotonicity}: $T$ is monotonically non-decreasing along rays from the source
    \item \textbf{Viscosity}: Unique viscosity solution that handles singularities automatically
\end{enumerate}

\textbf{Discretization on Triangular Meshes}:

For a vertex $v_i$ in a triangle $T = \{v_i, v_j, v_k\}$, the discrete update solves:
\begin{equation}
\max\left\{\frac{T_i - T_j}{\|v_i - v_j\|}, \frac{T_i - T_k}{\|v_i - v_k\|}\right\} = f_i
\end{equation}

subject to the causality constraint that the values $T_j, T_k$ are already determined.

\textbf{Ordered Heap Algorithm}:
\begin{enumerate}
    \item \textbf{States}: ALIVE (final value), NARROW\_BAND (candidate), FAR (not processed)
    \item \textbf{Initialization}: Mark sources as ALIVE, neighbors as NARROW\_BAND
    \item \textbf{Propagation}: Extract minimum from heap, update FAR neighbors
    \item \textbf{Termination}: When the heap empties or the maximum distance is reached
\end{enumerate}

\textbf{Theoretical Guarantees}:
\begin{itemize}
    \item \textbf{Convergence}: $O(h)$ for meshes of size $h$ on smooth domains
    \item \textbf{Complexity}: $O(n \log n)$ where $n$ is the number of vertices
    \item \textbf{Stability}: Conditioning controlled by the geometry of the mesh
    \item \textbf{Adaptability}: Automatic handling of variable velocities $f(x)$
\end{itemize}

\textbf{Applications to Geodesic Patches}:
To extract geodesic neighborhoods of radius $r$, one uses $f(x) \equiv 1$ (constant velocity) and terminates when $T(x) > r$. The vertices with $T(x) \leq r$ form the desired geodesic patch.

\subsubsection{Parallelization and GPU Optimization}\label{subsubsec:optimizacion_gpu}

CNNs on manifolds present unique challenges for parallelization due to the irregular geometry of the meshes.

\paragraph{Batching Strategies for Meshes}

Unlike regular images, meshes have variable topologies that complicate direct batching.

\begin{algorithm}[H]
\caption{Adaptive Batching for Meshes of Variable Sizes}
\begin{algorithmic}[1]
    \State \textbf{Input:} List of meshes $\{M_1, M_2, ..., M_B\}$ with $|V_i|$ vertices each
    \State Determine the maximum size: $N_{max} = \max_i |V_i|$
    \State \textbf{For each mesh} $M_i$:
    \State \quad \textbf{If} $|V_i| < N_{max}$:
    \State \quad \quad Add dummy vertices with zero features
    \State \quad \quad Create validity mask $mask_i \in \{0,1\}^{N_{max}}$
    \State \quad \textbf{Else}:
    \State \quad \quad Apply consistent subsampling to $N_{max}$ vertices
    \State Concatenate into batch tensor: $X \in \mathbb{R}^{B \times N_{max} \times F}$
    \State Apply CNNs in batch with attention to validity masks
    \State \textbf{Output:} Predictions masked by individual mesh
\end{algorithmic}
\end{algorithm}

\paragraph{Memory Optimization}

\textbf{Theoretical Strategies for Memory Optimization.} Geodesic CNNs present unique computational challenges due to the irregular geometry. The following theoretical strategies optimize resource usage:

\textbf{1. Gradient Checkpointing}:
\begin{itemize}
    \item \textbf{Principle}: Trade computation for memory, recomputing geodesic coordinates during backpropagation
    \item \textbf{Time complexity}: Factor $\times 2$ in computation, reduction $\times 1/N$ in memory
    \item \textbf{Applicability}: Particularly effective for costly differential operators
\end{itemize}

\textbf{2. Sequential Processing by Patches}:
\begin{itemize}
    \item \textbf{Spatial partitioning}: Divide the mesh into coherent geometric regions
    \item \textbf{Temporal locality}: Exploit cache coherence via geodesic ordering
    \item \textbf{Load-memory balance}: Adaptive chunk size according to local curvature
\end{itemize}

\textbf{3. Adaptive Mixed Precision}:
\begin{itemize}
    \item \textbf{Forward pass}: FP16 for regular geometric computations
    \item \textbf{Critical regions}: FP32 for areas of high curvature or singularities
    \item \textbf{Gradients}: Full precision for derivatives of differential operators
\end{itemize}

\textbf{4. Hierarchical Geodesic Cache}:
\begin{itemize}
    \item \textbf{Strategic precomputation}: Geodesic distances between reference points
    \item \textbf{Local interpolation}: Low-order approximations for frequent queries
    \item \textbf{Adaptive eviction}: Priority based on access frequency and recomputation cost
\end{itemize}

\textbf{Complexity Models}:
\begin{align}
\text{Base memory} &= O(nF + E) \\
\text{Memory with chunking} &= O(\text{chunk\_size} \cdot F + E) \\
\text{Theoretical speedup} &= \frac{\text{bandwidth}_{\text{memory}}}{\text{bandwidth}_{\text{compute}}} \times \text{reuse\_ratio}
\end{align}

where $n$ is the number of vertices, $F$ the number of features, and $E$ the edges of the mesh.

\subsection{Experimental Evaluation and Benchmarks}\label{subsec:evaluacion_experimental_variedades}

This subsection presents a systematic evaluation of CNNs on manifolds across multiple standard datasets, including ablation analysis and comparisons with Euclidean projection methods.

\subsubsection{Standard Benchmark Datasets}\label{subsubsec:datasets_benchmark}

\paragraph{Benchmarks for Shape Analysis}

\begin{definition}[Categories of Problems on Manifolds]
CNN methods on manifolds address the following classes of mathematical problems:

\begin{enumerate}
    \item \textbf{Shape classification}: Determine the equivalence class of $(M, g)$ under a transformation group $G$
    \item \textbf{Geometric correspondence}: Find a mapping $\phi: M_1 \to M_2$ that preserves geometric structure
    \item \textbf{Semantic segmentation}: Partition $M = \bigcup_{i=1}^k \Omega_i$ according to geometric/semantic criteria
    \item \textbf{Feature detection}: Identify regions $U \subset M$ with distinguishable geometric properties
\end{enumerate}

Each class requires specific mathematical formulations and evaluation metrics adapted to the geometry of the problem.
\end{definition}

\paragraph{Theoretical Framework for Rigorous Evaluation}

The mathematically rigorous evaluation of methods on manifolds requires special theoretical considerations:

\begin{definition}[Geometric-Invariant Evaluation Protocol]
\leavevmode
\begin{enumerate}
    \item \textbf{Canonical normalization}:
    \begin{itemize}
        \item Orientation: Alignment with the principal components of $\int_M x \otimes x \, d\text{Vol}(x)$
        \item Scale: Normalization $\text{Vol}_g(M) = 1$ via $g \mapsto g/\text{Vol}_g(M)$
        \item Center of mass: Translation for $\int_M x \, d\text{Vol}(x) = 0$
    \end{itemize}

    \item \textbf{Controlled geometric augmentation}:
    \begin{itemize}
        \item Isometries: $\phi \in \text{Isom}(M)$ preserving the metric
        \item Metric perturbations: $g \mapsto g + \epsilon h$ with $\|h\|_{C^k} \leq \delta$
        \item Topologically consistent resampling preserving genus and features
    \end{itemize}

    \item \textbf{Invariant metrics}:
    \begin{itemize}
        \item Geodesic error: $d_g(\phi(p), q)$ for correspondences
        \item Geometric IoU: $\frac{\text{Vol}(\Omega_1 \cap \Omega_2)}{\text{Vol}(\Omega_1 \cup \Omega_2)}$ for segmentation
        \item Spectral measures: Based on eigenvalues of the Laplacian
    \end{itemize}

    \item \textbf{Statistical validity}:
    \begin{itemize}
        \item Geometric bootstrap over manifolds
        \item Non-parametric tests for distributions on curved spaces
        \item Correction for multiple comparisons in Riemannian spaces
    \end{itemize}
\end{enumerate}
\end{definition}

\subsubsection{Comprehensive Comparative Results}\label{subsubsec:resultados_comparativos}

\paragraph{Main Benchmark: Shape Classification}

\begin{theorem}[Hierarchy of Representational Capacities]
The different approaches for CNNs on manifolds form a natural theoretical hierarchy:

\begin{equation}
\mathcal{F}_{\text{Euclidean}} \subset \mathcal{F}_{\text{Spectral}} \subset \mathcal{F}_{\text{Geodesic}} \subset \mathcal{F}_{\text{Gauge}}
\end{equation}

where $\mathcal{F}_X$ denotes the space of functions approximable by method $X$.
\end{theorem}

\begin{proof}[Sketch]
The inclusion is established via:
\begin{enumerate}
    \item \textbf{Euclidean $\subset$ Spectral}: Every Euclidean kernel can be expressed spectrally via the Fourier transform
    \item \textbf{Spectral $\subset$ Geodesic}: Spectral methods are particular cases with global geodesic coordinates
    \item \textbf{Geodesic $\subset$ Gauge}: Geodesic coordinates correspond to specific gauge choices
\end{enumerate}
Each inclusion is strict, as constructive counterexamples show.
\end{proof}

\begin{corollary}[Asymptotic Optimality]
In the limit of infinite data and sufficient complexity, gauge-equivariant methods achieve the Bayes error for classification problems on manifolds.
\end{corollary}

\begin{proof}
By the expressive hierarchy theorem (whose proof precedes this corollary), gauge-equivariant methods constitute the most expressive class in the hierarchy, containing as particular cases all equivariant estimators of lesser generality. By the Hunt-Stein theorem of statistical decision theory, for problems invariant under a compact group, the minimax equivariant estimator is admissible. The universality of gauge-equivariant methods guarantees that, with sufficient complexity (network depth and width), they can approximate any equivariant estimator. In the limit of infinite data, the consistency of the maximum likelihood estimator implies convergence to the Bayes classifier, which minimizes the expected risk.
\end{proof}

\paragraph{Theoretical Foundations of Significance}

The statistical superiority of geometric methods is grounded in deep theoretical principles:

\begin{theorem}[Principle of Geometric Statistical Efficiency]
For problems with intrinsic symmetries characterized by group $G$, equivariant estimators achieve lower asymptotic variance:
\begin{equation}
\text{Var}[\hat{\theta}_{\text{equivariant}}] \leq \frac{1}{|G|} \text{Var}[\hat{\theta}_{\text{euclidean}}]
\end{equation}
\end{theorem}

\begin{corollary}[Improved Convergence]
Geometric methods require $O(|G|)$ fewer samples to achieve the same accuracy as Euclidean approximations, where $|G|$ is the effective size of the symmetry group.
\end{corollary}

\begin{proof}
Direct consequence of the previous theorem. The convergence rate of an estimator scales as $O(1/\sqrt{n/\text{Var}})$, where $n$ is the sample size. By the geometric statistical efficiency theorem:
\[
\text{Var}[\hat{\theta}_{\text{equiv}}] \leq \frac{1}{|G|}\text{Var}[\hat{\theta}_{\text{eucl}}]
\]
To achieve error $\epsilon$, the Euclidean estimator needs $n_{\text{eucl}} = O(\text{Var}_{\text{eucl}}/\epsilon^2)$ samples, while the equivariant one needs $n_{\text{equiv}} = O(\text{Var}_{\text{equiv}}/\epsilon^2) \leq O(\text{Var}_{\text{eucl}}/(|G|\epsilon^2)) = n_{\text{eucl}}/|G|$. That is, the variance reduction by a factor $|G|$ is equivalent to multiplying the effective sample size by $|G|$.
\end{proof}

\subsubsection{Detailed Ablation Studies}\label{subsubsec:estudios_ablacion}

\paragraph{Impact of Architectural Components}

\begin{remark}[Analysis of Theoretical Contributions]\label{rem:analisis_contribuciones_teoricas}
The effectiveness of geodesic CNNs is grounded in the synergistic interaction of several mathematical components:

\begin{enumerate}
    \item \textbf{Coordinate Systems}:
    \begin{itemize}
        \item \textbf{Jacobian correction}: Essential for compensating distortions $J(\rho, \theta) = \sqrt{\det g}$
        \item \textbf{Geodesic coordinates}: Provide a natural parameterization that preserves the local metric structure
        \item \textbf{Alternative choices}: Spherical coordinates introduce bias due to non-constant curvature
    \end{itemize}

    \item \textbf{Intrinsic Geometric Features}:
    \begin{itemize}
        \item \textbf{External coordinates}: Basic positional information with invariance limitations
        \item \textbf{Local curvature}: Fundamental invariants $K$ (Gaussian) and $H$ (mean)
        \item \textbf{Spectral descriptors}: HKS = $\sum_k e^{-t\lambda_k} \phi_k^2(x)$ captures multi-scale geometry
        \item \textbf{Wavelet descriptors}: WKS provides optimal space-frequency localization
    \end{itemize}

    \item \textbf{Neighborhood Configuration}:
    \begin{itemize}
        \item \textbf{Optimal radius}: Trade-off between locality ($r \to 0$) and contextual information ($r \to \infty$)
        \item \textbf{Dynamic adaptation}: $r_{\text{eff}}(x) = r_0 / \sqrt{1 + |K(x)| r_0^2}$ according to local curvature
    \end{itemize}

    \item \textbf{Network Architecture}:
    \begin{itemize}
        \item \textbf{Depth}: Multiple layers allow composition of geometric transformations
        \item \textbf{Residual connections}: Stabilize gradients on curved manifolds
        \item \textbf{Geometric attention}: Weights information according to local geometric relevance
    \end{itemize}
\end{enumerate}
\end{remark}

\begin{remark}[Synergy of Components]
The effectiveness of geodesic CNNs emerges from the nonlinear interaction between components, not from the sum of individual contributions. Formally:
\begin{equation}
\text{Total efficacy} \neq \sum_{i} \text{Contribution}_i
\end{equation}
\end{remark}

\paragraph{Geometric Sensitivity Analysis}

\begin{figure}[h]
\centering
% Aquí iría una gráfica mostrando robustez a deformaciones
\begin{minipage}{0.48\textwidth}
\textbf{Theoretical Robustness Analysis}:
\begin{itemize}
    \item \textbf{Metric perturbations}: For $\|g - g_0\|_{C^k} \leq \epsilon$, error bounded by $O(\epsilon)$
    \item \textbf{Stochastic noise}: Stability theorem guarantees gradual degradation
    \item \textbf{Topological invariance}: Methods preserve robustness under homeomorphisms
    \item \textbf{Implicit regularization}: Laplacian smoothing improves numerical stability
\end{itemize}
\end{minipage}
\hfill
\begin{minipage}{0.48\textwidth}
\textbf{Theoretical Invariance Guarantees}:
\begin{itemize}
    \item \textbf{Isometries}: Exact invariance by construction for $\phi \in \text{Isom}(M)$
    \item \textbf{Homotheties}: Invariance under scaling $g \mapsto \lambda^2 g$ with $\lambda > 0$
    \item \textbf{Quasi-isometries}: Controlled stability for small deformations
    \item \textbf{Topological changes}: Outside the theoretical framework of differentiable manifolds
\end{itemize}
\end{minipage}
\caption[Geometric robustness analysis]{Geometric robustness analysis}
\label{fig:geometric_robustness}
\end{figure}

\subsubsection{Computational Efficiency Analysis}\label{subsubsec:eficiencia_computacional}

\paragraph{Detailed Performance Profiling}

\begin{proposition}[Computational Complexity of Geometric Approaches]\label{prop:complejidad_aproximaciones_geometricas}
The time complexity of different geometric approaches is characterized by:

\begin{align}
\text{Geodesic CNN}: &\quad O(n \log n + nk^2) \\
\text{MoNet}: &\quad O(n + E \cdot d^2) \\
\text{PTC}: &\quad O(n^{3/2} + nk^3) \\
\text{SpectralCNN}: &\quad O(n^3 + nk) \\
\text{PointNet++}: &\quad O(n \log n)
\end{align}

where $n$ is the number of vertices, $E$ the number of edges, $k$ the kernel size, and $d$ the dimension of pseudo-local coordinates.
\end{proposition}

\begin{proof}
Each bound is obtained by direct analysis of the operations:
\begin{itemize}
    \item \textbf{Geodesic CNN}: $O(n \log n)$ for computing geodesic distances via the Fast Marching Method, plus $O(nk^2)$ for the convolution with a kernel of size $k \times k$ at each vertex.
    \item \textbf{MoNet}: $O(E \cdot d^2)$ for the evaluation of $d$ Gaussian kernels on $d$ pseudo-local coordinates for each edge, plus $O(n)$ for the aggregation.
    \item \textbf{PTC}: $O(n^{3/2})$ for the computation of exact parallel transport (solving differential equations on the mesh), plus $O(nk^3)$ for the convolution with a $k \times k \times k$ kernel.
    \item \textbf{SpectralCNN}: $O(n^3)$ dominated by the diagonalization of the Laplacian, plus $O(nk)$ for the application of the spectral filter with $k$ coefficients.
    \item \textbf{PointNet++}: $O(n \log n)$ for the nearest neighbor search via kd-tree, with local aggregation in $O(n)$.
\end{itemize}
\end{proof}

\textbf{Fundamental Trade-offs.} There exists a theoretical hierarchy in the accuracy-efficiency trade-off:
\begin{enumerate}
    \item \textbf{Spectral methods}: High theoretical accuracy but complexity $O(n^3)$ due to diagonalization
    \item \textbf{Geodesic methods}: Optimal balance with complexity $O(n \log n)$ and invariance guarantees
    \item \textbf{Local approximations}: Efficiency $O(n)$ at the cost of loss of global information
    \item \textbf{Euclidean methods}: Maximum efficiency but without geometric guarantees
\end{enumerate}

\paragraph{Scalability Theory on Manifolds}

The scalability of different geometric approaches reveals fundamental principles:

\begin{proposition}[Hierarchy of Complexities]\label{prop:jerarquia_complejidades}
For manifolds discretized with $n$ vertices, the methods are classified according to:
\begin{enumerate}
    \item \textbf{Linear} $O(n)$: MoNet and methods based on pseudo-local coordinates
    \item \textbf{Quasi-linear} $O(n \log n)$: Geodesic CNN, Fast Marching
    \item \textbf{Superlinear} $O(n^{3/2})$: Exact parallel transport
    \item \textbf{Cubic} $O(n^3)$: Spectral methods with full diagonalization
\end{enumerate}
\end{proposition}

\begin{proof}
The classification follows from \autoref{prop:complejidad_aproximaciones_geometricas}. Methods based on pseudo-local coordinates (MoNet) only require evaluating local functions on each edge, giving $O(|E|) = O(n)$ for graphs of bounded degree. Exact geodesic distances require Fast Marching with complexity $O(n \log n)$ due to the use of a priority queue. Exact parallel transport on triangular meshes involves solving ODEs along geodesics, with complexity $O(n^{3/2})$ for regular meshes where the number of edges grows as $O(n)$ and each transport computation requires $O(\sqrt{n})$ steps. Finally, the full diagonalization of the Laplacian via standard algorithms (QR, divide and conquer) has complexity $O(n^3)$.
\end{proof}

\begin{remark}[Theoretical Efficiency Limits]
There exists a lower bound $\Omega(n \log n)$ for algorithms that require the construction of exact geodesic neighborhoods on manifolds of variable curvature. This bound is obtained by reduction to the sorting problem: determining the $k$ nearest neighbors in the geodesic metric requires comparisons that, in the general case of variable curvature, cannot be solved in sublinear time without preprocessing.
\end{remark}

\paragraph{Accuracy vs Efficiency Trade-offs}

\textbf{Theoretical Selection Principles.} The choice of the optimal method is based on fundamental mathematical principles:

\begin{enumerate}
    \item \textbf{Principle of parsimony}: Choose the method with the lowest complexity that satisfies the required theoretical guarantees
    \item \textbf{Approximation theorem}: Methods with greater expressiveness require less data to achieve a given accuracy
    \item \textbf{Bias-variance trade-off}: Geometric methods reduce variance at the cost of inductive bias
    \item \textbf{Invariance vs flexibility}: There is a fundamental trade-off between exact invariances and adaptability
    \item \textbf{Geometric regularity}: Spectral methods require smoothness; local methods tolerate irregularities
\end{enumerate}

\textbf{Theoretical criterion}: The improvement in the generalization bound $O(\sqrt{\log N/n})$ justifies the increase in computational complexity when $N$ (model capacity) decreases significantly.

This theoretical and practical foundation established on Riemannian manifolds provides the conceptual basis for even more general extensions, such as the gauge convolutions that we will explore in the following section.
\section{Gauge Convolutions}\label{sec:gauge_cnns}

The diversity of techniques specific to particular manifolds, examined in the previous section, raises a fundamental question: is there a unifying framework that transcends specific geometries? The answer emerges from \textbf{gauge theory}, providing general principles for the design of architectures invariant under local transformations.

Gauge convolutions represent a profound conceptual extension of the convolutional paradigm. While classical convolutions exploit global symmetries (translations), and geometric extensions exploit intrinsic symmetries of specific manifolds, \textbf{gauge convolutions} address local symmetries that may vary from point to point in space. This approach, inspired by fundamental physics, provides an elegant mathematical framework for the construction of neural networks with sophisticated invariance properties.

\subsection{From Global to Local Symmetries: Intuitive Motivation}\label{subsec:motivacion_simetrias_locales}

Before delving into the mathematical formalism, it is crucial to understand \textit{why} we need local symmetries and how they emerge naturally in real problems. This section provides the intuition and motivation necessary to justify the mathematical complexity that follows.

\begin{figure}[htbp]
	\centering
	\includegraphics[width=0.90\textwidth]{figures/gauge_freedom}
	\caption[Gauge freedom: choice of local frames]{Gauge freedom: the same surface admits different choices of local frames (Gauge 1 and Gauge 2). The gauge-equivariant convolution produces identical results regardless of the gauge choice, connected by pointwise transformations $g(x) \in SO(2)$.}
	\label{fig:gauge_freedom}
\end{figure}

\subsubsection{The problem of rigid symmetries}

Classical CNNs assume that the same transformations (translations) are valid throughout the domain. This assumption works well for flat images, but fails systematically when the data has intrinsic geometric structure.

\begin{example}[Rotations on a Sphere]
Consider an image of the Earth projected onto a sphere. If we apply a global rotation of $90^\circ$ around the north-south axis:
\begin{itemize}
    \item At the equator, the rotation preserves local shapes perfectly
    \item Near the poles, the same rotation causes severe distortion
    \item An image of a continent near the north pole looks completely different after the global rotation
\end{itemize}
This example illustrates that \textbf{the appropriate transformation depends on the position}.
\end{example}

\textbf{Failure of Global Equivariance.} In curved spaces or those with non-trivial topology, applying the same transformation everywhere can:
\begin{enumerate}
    \item \textbf{Break the natural geometric structure} of the domain
    \item \textbf{Introduce artificial distortions} that did not exist in the original data
    \item \textbf{Violate the physical principles} that govern the modeled phenomenon
\end{enumerate}

\subsubsection{Fundamental physical analogies}

Gauge theory has its roots in physics, where it describes phenomena in which the mathematical description depends on arbitrary choices of local coordinates or references.

\begin{example}[Electromagnetism: The Classical Paradigm]
In electromagnetism, the observable physical fields (such as the electric field $\mathbf{E}$ and magnetic field $\mathbf{B}$) do not depend on how we choose to describe the electromagnetic potential $A_\mu(x)$.

\textbf{Gauge Freedom}: We can transform $A_\mu(x) \mapsto A_\mu(x) + \partial_\mu \Lambda(x)$ where $\Lambda(x)$ is an arbitrary function that can vary from point to point.

\textbf{Physical Invariance}: The observable fields remain unaltered:
\begin{align}
\mathbf{E} &= -\nabla \phi - \frac{\partial \mathbf{A}}{\partial t} \\
\mathbf{B} &= \nabla \times \mathbf{A}
\end{align}
remain invariant under the gauge transformation.

\textbf{Connection with Deep Learning}: Just as physical laws should not depend on our choice of coordinates, neural representations should not depend on arbitrary choices of local reference systems.
\end{example}

\begin{example}[Navigation on Earth: Local Coordinate Systems]
Imagine you are developing a global navigation system:
\begin{itemize}
    \item In New York, it is natural to use coordinates that point toward ``geographic north''
    \item In Sydney, ``geographic north'' points in a completely different direction in three-dimensional space
    \item Nevertheless, the \textit{local relations} (left, right, forward, backward) are consistent in both cities
\end{itemize}

A \textbf{gauge equivariant} neural network can learn to recognize patterns that are invariant to these choices of local coordinate system, while maintaining the ability to exploit the real geometric structure of the problem.
\end{example}

\subsubsection{Conceptual progression: from simple to complex}

To build intuition, let us consider the natural progression from simple cases toward full generality:

\begin{enumerate}
    \item \textbf{Classical CNNs}: A single transformation (translation) valid globally
    \item \textbf{Group-CNNs}: A finite set of transformations valid globally
    \item \textbf{CNNs on Manifolds}: Transformations dictated by the specific geometry of the manifold
    \item \textbf{Gauge CNNs}: Complete freedom to choose local transformations, subject to consistency
\end{enumerate}

\textbf{Why not simply data augmentation?} One might ask: ``Why not simply augment the data with multiple local transformations?'' The answer is fundamental:

\begin{itemize}
    \item \textbf{Data augmentation}: Hopes to achieve invariance \textit{statistically} across multiple samples
    \item \textbf{Gauge equivariance}: Guarantees \textit{exact} invariance through the architecture
    \item \textbf{Efficiency}: Gauge equivariance achieves with a single forward pass what data augmentation requires multiple evaluations for
    \item \textbf{Consistency}: Guarantees that the network cannot learn to exploit ``accidental'' transformations that do not reflect the true structure of the problem
\end{itemize}

\subsection{Concrete Example: Gauge Equivariance in Action}\label{subsec:ejemplo_gauge_equivarianza}

To make these abstract concepts concrete, let us develop a simple but illustrative example that shows how gauge transformations work in practice.

\subsubsection{Scenario: Pattern recognition on a cylindrical surface}

Consider the problem of recognizing textures on a cylinder (such as a tree trunk). The data consist of an image ``wrapped'' around the cylinder.

\begin{example}[Local Coordinate Problem on a Cylinder]
Suppose we have a tree-bark texture photographed at different positions around the trunk:
\begin{itemize}
    \item \textbf{Position A}: Photo taken facing north
    \item \textbf{Position B}: Photo taken facing east
    \item \textbf{Position C}: Photo taken facing south
\end{itemize}

At each position, we establish local coordinates $(u,v)$ where:
\begin{itemize}
    \item $u$: ``horizontal'' direction relative to the camera orientation
    \item $v$: ``vertical'' direction (along the trunk axis)
\end{itemize}

\textbf{The problem}: The same physical bark patterns appear as completely different tensors in each local coordinate system.
\end{example}

\subsubsection{Gauge solution: coherent local transformations}

\begin{definition}[Simple Gauge Transformation]
For our cylinder, a \textbf{gauge transformation} is a function $g: S^1 \times \mathbb{R} \to SO(2)$ that assigns to each point $(\theta, z)$ of the cylinder a local rotation in the tangent plane:
\begin{equation}
g(\theta, z) = \begin{pmatrix}
\cos(\phi(\theta, z)) & -\sin(\phi(\theta, z)) \\
\sin(\phi(\theta, z)) & \cos(\phi(\theta, z))
\end{pmatrix}
\end{equation}
where $\phi(\theta, z)$ can vary smoothly over the surface.
\end{definition}

\begin{example}[Gauge Convolution on the Cylinder]\label{ex:convolucion_gauge_cilindro}
A gauge equivariant convolution in this scenario would be:
\begin{equation}
[f \star \psi]_\text{gauge}(\theta, z) = \int_{N(\theta,z)} \psi\left(g(\theta', z')^{-1} \cdot \text{connect}((\theta,z), (\theta',z'))\right) f(\theta', z') \, d\theta' dz'
\end{equation}

where:
\begin{itemize}
    \item $N(\theta,z)$ is a neighborhood of the point $(\theta,z)$
    \item $\text{connect}((\theta,z), (\theta',z'))$ describes how to ``connect'' the local coordinate systems
    \item $\psi$ is a kernel that operates in the space of ``orientation differences''
\end{itemize}

\textbf{Key property}: If we apply different choices of local orientation $g(\theta,z)$ at each point, the final result of the convolution remains invariant.
\end{example}

\subsubsection{Visualization of information transport}

\begin{figure}[htbp]
	\centering
	\includegraphics[width=0.95\textwidth]{figures/gauge_transport_cylinder.png}
	\caption[Gauge transport on a cylinder]{Gauge transport on a cylinder. Left: local coordinate systems (tangent vectors) at different points of the cylinder. Center: convolution kernels transported along a geodesic between neighboring points. Right: the concept of gauge invariance: different choices of local orientation (Gauge 1 vs. Gauge 2) produce the same detected pattern.}
	\label{fig:gauge_transport_cylinder}
\end{figure}

\textbf{Key Intuition: Separation of Shape and Orientation.} Gauge equivariance allows the network to learn about:
\begin{itemize}
    \item \textbf{Intrinsic shape}: Which patterns are present (independent of orientation)
    \item \textbf{Spatial relations}: How patterns relate geometrically
\end{itemize}
while completely ignoring:
\begin{itemize}
    \item \textbf{Absolute orientation}: Arbitrary choices of ``north'', ``up'', etc.
    \item \textbf{Coordinate systems}: Labeling conventions that do not affect the physics
\end{itemize}

\begin{proposition}[Gauge Equivariance on the Cylinder]\label{prop:gauge_equivarianza_cilindro}
Let $f: S^1 \times \mathbb{R} \to \mathbb{R}^d$ be a feature field over the cylinder and let $g: S^1 \times \mathbb{R} \to SO(2)$ be an arbitrary gauge transformation. The gauge convolution $[f \star \psi]_{\text{gauge}}$ defined in \autoref{ex:convolucion_gauge_cilindro} satisfies:
\begin{equation}
[(\rho(g) \cdot f) \star \psi]_{\text{gauge}} = \rho(g) \cdot [f \star \psi]_{\text{gauge}}
\end{equation}
where $(\rho(g) \cdot f)(\theta, z) = g(\theta,z) \cdot f(\theta, z)$ is the pointwise gauge action.
\end{proposition}

\begin{proof}
Substituting the definition:
\begin{align}
[(\rho(g) \cdot f) \star \psi]_{\text{gauge}}(\theta, z) &= \int_{N(\theta,z)} \psi\left(g(\theta',z')^{-1} \cdot c_{(\theta,z) \to (\theta',z')}\right) g(\theta',z') f(\theta',z') \, d\theta' dz'
\end{align}
where $c_{(\theta,z) \to (\theta',z')}$ denotes the connection (parallel transport) between the local frames. Since $g(\theta',z')^{-1} g(\theta',z') = \operatorname{Id}$, the gauge transformation of the field $f$ cancels with the inverse transformation in the argument of $\psi$, leaving the result transformed by $g(\theta,z)$ at the output point.
\end{proof}

\begin{remark}[Analogy with Gauge Theories in Physics]
The mathematical structure of this proposition is identical to that of gauge theories in particle physics. In electrodynamics, the vector potential $A_\mu$ (analogous to our field $f$) transforms under a gauge change $A_\mu \to A_\mu + \partial_\mu \chi$, but the field tensor $F_{\mu\nu} = \partial_\mu A_\nu - \partial_\nu A_\mu$ (analogous to the output of the convolution) is gauge-invariant. The gauge convolution of Cohen et al. implements exactly this principle in the context of neural networks.
\end{remark}

\subsection{Mathematical Foundations of Gauge Theory}\label{subsec:fundamentos_matematicos_gauge}

The fundamental distinction between global and local symmetries is formalized as follows:

\begin{definition}[Local vs. Global Transformation]
For a group $G$ and a function $f: \mathcal{M} \to \mathbb{R}^d$:

\textbf{Global Transformation}: $T_g[f](x) = \rho(g) f(g^{-1} \cdot x)$, where $g \in G$ is the same element for all points $x \in \mathcal{M}$.

\textbf{Local (Gauge) Transformation}: $T_{g(\cdot)}[f](x) = \rho(g(x)) f(\phi_{g(x)}(x))$, where $g: \mathcal{M} \to G$ assigns to each point $x$ a transformation $g(x) \in G$.
\end{definition}

The freedom to choose independent local transformations raises a fundamental problem: how do we compare information between points with different coordinate systems? The answer comes from the theory of principal bundles (\autoref{def:fibrado_principal_formal}, \autoref{ch:gdl}), which associates to each point of the space all its possible local orientations. In our cylinder example, the bundle is $P = (S^1 \times \mathbb{R}) \times SO(2)$, where each point $(\theta, z, R) \in P$ represents a position with local orientation $R$.

\textbf{Connections} (\autoref{def:conexion_fibrado}) solve the comparison problem by specifying how to transport orientations between neighboring points in a consistent way, and \textbf{parallel transport} (\autoref{def:transporte_paralelo_formal}) allows objects in different fibers to be compared. For a complete exposition of this formalism, see \autoref{sec:gauges_invarianzas}.

\begin{proposition}[Topological Obstruction to Global Gauges]\label{prop:obstruccion_gauge_global}
Let $\mathcal{M}$ be a differentiable manifold and $P \to \mathcal{M}$ a principal bundle with structure group $G$. The bundle admits a global section (that is, a consistent choice of gauge over all of $\mathcal{M}$) if and only if $P$ is trivial ($P \cong \mathcal{M} \times G$). In particular:
\begin{enumerate}
    \item The frame bundle of $S^2$ does not admit a global section (by the hairy ball theorem).
    \item Every principal bundle over a contractible manifold is trivial.
\end{enumerate}
\end{proposition}

\begin{proof}
A global section $s: \mathcal{M} \to P$ induces the trivialization $\mathcal{M} \times G \to P$ given by $(x, g) \mapsto s(x) \cdot g$, which is a bundle isomorphism by the free and transitive action of $G$ on each fiber. Conversely, a trivialization $P \cong \mathcal{M} \times G$ provides the section $s(x) = (x, e)$.

For (1), a section of the orthonormal frame bundle $\operatorname{SO}(\mathcal{M})$ of $S^2$ is equivalent to a continuous field of orthonormal bases, which in particular implies a nowhere-vanishing tangent vector field, contradicting the hairy ball theorem (the Euler characteristic $\chi(S^2) = 2 \neq 0$).

For (2), every contractible manifold is homotopy equivalent to a point, hence every principal bundle over it is trivial (since bundles are classified by homotopy classes of maps to the classifying space $BG$).
\end{proof}

\begin{remark}[Implication for Neural Networks]
\autoref{prop:obstruccion_gauge_global} explains why gauge equivariance is \textit{necessary} (not merely convenient) for neural networks on general surfaces: on the sphere $S^2$ and on most manifolds of interest, there is no consistent global coordinate system. Any global parametrization has singularities (such as the poles in spherical coordinates). Gauge equivariance allows the network to be independent of these local choices, avoiding artifacts at the singularities.
\end{remark}

\subsubsection{Toward practical implementation}

The gap between theory (infinite-dimensional bundles with smooth connections) and practice (discrete neural networks with finite parameters) is significant. The pioneering work of Cohen et al.\ showed how to find discretizations that preserve the essential theoretical properties. The key lies in the fact that gauge equivariance imposes \textit{linear} constraints on the network weights, reducing the parameter space to an affine subspace that can be solved analytically before training.

\subsection{Gauge Equivariant CNNs: The Model of Cohen et al.}\label{subsec:modelo_cohen_gauge}

The seminal work of Cohen et al.~\citep{cohen2019gaugeequivariantconvolutional} represents the first success in translating abstract gauge theory into a practical and efficient neural architecture. Their key contribution was to show how to implement exact gauge equivariance while maintaining the basic structure of convolutions.

\subsubsection{Design philosophy: from the abstract to the concrete}

The approach of Cohen et al. follows a clear design philosophy:
\begin{enumerate}
    \item \textbf{Principle}: Identify the fundamental geometric properties that the architecture must preserve
    \item \textbf{Discretization}: Find a finite representation that captures these properties
    \item \textbf{Implementation}: Translate into operations that can be executed efficiently
    \item \textbf{Verification}: Demonstrate that the implementation respects the theoretical properties
\end{enumerate}

The concepts of curvature (\autoref{def:curvatura_formal}) and parallel transport (\autoref{def:transporte_paralelo_formal}), fundamental to the construction of gauge equivariant convolutions, are developed formally in \autoref{sec:gauges_invarianzas}. Curvature measures the ``integrability defect'' of the connection (how much parallel transport depends on the chosen path), while parallel transport provides the mechanism for comparing elements in different fibers.

\subsubsection{The specific case: icosahedral CNNs}

To make these abstract ideas concrete, Cohen et al. focus on a specific case that combines theoretical richness with computational feasibility.

\begin{example}[Why the Icosahedron?]
The choice of the icosahedron as a domain is not arbitrary:

\textbf{Geometric properties}:
\begin{itemize}
    \item \textbf{Best spherical approximation}: Among the regular polyhedra, it is the one that best approximates the sphere
    \item \textbf{Rich symmetry}: Symmetry group $I_h$ with 120 elements
    \item \textbf{Geodesic regularity}: Approximately uniform distances between vertices
    \item \textbf{Canonical triangulation}: Well-defined mesh structure
\end{itemize}

\textbf{Computational properties}:
\begin{itemize}
    \item \textbf{Manageable size}: 12 vertices, 20 faces, 30 edges
    \item \textbf{Systematic subdivision}: Refinement by geodesic subdivision
    \item \textbf{Barycentric coordinates}: Smooth interpolation on each triangular face
    \item \textbf{Efficient implementation}: Optimized data structure
\end{itemize}
\end{example}

\begin{definition}[Icosahedral Discretization]\label{def:discretizacion_icosaedrica}
The icosahedral discretization of the sphere $S^2$ is obtained by means of:

\textbf{Geometric construction}:
\begin{enumerate}
    \item \textbf{Inscribed icosahedron}: 12 vertices located on $S^2$ with maximal symmetry
    \item \textbf{Radial projection}: Each point of the sphere is assigned to the nearest vertex
    \item \textbf{Geodesic refinement}: Subdivision of edges by great circles
\end{enumerate}

\textbf{Parametrization}: Each point in $S^2$ is represented as:
\begin{equation}
p = \sum_{i=1}^{12} w_i v_i, \quad \sum_{i=1}^{12} w_i = 1, \quad w_i \geq 0
\end{equation}
where $\{v_i\}$ are the icosahedral vertices and $\{w_i\}$ are barycentric coordinates.
\end{definition}

\begin{definition}[Local Gauge Transformation]\label{def:transformacion_gauge_local}
In the icosahedral discretization, the local gauge transformations are defined as:
\begin{equation}
h: S^2 \to SO(2)
\end{equation}
where $h(p) \in SO(2)$ specifies the transformation at the point $p$.

\textbf{Regularity constraint}: To guarantee differentiability, it is required that:
\begin{equation}
\|h(p) - h(q)\| \leq C d(p,q)^\alpha
\end{equation}
for some constant $C$ and Hölder exponent $\alpha \in (0,1]$.
\end{definition}

\begin{definition}[Gauge Equivariant Function]\label{def:funcion_gauge_equivariante}
A function $f: S^2 \to \mathbb{R}^n$ is said to be gauge equivariant if for every gauge transformation $h$:
\begin{equation}
f^h(p) = \rho(h(p)^{-1}) f(p)
\end{equation}
where $\rho: SO(2) \to GL(\mathbb{R}^n)$ is a representation of the group.

\textbf{Important special cases}:
\begin{itemize}
    \item \textbf{Scalars} ($n=1$): $\rho$ trivial, $f^h = f$
    \item \textbf{Vectors} ($n=2$): $\rho$ is the standard representation of $SO(2)$
    \item \textbf{Complex fields}: Multiplication by $e^{i\theta}$ for rotation by angle $\theta$
\end{itemize}
\end{definition}

\begin{theorem}[Gauge Equivariant Convolution]\label{thm:convolucion_gauge_equivariante}
For gauge equivariant functions $f, g: S^2 \to \mathbb{R}^n$, their gauge convolution is defined as:
\begin{equation}
(f *_G g)(p) = \int_{S^2} f(q) \cdot G(q, p) \cdot g(q) \, d\sigma(q)
\end{equation}
where $G(q,p)$ is the gauge propagator that encodes the parallel transport from $q$ to $p$.

If $f$ and $g$ are gauge equivariant, then $f *_G g$ is also gauge equivariant.
\end{theorem}

\begin{proof}
Let $\rho_p$ be a gauge change at the point $p$ (a change of local basis of the tangent space), so that gauge equivariant functions transform as $f(p) \mapsto \rho_p \cdot f(p)$. The gauge propagator $G(q,p)$ encodes the parallel transport from $q$ to $p$ and transforms as $G(q,p) \mapsto \rho_p \cdot G(q,p) \cdot \rho_q^{-1}$ (since it changes from the basis at $q$ to the basis at $p$). Then:
\begin{align}
(f *_G g)(p) &\mapsto \int_{S^2} (\rho_q f(q)) \cdot (\rho_p G(q,p) \rho_q^{-1}) \cdot (\rho_q g(q)) \, d\sigma(q) \\
&= \rho_p \int_{S^2} f(q) \cdot G(q,p) \cdot g(q) \, d\sigma(q) = \rho_p \cdot (f *_G g)(p)
\end{align}
where we used that $\rho_q^{-1} \rho_q = \operatorname{Id}$. Thus, $f *_G g$ transforms as $\rho_p \cdot (f *_G g)(p)$, which is the gauge equivariant transformation law.
\end{proof}

\textbf{Implementation of the gauge propagator}: In practice, $G(q,p)$ is implemented by means of:

\begin{algorithm}[H]
\caption{Computation of the Icosahedral Gauge Propagator}
\begin{algorithmic}[1]
    \State \textbf{Input:} Points $q, p \in S^2$, gauge connection $\omega$
    \State Find geodesic path $\gamma: q \to p$ in the discretization
    \State Initialize group element $g = I \in SO(2)$
    \For{each edge $e$ in the path $\gamma$}
        \State Update: $g \leftarrow g \cdot \exp(-\omega_e)$
    \EndFor
    \State \textbf{Output:} $G(q,p) = g$
\end{algorithmic}
\end{algorithm}

\subsubsection{Simplified theoretical construction}

\begin{theorem}[Simplified Construction of Gauge CNNs]
To make the ideas more accessible, we present a simplified construction that retains the essential elements:

\textbf{Input data}: A function $f: V \to \mathbb{R}^d$ defined on the icosahedral vertices $V = \{v_1, \ldots, v_{12}\}$.

\textbf{Local gauge convolution}: At each vertex $v_i$:
\begin{equation}
(f * \psi)(v_i) = \sum_{j \in \mathcal{N}(i)} \psi_{ij} \cdot \Gamma_{ij} \cdot f(v_j)
\end{equation}
where:
\begin{itemize}
    \item $\mathcal{N}(i)$ are the neighboring vertices of $v_i$
    \item $\psi_{ij}$ are learnable kernel weights
    \item $\Gamma_{ij} \in SO(2)$ is the gauge transport from $v_j$ to $v_i$
\end{itemize}

\textbf{Learning the connection}: The transport elements $\Gamma_{ij}$ are parametrized as:
\begin{equation}
\Gamma_{ij} = \exp(i\theta_{ij}), \quad \theta_{ij} \in [0, 2\pi)
\end{equation}
and the angles $\{\theta_{ij}\}$ are learnable parameters of the network.

\textbf{Gauge regularization}: To avoid degenerate solutions, a regularization term is added:
\begin{equation}
\mathcal{R}(\theta) = \sum_{\text{triangles } T} \left|\sum_{(i,j) \in \partial T} \theta_{ij}\right|^2
\end{equation}
which penalizes excessive curvature (non-trivial holonomy in small triangles). This regularization is a practical consequence of the Bianchi identity~\eqref{eq:bianchi}.
\end{theorem}

\subsection{General Framework of Principal Bundles}

Beyond the icosahedral example, gauge theory provides a general framework for the design of equivariant neural architectures in complex geometric contexts.

\subsubsection{Classification of gauge structures}

Different problems require different choices of principal bundle, determined by the intrinsic symmetries of the data.

\begin{theorem}[Classification of Bundles for Deep Learning]
Let $\mathcal{M}$ be a manifold that models the data domain and $G$ the relevant symmetry group. The principal bundles appropriate for gauge equivariant architectures are classified according to:

\begin{enumerate}
\item \textbf{Trivial bundles}: $P = \mathcal{M} \times G$ when the symmetries are globally consistent
\item \textbf{Frame bundles}: When invariance under local coordinate changes is required
\item \textbf{Spin bundles}: For data with orientation structure (common in particle physics)
\item \textbf{Gauge bundles}: For fields with internal symmetries (electromagnetism, Yang-Mills fields)
\end{enumerate}
\end{theorem}

\subsubsection{Universal approximation properties}

The rigorous formulation of gauge equivariant convolutions uses parallel transport (\autoref{def:transporte_paralelo_formal}) to compare fibers at different points of the base manifold, preserving the gauge structure.

A fundamental theoretical result establishes that gauge equivariant networks possess universal approximation properties in appropriate contexts.

\begin{theorem}[Universal Gauge Approximation]
Let $\mathcal{F}_{\text{gauge}}$ be the space of functions implementable by gauge equivariant networks with sufficient depth and width. Then:

\begin{enumerate}
\item $\mathcal{F}_{\text{gauge}}$ is dense in $L^2_{\text{gauge}}(P)$, the space of square-integrable gauge equivariant functions
\item The approximation rate is $O(n^{-2/(d_{\mathcal{M}} + d_G)})$ where $d_{\mathcal{M}}$ is the dimension of the base manifold and $d_G$ the dimension of the gauge group
\item The approximation constants depend only on the curvature of the bundle and the regularity properties of the connection
\end{enumerate}

See Section~\ref{sec:aproximacion_universal_gauge} of \autoref{ch:teoremas_complementarios} for the proof sketch.
\end{theorem}

\subsection{Deeper Mathematical Foundations}

The mathematical foundations of gauge theory (including the curvature of the connection (\autoref{def:curvatura_formal}), the Cartan structure equation (\autoref{thm:cartan_estructura}) and holonomy (\autoref{def:holonomia})) are developed formally in \autoref{sec:gauges_invarianzas} of \autoref{ch:gdl}. Here we focus on the applications of these concepts to geometry and physics.

\subsubsection{Mathematical applications in geometry and physics}

The theory of gauge convolutions provides unifying frameworks for modeling phenomena where local symmetries play a fundamental role.

\paragraph{Gauge field theories}
The fundamental gauge fields of theoretical physics are naturally modeled by principal bundles with specific structure groups.

\begin{example}[Quantum Electrodynamics]
Electromagnetism is described by the principal bundle $(P, \mathcal{M}, \pi, U(1))$ where:
\begin{itemize}
\item The base manifold $\mathcal{M}$ represents spacetime
\item The group $U(1)$ encodes the electromagnetic phase transformations
\item The connection $A$ is the electromagnetic vector potential
\item The curvature $F = dA$ is the electromagnetic field tensor
\end{itemize}
\end{example}

\paragraph{Differential geometry}
Gauge structures appear naturally in the study of Riemannian manifolds and their generalizations.

\begin{example}[Orthogonal Frame Bundle]
For a Riemannian manifold $(\mathcal{M}, g)$, the orthonormal frame bundle $O(\mathcal{M})$ is a principal bundle with structure group $O(n)$. Gauge convolutions in this context respect the local isometries of the Riemannian metric.
\end{example}

\paragraph{Algebraic topology}
The characteristic classes of principal bundles provide topological invariants that completely classify gauge structures up to isomorphism.

\begin{example}[Chern Classes]
For complex bundles, the Chern classes $c_k \in H^{2k}(\mathcal{M}, \mathbb{Z})$ provide topological obstructions to the existence of global sections and connections with prescribed curvature.
\end{example}

\subsection{Theoretical Extensions and Recent Developments}

\subsubsection{Gauge CNNs in higher dimensions}

The extensions of gauge theory to higher dimensions and more complex structures represent an active area of research.

\begin{definition}[Higher-Order Principal Bundle]
A \textbf{higher-order principal bundle} is a sequence of bundles:
\begin{equation}
G \to P_n \to P_{n-1} \to \cdots \to P_1 \to \mathcal{M}
\end{equation}
where each level introduces additional symmetries. These structures model physical systems with multiple symmetry scales.
\end{definition}

\paragraph{Applications in unified field theories}
Non-abelian gauge theories (such as quantum chromodynamics) require principal bundles with non-commutative groups:

\begin{example}[Lattice QCD]
Quantum chromodynamics on a lattice is modeled by:
\begin{itemize}
\item Base manifold: Discrete spacetime lattice
\item Gauge group: $SU(3)$ (color group)
\item Fields: Gauge connections on lattice links
\item Gauge structure: Describes the fundamental interactions between quarks and gluons
\end{itemize}
\end{example}

\subsubsection{Connection with generalized Lie algebras}

Recent developments explore the connection between gauge CNNs and more general algebraic structures.

\begin{theorem}[Representation by Lie Algebras]
Every gauge CNN can be represented by a generalized Lie algebra $\mathfrak{L}$ that acts on the feature space. The structure of $\mathfrak{L}$ completely determines the equivariance properties of the network.
\end{theorem}

This representation allows:
\begin{itemize}
\item Systematic classification of gauge structures by algebraic invariants
\item Complete characterization of equivariance properties via representation theory
\item Geometric analysis of gauge moduli spaces
\end{itemize}

\subsubsection{Quantum gauge CNNs}

An emerging frontier explores the extension of gauge principles to the domain of quantum computing.

\begin{definition}[Quantum Gauge CNN]
A \textbf{quantum gauge CNN} is a parametrized quantum circuit that preserves gauge symmetries through unitary operations:
\begin{equation}
U_{\text{gauge}}(\theta) = \exp\left(i \sum_j \theta_j H_j^{\text{gauge}}\right)
\end{equation}
where the $H_j^{\text{gauge}}$ are Hamiltonians that commute with gauge generators.
\end{definition}

The potential applications include:
\begin{itemize}
\item Simulation of quantum gauge systems
\item Data processing in curved Hilbert spaces
\item Quantum algorithms for optimization problems with symmetries
\end{itemize}

\subsubsection{Future perspectives}

The development of gauge convolutions represents only the beginning of a conceptual revolution in geometric deep learning. The promising directions include:

\paragraph{Adaptive gauge theory}
Development of variational methods for the optimal determination of bundle structures through generalized least-action principles, extending the classical Lagrangian formalism.

\paragraph{Geometric renormalization}
Integration of differential renormalization techniques for the systematic treatment of gauge singularities at multiple geometric scales.

\paragraph{Extensions to non-commutative geometries}
Generalization of the gauge formalism to non-commutative manifolds, where the principal bundle structure is replaced by projective modules over operator algebras.

Gauge theory, imported from fundamental physics into machine learning, promises to establish unifying principles for the design of neural architectures that respect the deep symmetries inherent in the structure of our physical and mathematical universe.

\section{Advanced Algebraic Extensions}\label{sec:algebraic_extensions}

After exploring gauge convolutions, which provide an elegant framework for handling local symmetries, we venture into even more abstract and powerful territories. Advanced algebraic extensions represent the most innovative frontiers of \gls{GDL}, where sophisticated mathematical structures such as Clifford algebras and topological methods open new possibilities for the processing of geometric data.

These extensions not only broaden the expressiveness of existing architectures, but also reveal deep connections between abstract algebra, algebraic topology, and machine learning. Their importance lies in providing unifying frameworks that can handle a much wider variety of geometric and structural data than was previously possible.

\subsection{Clifford and Geometric Algebras}

Clifford algebras, also known as geometric algebras, constitute a unified algebraic framework that generalizes complex numbers, quaternions, and other advanced number systems. Their power lies in the ability to represent geometric transformations, rotations, and reflections in a natural and computationally efficient way.

In the context of CNNs, Clifford algebras allow convolutions to be extended to domains where geometric transformations are inherently more complex than simple Euclidean translations. Unlike the complete treatment of Clifford algebras in \autoref{ch:attention} (\autoref{sec:clifford_attention}), where they are studied in the context of attention mechanisms, here we focus on their specific application to convolutional operations and their connection with topological structures.

\subsubsection{Foundations of Geometric Algebras}

\begin{definition}[Clifford Algebra]\label{def:clifford_cnn}
Let $V$ be a finite-dimensional vector space over $\mathbb{R}$ equipped with a quadratic form $Q$. The \textbf{Clifford algebra} $\text{Cl}(V,Q)$ is the unital associative algebra generated by $V$ subject to the relation:
\[v^2 = Q(v) \cdot 1\]
for all $v \in V$.
\end{definition}

In the context of \gls{GDL}, we typically work with Clifford algebras defined by a metric with signature $(p,q)$, denoted $\text{Cl}(p,q)$, where $p$ dimensions have positive norm and $q$ dimensions have negative norm.

\textbf{Relation to Known Structures.}\label{rem:clifford_structures} Clifford algebras unify several familiar algebraic structures:
\begin{itemize}
    \item $\text{Cl}(0,0) \cong \mathbb{R}$ (real numbers)
    \item $\text{Cl}(0,1) \cong \mathbb{C}$ (complex numbers)
    \item $\text{Cl}(0,2) \cong \mathbb{H}$ (Hamilton quaternions; see Section~\ref{sec:cuaterniones} of Appendix~\ref{ch:teoremas_complementarios} for their properties and their relation to $SU(2)$)
    \item $\text{Cl}(3,0) \cong \mathbb{R}^{8}$ (3D Euclidean geometric algebra)
\end{itemize}

\textbf{Fundamental Elements and Grading:}

\begin{definition}[Multivector and Grading]\label{def:multivector_cnn}
The elements of the Clifford algebra are called \textit{multivectors} and admit a unique decomposition by grades. For $\text{Cl}(p,q)$ with $n = p+q$:
\[M = \langle M \rangle_0 + \langle M \rangle_1 + \langle M \rangle_2 + \cdots + \langle M \rangle_n\]
where:
\begin{itemize}
    \item $\langle M \rangle_0 \in \mathbb{R}$: scalar
    \item $\langle M \rangle_1 \in \text{span}\{\mathbf{e}_1, \ldots, \mathbf{e}_n\}$: vector
    \item $\langle M \rangle_2 \in \text{span}\{\mathbf{e}_i \wedge \mathbf{e}_j : i < j\}$: bivector
    \item $\langle M \rangle_k$: $k$-vector with $\binom{n}{k}$ dimensions
    \item $\langle M \rangle_n$: pseudoscalar
\end{itemize}
\end{definition}

\begin{theorem}[Dimensionality of the Algebra]\label{thm:dim_clifford_cnn}
The Clifford algebra $\text{Cl}(p,q)$ is a vector space of dimension $2^{p+q}$, with canonical basis formed by all the wedge products of subsets of $\{\mathbf{e}_1, \ldots, \mathbf{e}_{p+q}\}$:
\[\dim(\text{Cl}(p,q)) = \sum_{k=0}^{p+q} \binom{p+q}{k} = 2^{p+q}\]
\end{theorem}

\begin{proof}
Let $n = p + q$ and $\{\mathbf{e}_1, \ldots, \mathbf{e}_n\}$ the orthonormal basis of $V$. The basis elements of the Clifford algebra are the ordered products $\mathbf{e}_{i_1} \mathbf{e}_{i_2} \cdots \mathbf{e}_{i_k}$ with $1 \leq i_1 < i_2 < \cdots < i_k \leq n$ and $0 \leq k \leq n$ (with the convention that $k = 0$ corresponds to the scalar $1$). These elements are linearly independent, since the relations $\mathbf{e}_i \mathbf{e}_j = -\mathbf{e}_j \mathbf{e}_i$ for $i \neq j$ allow every product to be ordered, and the relations $\mathbf{e}_i^2 = \pm 1$ eliminate squares. The number of subsets of $\{1, \ldots, n\}$ is $\sum_{k=0}^{n} \binom{n}{k} = 2^n$, which is the dimension of the algebra.
\end{proof}

\textbf{Geometric Product:}

The central concept of the Clifford algebra is the \textit{geometric product}, which unifies the inner product and the exterior product.

\begin{definition}[Geometric Product]\label{def:producto_geometrico_cnn}
For vectors $\mathbf{u}, \mathbf{v} \in V$, the \textbf{geometric product} is defined as:
\[\mathbf{u}\mathbf{v} = \mathbf{u} \cdot \mathbf{v} + \mathbf{u} \wedge \mathbf{v}\]
where:
\begin{itemize}
    \item $\mathbf{u} \cdot \mathbf{v} = \frac{1}{2}(\mathbf{u}\mathbf{v} + \mathbf{v}\mathbf{u})$ (inner product, scalar)
    \item $\mathbf{u} \wedge \mathbf{v} = \frac{1}{2}(\mathbf{u}\mathbf{v} - \mathbf{v}\mathbf{u})$ (exterior product, bivector)
\end{itemize}
\end{definition}

\begin{example}[Products in $\text{Cl}(3,0)$]\label{ex:productos_cl30}
Consider the three-dimensional Euclidean algebra $\text{Cl}(3,0)$ with basis $\{\mathbf{e}_1, \mathbf{e}_2, \mathbf{e}_3\}$ where $\mathbf{e}_i^2 = 1$.

\textbf{Products of the basis:}
\begin{align}
\mathbf{e}_1 \mathbf{e}_1 &= 1 \quad \text{(squares give a scalar)} \\
\mathbf{e}_1 \mathbf{e}_2 &= \mathbf{e}_1 \wedge \mathbf{e}_2 \quad \text{(orthogonal vectors)} \\
\mathbf{e}_2 \mathbf{e}_1 &= -\mathbf{e}_1 \wedge \mathbf{e}_2 = -\mathbf{e}_1 \mathbf{e}_2
\end{align}

\textbf{Product of bivector with vector:}
\[(\mathbf{e}_1 \wedge \mathbf{e}_2) \mathbf{e}_3 = \mathbf{e}_1 \mathbf{e}_2 \mathbf{e}_3 = \mathbf{e}_1 \wedge \mathbf{e}_2 \wedge \mathbf{e}_3\]
This is the unit pseudoscalar $I$ in $\text{Cl}(3,0)$.

\textbf{Bivector squared:}
\[(\mathbf{e}_1 \wedge \mathbf{e}_2)^2 = \mathbf{e}_1 \mathbf{e}_2 \mathbf{e}_1 \mathbf{e}_2 = -\mathbf{e}_1 \mathbf{e}_1 \mathbf{e}_2 \mathbf{e}_2 = -1\]
Bivectors behave like imaginary numbers.
\end{example}

\begin{example}[Rotors in 2D]\label{ex:rotores_2d}
In $\text{Cl}(2,0)$, a rotor representing a rotation by angle $\theta$ is:
\[R = \cos(\theta/2) + \sin(\theta/2) \mathbf{e}_1 \wedge \mathbf{e}_2 = e^{(\theta/2) \mathbf{e}_1 \wedge \mathbf{e}_2}\]

To rotate a vector $\mathbf{v}$:
\[\mathbf{v}' = R \mathbf{v} \widetilde{R}\]
where $\widetilde{R}$ is the reverse of $R$ (it reverses the order of multiplication in wedge products).

\textbf{Explicit verification for $\theta = \pi/2$:}
\[R = \frac{1}{\sqrt{2}}(1 + \mathbf{e}_1 \wedge \mathbf{e}_2), \quad \widetilde{R} = \frac{1}{\sqrt{2}}(1 - \mathbf{e}_1 \wedge \mathbf{e}_2)\]
\[R \mathbf{e}_1 \widetilde{R} = \frac{1}{2}(1 + \mathbf{e}_1 \mathbf{e}_2) \mathbf{e}_1 (1 - \mathbf{e}_1 \mathbf{e}_2) = \mathbf{e}_2\]
\end{example}

\begin{example}[Reflections]\label{ex:reflexiones_clifford}
The reflection of a vector $\mathbf{v}$ with respect to a hyperplane perpendicular to $\mathbf{n}$ (unit) is:
\[\mathbf{v}' = -\mathbf{n} \mathbf{v} \mathbf{n}\]

For $\mathbf{n} = \mathbf{e}_1$ and $\mathbf{v} = a\mathbf{e}_1 + b\mathbf{e}_2$:
\[\mathbf{v}' = -\mathbf{e}_1 (a\mathbf{e}_1 + b\mathbf{e}_2) \mathbf{e}_1 = -a\mathbf{e}_1 + b\mathbf{e}_2\]
As expected, it inverts the component perpendicular to the plane.
\end{example}

\subsubsection{Convolutions in Clifford Algebras}

Although GATr is explored extensively in \autoref{ch:attention} (\autoref{sec:clifford_attention}), here we examine how Clifford algebras allow convolutional operations to be defined in complex geometric domains.

\begin{definition}[Clifford-valued Convolution]\label{def:conv_clifford}
Let $f: \Omega \to \text{Cl}(p,q)$ be a signal valued in the Clifford algebra and $\psi: G \to \text{Cl}(p,q)$ a kernel defined over a group of transformations $G$. The \textbf{Clifford-valued convolution} is:
\[(f \star \psi)(x) = \int_G \rho_g^{-1}(\psi(g)) \cdot f(g^{-1} \cdot x) \, dg\]
where $\rho_g$ is the action of the group on the algebra and $\cdot$ is the geometric product.
\end{definition}

\textbf{Advantages over Scalar Convolutions.}\label{rem:ventajas_clifford_conv} Clifford-valued convolutions offer:
\begin{enumerate}
    \item \textbf{Rich representation}: Multivectors encode multiple geometric aspects simultaneously
    \item \textbf{Automatic equivariance}: The algebraic structure guarantees symmetry properties
    \item \textbf{Computational efficiency}: Low-level operations that can be optimized with SIMD
    \item \textbf{Interpretability}: The grades of the multivector have clear geometric meaning
\end{enumerate}

\begin{example}[Steerable CNN with Clifford Algebra]\label{ex:steerable_clifford}
Consider a 2D CNN equivariant to rotations using $\text{Cl}(2,0)$.

\textbf{Feature representation:}
A feature at position $\mathbf{x}$ is a multivector:
\[h(\mathbf{x}) = \alpha_0(\mathbf{x}) + \alpha_1(\mathbf{x}) \mathbf{e}_1 + \alpha_2(\mathbf{x}) \mathbf{e}_2 + \alpha_{12}(\mathbf{x}) \mathbf{e}_1 \wedge \mathbf{e}_2\]

\textbf{Rotation action:}
For rotation by angle $\theta$, the rotor $R = e^{(\theta/2)\mathbf{e}_1 \wedge \mathbf{e}_2}$ acts as:
\[h'(\mathbf{x}) = R h(R^{-1} \mathbf{x}) \widetilde{R}\]

\textbf{Equivariant kernel:}
The kernel $\psi: SO(2) \to \text{Cl}(2,0)$ is parametrized as:
\[\psi(R) = w_0 + w_1 R + w_2 R^2 + w_3 R^3\]
where the $w_i$ are learnable parameters.
\end{example}

\subsubsection{Connection with Steerable CNNs}

Steerable CNNs (see \autoref{subsubsec:implementacion_steerable_cnns}) can be understood as special cases of Clifford-valued CNNs.

\begin{remark}[Steerable-Clifford Connection]\label{thm:equiv_steerable_clifford}
Let $\mathcal{F}$ be a feature space with group action $G$ given by a representation $\rho: G \to GL(\mathcal{F})$. If $\rho$ admits a decomposition into irreducibles that are isomorphic to representations of $\text{Cl}(p,q)$, then the $G$-equivariant linear operations on $\mathcal{F}$ correspond to equivariant linear transformations in $\text{Cl}(p,q)$. The intuition is that the irreducible representations of $SO(n)$ are in correspondence with spaces of multivectors of fixed grade in $\text{Cl}(n,0)$, and an equivariant linear transformation preserves this grade structure, which corresponds to the block structure of steerable CNNs. The complete formalization of the conditions under which this correspondence is bijective remains an active area of research \citep{brehmer2023geometric}.
\end{remark}

\subsubsection{Applications to CNNs on Manifolds}

Clifford algebras provide a natural framework for CNNs on Riemannian manifolds.

\begin{definition}[Clifford Bundle]\label{def:clifford_bundle}
Let $(\mathcal{M}, g)$ be a Riemannian manifold. The \textbf{Clifford bundle} $\text{Cl}(T\mathcal{M})$ is the vector bundle whose fiber at $p \in \mathcal{M}$ is the Clifford algebra $\text{Cl}(T_p\mathcal{M}, g_p)$ of the tangent space with the Riemannian metric.
\end{definition}

\textbf{Sections as Features.}\label{rem:secciones_clifford} The features in a CNN over a manifold are sections of the Clifford bundle. A section $s: \mathcal{M} \to \text{Cl}(T\mathcal{M})$ assigns to each point $p$ a multivector $s(p) \in \text{Cl}(T_p\mathcal{M})$.

\begin{example}[CNN on a Sphere with Clifford]\label{ex:cnn_sphere_clifford}
For the sphere $S^2 \subset \mathbb{R}^3$, we use the ambient algebra $\text{Cl}(3,0)$.

\textbf{Embedding:}
Each point $\mathbf{p} \in S^2$ is represented as a vector $\mathbf{p} = p_1 \mathbf{e}_1 + p_2 \mathbf{e}_2 + p_3 \mathbf{e}_3$ with $\|\mathbf{p}\| = 1$.

\textbf{Tangent space:}
$T_\mathbf{p} S^2 = \{\mathbf{v} \in \mathbb{R}^3 : \mathbf{v} \cdot \mathbf{p} = 0\}$

\textbf{Tangential features:}
A feature is a multivector with vector part tangent to $S^2$:
\[h(\mathbf{p}) = \alpha(\mathbf{p}) + \mathbf{v}(\mathbf{p}) + \beta(\mathbf{p}) \mathbf{p} \wedge \mathbf{v}(\mathbf{p})\]
where $\mathbf{v}(\mathbf{p}) \in T_\mathbf{p} S^2$.

\textbf{Geodesic convolution:}
\[(\psi \star h)(\mathbf{p}) = \int_{S^2} \psi(\mathbf{p}, \mathbf{q}) \cdot h(\mathbf{q}) \, d\sigma(\mathbf{q})\]
where $\psi(\mathbf{p}, \mathbf{q})$ depends on the geodesic distance and relative orientation.
\end{example}

\subsubsection{Choice of Algebras: Euclidean, Projective, and Conformal}

The work of \citet{dehaan2024euclidean} provides a systematic analysis for the optimal selection of geometric algebras in different contexts.

\textbf{Euclidean Algebra $\mathcal{Cl}(3,0)$:}
\begin{itemize}
    \item \textbf{Dimensionality:} 8 dimensions
    \item \textbf{Advantages:} Computationally efficient, direct interpretation
    \item \textbf{Limitations:} Symmetry group restricted to $SO(3)$
    \item \textbf{Applications:} Cases with limited computational resources
\end{itemize}

\textbf{Projective Algebra $\mathcal{Cl}(4,1)$:}
\begin{itemize}
    \item \textbf{Dimensionality:} 32 dimensions
    \item \textbf{Representation:} Includes homogeneous coordinates
    \item \textbf{Limitations:} Insufficient expressiveness in the standard form
    \item \textbf{Improved version:} Modifications that increase the representational capacity
\end{itemize}

\textbf{Conformal Algebra $\mathcal{Cl}(4,1)$ with conformal metric:}
\begin{itemize}
    \item \textbf{Dimensionality:} 32 dimensions
    \item \textbf{Capabilities:} Natural representation of spheres, circles, inversions
    \item \textbf{Symmetry group:} Full conformal group
    \item \textbf{Performance:} Superior in empirical experiments
    \item \textbf{Cost:} Greater computational complexity
\end{itemize}

\textbf{Selection Criteria:}

The choice of the appropriate algebra must consider:
\begin{enumerate}
    \item \textbf{Required expressiveness:} Complexity of the necessary geometric transformations
    \item \textbf{Computational resources:} Available memory and computation time
    \item \textbf{Target precision:} Tolerance to errors in the specific application
    \item \textbf{Data characteristics:} Types of symmetries present in the domain
\end{enumerate}

\subsubsection{Metric Learning for Clifford Groups}

The work of \citet{ali2024metric} introduces a fundamental innovation: the adaptive learning of metrics in Clifford group equivariant neural networks (CGENNs).

\textbf{Motivation and Problem:}

Traditional equivariant architectures depend on predefined metrics (Euclidean, Minkowski) that may be suboptimal for specific domains. Metric learning allows the architecture to automatically adapt to the geometric characteristics of the data.

\textbf{Mathematical Formulation:}

For a Clifford algebra with learnable metric $G_\theta$:
\[G_\theta = \sum_{i=1}^n \lambda_i(\theta) \mathbf{v}_i \mathbf{v}_i^T\]
where $\lambda_i(\theta)$ are learnable eigenvalues and $\mathbf{v}_i$ eigenvectors that preserve the symmetric structure.

\textbf{Preservation of Equivariance:}

To maintain the equivariance properties of the Clifford group $g$:
\[G_\theta(g \cdot x) = \rho_g(G_\theta(x))\]
where $\rho_g$ is the representation of $g$ in the space of metrics.

\textbf{Integration into CGENNs:}

The equivariant layers incorporate the learnable metric:
\[h^{(l+1)} = \sigma\left(\sum_i W_i^{(l)} h^{(l)} \star G_\theta^{(l)}\right)\]
where $\star$ denotes the product in the Clifford algebra modified by the adaptive metric.

\textbf{Observed Advantages:}
\begin{itemize}
    \item \textbf{Adaptability:} Metrics that automatically adjust to the characteristics of the data
    \item \textbf{Improvement in precision:} Superior performance compared to fixed metrics
    \item \textbf{Generalization:} Better transfer to data with different geometric characteristics
\end{itemize}

\subsection{Topological Extensions}

Beyond algebraic structures, \gls{GDL} extends toward topological domains that capture structural relations more complex than traditional graphs. These topological methods represent an emerging paradigm that promises to revolutionize the processing of complex relational data.

Topological CNNs differ fundamentally from traditional CNNs in that they operate over spaces with complex combinatorial structure, where the neighborhood relations do not reduce to simple graphs but instead involve cells of multiple dimensions with boundary and incidence relations.

\subsubsection{Foundations of Algebraic Topology}

Before addressing deep topological architectures, we establish the necessary mathematical foundations.

We recall that a \textit{simplicial complex} $\mathcal{K}$ is a collection of finite sets closed downward and with proper intersection (\autoref{def:complejo_simplicial} of \autoref{ch:gdl}).\label{def:complejo_simplicial_cnn}

\begin{example}[Simplices in Low Dimensions]\label{ex:simplicios_bajos}
\leavevmode
\begin{itemize}
    \item \textbf{0-simplex}: A point $\{v\}$
    \item \textbf{1-simplex}: An edge $\{v_0, v_1\}$ with endpoints $\{v_0\}, \{v_1\}$
    \item \textbf{2-simplex}: A triangle $\{v_0, v_1, v_2\}$ with edges $\{v_0,v_1\}, \{v_0,v_2\}, \{v_1,v_2\}$
    \item \textbf{3-simplex}: A tetrahedron $\{v_0, v_1, v_2, v_3\}$ with 4 triangular faces
\end{itemize}
\end{example}

\begin{definition}[Cell Complexes]\label{def:complejo_celular}
A \textbf{cell complex} (or CW-complex) $X$ is a topological space constructed inductively by gluing cells $e^n$ (homeomorphic to open balls $\mathbb{D}^n$) by means of characteristic maps $\phi: \partial \mathbb{D}^n \to X^{n-1}$ that map the boundary of the cell to the $(n-1)$-dimensional skeleton.
\end{definition}

\textbf{Generalization of Simplicial Complexes.}\label{rem:celular_vs_simplicial} Cell complexes generalize simplicial complexes:
\begin{itemize}
    \item Cells of dimension $n$ can have any shape (not just simplices)
    \item Example: a cube is a 3-dimensional cell in a cell complex, but requires multiple tetrahedra in a simplicial representation
    \item Computational advantage: More compact representations for certain domains
\end{itemize}

\subsubsection{Boundary Operators and Chains}

The fundamental algebraic structure of topological complexes is captured by the boundary operators.

\begin{definition}[Chains and Boundary Operator]\label{def:cadenas_frontera}
Let $\mathcal{K}$ be a simplicial complex. The \textbf{space of $k$-chains} is:
\[C_k(\mathcal{K}) = \left\{\sum_{i} a_i \sigma_i : a_i \in \mathbb{R}, \sigma_i \text{ is a } k\text{-simplex of } \mathcal{K}\right\}\]

The \textbf{boundary operator} $\partial_k: C_k(\mathcal{K}) \to C_{k-1}(\mathcal{K})$ is defined on a $k$-simplex $\sigma = [v_0, v_1, \ldots, v_k]$ as:
\[\partial_k \sigma = \sum_{i=0}^k (-1)^i [v_0, \ldots, \hat{v}_i, \ldots, v_k]\]
where $\hat{v}_i$ indicates the omission of $v_i$, and it extends linearly.
\end{definition}

\begin{theorem}[Fundamental Property: $\partial^2 = 0$]\label{thm:partial_squared}
For all $k$, it holds that $\partial_{k-1} \circ \partial_k = 0$, that is, the boundary of a boundary is empty.
\end{theorem}

\begin{proof}
For a $k$-simplex $\sigma = [v_0, \ldots, v_k]$:
\begin{align}
\partial_{k-1}(\partial_k \sigma) &= \partial_{k-1}\left(\sum_{i=0}^k (-1)^i [v_0, \ldots, \hat{v}_i, \ldots, v_k]\right) \\
&= \sum_{i=0}^k (-1)^i \sum_{j<i} (-1)^j [v_0, \ldots, \hat{v}_j, \ldots, \hat{v}_i, \ldots, v_k] \\
&\quad + \sum_{i=0}^k (-1)^i \sum_{j>i} (-1)^{j-1} [v_0, \ldots, \hat{v}_i, \ldots, \hat{v}_j, \ldots, v_k]
\end{align}
Each term appears exactly twice with opposite signs, cancelling out.
\end{proof}

\begin{example}[Boundary Operator on a Triangle]\label{ex:frontera_triangulo}
For the 2-simplex $\sigma = [v_0, v_1, v_2]$:
\[\partial_2 [v_0, v_1, v_2] = [v_1, v_2] - [v_0, v_2] + [v_0, v_1]\]

Geometric interpretation: The boundary of the triangle consists of its three edges with coherent orientations.

Verification of $\partial^2 = 0$:
\begin{align}
\partial_1(\partial_2 [v_0, v_1, v_2]) &= \partial_1([v_1, v_2]) - \partial_1([v_0, v_2]) + \partial_1([v_0, v_1]) \\
&= ([v_2] - [v_1]) - ([v_2] - [v_0]) + ([v_1] - [v_0]) \\
&= 0
\end{align}
\end{example}

\begin{figure}[htbp]
	\centering
	\includegraphics[width=0.85\textwidth]{figures/simplicial_hierarchy}
	\caption[Simplicial hierarchy: multi-order message passing]{Message passing hierarchy in a simplicial complex: messages between nodes (0-chains), edges (1-chains) and faces (2-chains), connected by boundary $\partial_k$ and coboundary $\delta_k$ operators. The Hodge Laplacian $L_k$ captures the interaction between adjacent levels.}
	\label{fig:simplicial_hierarchy}
\end{figure}

\subsubsection{Hodge Laplacians}

Hodge Laplacians are fundamental operators that extend the graph Laplacian to higher-dimensional complexes (for the general definition and its properties, see \autoref{def:hodge_laplacian_gdl} and \autoref{prop:propiedades_hodge_gdl} of \autoref{ch:gdl}).

\begin{definition}[Coboundary Operator]\label{def:cofrontera}
The \textbf{coboundary operator} $\delta_k: C_k \to C_{k+1}$ is the formal adjoint of $\partial_{k+1}$:
\[\delta_k = \partial_{k+1}^T\]
with respect to the standard inner product on chain spaces.
\end{definition}

\begin{definition}[Hodge Laplacians]\label{def:hodge_laplacians_cnn}
The \textbf{Hodge Laplacian} of order $k$ is:
\[L_k = \partial_{k+1} \partial_{k+1}^T + \partial_k^T \partial_k = \delta_{k-1} \delta_{k-1}^T + \delta_k^T \delta_k\]

Decomposition into two terms:
\begin{itemize}
    \item $L_k^{\text{down}} = \partial_k^T \partial_k$: Lower Laplacian (flow from $(k-1)$-chains)
    \item $L_k^{\text{up}} = \partial_{k+1} \partial_{k+1}^T$: Upper Laplacian (flow toward $(k+1)$-chains)
\end{itemize}
\end{definition}

\begin{theorem}[Properties of the Hodge Laplacian]\label{thm:prop_hodge_laplacian}
The Laplacian $L_k$ satisfies:
\begin{enumerate}
    \item \textbf{Symmetry}: $L_k$ is symmetric and positive semidefinite
    \item \textbf{Kernel and homology}: $\ker(L_k) \cong H_k(\mathcal{K})$ (homology group of dimension $k$)
    \item \textbf{Hodge decomposition}:
    \[C_k = \text{Im}(\partial_{k+1}) \oplus \ker(L_k) \oplus \text{Im}(\partial_k^T)\]
\end{enumerate}
\end{theorem}

\begin{proof}
(1) \textit{Symmetry}: $L_k = \partial_k^T \partial_k + \partial_{k+1} \partial_{k+1}^T$ is a sum of matrices of the form $A^T A$, which are symmetric and positive semidefinite. To see that $L_k$ is positive semidefinite: $\langle L_k \mathbf{x}, \mathbf{x} \rangle = \|\partial_k \mathbf{x}\|^2 + \|\partial_{k+1}^T \mathbf{x}\|^2 \geq 0$.

(2) \textit{Kernel and homology}: $L_k \mathbf{x} = 0$ if and only if $\|\partial_k \mathbf{x}\|^2 + \|\partial_{k+1}^T \mathbf{x}\|^2 = 0$, that is, $\partial_k \mathbf{x} = 0$ and $\partial_{k+1}^T \mathbf{x} = 0$. The first condition says that $\mathbf{x}$ is a $k$-cycle ($\mathbf{x} \in \ker \partial_k$) and the second that $\mathbf{x}$ is orthogonal to the image of $\partial_{k+1}$ ($\mathbf{x} \perp \operatorname{Im} \partial_{k+1}$). Hence $\ker L_k = \ker \partial_k \cap (\operatorname{Im} \partial_{k+1})^\perp$, which is the harmonic representative of the homology class $H_k = \ker \partial_k / \operatorname{Im} \partial_{k+1}$.

(3) \textit{Hodge decomposition}: The three subspaces are mutually orthogonal: $\operatorname{Im}(\partial_{k+1}) \perp \operatorname{Im}(\partial_k^T)$ because $\partial_k \partial_{k+1} = 0$ (fundamental property of the boundary operators). The kernel $\ker L_k$ is the orthogonal complement of $\operatorname{Im}(\partial_{k+1}) \oplus \operatorname{Im}(\partial_k^T)$ in $C_k$.
\end{proof}

\begin{example}[Hodge Laplacian for a Triangular Mesh]\label{ex:hodge_triangle_mesh}
Consider a simple triangular mesh with 4 vertices forming two adjacent triangles.

\textbf{Structure:}
\begin{itemize}
    \item Vertices: $v_0, v_1, v_2, v_3$
    \item Edges: $e_1 = [v_0,v_1], e_2 = [v_1,v_2], e_3 = [v_0,v_2], e_4 = [v_1,v_3], e_5 = [v_2,v_3]$
    \item Triangles: $\tau_1 = [v_0,v_1,v_2], \tau_2 = [v_1,v_2,v_3]$
\end{itemize}

\textbf{Boundary matrix $\partial_1$ (edges $\to$ vertices):}
\[\partial_1 = \begin{pmatrix}
-1 & 0 & -1 & 0 & 0 \\
1 & -1 & 0 & -1 & 0 \\
0 & 1 & 1 & 0 & -1 \\
0 & 0 & 0 & 1 & 1
\end{pmatrix}\]

\textbf{Boundary matrix $\partial_2$ (triangles $\to$ edges):}

For $\tau_1 = [v_0,v_1,v_2]$: $\partial_2(\tau_1) = [v_1,v_2] - [v_0,v_2] + [v_0,v_1] = e_2 - e_3 + e_1$.
For $\tau_2 = [v_1,v_2,v_3]$: $\partial_2(\tau_2) = [v_2,v_3] - [v_1,v_3] + [v_1,v_2] = e_5 - e_4 + e_2$.
\[\partial_2 = \begin{pmatrix}
1 & 0 \\
1 & 1 \\
-1 & 0 \\
0 & -1 \\
0 & 1
\end{pmatrix}\]

It is verified that $\partial_1 \partial_2 = 0$, consistent with \autoref{thm:partial_squared}.

\textbf{Edge Laplacian (1-Laplacian):}
\[L_1 = \partial_1^T \partial_1 + \partial_2 \partial_2^T\]

This operator captures both the connectivity between edges via common vertices and via common triangles.
\end{example}

\subsubsection{Topological Deep Learning: Architectures}

The comprehensive work of \citet{hajij2023topological} establishes the theoretical and practical foundations of topological deep learning.

\begin{definition}[Combinatorial Complex]\label{def:complejo_combinatorial}
A \textbf{combinatorial complex} $\mathcal{K}$ is a collection of finite sets (cells) equipped with a partial inclusion relation, which generalizes both simplicial and cell complexes without imposing restrictive boundary conditions.
\end{definition}

This definition generalizes traditional relational data structures:
\begin{itemize}
    \item \textbf{Graphs:} Special case with cells of dimension 0 and 1
    \item \textbf{Hypergraphs:} Subset with hyperedges of variable dimension
    \item \textbf{Simplicial complexes:} Restricted case with specific boundary relations
    \item \textbf{Cell complexes:} Generalizations with cells of arbitrary shapes
\end{itemize}

\textbf{Combinatorial Complex Neural Networks (CCNNs):}

CCNNs operate by means of multi-level message passing between cells of different dimensions:
\[h_\sigma^{(l+1)} = \text{UPDATE}\left(h_\sigma^{(l)}, \text{AGGREGATE}\left(\{h_\tau^{(l)} : \tau \in \mathcal{N}(\sigma)\}\right)\right)\]
where $\mathcal{N}(\sigma)$ represents the neighborhood of the cell $\sigma$, which may include cells of various dimensions.

\begin{definition}[Neighborhoods in Complexes]\label{def:vecindarios_complejos}
For a $k$-cell $\sigma$ in a complex, we define:
\begin{itemize}
    \item \textbf{Lower neighborhood}: $\mathcal{N}_{\text{low}}(\sigma) = \{\tau : \tau \subseteq \sigma, \dim(\tau) = k-1\}$
    \item \textbf{Upper neighborhood}: $\mathcal{N}_{\text{up}}(\sigma) = \{\tau : \sigma \subseteq \tau, \dim(\tau) = k+1\}$
    \item \textbf{Co-dimensional neighborhood}: $\mathcal{N}_{\text{co}}(\sigma) = \{\tau : \dim(\tau) = k, \tau \cap \sigma \neq \emptyset\}$
\end{itemize}
\end{definition}

\textbf{Boundary and Co-boundary Operators:}

The topological structure is captured by means of linear operators derived from the algebraic boundary operators:
\begin{itemize}
    \item \textbf{Boundary operator} $\partial_k: C_k \to C_{k-1}$ (see \autoref{def:cadenas_frontera})
    \item \textbf{Co-boundary operator} $\delta_k = \partial_{k+1}^T: C_k \to C_{k+1}$
\end{itemize}

These operators define the generalized adjacency matrices that govern the message passing.

\textbf{Equivariance Properties:}

CCNNs maintain equivariance under:
\begin{enumerate}
    \item \textbf{Permutations}: Reordering of cells within the same dimension
    \item \textbf{Orientations}: Changes in the orientation of oriented cells
    \item \textbf{Topological transformations}: Homeomorphisms that preserve the structure
\end{enumerate}

\subsubsection{Simplicial Attention Networks}

The work of \citet{giusti2022simplicial} introduces attention mechanisms specifically designed for simplicial complexes, combining topological expressiveness with the flexibility of attention.

\textbf{Simplicial Complexes:}

\begin{definition}[Simplicial Complex]
A \textbf{simplicial complex} $\mathcal{K}$ is a collection of simplices that satisfies:
\begin{enumerate}
    \item If $\sigma \in \mathcal{K}$ and $\tau \subseteq \sigma$, then $\tau \in \mathcal{K}$ (downward closure)
    \item The intersection of two simplices is empty or a common simplex
\end{enumerate}
\end{definition}

\textbf{Simplicial Attention Mechanism:}

For a $k$-simplex $\sigma$ and its neighborhood $N(\sigma)$:
\[\alpha_{\sigma,\tau} = \text{softmax}\left(\text{LeakyReLU}\left(a^T[W_\sigma h_\sigma \| W_\tau h_\tau \| o_{\sigma,\tau}]\right)\right)\]
where:
\begin{itemize}
    \item $h_\sigma, h_\tau$: Representations of simplices
    \item $W_\sigma, W_\tau$: Learnable transformation matrices
    \item $a$: Learnable attention vector
    \item $o_{\sigma,\tau}$: Relative orientation information
\end{itemize}

\textbf{Attention with Orientation (Signed Attention):}

A distinctive feature is the incorporation of orientation information:
\[\alpha'_{\sigma,\tau} = s_{\sigma,\tau} \cdot \alpha_{\sigma,\tau}\]
where $s_{\sigma,\tau} \in \{-1, +1\}$ reflects the relative orientation between simplices.

This oriented attention preserves the cohomological properties of the complex and maintains equivariance under orientation changes:
\[\text{SAT}(\mathcal{O}(X)) = \mathcal{O}(\text{SAT}(X))\]
for an orientation change $\mathcal{O}$.

\textbf{Multi-dimensional Processing:}

SATs process information in multiple dimensions simultaneously:
\begin{itemize}
    \item \textbf{Parallel layers:} Independent processing by dimension
    \item \textbf{Inter-dimensional communication:} Flow of information via boundary operators
    \item \textbf{Hierarchical aggregation:} Topologically-aware pooling
\end{itemize}

\subsubsection{CW Complexes and Applications}

CW (Closure-finite, Weak topology) complexes provide an even more general framework than simplicial complexes, allowing cells of arbitrary shapes attached in a topologically coherent way.

\textbf{Definition and Structure:}

\begin{definition}[CW Complex]
A \textbf{CW complex} is a topological space $X$ together with a decomposition into cells $e^n_\alpha$ of dimension $n$, where each cell is attached to the lower-dimensional skeleton by means of a continuous characteristic map.
\end{definition}

\textbf{Advantages for Deep Learning:}
\begin{enumerate}
    \item \textbf{Geometric flexibility:} Cells of arbitrary shapes
    \item \textbf{Representational efficiency:} Fewer cells than equivalent simplicial representations
    \item \textbf{Computational compatibility:} Regular structure compatible with parallelization
\end{enumerate}

\textbf{Emerging Applications:}
\begin{itemize}
    \item \textbf{Image analysis:} Segmentation based on topological structure
    \item \textbf{Molecular modeling:} Representation of protein complexes
    \item \textbf{Neuroscience:} Analysis of brain connectomes
    \item \textbf{Social network analysis:} Modeling of hierarchical communities
\end{itemize}

\subsection{Connections between Algebra and Topology}

The true power of advanced algebraic extensions emerges from the deep connections between algebraic and topological structures. These connections are not merely theoretical, but provide unifying principles for the design of \gls{GDL} architectures.

\subsubsection{Homology and Cohomology}

Homology and cohomology provide algebraic invariants of topological spaces that are fundamental to \gls{GDL}.

We recall that the $k$-th homology group $H_k(\mathcal{K}) = Z_k / B_k$ measures the ``$k$-dimensional holes'' of the complex (\autoref{def:homologia_betti} of \autoref{ch:gdl}).\label{def:grupos_homologia_cnn}

\textbf{Geometric Interpretation.}\label{rem:interp_homologia} $H_k(\mathcal{K})$ measures the ``$k$-dimensional holes'' of the space:
\begin{itemize}
    \item $H_0$: Connected components
    \item $H_1$: Non-trivial cycles (tunnels, loops)
    \item $H_2$: Cavities (voids)
    \item $H_k$ ($k > 2$): Higher-dimensional holes
\end{itemize}

\begin{definition}[Betti Numbers]\label{def:numeros_betti_cnn}
The \textbf{$k$-th Betti number} is:
\[\beta_k = \dim(H_k(\mathcal{K})) = \text{rank}(Z_k) - \text{rank}(B_k)\]
\end{definition}

\begin{example}[Homology of the Sphere and Torus]\label{ex:homologia_esfera_toro}
\textbf{Sphere $S^2$:}
\begin{align}
H_0(S^2) &\cong \mathbb{R} \quad (\beta_0 = 1, \text{ connected}) \\
H_1(S^2) &= 0 \quad (\beta_1 = 0, \text{ no loops}) \\
H_2(S^2) &\cong \mathbb{R} \quad (\beta_2 = 1, \text{ cavity})
\end{align}

\textbf{Torus $T^2 = S^1 \times S^1$:}
\begin{align}
H_0(T^2) &\cong \mathbb{R} \quad (\beta_0 = 1) \\
H_1(T^2) &\cong \mathbb{R}^2 \quad (\beta_1 = 2, \text{ two independent loops}) \\
H_2(T^2) &\cong \mathbb{R} \quad (\beta_2 = 1)
\end{align}

The Euler characteristic:
\[\chi(S^2) = \beta_0 - \beta_1 + \beta_2 = 1 - 0 + 1 = 2\]
\[\chi(T^2) = \beta_0 - \beta_1 + \beta_2 = 1 - 2 + 1 = 0\]
\end{example}

\begin{definition}[De Rham Cohomology]\label{def:cohomologia_derham}
On differentiable manifolds, the \textbf{De Rham cohomology} $H^k_{\text{dR}}(\mathcal{M})$ is defined as:
\[H^k_{\text{dR}}(\mathcal{M}) = \frac{\ker(d: \Omega^k \to \Omega^{k+1})}{\text{Im}(d: \Omega^{k-1} \to \Omega^k)}\]
where $\Omega^k$ is the space of differential $k$-forms and $d$ is the exterior derivative.
\end{definition}

\begin{theorem}[De Rham Theorem]\label{thm:derham}
For a differentiable manifold $\mathcal{M}$:
\[H^k_{\text{dR}}(\mathcal{M}) \cong H^k(\mathcal{M}; \mathbb{R})\]
The De Rham cohomology is isomorphic to the singular cohomology with real coefficients.

See Section~\ref{sec:de_rham} of \autoref{ch:teoremas_complementarios} for the proof.
\end{theorem}

\subsubsection{Exterior Algebra and Connection with Clifford}

The exterior algebra provides the bridge between Clifford algebras and algebraic topology.

\begin{definition}[Exterior Algebra]\label{def:algebra_exterior}
Let $V$ be a vector space. The \textbf{exterior algebra} $\bigwedge V$ is the graded algebra:
\[\bigwedge V = \bigoplus_{k=0}^n \bigwedge^k V\]
where $\bigwedge^k V$ is the space of $k$-vectors with antisymmetric exterior product $\wedge$:
\[v \wedge w = -w \wedge v, \quad v \wedge v = 0\]
\end{definition}

\textbf{Clifford-Exterior Relation.}\label{rem:clifford_exterior} The Clifford algebra $\text{Cl}(V, Q)$ and the exterior algebra $\bigwedge V$ are related:
\begin{itemize}
    \item As vector spaces: $\text{Cl}(V,Q) \cong \bigwedge V$ (both have dimension $2^n$)
    \item The geometric product extends the exterior product: $uv = u \cdot v + u \wedge v$
    \item $\bigwedge V = \text{Cl}(V, 0)$ (Clifford algebra with null quadratic form)
\end{itemize}

\begin{example}[Basis of the Exterior Algebra $\bigwedge \mathbb{R}^3$]\label{ex:base_exterior_3d}
For $V = \mathbb{R}^3$ with basis $\{e_1, e_2, e_3\}$:
\begin{align}
\bigwedge^0 \mathbb{R}^3 &: \{1\} \quad \text{(scalars)} \\
\bigwedge^1 \mathbb{R}^3 &: \{e_1, e_2, e_3\} \quad \text{(vectors)} \\
\bigwedge^2 \mathbb{R}^3 &: \{e_1 \wedge e_2, e_1 \wedge e_3, e_2 \wedge e_3\} \quad \text{(bivectors)} \\
\bigwedge^3 \mathbb{R}^3 &: \{e_1 \wedge e_2 \wedge e_3\} \quad \text{(pseudoscalar)}
\end{align}
Total: $1 + 3 + 3 + 1 = 8$ dimensions.
\end{example}

\subsubsection{Persistent Homology}

Persistent homology provides a multiscale tool for the topological analysis of data.

\begin{definition}[Filtration]\label{def:filtracion_cnn}
A \textbf{filtration} of a simplicial complex is a nested sequence:
\[\emptyset = \mathcal{K}_0 \subseteq \mathcal{K}_1 \subseteq \mathcal{K}_2 \subseteq \cdots \subseteq \mathcal{K}_m = \mathcal{K}\]
where each $\mathcal{K}_i$ is a subcomplex of $\mathcal{K}$.
\end{definition}

\begin{definition}[Persistence Diagram]\label{def:diagrama_persistencia_cnn}
The \textbf{persistence diagram} records the pairs $(b, d)$ where:
\begin{itemize}
    \item $b$ (birth): Scale at which a homological cycle appears
    \item $d$ (death): Scale at which the cycle disappears
    \item Persistence: $\text{pers}(b,d) = d - b$
\end{itemize}
\end{definition}

\textbf{Applications in GDL.}\label{rem:persistencia_gdl} Persistent homology is used in:
\begin{itemize}
    \item \textbf{Featurization}: Persistence diagrams as features for CNNs
    \item \textbf{Topological pooling}: Preserving persistent structures
    \item \textbf{Regularization}: Penalizing undesired topologies
\end{itemize}

\subsubsection{Sheaf Theory}

Sheaf theory provides a framework for data distributed over topological spaces.

\begin{definition}[Cellular Sheaf]\label{def:sheaf_celular}
A \textbf{cellular sheaf} $\mathcal{F}$ over a cell complex $X$ assigns:
\begin{itemize}
    \item To each cell $\sigma \in X$ a vector space $\mathcal{F}(\sigma)$ (stalk)
    \item To each pair $\tau \subseteq \sigma$ a restriction map $\mathcal{F}_{\sigma \to \tau}: \mathcal{F}(\sigma) \to \mathcal{F}(\tau)$
\end{itemize}
satisfying composition properties.
\end{definition}

\textbf{Sheaf Neural Networks.}\label{rem:sheaf_nns} The Sheaf Neural Networks of \citet{hansen2021sheaf} and \citet{bodnar2023cellular} use sheaves to model:
\begin{itemize}
    \item \textbf{Heterogeneity}: Different feature spaces in each cell
    \item \textbf{Alignment}: Restriction maps as alignment transformations
    \item \textbf{Coherence}: Minimizing misalignment in the message passing
\end{itemize}

\begin{example}[Sheaf Laplacian]\label{ex:sheaf_laplacian}
For a sheaf $\mathcal{F}$ over a graph $G$, the \textbf{sheaf Laplacian} is:
\[L_{\mathcal{F}} = \delta^T \delta\]
where $\delta$ incorporates the restriction maps:
\[(\delta f)_e = \mathcal{F}_{t(e) \to e}(f_{t(e)}) - \mathcal{F}_{s(e) \to e}(f_{s(e)})\]
for an edge $e$ with source $s(e)$ and target $t(e)$.

It generalizes the standard graph Laplacian by allowing heterogeneous feature spaces.
\end{example}

\subsubsection{Learnable Topological Invariants}

Advanced architectures can learn to compute topological invariants such as:
\begin{itemize}
    \item \textbf{Betti numbers:} Dimensions of homology groups (\autoref{def:numeros_betti_cnn})
    \item \textbf{Euler characteristic:} Alternating sum of Betti numbers: $\chi = \sum_k (-1)^k \beta_k$
    \item \textbf{Characteristic classes:} Differentiable invariants of bundles
    \item \textbf{Persistence diagrams:} Multiscale representations of topology
\end{itemize}

These invariants provide robust and geometrically meaningful representations for learning.

\subsubsection{Categorical Perspectives}

Category theory provides a unifying language that connects Clifford algebras, topological complexes, and neural network architectures.

\textbf{Functors and Equivalences.}\label{rem:funtores_equivalencias} The functors between geometric categories formalize the transformations of representation:
\begin{itemize}
    \item \textbf{Chains functor}: $C_*: \textbf{Top} \to \textbf{Chain}$ (spaces to chain complexes)
    \item \textbf{Homology functor}: $H_*: \textbf{Chain} \to \textbf{GrVect}$ (chains to graded spaces)
    \item \textbf{Clifford functor}: Associates to each Riemannian manifold its Clifford bundle
\end{itemize}

The advanced algebraic extensions (Clifford, topological, categorical) face computational challenges of scalability and numerical stability in high-dimensional algebras, and require further theoretical development in expressive capabilities and convergence guarantees.

\section{Special Methods and Connections}\label{sec:special_methods}

This final section of the chapter explores specialized methods that complement the geometric convolutional techniques presented earlier, establishing connections with other learning paradigms and analyzing future perspectives of the field. The topics covered include advanced positional encoding, controlled symmetry breaking, applications in quantum computing, and scalability considerations that determine the practical viability of these methods.

\subsection{Positional Encoding and Fourier Features}

\subsubsection{Foundations of Positional Encoding}

Positional encoding represents a fundamental challenge in geometric learning: how to represent spatial information in a way that preserves both the local structure and the global relations. Traditional CNNs incorporate positional information implicitly through the neighborhood structure of convolutions, but in more general geometric contexts this approach proves insufficient.

The seminal work of \citet{tancik2020fourier} demonstrated that standard neural networks have difficulty learning functions that vary rapidly in low-dimensional domains. This problem arises because networks with \gls{ReLU} activations are naturally biased toward smooth functions, a limitation that manifests especially in tasks that require capturing high-frequency details.

\begin{theorem}[Frequency Bias in Neural Networks]
Let $f: \mathbb{R}^d \to \mathbb{R}$ be a target function with significant spectral components at high frequencies. A neural network $\mathcal{F}_\theta$ trained by gradient descent in the parameter space $\theta$ exhibits a bias toward the approximation of low-frequency components of $f$ during the first iterations of training.
\end{theorem}

This bias originates in the weight initialization and the nature of gradient descent, which tends to explore first the components of the function space of lower complexity. Mathematically, consider the Fourier expansion of the target function:
\begin{equation}
f(\mathbf{x}) = \sum_{\omega \in \mathbb{Z}^d} \hat{f}(\omega) e^{2\pi i \omega \cdot \mathbf{x}}
\end{equation}

During the first epochs of training, the network preferentially approximates the coefficients $\hat{f}(\omega)$ with small $\|\omega\|$, postponing the capture of high-frequency components.

\subsubsection{Random Fourier Features: Foundational Theory}

The pioneering work of Rahimi \& Recht (NIPS 2007) established the theoretical foundations of random Fourier features for approximating kernels in large-scale kernel machines. This approach has become a fundamental tool for positional encoding in geometric learning.

\begin{theorem}[Bochner's Theorem]\label{thm:bochner}
A continuous kernel $k: \mathbb{R}^d \times \mathbb{R}^d \to \mathbb{R}$ is translation invariant (i.e., $k(\mathbf{x}, \mathbf{y}) = k(\mathbf{x} - \mathbf{y})$) and positive definite if and only if there exists a finite non-negative probability measure $\mu$ such that:
\begin{equation}
k(\mathbf{x} - \mathbf{y}) = \int_{\mathbb{R}^d} e^{2\pi i \omega^T (\mathbf{x} - \mathbf{y})} d\mu(\omega)
\end{equation}
\end{theorem}

\begin{proof}
The forward direction follows from basic properties of the Fourier transform. For the reverse direction, the function $\phi(\mathbf{x}) = k(\mathbf{x}, \mathbf{0}) = k(\mathbf{x})$ is continuous and positive definite. By the classical Bochner theorem, $\phi$ admits a representation as the Fourier transform of a finite non-negative measure.
\end{proof}

Bochner's theorem provides the mathematical justification for approximating kernels by Monte Carlo sampling of their Fourier transform. If $p(\omega)$ is the density associated with $\mu$, we can approximate:
\begin{equation}
k(\mathbf{x}, \mathbf{y}) \approx \frac{1}{m} \sum_{j=1}^m \cos(2\pi \omega_j^T \mathbf{x}) \cos(2\pi \omega_j^T \mathbf{y}) + \sin(2\pi \omega_j^T \mathbf{x}) \sin(2\pi \omega_j^T \mathbf{y})
\end{equation}
where $\omega_j \sim p(\omega)$ are independently sampled frequencies.

\paragraph{Construction of the Feature Mapping}

We define the random Fourier features mapping as:
\begin{equation}
z(\mathbf{x}) = \sqrt{\frac{2}{m}} \begin{bmatrix} \cos(2\pi \omega_1^T \mathbf{x}) \\ \vdots \\ \cos(2\pi \omega_m^T \mathbf{x}) \\ \sin(2\pi \omega_1^T \mathbf{x}) \\ \vdots \\ \sin(2\pi \omega_m^T \mathbf{x}) \end{bmatrix} \in \mathbb{R}^{2m}
\end{equation}

This mapping satisfies the fundamental property:
\begin{equation}
\mathbb{E}_{\omega \sim p}[z(\mathbf{x})^T z(\mathbf{y})] = k(\mathbf{x}, \mathbf{y})
\end{equation}

\begin{theorem}[Convergence of Random Fourier Features]
Let $k$ be a Gaussian kernel $k(\mathbf{x}, \mathbf{y}) = \exp(-\|\mathbf{x} - \mathbf{y}\|^2 / 2\sigma^2)$ and $\hat{k}_m$ its approximation by $m$ random Fourier features. Then, for any $\epsilon, \delta > 0$, if:
\begin{equation}
m \geq \frac{4}{\epsilon^2} \log\left(\frac{2n}{\delta}\right)
\end{equation}
it holds that with probability at least $1 - \delta$:
\begin{equation}
\sup_{\mathbf{x}, \mathbf{y} \in \mathcal{X}} |k(\mathbf{x}, \mathbf{y}) - \hat{k}_m(\mathbf{x}, \mathbf{y})| \leq \epsilon
\end{equation}
where $\mathcal{X} \subset \mathbb{R}^d$ is a compact set of $n$ points.

See Section~\ref{sec:convergencia_rff} of \autoref{ch:teoremas_complementarios} for the complete proof.
\end{theorem}

This result guarantees uniform convergence of the approximation with rate $O(1/\sqrt{m})$, providing rigorous theoretical guarantees for the use of random Fourier features in positional encoding.

\subsubsection{Fourier Features as Universal Encoding}

The solution proposed by Tancik et al. consists of transforming the input coordinates by means of Fourier features before processing them with the neural network. For an input $\mathbf{x} \in \mathbb{R}^d$, the Fourier features transformation is defined as:

\begin{equation}
\gamma(\mathbf{x}) = \left[\sin(2\pi \mathbf{B}\mathbf{x}), \cos(2\pi \mathbf{B}\mathbf{x})\right]^T
\end{equation}

where $\mathbf{B} \in \mathbb{R}^{m \times d}$ is a frequency matrix sampled from a Gaussian distribution $\mathcal{N}(0, \sigma^2)$. The parameter $\sigma$ controls the frequency spectrum: small values favor smooth functions while large values allow high-frequency variations to be captured.

\begin{definition}[Fourier Features Mapping]
A Fourier features mapping $\gamma: \mathbb{R}^d \to \mathbb{R}^{2m}$ is a transformation that projects spatial coordinates to a feature space where periodic functions of different frequencies can be represented linearly.
\end{definition}

\paragraph{Spectral Analysis and Choice of $\sigma$}

The choice of the scale parameter $\sigma$ in the frequency distribution is crucial. For Gaussian RBF kernels $k(\mathbf{x}, \mathbf{y}) = \exp(-\|\mathbf{x} - \mathbf{y}\|^2 / 2\gamma^2)$, the optimal frequency distribution is $p(\omega) \propto \exp(-\|\omega\|^2 \gamma^2 / 2)$, that is, $\omega \sim \mathcal{N}(0, \gamma^{-2} \mathbf{I})$.

In practice, $\sigma$ is selected by cross-validation or following heuristics based on the characteristic scale of the data:
\begin{equation}
\sigma_{\text{optimal}} \approx \frac{1}{\text{median}(\{\|\mathbf{x}_i - \mathbf{x}_j\| : i, j \in [n]\})}
\end{equation}

\paragraph{Worked Example: Encoding in $\mathbb{R}^2$}

Consider spatial coordinates $\mathbf{x} = (x, y) \in [0, 1]^2$ and $m = 3$ frequency components. We sample:
\begin{align}
\omega_1 &= (2.1, -1.3)^T, \quad \omega_2 = (-0.8, 3.4)^T, \quad \omega_3 = (1.7, 2.2)^T
\end{align}

The Fourier encoding for a point $\mathbf{x}_0 = (0.5, 0.5)$ is:
\begin{align}
\gamma(\mathbf{x}_0) = [&\sin(2\pi \cdot 0.4), \sin(2\pi \cdot 1.3), \sin(2\pi \cdot 1.95), \\
&\cos(2\pi \cdot 0.4), \cos(2\pi \cdot 1.3), \cos(2\pi \cdot 1.95)]^T
\end{align}

This encoding captures information about the position at multiple frequency scales simultaneously.

\subsubsection{Connection with Harmonic Analysis on Manifolds}

This technique has profound implications for geometric CNNs. On Riemannian manifolds, the geodesic coordinates can be enriched with Fourier features to capture high-frequency local variations while preserving the global structure of the manifold.

Harmonic analysis on manifolds provides the theoretical framework for extending random Fourier features to non-Euclidean domains. Let $(\mathcal{M}, g)$ be a compact Riemannian manifold with Laplace-Beltrami operator $\Delta_{\mathcal{M}}$. The eigenfunctions $\{\phi_k\}$ of the Laplacian form a complete orthonormal basis of $L^2(\mathcal{M})$:
\begin{equation}
\Delta_{\mathcal{M}} \phi_k = -\lambda_k \phi_k, \quad \int_{\mathcal{M}} \phi_i \phi_j \, dV_g = \delta_{ij}
\end{equation}

\begin{proposition}[Fourier Features on Manifolds]
For a compact Riemannian manifold $\mathcal{M}$, a feature mapping based on eigenfunctions of the Laplacian:
\begin{equation}
\gamma_{\mathcal{M}}(p) = \sqrt{\frac{2}{m}} \begin{bmatrix} \phi_1(p) \\ \vdots \\ \phi_m(p) \end{bmatrix}
\end{equation}
approximates diffusion kernels $k_t(p, q) = \sum_{k=1}^\infty e^{-t\lambda_k} \phi_k(p)\phi_k(q)$ with error $O(e^{-t\lambda_{m+1}})$.
\end{proposition}

\begin{proof}
The diffusion kernel expands as $k_t(p,q) = \sum_{k=1}^{\infty} e^{-t\lambda_k} \phi_k(p) \phi_k(q)$, where $\{\lambda_k, \phi_k\}$ are the eigenpairs of the Laplace-Beltrami operator on $\mathcal{M}$ with $0 \leq \lambda_1 \leq \lambda_2 \leq \cdots$. The approximation with $m$ eigenfunctions truncates the series: $\hat{k}_t(p,q) = \sum_{k=1}^{m} e^{-t\lambda_k} \phi_k(p) \phi_k(q)$. The truncation error is:
\[
|k_t(p,q) - \hat{k}_t(p,q)| = \left|\sum_{k=m+1}^{\infty} e^{-t\lambda_k} \phi_k(p) \phi_k(q)\right| \leq e^{-t\lambda_{m+1}} \sum_{k=m+1}^{\infty} |\phi_k(p) \phi_k(q)|
\]
where we used $e^{-t\lambda_k} \leq e^{-t\lambda_{m+1}}$ for $k > m$. The convergence of the residual series is guaranteed by Weyl's law ($\lambda_k \sim k^{2/\dim \mathcal{M}}$) and the uniform bound on the eigenfunctions on compact manifolds.
\end{proof}

This result establishes that the eigenfunctions of the Laplacian on manifolds generalize the sinusoidal functions of the Euclidean case, providing a unified theory of positional encoding.

\subsubsection{Integration with Equivariant CNNs}

The combination of Fourier features with equivariant CNNs requires special care to preserve the symmetry properties. Let $g \in G$ be a transformation of the symmetry group and $\mathbf{x}$ a spatial coordinate. For the encoding to be equivariant, it must satisfy:

\begin{equation}
\gamma(g \cdot \mathbf{x}) = \rho(g) \gamma(\mathbf{x})
\end{equation}

where $\rho: G \to \text{GL}(\mathbb{R}^{2m})$ is a representation of the group in the feature space.

For rotation groups $\text{SO}(2)$, this is achieved using circular harmonics as a basis for the Fourier features. Specifically, for $G = \text{SO}(2)$, the frequencies $\omega$ must be selected so that $\omega \cdot (R\mathbf{x}) = (R^T\omega) \cdot \mathbf{x}$ for every rotation $R \in \text{SO}(2)$.

For rotations in 3D ($G = \text{SO}(3)$), the spherical harmonics $Y_\ell^m(\theta, \phi)$ provide the natural equivariant basis. The encoding is constructed as:
\begin{equation}
\gamma_{\text{SO}(3)}(\mathbf{x}) = [Y_0^0(\theta, \phi), Y_1^{-1}(\theta, \phi), Y_1^0(\theta, \phi), Y_1^1(\theta, \phi), \ldots]^T
\end{equation}
where $(\theta, \phi)$ are spherical coordinates associated with $\mathbf{x}$.

\paragraph{Equivariance for the Euclidean Group $E(3)$}

For the full Euclidean group $E(3) = \mathbb{R}^3 \rtimes \text{SO}(3)$ (rotations and translations), the construction is more complex. Features are used that separate magnitude information (invariant to rotations) and direction (equivariant):
\begin{equation}
\gamma_{E(3)}(\mathbf{x}) = [\|\mathbf{x}\|, Y_\ell^m(\hat{\mathbf{x}})]
\end{equation}
where $\hat{\mathbf{x}} = \mathbf{x}/\|\mathbf{x}\|$ is the unit direction.

\subsubsection{Applications in Differential Geometry}

In the context of manifolds, Fourier features can be adapted to respect the local Riemannian metric. Given a coordinate atlas $\{(U_i, \phi_i)\}$ for a manifold $\mathcal{M}$, the features are constructed as:

\begin{equation}
\gamma_{\mathcal{M}}(p) = \gamma(\phi_i(p)) \sqrt{|\det(g_{ij}(p))|}
\end{equation}

where $g_{ij}$ is the metric tensor in local coordinates. This modification ensures that the features are invariant under coordinate changes and reflect the intrinsic geometry of the manifold.

\paragraph{Gauge Theory and Fourier Features}

In principal bundles, Fourier features can incorporate gauge connection information. Let $(P, \mathcal{M}, \pi, G)$ be a principal bundle with connection $\omega$. The gauge-equivariant positional encoding is defined by means of:
\begin{equation}
\gamma_{\text{gauge}}(p) = \Gamma_\gamma^{-1} \gamma(\pi(p))
\end{equation}
where $\Gamma_\gamma$ is the parallel transport along a curve $\gamma$ from a base point to $\pi(p)$.

This construction guarantees that the features transform appropriately under local gauge transformations, preserving the geometric structure of the bundle.

\subsection{Controlled Symmetry Breaking}

\subsubsection{Theoretical Motivation}

Although equivariance is desirable for the generalization and data efficiency of models, in many real applications the symmetries are approximate or there is privileged information that breaks the natural symmetries. The work of \citet{vanderouderaa2024symmetry} formalizes the concept of controlled symmetry breaking as a regularization technique that allows models to adapt to deviations from perfect symmetry.

\begin{definition}[Controlled Symmetry Breaking]
A controlled symmetry breaking is an intentional modification of an equivariant architecture that introduces a specific directional bias while maintaining the ability to revert to equivariant behavior when the data require it.
\end{definition}

The mathematical formulation is based on the interpolation between equivariant and invariant operators. Let $\mathcal{F}_{\text{eq}}: X \to Y$ be an equivariant function and $\mathcal{F}_{\text{in}}: X \to Y$ an invariant function. The function with controlled breaking is defined as:

\begin{equation}
\mathcal{F}_\lambda = (1-\lambda) \mathcal{F}_{\text{eq}} + \lambda \mathcal{F}_{\text{in}}
\end{equation}

where $\lambda \in [0,1]$ is a parameter that controls the degree of symmetry breaking.

\paragraph{Group Theory Foundations}

To formalize controlled breaking mathematically, let $G$ be a group of symmetries acting on the input space $X$ and $\rho: G \to \text{GL}(Y)$ a representation in the output space. A function $f: X \to Y$ is $G$-equivariant if:
\begin{equation}
f(g \cdot x) = \rho(g) f(x), \quad \forall g \in G, x \in X
\end{equation}

Symmetry breaking introduces a parametrized family $f_\lambda$ that interpolates between full equivariance ($\lambda = 0$) and symmetry independence ($\lambda = 1$):
\begin{equation}
f_\lambda(g \cdot x) = (1-\lambda) \rho(g) f_\lambda(x) + \lambda h(g \cdot x)
\end{equation}
where $h: X \to Y$ is an auxiliary function without equivariance constraints.

\begin{proposition}[Smoothness of the Transition]
The family $\{f_\lambda\}_{\lambda \in [0,1]}$ defines a continuous curve in the function space. If $f_0$ is $G$-equivariant and $h$ is arbitrary, then:
\begin{equation}
\lim_{\lambda \to 0} \|f_\lambda(g \cdot x) - \rho(g) f_\lambda(x)\| = 0
\end{equation}
uniformly over compacts in $X \times G$.
\end{proposition}

\begin{proof}
Let $f_\lambda = (1-\lambda) f_0 + \lambda h$ with $f_0$ equivariant and $h$ arbitrary. For $g \in G$:
\begin{align}
\|f_\lambda(g \cdot x) - \rho(g) f_\lambda(x)\| &= \|(1-\lambda)(f_0(g \cdot x) - \rho(g) f_0(x)) + \lambda(h(g \cdot x) - \rho(g) h(x))\| \\
&= \lambda \|h(g \cdot x) - \rho(g) h(x)\|
\end{align}
where we used the exact equivariance of $f_0$. Over a compact $K \subset X \times G$, the continuous function $(x,g) \mapsto \|h(g \cdot x) - \rho(g) h(x)\|$ is bounded by some $M < \infty$. Thus, $\|f_\lambda(g \cdot x) - \rho(g) f_\lambda(x)\| \leq \lambda M \to 0$ as $\lambda \to 0$, uniformly over $K$.
\end{proof}

\subsubsection{Implementation in Geometric CNNs}

In the context of CNNs on manifolds, symmetry breaking can be implemented in several ways:

\textbf{1. Breaking in the Feature Space:} A directional bias is introduced in the intermediate representations:

\begin{equation}
\mathbf{h}_{l+1} = \sigma(\mathbf{W}_l * \mathbf{h}_l + \mathbf{b}_l + \lambda_l \mathbf{v}_{\text{bias}})
\end{equation}

where $\mathbf{v}_{\text{bias}}$ is a bias vector that breaks the symmetry in a controlled way.

\textbf{2. Breaking in the Convolutions:} The convolutional kernels include asymmetric components:

\begin{equation}
K_{\text{sb}}(\mathbf{x}, \mathbf{y}) = K_{\text{sym}}(\mathbf{x}, \mathbf{y}) + \lambda K_{\text{asym}}(\mathbf{x}, \mathbf{y})
\end{equation}

where $K_{\text{sym}}$ preserves the symmetries of the problem and $K_{\text{asym}}$ introduces a specific directionality.

\textbf{3. Adaptive Breaking:} The parameter $\lambda$ is learned during training, allowing the model to automatically determine the optimal degree of symmetry breaking for each layer or region of the input space.

The loss function for adaptive breaking incorporates a regularization term:
\begin{equation}
\mathcal{L}_{\text{total}} = \mathcal{L}_{\text{task}} + \alpha \sum_{l=1}^L \lambda_l^2
\end{equation}
where $\alpha > 0$ controls the bias toward equivariant solutions.

\paragraph{Variational Formulation}

From a Bayesian perspective, controlled breaking can be formulated as variational inference over the function space. We define a prior distribution $p(\lambda)$ that favors small values of $\lambda$:
\begin{equation}
p(\lambda) = \text{Beta}(\alpha, \beta), \quad \alpha > \beta > 0
\end{equation}

The posterior distribution $p(\lambda | \mathcal{D})$ is inferred from the data by means of:
\begin{equation}
p(\lambda | \mathcal{D}) \propto p(\mathcal{D} | \lambda) p(\lambda)
\end{equation}

This approach allows the uncertainty about the optimal degree of symmetry breaking to be quantified.

\subsubsection{Regularization Analysis}

Controlled symmetry breaking acts as a regularization mechanism that prevents overfitting to perfect symmetries that may not exist in real data. From an information-theoretic perspective, symmetry breaking introduces additional capacity to the model while maintaining an inductive bias toward symmetric solutions.

\begin{theorem}[Regularizing Effect of Symmetry Breaking]
Let $\mathcal{H}_{\text{eq}}$ be the space of equivariant hypotheses and $\mathcal{H}_{\lambda}$ the space with controlled breaking. Then:
\[\text{VC-dim}(\mathcal{H}_{\lambda}) = \text{VC-dim}(\mathcal{H}_{\text{eq}}) + O(\lambda \log(1/\lambda))\]
\end{theorem}

\begin{proof}[Proof sketch]
The Vapnik-Chervonenkis dimension measures the capacity of the hypothesis space. For $\lambda = 0$, we recover $\mathcal{H}_{\text{eq}}$. As $\lambda$ increases, additional directions are added in the function space. The logarithmic growth derives from the analysis of the covering of orbits under the action of $G$.
\end{proof}

This result shows that controlled breaking expands the capacity of the model logarithmically in the breaking parameter, providing additional flexibility without drastically compromising the generalization capacity.

\paragraph{Bias-Variance Trade-off}

Symmetry breaking modifies the classical bias-variance trade-off of statistical learning. For an estimator $\hat{f}_\lambda$ with symmetry breaking $\lambda$, the expected error decomposes as:
\begin{align}
\mathbb{E}[\|f - \hat{f}_\lambda\|^2] &= \underbrace{\|f - \mathbb{E}[\hat{f}_\lambda]\|^2}_{\text{Bias}^2} + \underbrace{\mathbb{E}[\|\hat{f}_\lambda - \mathbb{E}[\hat{f}_\lambda]\|^2]}_{\text{Variance}} \\
&= B(\lambda)^2 + V(\lambda)
\end{align}

where typically $B(\lambda)$ decreases and $V(\lambda)$ increases with $\lambda$, suggesting the existence of an optimal value $\lambda^*$.

\subsubsection{Practical Applications}

\paragraph{Natural Images with Directional Biases}

In photography, although objects can have intrinsic symmetry, factors such as directional lighting and gravity break the symmetry systematically. A CNN with controlled breaking can capture these effects:
\begin{equation}
\mathbf{h}^{(l)} = \text{ReLU}(W_{\text{eq}}^{(l)} * \mathbf{h}^{(l-1)} + \lambda_l \mathbf{b}_{\text{gravity}})
\end{equation}
where $\mathbf{b}_{\text{gravity}}$ encodes a vertical bias consistent with gravity.

\paragraph{Molecular Data with Approximate Symmetries}

In computational chemistry, molecules can have nominal symmetries that are broken by environmental or quantum effects. GNNs with controlled breaking model these deviations:
\begin{equation}
\mathbf{h}_v^{(l+1)} = \sum_{u \in \mathcal{N}(v)} [(1-\lambda) \phi_{\text{eq}}(\mathbf{h}_u^{(l)}, \mathbf{e}_{vu}) + \lambda \phi_{\text{asym}}(\mathbf{h}_u^{(l)}, \mathbf{e}_{vu}, \mathbf{c}_{\text{env}})]
\end{equation}
where $\mathbf{c}_{\text{env}}$ encodes environmental conditions that break the symmetry.

\subsection{Connections with Quantum Computing}

\subsubsection{Equivariant Quantum CNNs}

The intersection between geometric deep learning and quantum computing represents an emerging frontier with transformative potential. The work of \citet{west2024quantum} explores how the principles of equivariance extend naturally to the quantum regime, where quantum states and unitary operations respect fundamental symmetries of physics.

An equivariant quantum CNN operates in the Hilbert space $\mathcal{H} = (\mathbb{C}^2)^{\otimes n}$, where the quantum states $|\psi\rangle$ encode geometric information. Equivariance manifests through commutation with quantum symmetry operators:

\begin{equation}
[U_g, \mathcal{C}_\text{quantum}] = 0
\end{equation}

where $U_g$ is the unitary representation of the group element $g$ and $\mathcal{C}_\text{quantum}$ is the quantum convolutional operation.

\subsubsection{Embedding of Geometric Data}

The embedding of classical data in quantum states requires specialized techniques that preserve the geometric structure. For data on manifolds, the embedding is performed by means of:

\begin{equation}
|\psi(\mathbf{x})\rangle = \sum_{i=1}^{2^n} \alpha_i(\mathbf{x}) |i\rangle
\end{equation}

where the amplitudes $\alpha_i(\mathbf{x})$ encode the local coordinates in the manifold so that the geometric symmetries translate into symmetries of the quantum state.

\textbf{Fundamental Advantages of the Quantum Approach:}

\begin{itemize}
\item \textbf{Coherent Superposition:} The superposition principle allows the simultaneous evaluation of geometric configurations in the Hilbert space $\mathcal{H} = \bigotimes_{i=1}^n \mathbb{C}^{d_i}$.
\item \textbf{Geometric Entanglement:} Entangled states $|\psi\rangle \in \mathcal{H}_A \otimes \mathcal{H}_B$ naturally encode non-separable correlations between geometrically distant regions.
\item \textbf{Exponential Encoding:} The dimension of the Hilbert space grows exponentially, allowing dense representations of complex geometric structures using $\log_2(\dim \mathcal{H})$ qubits.
\end{itemize}

\subsubsection{Measurement and Decoherence}

The measurement process in quantum CNNs must be designed carefully to extract relevant geometric information while minimizing the loss of quantum coherence. Projective measurements are used in bases that respect the symmetries of the problem:

\begin{equation}
P_{\text{sym}} = \frac{1}{|G|} \sum_{g \in G} U_g P_0 U_g^\dagger
\end{equation}

where $P_0$ is a base projector and the sum over the group $G$ ensures that the measurement is invariant under the relevant symmetries.

\subsection{Scalability and Efficiency Analysis}

\subsubsection{The Equivariance vs. Scalability Dilemma}

A fundamental theoretical result established by \citet{park2024equivariancescale} examines the saturation of the advantages of equivariance in high-dimensionality regimes. This analysis reveals an intrinsic tension between geometric constraints and expressive capacity in hypothesis spaces of great complexity.

\begin{hypothesis}[Equivariance Saturation Hypothesis]
There exists a scale regime where the benefits of equivariance saturate and the additional capacity of the model can compensate for the loss of geometric inductive bias.
\end{hypothesis}

This saturation has theoretical foundations based on universal approximation theory and the Vapnik-Chervonenkis dimension:

\begin{itemize}
\item The symmetry orbits in high-dimensional spaces require exponential coverage
\item The Rademacher complexity of equivariant functions is related to the size of the group
\item There is a fundamental trade-off between expressive capacity and geometric constraints
\end{itemize}

\paragraph{Quantitative Analysis of Orbit Coverage}

For a finite group $G$ acting on $\mathbb{R}^d$, the number of distinct orbits in a ball of radius $R$ grows as:
\begin{equation}
|\{\text{Orbits in } B_R\}| \sim \frac{\text{Vol}(B_R)}{|G| \cdot \text{Vol}(\text{fundamental domain})}
\end{equation}

In high dimensions ($d \to \infty$), the volume $\text{Vol}(B_R) = \frac{\pi^{d/2}}{\Gamma(d/2 + 1)} R^d$ grows exponentially, requiring an exponential number of parameters to cover all orbits with fixed resolution.

\begin{proposition}[Parametric Cost of Equivariance]
Let $\mathcal{F}_{\text{eq}}$ be a family of $G$-equivariant functions and $\mathcal{F}_{\text{free}}$ functions without constraints. To achieve approximation error $\epsilon$ over a compact domain, the number of required parameters satisfies:
\begin{equation}
\frac{N_{\text{eq}}}{N_{\text{free}}} \geq \left(\frac{|G|}{\text{const}}\right)^{1/d}
\end{equation}
in the high-dimensional limit.
\end{proposition}

\begin{proof}[Proof sketch]
A $G$-equivariant function $f: \mathbb{R}^d \to \mathbb{R}^d$ satisfies $f(gx) = gf(x)$ for all $g \in G$, which imposes $|G|$ constraints for each point of the domain. These constraints reduce the effective degrees of freedom by a factor $|G|$, so that the space $\mathcal{F}_{\text{eq}}$ has dimension $\sim N/|G|$ per unit of volume of the domain. However, to \textit{approximate} arbitrary functions (not necessarily equivariant), an equivariant network needs to compensate for the imposed constraints. In high dimension $d$, the ``curse of dimensionality'' implies that the number of parameters needed to cover the compact domain grows exponentially in $d$, and the reduction by symmetry is attenuated as $(|G|/\text{const})^{1/d}$ by the packing properties in high dimension.
\end{proof}

This result suggests that in sufficiently high dimensions, equivariant architectures require significantly more parameters to achieve the same expressive capacity.

\subsubsection{Computational Complexity: Big-O Analysis}

We analyze the temporal and spatial complexity of different geometric architectures.

\paragraph{Graph Convolutional Networks (GCN)}

For a graph with $|V| = n$ vertices and $|E| = m$ edges:

\textbf{Temporal Complexity:}
\begin{itemize}
\item A GCN layer: $O(m \cdot d_{\text{in}} \cdot d_{\text{out}})$
\item Complete network ($L$ layers): $O(L \cdot m \cdot d^2)$ where $d = \max(d_{\text{in}}, d_{\text{out}})$
\end{itemize}

\textbf{Spatial Complexity:}
\begin{itemize}
\item Storage of representations: $O(L \cdot n \cdot d)$
\item Adjacency matrix (dense): $O(n^2)$
\item Adjacency matrix (sparse): $O(m)$
\end{itemize}

\paragraph{Graph Attention Networks (GAT)}

\textbf{Temporal Complexity:}
\begin{itemize}
\item Attention computation: $O(m \cdot d + n \cdot h \cdot d^2)$ where $h$ is the number of heads
\item Multi-head with $h$ heads: $O(h \cdot m \cdot d + n \cdot h \cdot d^2)$
\item Complete network: $O(L \cdot h \cdot m \cdot d^2)$
\end{itemize}

\textbf{Spatial Complexity:}
\begin{itemize}
\item Attention matrices: $O(h \cdot m)$
\item Cache of intermediate scores: $O(h \cdot n \cdot d)$
\end{itemize}

\paragraph{GraphGPS (Graph GPS Transformer)}

GraphGPS combines local \gls{MPNN} with global attention:

\textbf{Temporal Complexity:}
\begin{itemize}
\item Local \gls{MPNN}: $O(m \cdot d^2)$
\item Global attention: $O(n^2 \cdot d)$ (complete) or $O(n \cdot k \cdot d)$ (approximate with $k$ neighbors)
\item Spectral positional encoding: $O(n^3)$ (precomputation) or $O(n \cdot \log n \cdot d)$ (approximation)
\item Total per layer: $O(m \cdot d^2 + n \cdot k \cdot d + n \cdot \log n \cdot d)$
\end{itemize}

\textbf{Spatial Complexity:}
\begin{itemize}
\item Positional features: $O(n \cdot d_{\text{pos}})$
\item Attention: $O(n \cdot k)$ (approximate) or $O(n^2)$ (complete)
\end{itemize}

\subsubsection{Comparative Table of Architectures}

\begin{table}[h]
\centering
\begin{tabular}{l|c|c|c|c}
\toprule
\textbf{Architecture} & \textbf{Time/layer} & \textbf{Space} & \textbf{Expressiveness} & \textbf{Equivariance} \\
\midrule
GCN & $O(m \cdot d^2)$ & $O(n \cdot d)$ & Medium & Permutation \\
GAT & $O(h \cdot m \cdot d^2)$ & $O(h \cdot m)$ & High & Permutation \\
GraphGPS & $O(n^2 \cdot d)$ & $O(n^2)$ & Very high & Permutation \\
G-CNN & $O(|G| \cdot n \cdot d^2)$ & $O(|G| \cdot n \cdot d)$ & High & $G$-equivariant \\
Steerable CNN & $O(n \cdot d^2 \cdot r^3)$ & $O(n \cdot d)$ & High & SO(3) \\
\bottomrule
\end{tabular}
\caption[Computational complexity of geometric architectures]{Comparison of computational complexity and properties of different geometric architectures. Here $n$ = number of nodes, $m$ = number of edges, $d$ = feature dimension, $h$ = number of heads, $|G|$ = order of the group, $r$ = angular resolution.}
\label{tab:complejidad_arquitecturas_geometricas}
\end{table}

\subsubsection{Metrics of Geometric Efficiency}\label{subsubsec:metricas_eficiencia_geometrica}

To quantify the efficiency of geometric CNNs, the following metrics are proposed:

\textbf{1. Parametric Compression Index:}
\begin{equation}
\text{PCI} = \frac{\dim(\mathcal{H}_{\text{free}})}{\dim(\mathcal{H}_{G\text{-eq}})}
\end{equation}
where $\mathcal{H}_{\text{free}}$ and $\mathcal{H}_{G\text{-eq}}$ are the free and $G$-equivariant hypothesis spaces respectively.

For finite groups, $\text{PCI} = |G|$ approximately. For continuous groups such as $\text{SO}(3)$, the compression depends on the parametrization:
\begin{equation}
\text{PCI}_{\text{SO}(3)} \approx \frac{(2\ell_{\max} + 1)^2}{2\ell_{\max} + 1} = 2\ell_{\max} + 1
\end{equation}
where $\ell_{\max}$ is the maximum degree of spherical harmonics used.

\textbf{2. Geometric Generalization Gain:}
\begin{equation}
\text{GGG} = \frac{\mathcal{R}_{\text{emp}}[\mathcal{H}_{\text{free}}] - \mathcal{R}_{\text{emp}}[\mathcal{H}_{G\text{-eq}}]}{\sqrt{\frac{\log |G|}{n}}}
\end{equation}
where $\mathcal{R}_{\text{emp}}$ is the empirical risk and $n$ the sample size.

This metric normalizes the improvement in generalization by the size of the group and the sample size, providing a fair measure of the geometric benefit.

\textbf{3. Relative Algorithmic Complexity:}
\begin{equation}
\text{RAC} = \frac{\mathcal{C}(\mathcal{A}_{G\text{-eq}})}{\mathcal{C}(\mathcal{A}_{\text{standard}}) \cdot \mathcal{T}(\epsilon)}
\end{equation}
where $\mathcal{C}(\mathcal{A})$ is the algorithmic complexity of the method $\mathcal{A}$ and $\mathcal{T}(\epsilon)$ represents the convergence speedup factor to precision $\epsilon$.

\paragraph{Quantitative Example: Point Cloud Classification}

For the ModelNet40 dataset with $n = 10{,}000$ points per cloud:

\begin{itemize}
\item \textbf{PointNet (baseline)}:
  \begin{itemize}
  \item Complexity: $O(n \cdot d^2)$
  \item PCI: $1$ (no equivariance)
  \item Accuracy: $89.2\%$
  \end{itemize}

\item \textbf{PointNet++ with SO(3) equivariance}:
  \begin{itemize}
  \item Complexity: $O(n \cdot d^2 \cdot (2\ell_{\max} + 1)^2)$ with $\ell_{\max} = 2$
  \item PCI: $5$
  \item Accuracy: $93.7\%$ ($+4.5\%$)
  \item GGG: $\frac{4.5}{\sqrt{\log(5) / 10000}} \approx 352.6$
  \end{itemize}
\end{itemize}

This demonstrates that the generalization gain (GGG) is substantial, justifying the additional cost of equivariance.

\subsubsection{Efficient Scaling Strategies}

To address the scalability challenges, several strategies have been developed:

\textbf{1. Hierarchical Equivariance:} We define a sequence of groups $G_1 \supset G_2 \supset \cdots \supset G_L$ where each layer $l$ maintains $G_l$-equivariance, allowing controlled relaxation of constraints.

\textbf{Example:} In image processing:
\begin{itemize}
\item Early layers: $G_1 = p4m$ (symmetries of the full square)
\item Intermediate layers: $G_2 = p4$ (only $90^\circ$ rotations)
\item Final layers: $G_3 = \{e\}$ (no equivariance)
\end{itemize}

\textbf{2. $\epsilon$-Approximate Equivariance:} For $\epsilon > 0$, a function $f$ is $\epsilon$-equivariant if:
\begin{equation}
\sup_{g \in G} \|f(g \cdot x) - \rho(g)f(x)\| \leq \epsilon
\end{equation}

This provides theoretical flexibility with approximation guarantees. In practice, it is implemented by means of:
\begin{equation}
f_\epsilon(x) = \frac{1}{|S|} \sum_{g \in S} \rho(g)^{-1} f_{\text{approx}}(g \cdot x)
\end{equation}
where $S \subset G$ is a finite subset and $f_{\text{approx}}$ is not exactly equivariant.

\begin{proposition}[Error of $\epsilon$-Equivariance]
If $f_{\text{approx}}$ is $\epsilon$-equivariant and $|S| \geq \log(1/\delta) / \epsilon^2$, then with probability $1 - \delta$:
\begin{equation}
\|f_\epsilon(g \cdot x) - \rho(g) f_\epsilon(x)\| \leq 2\epsilon
\end{equation}
\end{proposition}

\begin{proof}
Let $f_\epsilon$ be an $\epsilon$-equivariant function, that is, $\|f_\epsilon(g \cdot x) - \rho(g) f_\epsilon(x)\| \leq \epsilon$ for all $g \in G$ and $x$ in the domain. Let $S = \{x_1, \ldots, x_{|S|}\}$ be a set of i.i.d. samples from the domain. By the triangle inequality, for any $g \in G$ and $x$ in the domain:
\[
\|f_\epsilon(g \cdot x) - \rho(g) f_\epsilon(x)\| \leq \|f_\epsilon(g \cdot x) - f_\epsilon(g \cdot x_i)\| + \|f_\epsilon(g \cdot x_i) - \rho(g) f_\epsilon(x_i)\| + \|\rho(g)(f_\epsilon(x_i) - f_\epsilon(x))\|
\]
where $x_i$ is the sample closest to $x$. The second term is bounded by $\epsilon$. For $|S| \geq \log(1/\delta)/\epsilon^2$, the concentration of measure guarantees that the discretization error (first and third terms) contributes at most an additional $\epsilon$ with probability $1 - \delta$, giving the total bound $2\epsilon$.
\end{proof}

\textbf{3. Measurable Conditional Equivariance:} It is defined by means of a measurable selection function $\chi: X \to \{0,1\}$ such that:
\begin{equation}
f(g \cdot x) = \chi(x) \rho(g)f(x) + (1-\chi(x))h(g \cdot x)
\end{equation}
where $h$ has no equivariance constraints.

The function $\chi$ is learned jointly with the model parameters, allowing the model to dynamically decide when to apply equivariance constraints.

\textbf{4. Low-Rank Approximation Techniques:}

To reduce the attention complexity in GraphGPS, low-rank factorizations are used:
\begin{equation}
\mathbf{A} = \mathbf{U} \mathbf{V}^T, \quad \mathbf{U}, \mathbf{V} \in \mathbb{R}^{n \times r}, \quad r \ll n
\end{equation}

This reduces the complexity from $O(n^2 \cdot d)$ to $O(n \cdot r \cdot d)$ with controlled error:
\begin{equation}
\|\mathbf{A} - \mathbf{U}\mathbf{V}^T\|_F \leq \sqrt{\sum_{i=r+1}^n \sigma_i^2}
\end{equation}
where $\sigma_i$ are singular values of $\mathbf{A}$.

\paragraph{Algorithmic Optimization: Case Study}

Consider the optimization of an equivariant attention layer in $\text{SO}(3)$:

\textbf{Naive Approach:}
\begin{itemize}
\item Discretize $\text{SO}(3)$ with $N_{\text{rot}}$ rotations
\item Complexity: $O(N_{\text{rot}} \cdot n^2 \cdot d)$
\item For $N_{\text{rot}} = 72$ ($5^\circ$ resolution): $\sim 10^6$ operations per point
\end{itemize}

\textbf{Optimized Approach (spherical harmonics):}
\begin{itemize}
\item Use spherical harmonics of degree $\leq \ell_{\max}$
\item Complexity: $O(n^2 \cdot d \cdot \ell_{\max}^2)$
\item For $\ell_{\max} = 3$: $\sim 10^4$ operations per point
\item \textbf{Speedup}: $100\times$
\end{itemize}

This example illustrates how the mathematical theory (representation theory) leads directly to efficient algorithms.

\subsection{Unifying Synthesis and Future Perspectives}

\subsubsection{Unified Theoretical Framework}

The methods presented in this chapter can be unified under a general theoretical framework that combines:

\begin{itemize}
\item \textbf{Algebraic Structure:} Symmetry groups, Lie algebras, and representations
\item \textbf{Differential Geometry:} Manifolds, connections, and parallel transport
\item \textbf{Harmonic Analysis:} Fourier transforms, spherical harmonics, and wavelets
\item \textbf{Approximation Theory:} Representation capacity and convergence
\end{itemize}

\begin{definition}[Unified Geometric CNN]
A unified geometric CNN is a triple $(\mathcal{M}, G, \Phi)$ where:
\begin{itemize}
\item $\mathcal{M}$ is a Riemannian manifold (possibly with group structure)
\item $G$ is a group of symmetries that acts on $\mathcal{M}$
\item $\Phi$ is a family of equivariant differential operators that generalize the convolution
\end{itemize}
\end{definition}

\subsubsection{Future Research Directions}

\textbf{1. Geometric Approximation Theory:} Development of universal approximation theorems for CNNs on non-compact manifolds and spaces with singularities.

\textbf{2. Optimization on Manifolds:} Optimization algorithms that respect the geometry of the parameter space, especially for models with equivariance constraints.

\textbf{3. Symmetry Learning:} Methods that can automatically discover the relevant symmetries in the data without prior knowledge.

\textbf{4. Quantum Scalability:} Development of quantum CNNs that can process large-scale classical data while maintaining quantum advantage.

\textbf{5. Scientific Applications:} Extension of geometric methods to problems in physics, chemistry, and molecular biology where symmetries are fundamental.

\subsubsection{Connection with Geometric Attention Models}

The developments presented in this chapter establish essential foundations that connect directly with the geometric attention mechanisms explored in \autoref{ch:attention}. This connection is not accidental but reflects a deep convergence in the paradigm of Geometric Deep Learning.

\paragraph{Conceptual Convergence CNN-Attention}

The fundamental concepts developed in the preceding sections find direct parallels in geometric attention:

\begin{itemize}
\item \textbf{\gls{equivarianza} (Section~\ref{subsec:equivarianza_invarianza})}: The equivariant attention mechanisms of \autoref{ch:attention} extend the equivariance constraints of CNNs by means of adaptive weights that respect symmetries.

\item \textbf{\gls{analisis_armonico} (\autoref{sec:analisis_armonico})}: The spectral techniques developed here provide the mathematical basis for the spectral positional encoding in Graph Attention Networks (\autoref{sec:graph_attention}).

\item \textbf{Group convolutions (\autoref{sec:group_convolutions})}: G-CNNs anticipate the formulation of attention over homogeneous spaces.

\item \textbf{Gauge equivariance (\autoref{sec:gauge_cnns})}: The principal bundles introduced here reappear in Geometric Algebra Transformers (\autoref{sec:clifford_attention}), where gauge transformations are realized by means of geometric products and Clifford operations.
\end{itemize}

\paragraph{Unifying Framework: Geometric Message Passing}

Geometric CNNs can be reinterpreted as special cases of the \gls{marco_unificador} of message passing that is formalized completely in \autoref{ch:attention}:

\begin{definition}[CNN as Structured Message Passing]\label{def:cnn_as_structured_mp}
A geometric CNN over domain $\mathcal{D}$ with symmetry group $G$ is a Message Passing Neural Network where:
\begin{itemize}
    \item The underlying graph reflects the geometric structure of $\mathcal{D}$
    \item The message functions preserve the equivariance to $G$
    \item The aggregation respects the local geometry (convolution → locally weighted sum)
    \item The weights are shared according to the orbits of $G$
\end{itemize}
\end{definition}

This perspective, developed completely in \autoref{subsubsec:unified_framework_advanced}, shows how the structural rigidity of CNNs generalizes toward the adaptive flexibility of attention.

\paragraph{Limitations that Motivate Attention}

The limitations identified in traditional CNNs (Section~\ref{subsec:fundamentos_limitaciones_cnn}) directly motivate the developments of \autoref{ch:attention}:

\begin{itemize}
\item \textbf{Long-range dependencies}: The inherent locality of convolutions (limited receptive field) is overcome by means of global attention mechanisms while preserving equivariance. Graph Transformers and Geometric Transformers implement global attention over graphs and manifolds respectively, allowing arbitrarily distant dependencies to be captured without sacrificing geometric properties.

\item \textbf{Irregular geometries}: Complex manifolds with variable curvature and non-trivial topology require adaptive attention patterns that generalize the fixed convolutional kernels. In \autoref{ch:attention}, attention kernels are developed that adapt to the local curvature by means of parallel transport and Riemannian metrics.

\item \textbf{Dimensionality and scalability}: Attention-based methods allow the treatment of high-dimensional spaces by means of low-rank approximation techniques and matrix factorizations, as analyzed in §\ref{subsubsec:metricas_eficiencia_geometrica}. These techniques are implemented in Linformer, Performer, and Flash Attention, demonstrating that scalability does not require sacrificing geometric expressiveness.

\item \textbf{Data heterogeneity}: CNNs assume homogeneous data in regular domains. Attention mechanisms, especially in Graph Attention Networks and Heterogeneous Graph Transformers, naturally handle heterogeneous data with multiple types of nodes and edges, each with its own feature space and symmetries.
\end{itemize}

\paragraph{Theoretical Bridge: From Convolution to Attention}

The mathematical transition from convolution to attention can be formalized by means of:

\begin{theorem}[Attention as Generalization of Convolution]
A geometric convolution operation $f * g$ on a manifold $\mathcal{M}$ can be expressed as a limit case of attention:
\begin{equation}
(f * g)(x) = \lim_{\beta \to \infty} \int_{\mathcal{M}} \frac{e^{-\beta d(x,y)^2}}{\int_{\mathcal{M}'} e^{-\beta d(x,y')^2} dy'} g(y) f(y) \, dy
\end{equation}
where $d$ is the Riemannian distance and the limit $\beta \to \infty$ recovers locally compact support.
\end{theorem}

\begin{proof}[Proof sketch]
The Gaussian kernel $K_\beta(x,y) = e^{-\beta d(x,y)^2}$ defines a family of distributions parametrized by $\beta > 0$. For $\beta \to \infty$, $K_\beta(x,y) \to 0$ exponentially fast for all $y \neq x$, while $K_\beta(x,x) = 1$. The normalized distribution $\frac{K_\beta(x,y)}{\int K_\beta(x,y') dy'}$ converges weakly to the Dirac delta $\delta_x(y)$, recovering the locality of the convolution. For finite $\beta$, the attention weights distribute mass in a geodesic neighborhood of radius $\sim 1/\sqrt{\beta}$, producing a ``smoothed'' convolution whose support narrows with $\beta$. See \autoref{ch:attention} for the complete proof.
\end{proof}

This result, proved formally in \autoref{ch:attention}, establishes that attention with Gaussian kernel $K(x,y) = \exp(-\beta d(x,y)^2)$ interpolates continuously between global attention ($\beta \to 0$) and local convolution ($\beta \to \infty$).

\subsubsection{Interdisciplinary Connections}

The field of geometric deep learning is evolving toward a truly interdisciplinary discipline that connects:

\begin{itemize}
\item \textbf{Mathematics:} Differential geometry, representation theory, harmonic analysis, algebraic topology
\item \textbf{Theoretical Physics:} Conformal field theory, gauge theory, geometry of curved spacetimes
\item \textbf{Applied Mathematics:} Approximation theory, functional analysis, dynamical systems
\item \textbf{Logic and Foundations:} Category theory, homological algebra, non-commutative geometry
\end{itemize}

The convergence of these fields suggests that geometric deep learning is not simply a computational tool, but a unifying mathematical language for the extraction of structure from complex data.
